%============================================================
%  CHAPTER 1 — INTRODUCTION TO GRAPHS
%  Entry file: chapter opener + \input's only.
%
%  Folder layout
%     lessons/               Part A: what a graph is, vocabulary, storage, a first traversal, modelling
%     p1_degrees/            theory, 2 solved problems, exercises
%     p2_reachability/
%     p3_grids/
%     p4_state_graphs/
%     p5_functional/
%     summary.tex, problem_set.tex, problem_set_solutions.tex
%============================================================
\chapter{Introduction to Graphs}
\label{ch:introduction}

\begin{tcolorbox}[enhanced, colback=white, colframe=secondaryblue, boxrule=1pt, arc=2mm,
                  left=10pt, right=10pt]
\textbf{\large\textcolor{primaryblue}{Learn to see graphs everywhere: in maps and friendships, but also in grids, numbers, words and game states. Once a problem is drawn as vertices and edges, a handful of simple ideas --- degrees, reachability, flood fill, state search and following arrows --- already solve a surprising number of contest problems.}}
\end{tcolorbox}

\chapterfacts{Beginner}{1400}{5}{11}{28}

\subsection*{How This Chapter Is Organised}

\begin{diagrambox}
\begin{tikzpicture}[node distance=3mm and 3mm,
    th/.style={rectangle, rounded corners=3pt, draw=theoryborder, fill=theorybg, line width=1pt,
               font=\small\bfseries, align=center, minimum height=10mm, text width=13.2cm},
    pt/.style={rectangle, rounded corners=3pt, draw=patternborder, fill=patternbg, line width=1pt,
               font=\scriptsize, align=center, minimum height=12mm, text width=2.35cm},
    pr/.style={rectangle, rounded corners=2pt, draw=primaryblue, fill=white,
               font=\scriptsize, align=center, text width=2.35cm, minimum height=6mm}]
    \node[th] (T) {Part A --- What a Graph Is and How to Use One (\S\ref{sec:ch1-lessons})\\[-1pt]
        {\scriptsize\mdseries vertices \& edges $\cdot$ vocabulary $\cdot$ storing a graph $\cdot$ a first traversal $\cdot$ modelling}};
    \node[pt, below=6mm of T.south west, anchor=north west, xshift=1mm] (P1) {\textbf{P1 Degrees}\\\textbf{\& Counting}\\no traversal needed};
    \node[pt, right=of P1] (P2) {\textbf{P2 Reach-}\\\textbf{ability}\\who can reach whom};
    \node[pt, right=of P2] (P3) {\textbf{P3 Grids}\\\textbf{as Graphs}\\cells \& neighbours};
    \node[pt, right=of P3] (P4) {\textbf{P4 Implicit}\\\textbf{State Graphs}\\numbers \& states};
    \node[pt, right=of P4] (P5) {\textbf{P5 Functional}\\\textbf{Graphs}\\one arrow out};
    \foreach \p/\a/\b in {P1/Find the Town Judge/Maximal Network Rank, P2/Find if Path Exists/Keys and Rooms,
                          P3/Counting Rooms/Max Area of Island, P4/Two Buttons/Open the Lock,
                          P5/Badge/Longest Cycle in a Graph}{
        \node[pr, below=of \p] (\p a) {\a};
        \node[pr, below=2mm of \p a] (\p b) {\b};
    }
    \node[pr, below=2mm of P5b] (P5c) {Find the Duplicate};
    \foreach \p in {P1,P2,P3,P4,P5} \draw[flow] (T.south -| \p.north) -- (\p.north);
    \node[th, anchor=north, at={([yshift=-5mm]P3b.south |- P5c.south)}, draw=successgreen, fill=solutionbg] (S) {Summary, Cheat Sheet, Practice Set \& Solution Sketches};
    \node[font=\small\bfseries\color{patternborder}, rotate=90, anchor=south] at ([xshift=-3mm]P1a.west |- P1a.south) {Part B};
\end{tikzpicture}
\end{diagrambox}

\subsection*{Before You Start}

You need basic C++ (vectors, queues, recursion) and the idea of time complexity. No graph theory is assumed: every term is defined in Part~A. Depth-first and breadth-first search are introduced informally in Lesson~\ref{ssec:ch1-first-traversal} and studied properly in Chapters~2 and~3.

\subsection*{After This Chapter You Will Be Able To}
\begin{enumerate}[leftmargin=*]
    \item \textbf{Speak graph}: vertices, edges, degrees, paths, cycles, connectivity, trees, DAGs, bipartite and functional graphs.
    \item \textbf{Store} a graph as an edge list, adjacency matrix or adjacency list, and read contest input into it correctly.
    \item \textbf{Answer} degree questions without any traversal, and reachability questions with a simple DFS or BFS.
    \item \textbf{Turn} grids, numbers, locks and other state spaces into graphs, and bound their size before coding.
    \item \textbf{Follow arrows} in functional graphs: find where a walk enters its cycle, and measure cycles.
\end{enumerate}

\chaptercontents

%============================================================
\section{Part A --- What a Graph Is and How to Use One}
\label{sec:ch1-lessons}
%============================================================
%============================================================
% CH1 · LESSON 1 — VERTICES AND EDGES
%============================================================
\subsection{Vertices and Edges}
\label{ssec:ch1-vertices-edges}

\subsubsection*{The one-paragraph answer}

A graph is the simplest possible description of \emph{things and the connections between them}. The things are called \textbf{vertices} (or nodes), the connections are called \textbf{edges}. Cities and roads, people and friendships, web pages and links, courses and prerequisites, cells of a maze and the doors between them --- all of these are graphs. The power of the idea is that once you have drawn the dots and the lines, the \emph{meaning} of the dots no longer matters: the same few algorithms answer ``can I get from here to there?'', ``how many steps does it take?'', ``which groups fall apart?'' for roads, friendships and mazes alike.

\begin{definitionbox}[Definition (Graph)]
A \textbf{graph} $G=(V,E)$ consists of
\begin{itemize}[leftmargin=*]
    \item a finite set $V$ of \textbf{vertices}; we write $n=\nV$ for their number;
    \item a set $E$ of \textbf{edges}; each edge connects two vertices $u$ and $v$, its \textbf{endpoints}. We write $m=\nE$.
\end{itemize}
Two vertices joined by an edge are \textbf{adjacent} (or \textbf{neighbours}); the edge is \textbf{incident} to both of them. In contest input the vertices are almost always the integers $1,2,\dots,n$ and the edges arrive as a list of pairs.
\end{definitionbox}

\subsubsection*{A first picture: five towns}

Five towns A, B, C, D, E and the roads between them. The geography is gone; only ``who is connected to whom'' remains.

\begin{diagrambox}
\begin{tikzpicture}[x=1.3cm, y=1.1cm]
    % left: the "map"
    \begin{scope}
        \fill[insightbg, rounded corners=6pt] (-0.6,-1.7) rectangle (3.6,1.5);
        \node[font=\scriptsize\bfseries\color{darkgray}] at (1.5,1.25) {a map};
        \foreach \p/\n in {(0,0.6)/A, (1.3,0.9)/B, (2.9,0.5)/C, (0.6,-1.1)/D, (2.4,-1.0)/E}
            \node[rectangle, draw=darkgray, fill=white, font=\scriptsize\bfseries, inner sep=2pt] (m\n) at \p {\n};
        \draw[accentorange, line width=1.6pt, rounded corners=4pt] (mA) -- (0.7,0.9) -- (mB);
        \draw[accentorange, line width=1.6pt] (mB) .. controls (2.2,1.1) .. (mC);
        \draw[accentorange, line width=1.6pt] (mA) .. controls (0.1,-0.5) .. (mD);
        \draw[accentorange, line width=1.6pt] (mD) -- (mE);
        \draw[accentorange, line width=1.6pt] (mB) .. controls (1.6,-0.2) .. (mE);
        \draw[accentorange, line width=1.6pt] (mC) -- (mE);
    \end{scope}
    \draw[flow, line width=1.2pt] (3.9,0) -- node[above, font=\scriptsize] {keep only} node[below, font=\scriptsize] {connections} (5.3,0);
    % right: the graph
    \begin{scope}[xshift=7.2cm]
        \node[gv] (A) at (0,0.6) {A};
        \node[gv] (B) at (1.3,1.0) {B};
        \node[gv] (C) at (2.6,0.5) {C};
        \node[gv] (D) at (0.4,-1.1) {D};
        \node[gv] (E) at (2.0,-1.0) {E};
        \foreach \a/\b in {A/B, B/C, A/D, D/E, B/E, C/E} \draw[ge] (\a) -- (\b);
        \node[font=\scriptsize\bfseries\color{darkgray}] at (1.3,1.6) {a graph: $n=5$, $m=6$};
    \end{scope}
\end{tikzpicture}
\figcap{The road map becomes a graph with $V=\{A,B,C,D,E\}$ and $E=\{AB, BC, AD, DE, BE, CE\}$. Lengths, curves and positions are forgotten: the picture on the right could be redrawn in any layout and it would still be the same graph.}
\end{diagrambox}

\begin{insightbox}[The drawing is not the graph]
A graph is only the two sets $V$ and $E$. Where you put the dots and whether the lines cross is your choice; two very different-looking pictures can show the same graph. Draw the version that makes the structure obvious --- this is a skill, and every chapter of this book practises it.
\end{insightbox}

%------------------------------------------------------------
\subsubsection*{Undirected and directed edges}

A road usable in both directions is an \textbf{undirected} edge $\{u,v\}$: it can be walked from $u$ to $v$ and from $v$ to $u$. A one-way street, a ``follows'' relation on a social network or ``course $u$ must be taken before course $v$'' is a \textbf{directed} edge (an \textbf{arc}) $u\to v$: an ordered pair $(u,v)$ with a \textbf{tail} $u$ and a \textbf{head} $v$. A graph whose edges are all directed is a \textbf{directed graph} (\textbf{digraph}).

\begin{diagrambox}
\begin{tikzpicture}[x=1.25cm, y=1.1cm]
    \begin{scope}
        \node[gv] (1) at (0,1) {1};
        \node[gv] (2) at (1.6,1) {2};
        \node[gv] (3) at (0,-0.3) {3};
        \node[gv] (4) at (1.6,-0.3) {4};
        \foreach \a/\b in {1/2, 1/3, 2/3, 2/4, 3/4} \draw[ge] (\a) -- (\b);
        \node[font=\small\bfseries] at (0.8,-1.15) {undirected};
        \node[font=\scriptsize\color{darkgray}] at (0.8,-1.6) {friendship, two-way road};
    \end{scope}
    \begin{scope}[xshift=4.2cm]
        \node[gv] (1) at (0,1) {1};
        \node[gv] (2) at (1.6,1) {2};
        \node[gv] (3) at (0,-0.3) {3};
        \node[gv] (4) at (1.6,-0.3) {4};
        \foreach \a/\b in {1/2, 3/1, 2/3, 2/4, 4/3} \draw[gde] (\a) -- (\b);
        \node[font=\small\bfseries] at (0.8,-1.15) {directed};
        \node[font=\scriptsize\color{darkgray}] at (0.8,-1.6) {follows, one-way street};
    \end{scope}
    \begin{scope}[xshift=8.4cm]
        \node[gv] (1) at (0,1) {1};
        \node[gv] (2) at (1.6,1) {2};
        \draw[ge] (1) -- (2);
        \node[font=\scriptsize] at (0.8,0.35) {is stored as};
        \node[gv] (1b) at (0,-0.3) {1};
        \node[gv] (2b) at (1.6,-0.3) {2};
        \draw[gde] (1b) to[bend left=25] (2b);
        \draw[gde] (2b) to[bend left=25] (1b);
        \node[font=\small\bfseries] at (0.8,-1.15) {in code};
        \node[font=\scriptsize\color{darkgray}] at (0.8,-1.6) {one undirected = two arcs};
    \end{scope}
\end{tikzpicture}
\figcap{Left: from 2 you can go to 1, 3 or 4. Middle: from 2 you can go to 3 and 4 but \emph{not} to 1 (the arc is $1\to 2$). Right: programs store an undirected edge as two opposite arcs (Lesson~\ref{ssec:ch1-storing}).}
\end{diagrambox}

\begin{warningbox}[Read the statement for the direction]
``There is a road between $u$ and $v$'' is undirected; ``$u$ trusts $v$'', ``$u$ points to $v$'', ``from room $u$ you can enter room $v$'', ``$a_i$ is the next element after $i$'' are directed. Treating a directed edge as undirected (or forgetting the reverse copy of an undirected edge) is one of the most frequent wrong answers in graph problems.
\end{warningbox}

%------------------------------------------------------------
\subsubsection*{Weighted edges}

Often an edge carries a number: a road length, a ticket price, a travel time, a capacity. Such a graph is \textbf{weighted}, and we write $w(u,v)$ for the weight of edge $uv$. An \textbf{unweighted} graph behaves as if every edge had weight 1, so ``shortest'' then means ``fewest edges''.

\begin{diagrambox}
\begin{tikzpicture}[x=1.5cm, y=1.1cm]
    \node[gvsrc] (s) at (0,0) {S};
    \node[gv] (a) at (1.4,1) {A};
    \node[gv] (b) at (1.4,-1) {B};
    \node[gvdst] (t) at (2.8,0) {T};
    \draw[ge] (s) -- node[gw, above left] {2} (a);
    \draw[ge] (s) -- node[gw, below left] {5} (b);
    \draw[ge] (a) -- node[gw, above right] {7} (t);
    \draw[ge] (b) -- node[gw, below right] {1} (t);
    \draw[ge] (a) -- node[gw, right] {1} (b);
    \begin{scope}[xshift=5.5cm]
        \node[gvsrc] (s) at (0,0) {S};
        \node[gv] (a) at (1.4,1) {A};
        \node[gv] (b) at (1.4,-1) {B};
        \node[gvdst] (t) at (2.8,0) {T};
        \draw[gpath] (s) -- (a); \draw[gpath] (a) -- (b); \draw[gpath] (b) -- (t);
        \draw[gfaded] (s) -- (b); \draw[gfaded] (a) -- (t);
        \node[gw] at (0.6,0.65) {2}; \node[gw] at (1.62,0) {1}; \node[gw] at (2.2,-0.65) {1};
        \node[gdist, right=1mm of t] {$2+1+1=4$};
    \end{scope}
\end{tikzpicture}
\figcap{A weighted graph. The cheapest route from \legendchip{gvsrc}{S} to \legendchip{gvdst}{T} is S--A--B--T of total weight 4, although it uses three edges; the routes S--B--T (weight 6) and S--A--T (weight 9) use only two. Weights change the answer --- shortest paths with weights are the subject of Chapter~7.}
\end{diagrambox}

%------------------------------------------------------------
\subsubsection*{Simple graphs, multigraphs and self-loops}

Two more kinds of edges can appear in real input:
\begin{itemize}[leftmargin=*]
    \item a \textbf{self-loop} is an edge from a vertex to itself, $u$--$u$ (``room 3 has a door back into room 3'');
    \item \textbf{parallel edges} (multi-edges) are two or more edges between the same pair of vertices (two different roads between the same towns, maybe of different lengths).
\end{itemize}
A graph with neither is \textbf{simple}; a graph that allows them is a \textbf{multigraph}.

\begin{diagrambox}
\begin{tikzpicture}[x=1.25cm, y=1.1cm]
    \begin{scope}
        \node[gv] (1) at (0,1) {1};
        \node[gv] (2) at (1.6,1) {2};
        \node[gv] (3) at (1.6,-0.3) {3};
        \node[gv] (4) at (0,-0.3) {4};
        \foreach \a/\b in {1/2, 2/3, 3/4, 4/1} \draw[ge] (\a) -- (\b);
        \node[font=\small\bfseries] at (0.8,-1.15) {simple};
    \end{scope}
    \begin{scope}[xshift=5cm]
        \node[gv] (1) at (0,1) {1};
        \node[gv] (2) at (1.6,1) {2};
        \node[gvcur] (3) at (1.6,-0.3) {3};
        \node[gv] (4) at (0,-0.3) {4};
        \draw[ge] (1) -- (2);
        \draw[ge, draw=recursionred, line width=1.3pt] (2) to[bend left=22] (3);
        \draw[ge, draw=recursionred, line width=1.3pt] (2) to[bend right=22] (3);
        \draw[ge] (3) -- (4); \draw[ge] (4) -- (1);
        \draw[ge, draw=recursionred, line width=1.3pt] (3) to[out=-20, in=-80, looseness=7] (3);
        \node[font=\scriptsize\color{recursionred}, anchor=west] at (2.0,1.0) {parallel edges 2--3};
        \node[font=\scriptsize\color{recursionred}, anchor=west] at (2.3,-0.85) {self-loop 3--3};
        \node[font=\small\bfseries] at (0.8,-1.15) {multigraph};
    \end{scope}
\end{tikzpicture}
\figcap{Right: the input \texttt{1 2 / 2 3 / 3 2 / 3 3 / 3 4 / 4 1} ($n=4$, $m=6$). The edges \texttt{2 3} and \texttt{3 2} are the same undirected pair, so they are parallel; \texttt{3 3} is a self-loop.}
\end{diagrambox}

Most contest statements promise ``no self-loops and no multiple edges''. When they do \emph{not}, the promise is missing for a reason: the tests will contain them. Most algorithms in this book work on multigraphs unchanged, but counting problems (``how many neighbours does $v$ have?'') and adjacency matrices (one cell per pair) do not. The following program detects both kinds; it is a good first exercise in reading a graph.

\begin{lstlisting}[style=cpp]
#include <bits/stdc++.h>
using namespace std;

// Reads an undirected graph and reports whether it is simple.
// Input: n m, then m lines "u v" (1-indexed).
int main() {
    int n, m;
    if (!(cin >> n >> m)) return 0;
    set<pair<int,int>> seen;          // each edge stored as (smaller, larger)
    int loops = 0, parallel = 0;
    for (int i = 0; i < m; i++) {
        int u, v;
        cin >> u >> v;
        if (u == v) { loops++; continue; }        // self-loop u-u
        if (u > v) swap(u, v);                    // {u,v} == {v,u}
        if (!seen.insert(make_pair(u, v)).second) // already there?
            parallel++;                           // a second copy
    }
    cout << "self-loops: " << loops << "\n";
    cout << "parallel edges: " << parallel << "\n";
    cout << (loops == 0 && parallel == 0 ? "simple graph" : "multigraph") << "\n";
    return 0;
}
\end{lstlisting}

\begin{lstlisting}[style=trace]
input : 4 6   1 2   2 3   3 2   3 3   3 4   4 1
output: self-loops: 1
        parallel edges: 1
        multigraph
\end{lstlisting}

The line \texttt{if (u > v) swap(u, v)} deserves a name: \textbf{normalising} an undirected edge so that $\{u,v\}$ and $\{v,u\}$ look identical. You will use it whenever undirected edges go into a \texttt{set} or a \texttt{map}.

%------------------------------------------------------------
\subsubsection*{Graphs in real life and in contests}

\begin{table}[H]
\centering\small
\begin{tabularx}{\linewidth}{@{}l l l X@{}}
\toprule
\textbf{Situation} & \textbf{Vertices} & \textbf{Edges} & \textbf{Kind and typical question} \\
\midrule
Road network & towns, crossings & roads & undirected, weighted: shortest route (Ch.~7) \\
Social network & people & friendships / follows & undirected / directed: groups (Ch.~4) \\
The web & pages & hyperlinks & directed: which pages are reachable \\
Course plan & courses & ``$u$ before $v$'' & directed, no cycles: a valid order (Ch.~5) \\
Maze, grid map & free cells & shared sides & undirected, implicit: flood fill (Pattern~3) \\
Puzzle, lock & configurations & one move & implicit: fewest moves (Pattern~4) \\
Permutation $p$ & indices $i$ & $i\to p_i$ & one outgoing arc per vertex: cycles (Pattern~5) \\
Company & employees & boss $\to$ subordinate & a rooted tree (Ch.~13) \\
\bottomrule
\end{tabularx}
\end{table}

\begin{diagrambox}
\begin{tikzpicture}[x=1cm, y=1cm]
    % social
    \begin{scope}
        \node[gvA, minimum size=8.5mm] (a) at (0,0.8) {\scriptsize Ann};
        \node[gvA, minimum size=8.5mm] (b) at (1.4,1.2) {\scriptsize Bob};
        \node[gvA, minimum size=8.5mm] (c) at (1.7,-0.1) {\scriptsize Cy};
        \node[gvB, minimum size=8.5mm] (d) at (0.1,-0.7) {\scriptsize Dee};
        \node[gvB, minimum size=8.5mm] (e) at (1.5,-1.4) {\scriptsize Eve};
        \draw[ge] (a) -- (b); \draw[ge] (b) -- (c); \draw[ge] (a) -- (c);
        \draw[ge] (d) -- (e);
        \node[font=\scriptsize\bfseries] at (0.8,-2.2) {friends: two groups};
    \end{scope}
    % prerequisites
    \begin{scope}[xshift=4.2cm]
        \node[gv, minimum size=9mm] (p) at (0,1) {\scriptsize Prog};
        \node[gv, minimum size=9mm] (d) at (1.9,1) {\scriptsize DS};
        \node[gv, minimum size=9mm] (m) at (0,-0.8) {\scriptsize Math};
        \node[gv, minimum size=9mm] (a) at (1.9,-0.8) {\scriptsize Algo};
        \draw[gde] (p) -- (d); \draw[gde] (d) -- (a); \draw[gde] (m) -- (a);
        \node[font=\scriptsize\bfseries] at (0.95,-2.2) {prerequisites: a DAG};
    \end{scope}
    % maze
    \begin{scope}[xshift=8.6cm, yshift=-1.5cm]
        \foreach \r/\c/\t in {0/0/gcell, 0/1/gcell, 0/2/gwall, 1/0/gwall, 1/1/gcell, 1/2/gcell, 2/0/gcell, 2/1/gcell, 2/2/gwall}
            \node[\t, minimum size=7mm] at (\c*0.7, 2.8-\r*0.7) {};
        \draw[primaryblue, line width=1.2pt] (0,2.8) -- (0.7,2.8) -- (0.7,2.1) -- (1.4,2.1);
        \draw[primaryblue, line width=1.2pt] (0.7,2.1) -- (0.7,1.4) -- (0,1.4);
        \foreach \x/\y in {0/2.8, 0.7/2.8, 0.7/2.1, 1.4/2.1, 0.7/1.4, 0/1.4}
            \fill[primaryblue] (\x,\y) circle (2.2pt);
        \node[font=\scriptsize\bfseries] at (0.7,0.8) {maze: cells + doors};
    \end{scope}
\end{tikzpicture}
\figcap{Three faces of the same object. Friendships (undirected; the two colours show the two separate groups), prerequisites (directed, no cycle), and a maze whose free cells become vertices joined when they share a side (dark cells are walls).}
\end{diagrambox}

%------------------------------------------------------------
\subsubsection*{How a graph arrives in contest input}

The standard format is ``$n$ $m$'' on the first line, followed by $m$ lines with one edge each (and a weight, if the graph is weighted). LeetCode passes the same information as \texttt{n} and a vector \texttt{edges} of pairs. Learn to \emph{draw} such input on paper in under a minute; it is the first thing to do with every sample.

\begin{diagrambox}
\begin{tikzpicture}
    \node[draw=darkgray, fill=lightgray!60, rounded corners=2pt, font=\ttfamily\small, align=left, inner sep=5pt] (in) at (0,0)
        {6 5\\1 2\\2 3\\1 3\\4 5\\5 6};
    \node[font=\scriptsize\color{darkgray}, above=1mm of in] {input};
    \draw[flow, line width=1.2pt] (0.9,0) -- (2.1,0);
    \begin{scope}[xshift=3.0cm, yshift=-0.2cm]
        \node[gv] (1) at (0,0.9) {1};
        \node[gv] (2) at (1.3,0.9) {2};
        \node[gv] (3) at (0.65,-0.3) {3};
        \node[gv] (4) at (2.7,0.9) {4};
        \node[gv] (5) at (4.0,0.9) {5};
        \node[gv] (6) at (3.35,-0.3) {6};
        \foreach \a/\b in {1/2, 2/3, 1/3, 4/5, 5/6} \draw[ge] (\a) -- (\b);
    \end{scope}
    \node[note, text width=4cm] at (10.0,0.3) {$n=6$ vertices, $m=5$ edges. The graph falls into two pieces: $\{1,2,3\}$ and $\{4,5,6\}$.};
\end{tikzpicture}
\figcap{Read the first line, draw $n$ dots, then one line per edge. Nothing says the graph is in one piece --- here it is not.}
\end{diagrambox}

%------------------------------------------------------------
\subsubsection*{Drawing conventions used in this book}

Every diagram in the book uses the same visual vocabulary, so a picture can be read without its caption once you know it. Algorithms colour vertices as they progress; the colours always mean the same thing.

\begin{diagrambox}
\begin{tikzpicture}[x=1cm, y=1cm]
    \foreach \s/\t/\d [count=\i from 0] in {gv/plain/a vertex, gvsrc/source/start here, gvdst/target/answer, gvcur/current/being processed,
                                         gvseen/discovered/queued or stacked, gvdone/finished/fully explored, gvdead/removed/ignored, gvsel/selected/chosen set}{
        \pgfmathsetmacro{\xx}{mod(\i,4)*3.6}
        \pgfmathsetmacro{\yy}{-floor(\i/4)*1.3}
        \node[\s] at (\xx,\yy) {v};
        \node[anchor=west, font=\scriptsize\bfseries] at (\xx+0.5,\yy+0.15) {\t};
        \node[anchor=west, font=\scriptsize\color{darkgray}] at (\xx+0.5,\yy-0.2) {\d};
    }
    \node[gvA] at (0,-2.6) {v}; \node[anchor=west, font=\scriptsize\bfseries] at (0.5,-2.45) {side A};
    \node[anchor=west, font=\scriptsize\color{darkgray}] at (0.5,-2.8) {colour class 1};
    \node[gvB] at (3.6,-2.6) {v}; \node[anchor=west, font=\scriptsize\bfseries] at (4.1,-2.45) {side B};
    \node[anchor=west, font=\scriptsize\color{darkgray}] at (4.1,-2.8) {colour class 2};
    \node[gcell] at (7.2,-2.6) {}; \node[anchor=west, font=\scriptsize\bfseries] at (7.7,-2.45) {grid cell};
    \node[anchor=west, font=\scriptsize\color{darkgray}] at (7.7,-2.8) {free square};
    \node[gwall] at (10.8,-2.6) {}; \node[anchor=west, font=\scriptsize\bfseries] at (11.3,-2.45) {wall};
    \node[anchor=west, font=\scriptsize\color{darkgray}] at (11.3,-2.8) {blocked};
\end{tikzpicture}
\figcap{Vertex and cell styles. A traversal turns vertices from plain to \emph{discovered} (orange) to \emph{current} (yellow) to \emph{finished} (blue).}
\end{diagrambox}

\begin{diagrambox}
\begin{tikzpicture}[x=1cm, y=1cm]
    \foreach \s/\t/\d [count=\i from 0] in {ge/edge/undirected, gde/arc/directed, gtree/tree edge/used by a search, gpath/path/chosen route,
                                         gback/back edge/closes a cycle, gfaded/faded/not relevant now, gcut/cut/separates parts}{
        \pgfmathsetmacro{\xx}{mod(\i,4)*3.6}
        \pgfmathsetmacro{\yy}{-floor(\i/4)*1.1}
        \draw[\s] (\xx-0.4,\yy) -- (\xx+0.5,\yy);
        \node[anchor=west, font=\scriptsize\bfseries] at (\xx+0.6,\yy+0.15) {\t};
        \node[anchor=west, font=\scriptsize\color{darkgray}] at (\xx+0.6,\yy-0.2) {\d};
    }
    \node[gw] at (10.8,-1.1) {5};
    \node[anchor=west, font=\scriptsize\bfseries] at (11.1,-0.95) {weight};
    \node[anchor=west, font=\scriptsize\color{darkgray}] at (11.1,-1.3) {$w(u,v)=5$};
\end{tikzpicture}
\figcap{Edge styles and labels. A small green number next to a vertex, like \tikz[baseline=-0.6ex]\node[gdist]{3};, is a distance computed by an algorithm.}
\end{diagrambox}

%------------------------------------------------------------
\subsubsection*{Check your understanding}

\begin{enumerate}[leftmargin=*]
    \item Draw the graph given by \texttt{n = 4, edges = [[1,2],[2,3],[3,1]]}. How many vertices have no neighbour?
    \item ``Each of the $n$ employees except the director has exactly one direct boss.'' What are the vertices and edges, and should the edges be directed?
    \item Input \texttt{3 4 / 1 2 / 2 1 / 2 3 / 3 3} describes an undirected graph. Is it simple? If not, what makes it a multigraph?
    \item A map has 6 towns and every pair of towns is joined by a two-way road. How many edges? How many arcs if every road is replaced by two one-way roads?
    \item In a weighted graph, is the path with the fewest edges always the cheapest path? Give a reason or a counter-example from this lesson.
\end{enumerate}

\begin{solutionbox}[Answers]
\begin{enumerate}[leftmargin=*]
    \item A triangle 1--2--3 plus the isolated vertex 4: \textbf{one} vertex (4) has no neighbour. Isolated vertices are easy to forget when $n$ is given separately from the edges.
    \item Vertices = employees, one edge per employee--boss pair, $n-1$ edges. Direct them (boss $\to$ employee, or employee $\to$ boss) because ``is the boss of'' is not symmetric. This graph is a \emph{rooted tree} (Lesson~\ref{ssec:ch1-vocabulary}).
    \item Not simple: \texttt{1 2} and \texttt{2 1} are parallel edges and \texttt{3 3} is a self-loop. The program above prints \texttt{self-loops: 1}, \texttt{parallel edges: 1}.
    \item $\binom{6}{2}=15$ undirected edges; $6\cdot 5=30$ arcs.
    \item No. In the weighted figure above, S--A--T and S--B--T have 2 edges and weights 9 and 6, while S--A--B--T has 3 edges and weight 4.
\end{enumerate}
\end{solutionbox}

%============================================================
% CH1 · LESSON 2 — THE VOCABULARY OF GRAPHS
%============================================================
\subsection{The Vocabulary of Graphs}
\label{ssec:ch1-vocabulary}

Graph theory has a lot of words, but each one names a simple picture. This lesson is a picture dictionary: every term comes with a drawing, and the terms are ordered so that each one uses only the earlier ones. You do not have to memorise it in one go --- come back whenever a later chapter uses a word you are unsure of.

%------------------------------------------------------------
\subsubsection*{Degree}

The \textbf{degree} $\deg(v)$ of a vertex in an undirected graph is the number of edges touching it (a self-loop counts twice, because both of its ends touch $v$). In a simple graph this is the number of neighbours. A vertex of degree 0 is \textbf{isolated}; a vertex of degree 1 is a \textbf{leaf}.

We use one running example, the ``bow tie'' $G_1$ with $n=6$ and $m=7$:

\begin{diagrambox}
\begin{tikzpicture}[x=1.25cm, y=1cm]
    \node[gv] (1) at (0,0.8) {1};
    \node[gv] (2) at (0,-0.8) {2};
    \node[gvsel] (3) at (1.3,0) {3};
    \node[gvsel] (4) at (2.8,0) {4};
    \node[gv] (5) at (4.1,0.8) {5};
    \node[gv] (6) at (4.1,-0.8) {6};
    \foreach \a/\b in {1/2, 1/3, 2/3, 3/4, 4/5, 4/6, 5/6} \draw[ge] (\a) -- (\b);
    \foreach \v/\d/\pos in {1/2/left, 2/2/left, 3/3/above, 4/3/above, 5/2/right, 6/2/right}
        \node[gdist, \pos=1.5mm of \v] {deg $\d$};
\end{tikzpicture}
\figcap{$G_1$: edges 12, 13, 23, 34, 45, 46, 56. The two \legendchip{gvsel}{selected} vertices 3 and 4 have degree 3; all others have degree 2. Sum of degrees: $2+2+3+3+2+2=14=2\cdot 7$.}
\end{diagrambox}

In a \textbf{directed} graph each vertex has an \textbf{in-degree} $\indeg{v}$ (arcs entering $v$) and an \textbf{out-degree} $\outdeg{v}$ (arcs leaving $v$).

\begin{diagrambox}
\begin{tikzpicture}[x=1.4cm, y=1.1cm]
    \node[gv] (1) at (0,1) {1};
    \node[gv] (2) at (2,1) {2};
    \node[gv] (3) at (0,-0.6) {3};
    \node[gv] (4) at (2,-0.6) {4};
    \draw[gde] (1) -- (2);
    \draw[gde] (1) to[bend right=25] (3);
    \draw[gde] (3) to[bend right=25] (1);
    \draw[gde] (2) -- (3);
    \draw[gde] (3) -- (4);
    \draw[gde] (4) -- (2);
    \foreach \v/\i/\o/\pos in {1/1/2/left, 2/2/1/right, 3/2/2/left, 4/1/1/right}
        \node[gdist, \pos=1.5mm of \v, align=center] {in \i\\out \o};
    \node[note, text width=4.6cm, anchor=west] at (3.3,0.2) {arcs $1{\to}2$, $1{\to}3$, $2{\to}3$, $3{\to}1$, $3{\to}4$, $4{\to}2$\\[2pt]$\sum \text{in} = 1{+}2{+}2{+}1 = 6$\\$\sum \text{out} = 2{+}1{+}2{+}1 = 6 = m$};
\end{tikzpicture}
\figcap{Every arc has exactly one tail and one head, so the in-degrees and the out-degrees each add up to $m$.}
\end{diagrambox}

Degrees are the cheapest information a graph offers: one pass over the edge list computes all of them, without any search. Pattern~1 of this chapter solves whole problems with nothing else.

\begin{lstlisting}[style=cpp]
#include <bits/stdc++.h>
using namespace std;

// Degrees of an undirected graph + a check of the handshake lemma.
// Input: n m, then m lines "u v" (1-indexed).
int main() {
    int n, m;
    if (!(cin >> n >> m)) return 0;
    vector<int> deg(n + 1, 0);
    for (int i = 0; i < m; i++) {
        int u, v;
        cin >> u >> v;
        deg[u]++;            // the edge touches u ...
        deg[v]++;            // ... and v (a self-loop u==v adds 2 to u)
    }
    long long sum = 0;
    int odd = 0;
    for (int v = 1; v <= n; v++) {
        cout << "deg(" << v << ") = " << deg[v] << "\n";
        sum += deg[v];
        if (deg[v] % 2 == 1) odd++;
    }
    cout << "sum of degrees = " << sum << ", 2m = " << 2LL * m << "\n";
    cout << "odd-degree vertices = " << odd << "\n";
    return 0;
}
\end{lstlisting}

\begin{lstlisting}[style=trace]
input : 6 7   1 2  1 3  2 3  3 4  4 5  4 6  5 6        (the bow tie G1)
output: deg(1) = 2   deg(2) = 2   deg(3) = 3   deg(4) = 3   deg(5) = 2   deg(6) = 2
        sum of degrees = 14, 2m = 14
        odd-degree vertices = 2
\end{lstlisting}

%------------------------------------------------------------
\subsubsection*{The handshake lemma}

\begin{theorembox}[Handshake lemma]
In every undirected graph (self-loops and parallel edges allowed)
\[
    \sum_{v\in V} \deg(v) \;=\; 2\,|E|.
\]
\textbf{Proof.} Cut every edge in the middle into two \emph{half-edges}, one attached to each endpoint (a self-loop gives two half-edges at the same vertex). The degree of $v$ is exactly the number of half-edges at $v$, so the left side counts all half-edges. Every edge produced exactly two of them, so there are $2|E|$ half-edges. \hfill$\square$
\end{theorembox}

\begin{diagrambox}
\begin{tikzpicture}[x=1.25cm, y=1cm]
    \node[gv] (1) at (0,0.8) {1};
    \node[gv] (2) at (0,-0.8) {2};
    \node[gv] (3) at (1.3,0) {3};
    \node[gv] (4) at (2.8,0) {4};
    \node[gv] (5) at (4.1,0.8) {5};
    \node[gv] (6) at (4.1,-0.8) {6};
    \foreach \a/\b in {1/2, 1/3, 2/3, 3/4, 4/5, 4/6, 5/6}{
        \draw[draw=primaryblue, line width=1.4pt] (\a) -- ($(\a)!0.5!(\b)$);
        \draw[draw=accentorange, line width=1.4pt] ($(\a)!0.5!(\b)$) -- (\b);
        \fill[white] ($(\a)!0.5!(\b)$) circle (2.2pt);
    }
    \node[note, text width=4.2cm, anchor=west] at (5,0) {7 edges $\Rightarrow$ 14 half-edges.\\Each vertex owns as many half-edges as its degree.};
\end{tikzpicture}
\figcap{The proof in one picture: every edge of $G_1$ is cut into a blue and an orange half; counting halves by vertex gives $\sum\deg(v)$, counting them by edge gives $2m$.}
\end{diagrambox}

Three consequences used again and again:
\begin{enumerate}[leftmargin=*]
    \item \textbf{The number of odd-degree vertices is even.} The sum of all degrees is even; the even degrees contribute an even amount, so the odd degrees must add up to an even number, which needs an even count of them. ($G_1$ has two: vertices 3 and 4.)
    \item \textbf{Average degree $= 2m/n$.} A graph with $n=2\cdot10^5$ vertices and $m=2\cdot10^5$ edges has average degree 2 --- but a single vertex may still have degree $2\cdot10^5$ (a ``star''). Algorithms whose cost is $\sum_v \deg(v)$ are $\bigO{m}$; algorithms whose cost is $n\cdot\max\deg$ can be $\bigO{n^2}$.
    \item \textbf{Iterating over all adjacency lists costs $\bigO{n+m}$}, not $\bigO{n\cdot m}$: the lists have total length $\sum\deg(v)=2m$. This one fact is behind the running time of every traversal in this book.
\end{enumerate}

For directed graphs the same argument gives $\sum_v \indeg{v} = \sum_v \outdeg{v} = |E|$.

%------------------------------------------------------------
\subsubsection*{Walks, paths and cycles}

\begin{definitionbox}[Walk, path, cycle]
\begin{itemize}[leftmargin=*]
    \item A \textbf{walk} is a sequence of vertices $v_0, v_1, \dots, v_k$ in which consecutive vertices are joined by an edge (in a digraph: by an arc $v_{i}\to v_{i+1}$). Its \textbf{length} is the number of edges, $k$. Vertices and edges may repeat.
    \item A \textbf{path} (or \textbf{simple path}) is a walk with no repeated vertex.
    \item A \textbf{cycle} is a walk of length $\ge 3$ (in a digraph, $\ge 1$) that returns to its start, $v_k=v_0$, and repeats no other vertex.
\end{itemize}
\end{definitionbox}

\begin{diagrambox}
\begin{tikzpicture}[x=0.82cm, y=0.85cm]
    \foreach \sx/\title/\sub [count=\k] in {0/walk/{1, 2, 3, 1, 2}, 4.8/path/{1, 3, 4, 5}, 9.6/cycle/{4, 5, 6, 4}}{
        \begin{scope}[xshift=\sx cm]
            \node[gv] (a\k1) at (0,0.8) {1};
            \node[gv] (a\k2) at (0,-0.8) {2};
            \node[gv] (a\k3) at (1.3,0) {3};
            \node[gv] (a\k4) at (2.7,0) {4};
            \node[gv] (a\k5) at (4,0.8) {5};
            \node[gv] (a\k6) at (4,-0.8) {6};
            \foreach \a/\b in {1/2, 1/3, 2/3, 3/4, 4/5, 4/6, 5/6} \draw[gfaded] (a\k\a) -- (a\k\b);
            \node[font=\small\bfseries] at (2,-1.6) {\title};
            \node[font=\scriptsize\color{darkgray}] at (2,-2.05) {\sub};
        \end{scope}
    }
    \draw[gpath] (a11) -- (a12); \draw[gpath] (a12) -- (a13); \draw[gpath] (a13) -- (a11);
    \node[font=\scriptsize\color{accentorange}, anchor=west] at (-0.4,1.5) {edge 12 used twice};
    \draw[gpath] (a21) -- (a23); \draw[gpath] (a23) -- (a24); \draw[gpath] (a24) -- (a25);
    \draw[gpath] (a34) -- (a35); \draw[gpath] (a35) -- (a36); \draw[gpath] (a36) -- (a34);
\end{tikzpicture}
\figcap{On $G_1$: the walk 1, 2, 3, 1, 2 (length 4) repeats vertices; the path 1, 3, 4, 5 (length 3) does not; the cycle 4, 5, 6, 4 (length 3) returns to its start.}
\end{diagrambox}

\begin{insightbox}[Every walk contains a path]
If a walk from $u$ to $v$ visits some vertex twice, cut out the loop between the two visits; the result is a shorter walk from $u$ to $v$. Repeating this gives a path. So ``$v$ can be reached from $u$'' means the same whether we think of walks or paths --- and the \emph{shortest} walk is always a path.
\end{insightbox}

%------------------------------------------------------------
\subsubsection*{Connectivity and components}

An undirected graph is \textbf{connected} if every vertex can reach every other vertex. Otherwise it falls into \textbf{connected components}: maximal groups of vertices that can all reach each other. ``Maximal'' means you cannot add another vertex to the group and keep it connected. Every vertex belongs to exactly one component; an isolated vertex is a component on its own.

\begin{diagrambox}
\begin{tikzpicture}[x=1.1cm, y=1cm]
    \node[gvA] (1) at (0,0.7) {1};
    \node[gvA] (2) at (1.2,0.7) {2};
    \node[gvA] (3) at (0.6,-0.5) {3};
    \node[gvB] (4) at (2.8,0.7) {4};
    \node[gvB] (5) at (2.8,-0.6) {5};
    \node[gvsel] (6) at (4.4,0.1) {6};
    \node[gvseen] (7) at (5.9,0.7) {7};
    \node[gvseen] (8) at (7.1,0.7) {8};
    \node[gvseen] (9) at (6.5,-0.5) {9};
    \foreach \a/\b in {1/2, 2/3, 3/1, 4/5, 7/8, 8/9} \draw[ge] (\a) -- (\b);
    \foreach \x/\t in {0.6/{$\{1,2,3\}$}, 2.8/{$\{4,5\}$}, 4.4/{$\{6\}$}, 6.5/{$\{7,8,9\}$}}
        \node[font=\scriptsize] at (\x,-1.35) {\t};
\end{tikzpicture}
\figcap{A graph with $n=9$, $m=6$ and four components (one colour each). Vertex 6 is isolated: it is its own component. Counting and labelling components is Chapter~4; doing it with DSU is Chapter~6.}
\end{diagrambox}

For directed graphs ``connected'' splits into two notions: $v$ is \textbf{reachable} from $u$ if there is a directed path $u\leadsto v$, and the graph is \textbf{strongly connected} if every vertex is reachable from every other. A one-way road from A to B makes B reachable from A but not the other way round.

%------------------------------------------------------------
\subsubsection*{Trees and forests}

\begin{definitionbox}[Tree, forest]
A \textbf{tree} is a connected undirected graph with no cycle. A \textbf{forest} is a graph with no cycle (its components are trees). For a graph with $n$ vertices the following are equivalent --- any two of the first three imply the third:
\begin{itemize}[leftmargin=*]
    \item it is connected;\quad it has no cycle;\quad it has exactly $n-1$ edges;
    \item between any two vertices there is \textbf{exactly one} path.
\end{itemize}
A forest with $n$ vertices and $k$ trees has exactly $n-k$ edges.
\end{definitionbox}

\begin{diagrambox}
\begin{tikzpicture}[x=0.95cm, y=0.95cm]
    \begin{scope}
        \node[gv] (1) at (1,1) {1};
        \node[gv] (2) at (0,0) {2};
        \node[gv] (3) at (2,0) {3};
        \node[gv] (4) at (-0.5,-1.1) {4};
        \node[gv] (5) at (0.5,-1.1) {5};
        \node[gv] (6) at (2,-1.1) {6};
        \foreach \a/\b in {1/2, 1/3, 2/4, 2/5, 3/6} \draw[gtree] (\a) -- (\b);
        \node[font=\small\bfseries] at (0.75,-1.9) {tree};
        \node[font=\scriptsize\color{darkgray}] at (0.75,-2.3) {$n=6$, $m=5$};
    \end{scope}
    \begin{scope}[xshift=4.4cm]
        \node[gv] (1) at (1,1) {1};
        \node[gv] (2) at (0,0) {2};
        \node[gv] (3) at (2,0) {3};
        \node[gv] (4) at (-0.5,-1.1) {4};
        \node[gv] (5) at (0.5,-1.1) {5};
        \node[gv] (6) at (2,-1.1) {6};
        \foreach \a/\b in {1/2, 1/3, 2/4, 2/5, 3/6} \draw[gtree] (\a) -- (\b);
        \draw[gback, -] (5) -- (6);
        \node[font=\small\bfseries] at (0.75,-1.9) {add an edge};
        \node[font=\scriptsize\color{darkgray}] at (0.75,-2.3) {cycle 5--2--1--3--6--5};
    \end{scope}
    \begin{scope}[xshift=8.8cm]
        \node[gv] (1) at (1,1) {1};
        \node[gv] (2) at (0,0) {2};
        \node[gv] (3) at (2,0) {3};
        \node[gv] (4) at (-0.5,-1.1) {4};
        \node[gv] (5) at (0.5,-1.1) {5};
        \node[gv] (6) at (2,-1.1) {6};
        \foreach \a/\b in {1/2, 2/4, 2/5, 3/6} \draw[gtree] (\a) -- (\b);
        \draw[gcut] (1.1,0.2) -- (1.9,0.9);
        \node[font=\small\bfseries] at (0.75,-1.9) {remove an edge};
        \node[font=\scriptsize\color{darkgray}] at (0.75,-2.3) {forest of 2 trees, $m=4$};
    \end{scope}
\end{tikzpicture}
\figcap{A tree is ``just connected'': adding any edge closes exactly one cycle (dashed), removing any edge disconnects it. That is why $n-1$ edges is the magic number.}
\end{diagrambox}

When one vertex is chosen as the \textbf{root}, every other vertex gets a \textbf{parent} (its neighbour towards the root) and \textbf{children}; the tree becomes a \textbf{rooted tree}. Company hierarchies, file systems and DFS trees are rooted trees. Trees are so important that Chapter~13 is devoted to them.

%------------------------------------------------------------
\subsubsection*{DAGs}

A directed graph with no directed cycle is a \textbf{DAG} (directed acyclic graph). Prerequisites, task dependencies and ``$i<j$'' relations are DAGs. Their vertices can be listed so that every arc points forward --- a \textbf{topological order} (Chapter~5) --- and that is exactly the order in which DP fills its table.

\begin{diagrambox}
\begin{tikzpicture}[x=1.15cm, y=1cm]
    \begin{scope}
        \node[gv] (a) at (0,0.8) {a};
        \node[gv] (b) at (1.5,0.8) {b};
        \node[gv] (c) at (0,-0.6) {c};
        \node[gv] (d) at (1.5,-0.6) {d};
        \node[gv] (e) at (3,0.1) {e};
        \draw[gde] (a) -- (b); \draw[gde] (a) -- (c); \draw[gde] (c) -- (d);
        \draw[gde] (b) -- (d); \draw[gde] (b) -- (e); \draw[gde] (d) -- (e);
        \node[font=\small\bfseries] at (1.5,-1.4) {a DAG};
    \end{scope}
    \begin{scope}[xshift=4.8cm]
        \foreach \v/\x in {a/0, c/1.1, b/2.2, d/3.3, e/4.4} \node[gv] (\v) at (\x,0.1) {\v};
        \draw[gde] (a) -- (c); \draw[gde] (b) -- (d); \draw[gde] (d) -- (e);
        \draw[gde] (a) to[bend left=35] (b);
        \draw[gde] (b) to[bend left=35] (e);
        \draw[gde] (c) to[bend right=35] (d);
        \node[font=\small\bfseries] at (2.2,-1.4) {same DAG, topological order};
        \node[font=\scriptsize\color{darkgray}] at (2.2,-1.8) {every arc points to the right};
    \end{scope}
    \begin{scope}[xshift=11.5cm]
        \node[gv] (x) at (0,0.8) {x};
        \node[gv] (y) at (1.4,0.8) {y};
        \node[gv] (z) at (0.7,-0.5) {z};
        \draw[gde] (x) -- (y); \draw[gde] (y) -- (z); \draw[gback] (z) -- (x);
        \node[font=\small\bfseries] at (0.7,-1.4) {not a DAG};
    \end{scope}
\end{tikzpicture}
\figcap{A DAG can be laid out on a line with all arcs going right. The cycle $x\to y\to z\to x$ makes such a layout impossible: whichever of the three comes first has an arc coming \emph{into} it from the right.}
\end{diagrambox}

%------------------------------------------------------------
\subsubsection*{Bipartite graphs}

A graph is \textbf{bipartite} if its vertices can be coloured with two colours so that every edge joins two \emph{different} colours --- for example, students and the clubs they belong to, or the black and white squares of a chessboard (a knight always changes colour). A graph is bipartite exactly when it has \textbf{no cycle of odd length} (Chapter~4).

\begin{diagrambox}
\begin{tikzpicture}[x=1.1cm, y=1cm]
    \begin{scope}
        \node[gvA] (1) at (0,1) {1};
        \node[gvB] (2) at (1.3,1) {2};
        \node[gvA] (3) at (1.3,-0.3) {3};
        \node[gvB] (4) at (0,-0.3) {4};
        \node[gvB] (5) at (2.6,-0.3) {5};
        \foreach \a/\b in {1/2, 2/3, 3/4, 4/1, 3/5} \draw[ge] (\a) -- (\b);
        \node[font=\small\bfseries] at (1.3,-1.2) {bipartite};
    \end{scope}
    \begin{scope}[xshift=4.5cm]
        \node[gvA] (1) at (0,1) {1};
        \node[gvA] (3) at (0,0.1) {3};
        \node[gvB] (2) at (2,1.2) {2};
        \node[gvB] (4) at (2,0.35) {4};
        \node[gvB] (5) at (2,-0.5) {5};
        \foreach \a/\b in {1/2, 3/2, 3/4, 1/4, 3/5} \draw[ge] (\a) -- (\b);
        \node[font=\small\bfseries] at (1,-1.2) {same graph, two sides};
    \end{scope}
    \begin{scope}[xshift=9.2cm]
        \node[gvA] (a) at (0,1) {a};
        \node[gvB] (b) at (1.4,1) {b};
        \node[gvdead] (c) at (0.7,-0.2) {c};
        \draw[ge] (a) -- (b); \draw[ge] (b) -- (c); \draw[gcut] (c) -- (a);
        \node[font=\scriptsize\color{warningred}, align=center] at (2.5,-0.2) {c needs a colour\\different from both};
        \node[font=\small\bfseries] at (0.7,-1.2) {odd cycle: not bipartite};
    \end{scope}
\end{tikzpicture}
\figcap{Left and middle: a bipartite graph redrawn with its two sides as columns --- every edge crosses. Right: a triangle cannot be 2-coloured.}
\end{diagrambox}

%------------------------------------------------------------
\subsubsection*{Complete graphs and complements}

The \textbf{complete graph} $K_n$ has an edge between \emph{every} pair of its $n$ vertices, so $m=\binom{n}{2}=\frac{n(n-1)}{2}$, the maximum for a simple undirected graph. The \textbf{complement} $\overline{G}$ of a simple graph $G$ has the same vertices and exactly the edges that $G$ does \emph{not} have; together, $G$ and $\overline{G}$ form $K_n$.

\begin{diagrambox}
\begin{tikzpicture}[x=1cm, y=1cm]
    \foreach \k/\sx in {3/0, 4/2.5, 5/5.3}{
        \begin{scope}[xshift=\sx cm]
            \foreach \i in {1,...,\k} \node[gv, minimum size=5.5mm, font=\scriptsize\bfseries] (k\k-\i) at ({90+360/\k*(\i-1)}:0.9) {\i};
            \foreach \i in {1,...,\k} \foreach \j in {1,...,\k} {
                \ifnum\i<\j \draw[ge] (k\k-\i) -- (k\k-\j); \fi }
            \pgfmathtruncatemacro{\mm}{\k*(\k-1)/2}
            \node[font=\scriptsize\bfseries] at (0,-1.45) {$K_{\k}$: $m=\mm$};
        \end{scope}
    }
    \begin{scope}[xshift=8.4cm, yshift=0.3cm]
        \foreach \i/\x in {1/0, 2/0.9, 3/1.8, 4/2.7} \node[gv, minimum size=5.5mm, font=\scriptsize\bfseries] (p\i) at (\x,0.5) {\i};
        \draw[ge] (p1) -- (p2); \draw[ge] (p2) -- (p3); \draw[ge] (p3) -- (p4);
        \foreach \i/\x in {3/0, 1/0.9, 4/1.8, 2/2.7} \node[gv, minimum size=5.5mm, font=\scriptsize\bfseries] (q\i) at (\x,-0.9) {\i};
        \draw[gpath, line width=1.2pt] (q3) -- (q1); \draw[gpath, line width=1.2pt] (q1) -- (q4); \draw[gpath, line width=1.2pt] (q4) -- (q2);
        \node[font=\scriptsize, anchor=west] at (3.0,0.5) {$G$};
        \node[font=\scriptsize, anchor=west] at (3.0,-0.9) {$\overline{G}$};
    \end{scope}
\end{tikzpicture}
\figcap{$K_3$, $K_4$, $K_5$ with 3, 6 and 10 edges. Right: the path $G=1$--$2$--$3$--$4$ (3 edges) and its complement (the other $6-3=3$ pairs: 13, 14, 24), redrawn: it is again a path, 3--1--4--2.}
\end{diagrambox}

Complements matter in contests because a statement may describe the edges a graph does \emph{not} have (``every two cities are connected except for these $m$ pairs''). If $n$ is large, the complement is dense and cannot be built explicitly --- a classic trap.

%------------------------------------------------------------
\subsubsection*{Sparse and dense graphs}

A simple graph has between $0$ and $\binom n2\approx n^2/2$ edges. When $m$ is of the order of $n$ (trees, road maps, grids, typical input with $n,m\le 2\cdot10^5$) the graph is \textbf{sparse}; when $m$ is of the order of $n^2$ it is \textbf{dense}. The distinction decides how to store the graph (Lesson~\ref{ssec:ch1-storing}) and which algorithm is fastest.

\begin{table}[H]
\centering\small
\begin{tabular}{@{}l l l l@{}}
\toprule
\textbf{Graph} & \textbf{$m$ as a function of $n$} & \textbf{$n=1000$} & \textbf{kind} \\
\midrule
tree, path, functional graph & $n-1$ or $n$ & $\approx 10^3$ & sparse \\
grid $r\times c$ (4 neighbours) & $\approx 2n$ & $\approx 2\cdot 10^3$ & sparse \\
typical contest input & given, usually $\le 2\cdot 10^5$ & --- & sparse \\
complete graph $K_n$ & $n(n-1)/2$ & $499\,500$ & dense \\
complete digraph & $n(n-1)$ & $999\,000$ & dense \\
\bottomrule
\end{tabular}
\end{table}

%------------------------------------------------------------
\subsubsection*{Functional graphs: exactly one arrow out}

A directed graph in which \textbf{every vertex has out-degree exactly 1} is a \textbf{functional graph}: it is the picture of a function $f:\{1..n\}\to\{1..n\}$, with an arc $v\to f(v)$. Permutations, ``each student blames one other student'', ``teleporter $i$ sends you to $a_i$'' are all functional graphs. It has exactly $n$ arcs, and its shape is always the same: \emph{every component is one cycle with trees hanging into it}.

\begin{diagrambox}
\begin{tikzpicture}[x=1.15cm, y=1cm]
    \node[gvsel] (2) at (2,0.8) {2};
    \node[gvsel] (3) at (3.2,0) {3};
    \node[gvsel] (4) at (2,-0.8) {4};
    \node[gv] (1) at (0.8,1.5) {1};
    \node[gv] (5) at (0.8,-0.8) {5};
    \node[gv] (6) at (-0.4,-0.8) {6};
    \node[gvsrc] (7) at (-1.6,-0.8) {7};
    \draw[gde] (2) to[bend left=20] (3);
    \draw[gde] (3) to[bend left=20] (4);
    \draw[gde] (4) to[bend left=20] (2);
    \draw[gde] (1) -- (2); \draw[gde] (5) -- (4); \draw[gde] (6) -- (5); \draw[gde] (7) -- (6);
    \node[gvsel] (8) at (5,0.4) {8};
    \node[gvsel] (9) at (6.4,0.4) {9};
    \draw[gde] (8) to[bend left=30] (9);
    \draw[gde] (9) to[bend left=30] (8);
    \node[note, text width=4cm] at (5.7,-0.9) {$f = (2,3,4,2,4,5,6,9,8)$\\for $v = 1..9$};
\end{tikzpicture}
\figcap{$f(1)=2, f(2)=3, f(3)=4, f(4)=2, f(5)=4, f(6)=5, f(7)=6, f(8)=9, f(9)=8$. Two components; the \legendchip{gvsel}{cycle} vertices are $\{2,3,4\}$ and $\{8,9\}$. Starting at 7 and following arrows: 7, 6, 5, 4, 2, 3, 4, 2, 3, 4, \ldots{} --- a tail of length 3, then the cycle forever (the shape of the Greek letter $\rho$).}
\end{diagrambox}

\begin{insightbox}[Why there is always a cycle]
Follow the arrows from any vertex. There are only $n$ vertices, so within $n$ steps some vertex repeats; from then on the walk repeats forever. Pattern~5 turns this observation into algorithms, and Chapter~11 (functional graphs and binary lifting) jumps $10^{18}$ steps along such walks.
\end{insightbox}

%------------------------------------------------------------
\subsubsection*{The dictionary at a glance}

\begin{table}[H]
\centering\small
\begin{tabularx}{\linewidth}{@{}l X l@{}}
\toprule
\textbf{Term} & \textbf{Meaning} & \textbf{Typical clue in a statement} \\
\midrule
degree & number of incident edges & ``how many friends'', ``number of roads at'' \\
path / cycle & no repeated vertex / returns to start & ``route'', ``loop'', ``come back'' \\
component & maximal mutually reachable group & ``groups'', ``islands'', ``clusters'' \\
tree & connected, $n-1$ edges, no cycle & ``$n-1$ roads'', ``unique path'' \\
DAG & directed, no directed cycle & ``prerequisites'', ``before'', ``$i<j$'' \\
bipartite & 2-colourable, no odd cycle & ``two teams'', ``enemies in different groups'' \\
complete / complement & all pairs / the missing pairs & ``every pair except'' \\
functional graph & out-degree exactly 1 & ``$a_i$'', ``next'', ``teleports to'' \\
\bottomrule
\end{tabularx}
\end{table}

%------------------------------------------------------------
\subsubsection*{Check your understanding}

\begin{enumerate}[leftmargin=*]
    \item Is there a simple graph with degrees $3,3,3,2,2$?
    \item Nine people at a party: can each shake hands with exactly three others?
    \item A connected graph has $n=10$ vertices and $m=9$ edges. What is it? What if it had $m=10$ edges?
    \item A forest has 20 vertices and 16 edges. How many trees does it have?
    \item Is the cycle on 6 vertices bipartite? The cycle on 7 vertices?
    \item How many edges does the complement of a tree on $n$ vertices have?
    \item A functional graph has $n$ vertices. How many arcs does it have, and can it be a DAG?
\end{enumerate}

\begin{solutionbox}[Answers]
\begin{enumerate}[leftmargin=*]
    \item No: the degrees add up to 13, which is odd, but $\sum\deg = 2m$ is even (equivalently: three odd degrees, an odd number).
    \item No, for the same reason: $9\cdot 3=27$ handshake-ends cannot be split into pairs.
    \item With 9 edges it is a tree. With 10 edges it has exactly one cycle (a tree plus one edge; sometimes called a \emph{unicyclic} graph).
    \item $n-k=m$ gives $k=20-16=4$ trees.
    \item $C_6$ is bipartite (colour alternately); $C_7$ is not --- it is an odd cycle.
    \item $\binom n2-(n-1)=\frac{(n-1)(n-2)}{2}$.
    \item Exactly $n$ arcs (one per vertex). It is never a DAG: following arrows must eventually repeat a vertex, so every functional graph contains a cycle.
\end{enumerate}
\end{solutionbox}

%============================================================
% CH1 · LESSON 3 — STORING A GRAPH
%============================================================
\subsection{Storing a Graph}
\label{ssec:ch1-storing}

A drawing is for humans; a program needs the graph in memory, in a form that answers the questions the algorithm asks. Almost every algorithm asks one question over and over: \textbf{``who are the neighbours of $u$?''} The structure that answers it in time proportional to the answer --- the \emph{adjacency list} --- is what you will use in 95\% of problems. This lesson shows it together with the three alternatives, so that you know when each one is the right tool.

Throughout the lesson we store the same weighted undirected graph $G_2$:

\begin{diagrambox}
\begin{tikzpicture}[x=1.4cm, y=1.1cm]
    \node[gv] (1) at (0,0.8) {1};
    \node[gv] (2) at (0,-0.8) {2};
    \node[gv] (3) at (1.4,0) {3};
    \node[gv] (4) at (2.9,0) {4};
    \node[gv] (5) at (4.4,0) {5};
    \draw[ge] (1) -- node[gw, left] {4} (2);
    \draw[ge] (1) -- node[gw, above right] {1} (3);
    \draw[ge] (2) -- node[gw, below right] {2} (3);
    \draw[ge] (3) -- node[gw, above] {7} (4);
    \draw[ge] (4) -- node[gw, above] {3} (5);
    \node[draw=darkgray, fill=lightgray!60, rounded corners=2pt, font=\ttfamily\small, align=left, inner sep=4pt, anchor=west] at (5.4,0)
        {5 5\\1 2 4\\1 3 1\\2 3 2\\3 4 7\\4 5 3};
    \node[font=\scriptsize\color{darkgray}] at (6.0,1.35) {input};
\end{tikzpicture}
\figcap{$G_2$: $n=5$ vertices, $m=5$ weighted edges, given as ``$u$ $v$ $w$'' lines.}
\end{diagrambox}

%------------------------------------------------------------
\subsubsection*{1. Edge list --- just keep the input}

The simplest representation is the input itself: an array of $m$ triples $(u,v,w)$.

\begin{diagrambox}
\begin{tikzpicture}
    \dpindex{ei}{1.4}{$i$}{0}{4}
    \dparray{eu}{0.7}{$u$}{1/cell, 1/cell, 2/cell, 3/cell, 4/cell}
    \dparray{ev}{0}{$v$}{2/cell, 3/cell, 3/cell, 4/cell, 5/cell}
    \dparray{ew}{-0.7}{$w$}{4/compcell, 1/compcell, 2/compcell, 7/compcell, 3/compcell}
    \node[note, anchor=west, text width=5.2cm] at (5.2,0.35) {Good for: algorithms that look at every edge as a whole --- sorting edges by weight (Kruskal, Ch.~10), relaxing all edges (Bellman--Ford, Ch.~8), DSU (Ch.~6), counting degrees (Lesson~\ref{ssec:ch1-vocabulary}).\\[2pt] Bad for: ``neighbours of $u$'' costs $\bigO{m}$.};
\end{tikzpicture}
\figcap{The edge list of $G_2$: three parallel arrays (or one \texttt{vector<array<int,3>>}).}
\end{diagrambox}

\begin{lstlisting}[style=snippet]
struct Edge { int u, v, w; };
vector<Edge> edges(m);
for (auto &e : edges) cin >> e.u >> e.v >> e.w;
sort(edges.begin(), edges.end(), [](const Edge &a, const Edge &b) { return a.w < b.w; });
\end{lstlisting}

%------------------------------------------------------------
\subsubsection*{2. Adjacency matrix --- one cell per pair}

An $n\times n$ table with \texttt{g[u][v]} $=$ weight of the edge $u$--$v$ (or 1/0 for ``edge / no edge'' in an unweighted graph). An undirected graph gives a \emph{symmetric} matrix.

\begin{diagrambox}
\begin{tikzpicture}[x=0.8cm, y=0.65cm]
    \foreach \j in {1,...,5} \node[hdr, minimum width=7mm, minimum height=5.5mm] at (\j,0.9) {\j};
    \foreach \i in {1,...,5} \node[hdr, minimum width=7mm, minimum height=5.5mm] at (0,1-\i-0.1) {\i};
    \foreach \row [count=\i] in {{0,4,1,0,0},{4,0,2,0,0},{1,2,0,7,0},{0,0,7,0,3},{0,0,0,3,0}}{
        \foreach \val [count=\j] in \row {
            \ifnum\val=0
                \node[emptycell, minimum width=7mm, minimum height=5.5mm] at (\j,1-\i-0.1) {$\cdot$};
            \else
                \node[compcell, minimum width=7mm, minimum height=5.5mm, font=\small\bfseries] at (\j,1-\i-0.1) {\val};
            \fi
        }
    }
    \node[font=\scriptsize\color{darkgray}] at (-0.1,0.9) {$u\backslash v$};
    \draw[accentorange, line width=1.2pt, rounded corners=2pt] (0.55,-2.45) rectangle (5.45,-1.75);
    \node[note, anchor=west, text width=5.4cm] at (6.2,-1.6) {Row 3 lists the neighbours of vertex 3 --- but you must scan \emph{all} $n$ cells to find the three non-empty ones.\\[3pt]
        Good for: $n\le$ a few thousand, dense graphs, ``is $u$--$v$ an edge?'' in $\bigO{1}$, Floyd--Warshall (Ch.~9).\\[2pt]
        Bad for: memory $n^2$; neighbours cost $\bigO{n}$; parallel edges collapse into one cell.};
\end{tikzpicture}
\figcap{Adjacency matrix of $G_2$ (dots are ``no edge''). It is symmetric about the diagonal because $G_2$ is undirected.}
\end{diagrambox}

\begin{lstlisting}[style=snippet]
vector<vector<int>> g(n + 1, vector<int>(n + 1, 0));      // n <= ~5000 only!
for (int i = 0; i < m; i++) {
    int u, v, w; cin >> u >> v >> w;
    g[u][v] = g[v][u] = w;          // parallel edges: keep min(w) if needed
}
\end{lstlisting}

%------------------------------------------------------------
\subsubsection*{3. Adjacency list --- the default}

For each vertex keep a dynamic array of its neighbours: \texttt{vector<vector<int>> adj(n + 1)}. An undirected edge $u$--$v$ is pushed twice, as $v$ into \texttt{adj[u]} and as $u$ into \texttt{adj[v]}. For a weighted graph each entry is a pair (neighbour, weight): \texttt{vector<vector<pair<int,int>>>}.

\begin{diagrambox}
\begin{tikzpicture}[x=1cm, y=0.7cm,
    hd/.style={hdr, minimum width=8mm, minimum height=6mm},
    it/.style={cell, minimum width=12mm, minimum height=6mm, font=\small},
    it3/.style={it, fill=currentfill, draw=accentorange, line width=1pt}]
    \foreach \u/\lst in {1/{(2,4),(3,1)}, 2/{(1,4),(3,2)}, 3/{(1,1),(2,2),(4,7)}, 4/{(3,7),(5,3)}, 5/{(4,3)}}{
        \node[hd] (h\u) at (0,-\u) {\u};
        \foreach \e [count=\k] in \lst {
            \ifnum\u=3 \node[it3] (c\u\k) at (0.6+1.35*\k,-\u) {\e};
            \else \node[it] (c\u\k) at (0.6+1.35*\k,-\u) {\e}; \fi
        }
        \draw[flow] (h\u) -- (c\u1);
    }
    \node[font=\scriptsize\bfseries\color{darkgray}] at (0,-0.25) {adj};
    \node[note, anchor=west, text width=5.2cm] at (5.2,-3) {Entry (neighbour, weight). Vertex 3 has degree 3, so its list has 3 entries and ``neighbours of 3'' costs exactly 3 steps.\\[3pt]
        Total length $=\sum\deg(v)=2m=10$ entries. Memory $\bigO{n+m}$.};
\end{tikzpicture}
\figcap{The weighted adjacency list of $G_2$, in input order. Each undirected edge appears twice, once in each endpoint's list (e.g.\ $(3,7)$ in list 4 and $(4,7)$ in list 3).}
\end{diagrambox}

\begin{insightbox}[Why the adjacency list is the default]
A traversal looks at the neighbours of each vertex once. With adjacency lists that costs $\sum_v \deg(v) = 2m$ steps in total (handshake lemma), so the whole traversal is $\bigO{n+m}$. With a matrix the same traversal scans $n$ cells per vertex: $\bigO{n^2}$, even if the graph has only a few edges.
\end{insightbox}

%------------------------------------------------------------
\subsubsection*{4. Forward star and CSR --- adjacency lists in flat arrays}

A \texttt{vector<vector<int>>} allocates one small array per vertex. Two classic layouts store the same lists in a few flat arrays instead --- less memory, fewer allocations, faster on huge inputs ($n,m\approx 10^6$). You will meet both in other people's code.

\textbf{Forward star} (``linked list in arrays'', \texttt{head}/\texttt{nxt}/\texttt{to}). Edges are numbered $0,1,2,\dots$ in the order they are added; \texttt{head[u]} is the last edge added at $u$, and \texttt{nxt[e]} is the edge added at the same vertex before $e$ (or $-1$). Walking \texttt{e = head[u], nxt[e], \ldots} lists $u$'s neighbours in \emph{reverse} input order.

\begin{diagrambox}
\begin{tikzpicture}
    \dpindex{fh}{2.2}{$u$}{1}{5}
    \dparray{hd}{1.5}{head}{2/cell, 4/cell, 6/curcell, 8/cell, 9/cell}
    \begin{scope}[yshift=-0.4cm]
    \dpindex{fe}{0.6}{$e$}{0}{9}
    \dparray{to}{-0.1}{to}{2/cell, 1/cell, 3/cell, 1/cell, 3/cell, 2/cell, 4/curcell, 3/cell, 5/cell, 4/cell}
    \dparray{nx}{-0.8}{nxt}{-1/cell, -1/cell, 0/cell, -1/cell, 1/cell, 3/cell, 5/curcell, -1/cell, 7/cell, -1/cell}
    \end{scope}
    \draw[dep] (hd2.south) to[out=-90, in=90] (to6.north);
    \draw[dep] (nx6.south) to[out=-120, in=-60] (nx5.south);
    \draw[dep] (nx5.south) to[out=-120, in=-60] (nx3.south);
    \node[note, anchor=west, text width=9cm] at (0,-2.75) {Neighbours of 3: \texttt{head[3]}$=6$ gives \texttt{to[6]}$=4$; \texttt{nxt[6]}$=5$ gives \texttt{to[5]}$=2$; \texttt{nxt[5]}$=3$ gives \texttt{to[3]}$=1$; \texttt{nxt[3]}$=-1$: stop. Order: 4, 2, 1.};
\end{tikzpicture}
\figcap{Forward star of $G_2$ (weights omitted; a fourth array \texttt{wt[e]} would hold them). Edge $2i$ is $u_i\to v_i$ and edge $2i+1$ is $v_i\to u_i$, so \texttt{e \^{} 1} is always the reverse edge --- a property flow algorithms rely on.}
\end{diagrambox}

\textbf{CSR} (compressed sparse row). Count the degrees, take prefix sums to get \texttt{start[u]}, then drop each neighbour into its slot. The neighbours of $u$ are \texttt{nbr[start[u]]} \ldots{} \texttt{nbr[start[u+1]-1]}, contiguous in memory and in input order.

\begin{diagrambox}
\begin{tikzpicture}
    \dpindex{ci}{1.4}{$u$}{1}{6}
    \dparray{st}{0.7}{start}{0/cell, 2/cell, 4/curcell, 7/curcell, 9/cell, 10/cell}
    \begin{scope}[yshift=-0.3cm]
    \dpindex{ck}{-0.1}{$k$}{0}{9}
    \dparray{nb}{-0.8}{nbr}{2/cell, 3/cell, 1/cell, 3/cell, 1/curcell, 2/curcell, 4/curcell, 3/cell, 5/cell, 4/cell}
    \end{scope}
    \draw[accentorange, line width=1.3pt, decorate, decoration={brace, amplitude=4pt, mirror}] (3.7,-1.55) -- (6.4,-1.55)
        node[midway, below=4pt, font=\scriptsize\color{accentorange}] {neighbours of 3 = slots $4..6$};
    \draw[dep] (st2.south) to[out=-80, in=100] (nb4.north);
    \draw[dep] (st3.south) to[out=-80, in=100] (nb7.north);
\end{tikzpicture}
\figcap{CSR of $G_2$: \texttt{start} has $n+1$ entries ($u=1..6$ here, with \texttt{start[6]}$=2m=10$ as a sentinel); the slice for vertex 3 runs from \texttt{start[3]}$=4$ to \texttt{start[4]}$-1=6$.}
\end{diagrambox}

\begin{lstlisting}[style=cpp]
#include <bits/stdc++.h>
using namespace std;

// Forward star (head / nxt / to arrays) and CSR (start / nbr arrays)
// built from the same undirected edge list.
const int MAXN = 200005, MAXM = 400010;   // MAXM = 2 * (max edges) + slack
int head[MAXN], nxt[MAXM], to_[MAXM], cnt = 0;

void addEdge(int u, int v) {           // directed u -> v
    to_[cnt] = v;
    nxt[cnt] = head[u];                // new edge points to the old first edge
    head[u] = cnt++;                   // and becomes the first edge of u
}

int main() {
    int n, m;
    if (!(cin >> n >> m)) return 0;
    vector<int> eu(m), ev(m);
    for (int i = 0; i < m; i++) cin >> eu[i] >> ev[i];

    // ---- forward star ----
    fill(head, head + n + 1, -1);
    for (int i = 0; i < m; i++) { addEdge(eu[i], ev[i]); addEdge(ev[i], eu[i]); }
    for (int u = 1; u <= n; u++) {
        cout << "star " << u << ":";
        for (int e = head[u]; e != -1; e = nxt[e]) cout << " " << to_[e];
        cout << "\n";
    }

    // ---- CSR: count, prefix sums, fill ----
    vector<int> start(n + 2, 0);
    for (int i = 0; i < m; i++) { start[eu[i] + 1]++; start[ev[i] + 1]++; }
    for (int u = 1; u <= n + 1; u++) start[u] += start[u - 1];
    vector<int> pos(start), nbr(2 * m);
    for (int i = 0; i < m; i++) { nbr[pos[eu[i]]++] = ev[i]; nbr[pos[ev[i]]++] = eu[i]; }
    for (int u = 1; u <= n; u++) {
        cout << "csr  " << u << ":";
        for (int k = start[u]; k < start[u + 1]; k++) cout << " " << nbr[k];
        cout << "\n";
    }
    return 0;
}
\end{lstlisting}

\begin{lstlisting}[style=trace]
input : 5 5   1 2  1 3  2 3  3 4  4 5          (G2 without weights)
output: star 1: 3 2      csr  1: 2 3
        star 2: 3 1      csr  2: 1 3
        star 3: 4 2 1    csr  3: 1 2 4
        star 4: 5 3      csr  4: 3 5
        star 5: 4        csr  5: 4
\end{lstlisting}
(The program prints the ``star'' block first and the ``csr'' block after it; they are shown side by side here.)

%------------------------------------------------------------
\subsubsection*{Choosing a representation}

\begin{table}[H]
\centering\small
\begin{tabularx}{\linewidth}{@{}X c c c@{}}
\toprule
\textbf{Operation} & \textbf{Edge list} & \textbf{Matrix} & \textbf{Adjacency list / CSR} \\
\midrule
memory & $\bigO{m}$ & $\bigO{n^2}$ & $\bigO{n+m}$ \\
build from input & $\bigO{m}$ & $\bigO{n^2+m}$ & $\bigO{n+m}$ \\
list the neighbours of $u$ & $\bigO{m}$ & $\bigO{n}$ & $\bigO{\deg u}$ \\
is $u$--$v$ an edge? & $\bigO{m}$ & $\bigO{1}$ & $\bigO{\deg u}$ \\
visit every edge from every vertex (a traversal) & --- & $\bigO{n^2}$ & $\bigO{n+m}$ \\
sort edges by weight & $\bigO{m\log m}$ & --- & --- \\
\bottomrule
\end{tabularx}
\end{table}

\begin{warningbox}[{$n = 2\cdot 10^5$: the matrix is impossible}]
An \texttt{int} matrix for $n=2\cdot10^5$ needs $n^2\cdot 4 = 1.6\cdot10^{11}$ bytes $=160$ GB; even one \emph{bit} per cell is 5 GB. A typical memory limit is 256 MB. The adjacency list of the same graph with $m=2\cdot10^5$ edges stores $n+2m=6\cdot10^5$ integers: a few megabytes. Rule of thumb: \textbf{use a matrix only when $n\le 5000$} (an \texttt{int} matrix of $5000^2$ is 100 MB --- already borderline).
\end{warningbox}

%------------------------------------------------------------
\subsubsection*{Reading input correctly}

Four details cause most ``my graph code is right but I get WA/RE'' moments:

\begin{enumerate}[leftmargin=*]
    \item \textbf{0- or 1-indexed?} Codeforces, CSES and AtCoder usually number vertices $1..n$: allocate \texttt{n + 1} and ignore index 0. LeetCode usually numbers them $0..n-1$: allocate \texttt{n}. Mixing them up gives an out-of-bounds access on vertex $n$. Alternatively read \texttt{u}, \texttt{v} and immediately do \texttt{u-{}-, v-{}-}, then work 0-indexed everywhere.
    \item \textbf{Undirected = two directed edges.} Push both directions. Forgetting \texttt{adj[v].push\_back(u)} makes half of the graph invisible; pushing both directions for a directed graph invents edges that do not exist.
    \item \textbf{Multiple test cases.} A global \texttt{adj} must be \emph{cleared} for vertices $1..n$ of every test --- or declared inside the test loop. And clear only what the test uses: resetting a $2\cdot10^5$ array for each of $10^4$ tiny tests is $2\cdot10^9$ operations (\tle). Statements that say ``the sum of $n$ over all tests does not exceed $2\cdot10^5$'' are telling you to allocate per test.
    \item \textbf{Fast I/O.} With $m$ up to $10^6$, \texttt{ios::sync\_with\_stdio(false); cin.tie(nullptr);} is mandatory.
\end{enumerate}

\begin{diagrambox}
\begin{tikzpicture}[x=1cm, y=1cm]
    \node[draw=darkgray, fill=lightgray!60, font=\ttfamily\small, align=left, inner sep=3pt, text width=1.1cm] (t0) at (0,1.6) {2};
    \node[draw=primaryblue, fill=lightgray!60, font=\ttfamily\small, align=left, inner sep=3pt, text width=1.1cm, below=0pt of t0] (ta)
        {5 5\\1 2 4\\1 3 1\\2 3 2\\3 4 7\\4 5 3};
    \node[draw=accentorange, fill=lightgray!60, font=\ttfamily\small, align=left, inner sep=3pt, text width=1.1cm, below=0pt of ta] (tb)
        {3 1\\1 3 5};
    \draw[decorate, decoration={brace, amplitude=4pt}, primaryblue, thick] ([xshift=2pt]ta.north east) -- ([xshift=2pt]ta.south east) node[midway, right=5pt, font=\scriptsize, text=primaryblue] {test 1: $G_2$};
    \draw[decorate, decoration={brace, amplitude=4pt}, accentorange, thick] ([xshift=2pt]tb.north east) -- ([xshift=2pt]tb.south east) node[midway, right=5pt, font=\scriptsize, text=accentorange] {test 2};
    \node[font=\scriptsize\color{darkgray}, left=2pt of t0] {$T$};
    \begin{scope}[xshift=4.6cm, yshift=-0.9cm]
        \node[gv] (1) at (0,0) {1};
        \node[gv] (2) at (1.2,0.9) {2};
        \node[gv] (3) at (1.7,-0.2) {3};
        \draw[ge] (1) -- node[gw, below] {5} (3);
        \node[font=\scriptsize\bfseries, text=accentorange] at (0.85,-1.0) {test 2: $n=3$, $m=1$};
        \node[note, anchor=west, text width=4.6cm] at (2.6,0.4) {Vertex 2 is isolated. If \texttt{adj} were not reset, list 1 would still contain 2 and 3 from test 1 --- a wrong answer that passes test 1 and fails test 2.};
    \end{scope}
\end{tikzpicture}
\figcap{Two tests in one input. Each test must start with empty adjacency lists.}
\end{diagrambox}

Here is a complete reading template that gets all four details right. Copy it, then replace the printing loop with your algorithm.

\begin{lstlisting}[style=cpp]
#include <bits/stdc++.h>
using namespace std;

// Complete reading template: T test cases, each a weighted graph.
// Input:  T
//         n m            (vertices 1..n, m edges)
//         u v w   x m    (undirected edge u-v of weight w)
// Output: for every vertex, its neighbours with weights, and its degree.
int main() {
    ios::sync_with_stdio(false);
    cin.tie(nullptr);
    int T;
    if (!(cin >> T)) return 0;
    while (T--) {
        int n, m;
        cin >> n >> m;
        // fresh containers for every test: nothing leaks between tests
        vector<vector<pair<int,int>>> adj(n + 1);   // adj[u] = {(v, w), ...}
        for (int i = 0; i < m; i++) {
            int u, v, w;
            cin >> u >> v >> w;
            adj[u].push_back(make_pair(v, w));      // u -> v
            adj[v].push_back(make_pair(u, w));      // v -> u (undirected!)
        }
        for (int u = 1; u <= n; u++) {
            cout << u << " (deg " << adj[u].size() << "):";
            for (size_t k = 0; k < adj[u].size(); k++)
                cout << " " << adj[u][k].first << "/" << adj[u][k].second;
            cout << "\n";
        }
        cout << "--\n";
    }
    return 0;
}
\end{lstlisting}

\begin{lstlisting}[style=trace]
output for the input in the figure:
1 (deg 2): 2/4 3/1            1 (deg 1): 3/5
2 (deg 2): 1/4 3/2            2 (deg 0):
3 (deg 3): 1/1 2/2 4/7        3 (deg 1): 1/5
4 (deg 2): 3/7 5/3            --
5 (deg 1): 4/3
--
\end{lstlisting}
(Test 1 on the left, test 2 on the right; the program prints them one after the other.)

\begin{tricksbox}[For a directed graph, and for LeetCode]
For a directed graph delete the second \texttt{push\_back}. On LeetCode the graph arrives as \texttt{int n, vector<vector<int>>\& edges}; the same loop becomes \texttt{for (auto \&e : edges) \{ adj[e[0]].push\_back(e[1]); adj[e[1]].push\_back(e[0]); \}} with \texttt{adj(n)} because vertices are $0..n-1$.
\end{tricksbox}

%------------------------------------------------------------
\subsubsection*{Reading a grid}

Many graphs arrive as a picture: $R$ lines of $C$ characters, where \texttt{.} is free and \texttt{\#} is a wall. We do \emph{not} build an adjacency list for a grid. The cell $(r,c)$ is the vertex, and its neighbours are computed on the fly with two small direction arrays. When an algorithm needs a single integer per vertex (an array index, a DSU element), use $\mathit{id}=r\cdot C+c$.

\begin{diagrambox}
\begin{tikzpicture}[x=0.75cm, y=0.75cm]
    \foreach \row [count=\r from 0] in {{.,.,\#,.},{.,.,.,.},{\#,.,.,\#}}{
        \foreach \ch [count=\c from 0] in \row {
            \pgfmathtruncatemacro{\id}{\r*4+\c}
            \if\ch\#
                \node[gwall] at (\c,-\r) {};
            \else
                \ifnum\id=5 \node[gcell, fill=currentfill] at (\c,-\r) {\scriptsize\id};
                \else \node[gcell] at (\c,-\r) {\scriptsize\id}; \fi
            \fi
        }
    }
    \foreach \c in {0,...,3} \node[font=\scriptsize\color{darkgray}] at (\c,0.75) {$c{=}\c$};
    \foreach \r in {0,...,2} \node[font=\scriptsize\color{darkgray}] at (-0.95,-\r) {$r{=}\r$};
    \foreach \dx/\dy in {0/1, 0/-1, -1/0, 1/0} \draw[dep, line width=1.3pt] (1+0.25*\dx,-1+0.25*\dy) -- (1+0.7*\dx,-1+0.7*\dy);
    \begin{scope}[xshift=4.4cm]
        \node[draw=darkgray, fill=lightgray!60, rounded corners=2pt, font=\ttfamily\small, align=left, inner sep=4pt] at (0.4,-1) {3 4\\..\#.\\....\\\#..\#};
        \node[font=\scriptsize\color{darkgray}] at (0.4,0.45) {input};
    \end{scope}
    \begin{scope}[xshift=7.6cm]
        \foreach \row [count=\r from 0] in {{2,2,\#,1},{2,4,3,2},{\#,2,2,\#}}{
            \foreach \ch [count=\c from 0] in \row {
                \if\ch\# \node[gwall] at (\c,-\r) {};
                \else \node[gcell, font=\small\bfseries] at (\c,-\r) {\ch}; \fi
            }
        }
        \node[font=\scriptsize\color{darkgray}] at (1.5,0.75) {free neighbours of each cell};
    \end{scope}
\end{tikzpicture}
\figcap{Left: cell ids $r\cdot C+c$ for $C=4$; the yellow cell $(1,1)$ has id 5 and four candidate neighbours (up, down, left, right). Right: the output of the program below --- how many of the four candidates are inside the grid and free.}
\end{diagrambox}

\begin{lstlisting}[style=cpp]
#include <bits/stdc++.h>
using namespace std;

// Reading a grid: R rows of C characters, '.' = free, '#' = wall.
// Prints, for every free cell, how many free 4-neighbours it has.
int main() {
    int R, C;
    if (!(cin >> R >> C)) return 0;
    vector<string> g(R);
    for (int r = 0; r < R; r++) cin >> g[r];         // one row = one string
    const int dr[4] = {-1, 1, 0, 0};
    const int dc[4] = {0, 0, -1, 1};
    for (int r = 0; r < R; r++) {
        for (int c = 0; c < C; c++) {
            if (g[r][c] == '#') { cout << '#'; continue; }
            int k = 0;
            for (int d = 0; d < 4; d++) {
                int nr = r + dr[d], nc = c + dc[d];
                if (nr < 0 || nr >= R || nc < 0 || nc >= C) continue;  // off the board
                if (g[nr][nc] == '#') continue;                        // wall
                k++;
            }
            cout << k;
        }
        cout << "\n";
    }
    return 0;
}
\end{lstlisting}

\begin{lstlisting}[style=trace]
input : 3 4  ..#.  ....  #..#
output: 22#1
        2432
        #22#
\end{lstlisting}

The bounds check \emph{before} the wall check matters: \texttt{g[-1][c]} is undefined behaviour. Grids are the subject of Pattern~3.

%------------------------------------------------------------
\subsubsection*{Check your understanding}

\begin{enumerate}[leftmargin=*]
    \item $n=1000$, $m=4000$, and the algorithm asks ``is there an edge $u$--$v$?'' $10^6$ times. Which representation?
    \item $n=10^5$, $m=2\cdot10^5$, undirected. How many entries do the adjacency lists hold in total?
    \item In the forward-star figure, list the neighbours of vertex 2 by following \texttt{head} and \texttt{nxt}.
    \item A LeetCode problem gives \texttt{n = 4, edges = [[0,1],[1,2]]}. What goes wrong with \texttt{vector<vector<int>> adj(n)} if you then write \texttt{adj[u+1]}?
    \item Your solution passes the first test case but fails the second, and the two test cases are each correct when run alone. What is the most likely bug?
\end{enumerate}

\begin{solutionbox}[Answers]
\begin{enumerate}[leftmargin=*]
    \item An adjacency matrix ($10^6$ cells, $\bigO{1}$ per query). Alternatively an adjacency list with each list sorted and binary search, or a hash set of normalised pairs.
    \item $2m = 4\cdot10^5$ entries (each undirected edge is stored twice), plus $n+1$ list headers.
    \item \texttt{head[2]}$=4$: \texttt{to[4]}$=3$; \texttt{nxt[4]}$=1$: \texttt{to[1]}$=1$; \texttt{nxt[1]}$=-1$: stop. Neighbours 3, 1 --- matching ``star 2: 3 1'' in the output.
    \item Vertex 3 becomes index 4, which is out of bounds for a vector of size 4: undefined behaviour (crash or silent corruption). Stay 0-indexed on LeetCode.
    \item State that survives between tests: an adjacency list, \texttt{visited} array or counter that was not reset.
\end{enumerate}
\end{solutionbox}

%============================================================
% CH1 · LESSON 4 — A FIRST TRAVERSAL
%============================================================
\subsection{A First Traversal}
\label{ssec:ch1-first-traversal}

The most basic question you can ask a graph is \textbf{reachability}: \emph{starting from vertex $s$, which vertices can I get to?} Two strategies answer it, and between them they are the backbone of half of this book:
\begin{itemize}[leftmargin=*]
    \item \textbf{Depth-first search (DFS)} --- walk as far as you can, and back up only when you are stuck. Like exploring a maze with one hand on the wall.
    \item \textbf{Breadth-first search (BFS)} --- explore in rings: first everything one step away, then everything two steps away, and so on. Like a rumour spreading, or ink spreading in water.
\end{itemize}
Both need one extra ingredient, a \texttt{visited} array, and both run in $\bigO{n+m}$. This lesson runs them by hand on the same small graph $G_3$; Chapter~2 (DFS) and Chapter~3 (BFS) study each one in depth.

% local drawing of G3: a pic whose vertex / edge styles are overridden per snapshot
\tikzset{
    g1l-s1/.style=gv, g1l-s2/.style=gv, g1l-s3/.style=gv, g1l-s4/.style=gv,
    g1l-s5/.style=gv, g1l-s6/.style=gv, g1l-s7/.style=gv,
    g1l-e12/.style=ge, g1l-e14/.style=ge, g1l-e23/.style=ge, g1l-e25/.style=ge, g1l-e45/.style=ge, g1l-e56/.style=ge,
    g1l-g3/.pic={
        \node[g1l-s1] (v1) at (0,0) {1};
        \node[g1l-s2] (v2) at (1.4,0.9) {2};
        \node[g1l-s4] (v4) at (1.4,-0.9) {4};
        \node[g1l-s3] (v3) at (2.9,0.9) {3};
        \node[g1l-s5] (v5) at (2.9,-0.9) {5};
        \node[g1l-s6] (v6) at (4.3,-0.9) {6};
        \node[g1l-s7] (v7) at (4.3,0.9) {7};
        \draw[g1l-e12] (v1) -- (v2); \draw[g1l-e14] (v1) -- (v4); \draw[g1l-e23] (v2) -- (v3);
        \draw[g1l-e25] (v2) -- (v5); \draw[g1l-e45] (v4) -- (v5); \draw[g1l-e56] (v5) -- (v6);
    }
}

\begin{diagrambox}
\begin{tikzpicture}[x=1.15cm, y=1cm]
    \pic {g1l-g3};
    \node[draw=darkgray, fill=lightgray!60, rounded corners=2pt, font=\ttfamily\small, align=left, inner sep=4pt, anchor=west] at (6.0,0)
        {7 6\\1 2\\1 4\\2 3\\2 5\\4 5\\5 6\\1};
    \node[font=\scriptsize\color{darkgray}, anchor=west] at (6.0,2.0) {input ($s$ last)};
    \node[note, anchor=west, text width=4.2cm] at (7.4,0) {adj[1] = 2 4\\adj[2] = 1 3 5\\adj[3] = 2\\adj[4] = 1 5\\adj[5] = 2 4 6\\adj[6] = 5\\adj[7] = (empty)};
\end{tikzpicture}
\figcap{$G_3$: $n=7$, $m=6$, start $s=1$. Vertex 7 is isolated, so the answer to ``what can 1 reach?'' is $\{1,2,3,4,5,6\}$. Neighbours are visited in the order of the adjacency lists (input order), which fixes the exact order of both searches.}
\end{diagrambox}

%------------------------------------------------------------
\subsubsection*{Why a visited array is essential}

The naive idea ``from $u$, go to every neighbour, and from there to every neighbour, \ldots'' never stops on a graph with a cycle: from 1 go to 2, from 2 back to 1, from 1 to 2 again, \ldots{} Even on a graph \emph{without} cycles it can explore the same vertex many times. The fix is one boolean per vertex: \textbf{mark a vertex the first time you reach it, and never enter a marked vertex again.}

\begin{diagrambox}
\begin{tikzpicture}[x=1cm, y=1cm]
    \begin{scope}
        \node[gv] (a) at (0,0.6) {1};
        \node[gv] (b) at (1.4,0.6) {2};
        \node[gv] (c) at (0.7,-0.6) {3};
        \draw[ge] (a) -- (b); \draw[ge] (b) -- (c); \draw[ge] (c) -- (a);
        \node[font=\scriptsize\bfseries] at (0.7,-1.5) {a triangle};
    \end{scope}
    \begin{scope}[xshift=3.4cm, yshift=0.9cm]
        \node[gvcur, minimum size=6mm] (r) at (1.8,0) {1};
        \node[gv, minimum size=6mm] (r2) at (0.9,-0.8) {2};
        \node[gv, minimum size=6mm] (r3) at (2.7,-0.8) {3};
        \node[gvB, minimum size=6mm] (r21) at (0.4,-1.6) {1};
        \node[gv, minimum size=6mm] (r23) at (1.4,-1.6) {3};
        \node[gvB, minimum size=6mm] (r31) at (2.2,-1.6) {1};
        \node[gv, minimum size=6mm] (r32) at (3.2,-1.6) {2};
        \foreach \a/\b in {r/r2, r/r3, r2/r21, r2/r23, r3/r31, r3/r32} \draw[faded] (\a) -- (\b);
        \node[font=\scriptsize\color{recursionred}] at (1.8,-2.25) {$\cdots$ forever};
        \node[font=\scriptsize\bfseries] at (1.8,-2.75) {without \texttt{visited}};
    \end{scope}
    \begin{scope}[xshift=8.6cm, yshift=0.9cm]
        \node[gvdone, minimum size=6mm] (r) at (0.9,0) {1};
        \node[gvdone, minimum size=6mm] (r2) at (0.9,-0.8) {2};
        \node[gvdone, minimum size=6mm] (r3) at (0.9,-1.6) {3};
        \draw[gtree] (r) -- (r2); \draw[gtree] (r2) -- (r3);
        \node[font=\scriptsize, anchor=west, text width=3cm] at (1.4,-0.8) {3 calls in total; every other edge leads to a marked vertex and is skipped};
        \node[font=\scriptsize\bfseries] at (1.8,-2.75) {with \texttt{visited}};
    \end{scope}
\end{tikzpicture}
\figcap{On a triangle the unmarked search revisits vertex 1 (red) again and again --- an infinite recursion. With marks, each vertex is entered once.}
\end{diagrambox}

%------------------------------------------------------------
\subsubsection*{Depth-first search (recursive)}

\begin{algorithmbox}[DFS: go deep first]
\texttt{dfs(u)}: mark $u$ as visited; for each neighbour $v$ of $u$: if $v$ is not visited, call \texttt{dfs(v)}.

When \texttt{dfs(v)} returns, the search is \emph{back at $u$} and tries $u$'s next neighbour. The chain of calls currently open (the recursion stack) is always a path from $s$ to the current vertex.
\end{algorithmbox}

\begin{lstlisting}[style=cpp]
#include <bits/stdc++.h>
using namespace std;

// Recursive DFS: which vertices can be reached from s?
// Input: n m, then m undirected edges "u v", then s.
vector<vector<int>> adj;
vector<bool> visited;
vector<int> order;                 // the order in which vertices are entered

void dfs(int u) {
    visited[u] = true;             // mark FIRST, then explore
    order.push_back(u);
    for (int v : adj[u])
        if (!visited[v])
            dfs(v);                // go deeper; come back when v is finished
}

int main() {
    int n, m;
    if (!(cin >> n >> m)) return 0;
    adj.assign(n + 1, vector<int>());
    for (int i = 0; i < m; i++) {
        int u, v;
        cin >> u >> v;
        adj[u].push_back(v);
        adj[v].push_back(u);
    }
    int s;
    cin >> s;
    visited.assign(n + 1, false);
    dfs(s);
    cout << "DFS order:";
    for (int v : order) cout << " " << v;
    cout << "\nunreachable:";
    for (int v = 1; v <= n; v++)
        if (!visited[v]) cout << " " << v;
    cout << "\n";
    return 0;
}
\end{lstlisting}

\begin{lstlisting}[style=trace]
input : the G3 input above
output: DFS order: 1 2 3 5 4 6
        unreachable: 7
\end{lstlisting}

\begin{tracebox}[DFS on $G_3$ from $s=1$, step by step]
Colours: \legendchip{gvcur}{current vertex}, \legendchip{gvseen}{on the stack (entered, not finished)}, \legendchip{gvdone}{finished}; thick blue edges are the \emph{tree edges} along which DFS entered a new vertex.

\begin{center}
\begin{tikzpicture}[x=0.62cm, y=0.62cm, every node/.append style={transform shape=false}]
    \tikzset{gv/.append style={minimum size=6mm, font=\scriptsize\bfseries}}
    % 1
    \begin{scope}[g1l-s1/.style=gvcur]
        \pic {g1l-g3};
        \node[font=\scriptsize, anchor=west] at (-0.3,-1.9) {\textbf{1.} enter 1};
        \node[font=\scriptsize\ttfamily, anchor=west] at (-0.3,-2.6) {stack: 1};
    \end{scope}
    % 2
    \begin{scope}[xshift=4.6cm, g1l-s1/.style=gvseen, g1l-s2/.style=gvcur, g1l-e12/.style=gtree]
        \pic {g1l-g3};
        \node[font=\scriptsize, anchor=west] at (-0.3,-1.9) {\textbf{2.} 1$\to$2 new: enter 2};
        \node[font=\scriptsize\ttfamily, anchor=west] at (-0.3,-2.6) {stack: 1 2};
    \end{scope}
    % 3
    \begin{scope}[xshift=9.2cm, g1l-s1/.style=gvseen, g1l-s2/.style=gvseen, g1l-s3/.style=gvcur,
                  g1l-e12/.style=gtree, g1l-e23/.style=gtree]
        \pic {g1l-g3};
        \node[font=\scriptsize, anchor=west] at (-0.3,-1.9) {\textbf{3.} 2$\to$1 seen; 2$\to$3: enter 3};
        \node[font=\scriptsize\ttfamily, anchor=west] at (-0.3,-2.6) {stack: 1 2 3};
    \end{scope}
    % 4
    \begin{scope}[yshift=-4.4cm, g1l-s1/.style=gvseen, g1l-s2/.style=gvseen, g1l-s3/.style=gvdone, g1l-s5/.style=gvcur,
                  g1l-e12/.style=gtree, g1l-e23/.style=gtree, g1l-e25/.style=gtree]
        \pic {g1l-g3};
        \node[font=\scriptsize, anchor=west] at (-0.3,-1.9) {\textbf{4.} 3 done, back to 2; enter 5};
        \node[font=\scriptsize\ttfamily, anchor=west] at (-0.3,-2.6) {stack: 1 2 5};
    \end{scope}
    % 5
    \begin{scope}[xshift=4.6cm, yshift=-4.4cm, g1l-s1/.style=gvseen, g1l-s2/.style=gvseen, g1l-s3/.style=gvdone, g1l-s5/.style=gvseen, g1l-s4/.style=gvcur,
                  g1l-e12/.style=gtree, g1l-e23/.style=gtree, g1l-e25/.style=gtree, g1l-e45/.style=gtree]
        \pic {g1l-g3};
        \node[font=\scriptsize, anchor=west] at (-0.3,-1.9) {\textbf{5.} 5$\to$2 seen; 5$\to$4: enter 4};
        \node[font=\scriptsize\ttfamily, anchor=west] at (-0.3,-2.6) {stack: 1 2 5 4};
    \end{scope}
    % 6
    \begin{scope}[xshift=9.2cm, yshift=-4.4cm, g1l-s1/.style=gvseen, g1l-s2/.style=gvseen, g1l-s3/.style=gvdone, g1l-s5/.style=gvseen, g1l-s4/.style=gvdone, g1l-s6/.style=gvcur,
                  g1l-e12/.style=gtree, g1l-e23/.style=gtree, g1l-e25/.style=gtree, g1l-e45/.style=gtree, g1l-e56/.style=gtree]
        \pic {g1l-g3};
        \node[font=\scriptsize, anchor=west] at (-0.3,-1.9) {\textbf{6.} 4: both seen, done; enter 6};
        \node[font=\scriptsize\ttfamily, anchor=west] at (-0.3,-2.6) {stack: 1 2 5 6};
    \end{scope}
\end{tikzpicture}
\end{center}

After step 6: vertex 6 sees only 5 (visited) and finishes; then 5, 2 finish; back at 1, the neighbour 4 is already visited, and 1 finishes. Vertex 7 is never touched. Entry order \textbf{1, 2, 3, 5, 4, 6}; maximum stack depth 4 (\texttt{1 2 5 4}).
\end{tracebox}

Notice what happened at step~5: DFS reached 4 through the long way round (1--2--5--4), not via the direct edge 1--4. DFS finds \emph{a} path to every reachable vertex, but not necessarily the shortest one.

%------------------------------------------------------------
\subsubsection*{Breadth-first search}

\begin{algorithmbox}[BFS: explore in rings]
Put $s$ in a queue with $\dist{s}=0$. Repeatedly take the vertex $u$ at the \emph{front} of the queue; for each neighbour $v$ that has not been discovered, set $\dist{v}=\dist{u}+1$ and push $v$ at the \emph{back}.

Because the queue is first-in first-out, all vertices at distance $d$ leave the queue before any vertex at distance $d+1$. So \texttt{dist[v]} is the \textbf{number of edges on a shortest path} from $s$ to $v$ --- a bonus DFS does not give.
\end{algorithmbox}

\begin{lstlisting}[style=cpp]
#include <bits/stdc++.h>
using namespace std;

// BFS: reachability from s plus the number of edges on a shortest path.
// Input: n m, then m undirected edges "u v", then s.
int main() {
    int n, m;
    if (!(cin >> n >> m)) return 0;
    vector<vector<int>> adj(n + 1);
    for (int i = 0; i < m; i++) {
        int u, v;
        cin >> u >> v;
        adj[u].push_back(v);
        adj[v].push_back(u);
    }
    int s;
    cin >> s;
    vector<int> dist(n + 1, -1);   // -1 = not discovered yet (doubles as visited)
    queue<int> q;
    dist[s] = 0;
    q.push(s);
    cout << "BFS order:";
    while (!q.empty()) {
        int u = q.front();
        q.pop();
        cout << " " << u;
        for (int v : adj[u]) {
            if (dist[v] == -1) {           // first time we see v
                dist[v] = dist[u] + 1;     // mark when PUSHED, not when popped
                q.push(v);
            }
        }
    }
    cout << "\ndist:";
    for (int v = 1; v <= n; v++) cout << " " << dist[v];
    cout << "\n";
    return 0;
}
\end{lstlisting}

\begin{lstlisting}[style=trace]
input : the G3 input above
output: BFS order: 1 2 4 3 5 6
        dist: 0 1 2 1 2 3 -1
\end{lstlisting}

\begin{tracebox}[BFS on $G_3$ from $s=1$, step by step]
Colours: \legendchip{gvcur}{popped now}, \legendchip{gvseen}{in the queue}, \legendchip{gvdone}{popped earlier}; green numbers are \texttt{dist}; thick blue edges are the edges along which a vertex was discovered.

\begin{center}
\begin{tikzpicture}[x=0.62cm, y=0.62cm]
    \tikzset{gv/.append style={minimum size=6mm, font=\scriptsize\bfseries}}
    % 1
    \begin{scope}[g1l-s1/.style=gvcur, g1l-s2/.style=gvseen, g1l-s4/.style=gvseen, g1l-e12/.style=gtree, g1l-e14/.style=gtree]
        \pic {g1l-g3};
        \foreach \v/\d in {1/0, 2/1, 4/1} \node[gdist, above left=-1pt of v\v] {\d};
        \node[font=\scriptsize, anchor=west] at (-0.3,-1.9) {\textbf{1.} pop 1: push 2, 4};
        \node[font=\scriptsize\ttfamily, anchor=west] at (-0.3,-2.6) {queue: 2 4};
    \end{scope}
    % 2
    \begin{scope}[xshift=4.6cm, g1l-s1/.style=gvdone, g1l-s2/.style=gvcur, g1l-s4/.style=gvseen, g1l-s3/.style=gvseen, g1l-s5/.style=gvseen,
                  g1l-e12/.style=gtree, g1l-e14/.style=gtree, g1l-e23/.style=gtree, g1l-e25/.style=gtree]
        \pic {g1l-g3};
        \foreach \v/\d in {1/0, 2/1, 4/1, 3/2, 5/2} \node[gdist, above left=-1pt of v\v] {\d};
        \node[font=\scriptsize, anchor=west] at (-0.3,-1.9) {\textbf{2.} pop 2: 1 seen; push 3, 5};
        \node[font=\scriptsize\ttfamily, anchor=west] at (-0.3,-2.6) {queue: 4 3 5};
    \end{scope}
    % 3
    \begin{scope}[xshift=9.2cm, g1l-s1/.style=gvdone, g1l-s2/.style=gvdone, g1l-s4/.style=gvcur, g1l-s3/.style=gvseen, g1l-s5/.style=gvseen,
                  g1l-e12/.style=gtree, g1l-e14/.style=gtree, g1l-e23/.style=gtree, g1l-e25/.style=gtree]
        \pic {g1l-g3};
        \foreach \v/\d in {1/0, 2/1, 4/1, 3/2, 5/2} \node[gdist, above left=-1pt of v\v] {\d};
        \node[font=\scriptsize, anchor=west] at (-0.3,-1.9) {\textbf{3.} pop 4: 1, 5 seen};
        \node[font=\scriptsize\ttfamily, anchor=west] at (-0.3,-2.6) {queue: 3 5};
    \end{scope}
    % 4
    \begin{scope}[yshift=-4.4cm, g1l-s1/.style=gvdone, g1l-s2/.style=gvdone, g1l-s4/.style=gvdone, g1l-s3/.style=gvcur, g1l-s5/.style=gvseen,
                  g1l-e12/.style=gtree, g1l-e14/.style=gtree, g1l-e23/.style=gtree, g1l-e25/.style=gtree]
        \pic {g1l-g3};
        \foreach \v/\d in {1/0, 2/1, 4/1, 3/2, 5/2} \node[gdist, above left=-1pt of v\v] {\d};
        \node[font=\scriptsize, anchor=west] at (-0.3,-1.9) {\textbf{4.} pop 3: 2 seen};
        \node[font=\scriptsize\ttfamily, anchor=west] at (-0.3,-2.6) {queue: 5};
    \end{scope}
    % 5
    \begin{scope}[xshift=4.6cm, yshift=-4.4cm, g1l-s1/.style=gvdone, g1l-s2/.style=gvdone, g1l-s4/.style=gvdone, g1l-s3/.style=gvdone, g1l-s5/.style=gvcur, g1l-s6/.style=gvseen,
                  g1l-e12/.style=gtree, g1l-e14/.style=gtree, g1l-e23/.style=gtree, g1l-e25/.style=gtree, g1l-e56/.style=gtree]
        \pic {g1l-g3};
        \foreach \v/\d in {1/0, 2/1, 4/1, 3/2, 5/2, 6/3} \node[gdist, above left=-1pt of v\v] {\d};
        \node[font=\scriptsize, anchor=west] at (-0.3,-1.9) {\textbf{5.} pop 5: 2, 4 seen; push 6};
        \node[font=\scriptsize\ttfamily, anchor=west] at (-0.3,-2.6) {queue: 6};
    \end{scope}
    % 6
    \begin{scope}[xshift=9.2cm, yshift=-4.4cm, g1l-s1/.style=gvdone, g1l-s2/.style=gvdone, g1l-s4/.style=gvdone, g1l-s3/.style=gvdone, g1l-s5/.style=gvdone, g1l-s6/.style=gvcur, g1l-s7/.style=gvdead,
                  g1l-e12/.style=gtree, g1l-e14/.style=gtree, g1l-e23/.style=gtree, g1l-e25/.style=gtree, g1l-e56/.style=gtree]
        \pic {g1l-g3};
        \foreach \v/\d in {1/0, 2/1, 4/1, 3/2, 5/2, 6/3} \node[gdist, above left=-1pt of v\v] {\d};
        \node[font=\scriptsize, anchor=west] at (-0.3,-1.9) {\textbf{6.} pop 6: 5 seen};
        \node[font=\scriptsize\ttfamily, anchor=west] at (-0.3,-2.6) {queue: (empty) --- stop};
    \end{scope}
\end{tikzpicture}
\end{center}

Pop order \textbf{1, 2, 4, 3, 5, 6}; final \texttt{dist} $=(0,1,2,1,2,3,-1)$ for vertices $1..7$. Vertex 7 keeps $-1$: unreachable.
\end{tracebox}

\begin{diagrambox}
\begin{tikzpicture}[x=1.3cm, y=1cm]
    \foreach \d/\x in {0/0, 1/1.6, 2/3.2, 3/4.8} {
        \fill[primaryblue!\the\numexpr 6+5*\d\relax, rounded corners=4pt] (\x-0.55,-1.5) rectangle (\x+0.55,1.5);
        \node[font=\scriptsize\bfseries\color{deepblue}] at (\x,1.8) {dist $\d$};
    }
    \node[gvsrc] (1) at (0,0) {1};
    \node[gv] (2) at (1.6,0.7) {2};
    \node[gv] (4) at (1.6,-0.7) {4};
    \node[gv] (3) at (3.2,0.7) {3};
    \node[gv] (5) at (3.2,-0.7) {5};
    \node[gv] (6) at (4.8,-0.7) {6};
    \draw[gtree] (1) -- (2); \draw[gtree] (1) -- (4); \draw[gtree] (2) -- (3); \draw[gtree] (2) -- (5); \draw[gtree] (5) -- (6);
    \draw[gfaded, dashed] (4) -- (5);
    \node[gvdead] (7) at (6.4,0) {7};
    \node[font=\scriptsize\color{darkgray}] at (6.4,-0.7) {never reached};
\end{tikzpicture}
\figcap{BFS layers of $G_3$. Every edge joins two vertices whose distances differ by at most 1; the dashed edge 4--5 joins layer 1 to layer 2 but was not used to discover anyone.}
\end{diagrambox}

%------------------------------------------------------------
\subsubsection*{Same graph, two different trees}

Each search discovers every reachable vertex through exactly one edge --- its \textbf{tree edge}. The tree edges form a tree (the \textbf{DFS tree} or \textbf{BFS tree}) rooted at $s$ that spans everything reachable. The two trees differ:

\begin{diagrambox}
\begin{tikzpicture}[x=1cm, y=0.9cm]
    \begin{scope}
        \node[gvsrc] (1) at (0,0) {1};
        \node[gv] (2) at (0,-1) {2};
        \node[gv] (3) at (-0.9,-2) {3};
        \node[gv] (5) at (0.9,-2) {5};
        \node[gv] (4) at (0.3,-3) {4};
        \node[gv] (6) at (1.5,-3) {6};
        \draw[gtree] (1) -- (2); \draw[gtree] (2) -- (3); \draw[gtree] (2) -- (5); \draw[gtree] (5) -- (4); \draw[gtree] (5) -- (6);
        \draw[gback, -] (4) to[bend left=50] (1);
        \node[font=\small\bfseries] at (0.3,-3.9) {DFS tree: depth 3};
        \node[font=\scriptsize\color{darkgray}] at (0.3,-4.35) {non-tree edge 1--4 (dashed)};
    \end{scope}
    \begin{scope}[xshift=6cm]
        \node[gvsrc] (1) at (0,0) {1};
        \node[gv] (2) at (-0.8,-1) {2};
        \node[gv] (4) at (0.8,-1) {4};
        \node[gv] (3) at (-1.4,-2) {3};
        \node[gv] (5) at (-0.2,-2) {5};
        \node[gv] (6) at (-0.2,-3) {6};
        \draw[gtree] (1) -- (2); \draw[gtree] (1) -- (4); \draw[gtree] (2) -- (3); \draw[gtree] (2) -- (5); \draw[gtree] (5) -- (6);
        \draw[gback, -] (4) -- (5);
        \node[font=\small\bfseries] at (0,-3.9) {BFS tree: depth = distance};
        \node[font=\scriptsize\color{darkgray}] at (0,-4.35) {non-tree edge 4--5 (dashed)};
    \end{scope}
\end{tikzpicture}
\figcap{Both trees contain the same six vertices (everything reachable from 1) and $6-1=5$ edges. In the BFS tree the depth of a vertex equals its distance from $s$; in the DFS tree vertex 4 sits at depth 3 although $\dist{4}=1$.}
\end{diagrambox}

\begin{comparisonbox}[DFS vs.\ BFS at a glance]
\small
\begin{tabularx}{\linewidth}{@{}>{\raggedright\arraybackslash}p{4.2cm} X X@{}}
\toprule
 & \textbf{DFS} & \textbf{BFS} \\
\midrule
data structure & recursion (a stack) & a queue \\
order on $G_3$ & 1, 2, 3, 5, 4, 6 & 1, 2, 4, 3, 5, 6 \\
reachability & yes & yes \\
shortest paths (unweighted) & no & \textbf{yes}: \texttt{dist[]} \\
code length & shortest & a few lines more \\
deep graphs ($10^6$ vertices on a path) & recursion may overflow & safe \\
strengths later in the book & cycles, bridges, SCC, trees (Ch.~2) & distances, grids, states (Ch.~3) \\
\bottomrule
\end{tabularx}
\end{comparisonbox}

%------------------------------------------------------------
\subsubsection*{Why both run in $\bigO{n+m}$}

\begin{invariantbox}[Each vertex is expanded once, each edge is looked at twice]
A vertex is marked the moment it is first reached and is never entered (DFS) or pushed (BFS) again. So each vertex $u$ is \emph{expanded} --- its adjacency list scanned --- at most once, at a cost of $\deg(u)+\bigO{1}$. Summing over all vertices, the total work is
\[
    \sum_{u}\big(\deg(u)+1\big) \;=\; 2m + n \;=\; \bigO{n+m}
\]
by the handshake lemma (for a directed graph the sum of out-degrees is $m$). On $G_3$ both searches perform exactly $2m=12$ neighbour checks.
\end{invariantbox}

\begin{warningbox}[In BFS, mark when you push, not when you pop]
If you set \texttt{dist[v]} only when $v$ is popped, the same vertex can be pushed by several neighbours before it is first popped. The answers stay correct, but the queue can grow to $\bigO{m}$ entries and on dense graphs the running time explodes. Mark at push time, as in the code above.
\end{warningbox}

%------------------------------------------------------------
\subsubsection*{The recursion-depth trap}

The recursive DFS keeps one stack frame for every vertex on the current path. On a path graph 1--2--3--$\cdots$--$n$ with $n=10^6$ the recursion goes $10^6$ levels deep. Codeforces gives a large stack (256 MB) and this is usually fine there, but on a default 1--8 MB stack (Windows at home, many other judges, LeetCode) it is a \stackoverflow{} --- reported as a mysterious runtime error.

\begin{diagrambox}
\begin{tikzpicture}[x=1cm, y=1cm]
    \foreach \i/\l in {0/1, 1/2, 2/3, 3/4} \node[gvseen] (p\i) at (\i*1.3,0) {\l};
    \node (dots) at (5.2,0) {$\cdots$};
    \node[gvcur, minimum size=9mm, font=\scriptsize\bfseries] (pn) at (6.6,0) {$10^6$};
    \foreach \a/\b in {p0/p1, p1/p2, p2/p3, p3/dots, dots/pn} \draw[gtree] (\a) -- (\b);
    \foreach \i/\h in {0/0.35, 1/0.7, 2/1.05, 3/1.4}
        \fill[recursionred!25, draw=recursionred] (8.2,-0.6+\h-0.35) rectangle (10.2,-0.6+\h);
    \node[font=\scriptsize] at (9.2,1.15) {$\vdots$};
    \fill[recursionred!25, draw=recursionred] (8.2,1.35) rectangle (10.2,1.7);
    \node[font=\scriptsize\ttfamily] at (9.2,-0.42) {dfs(1)};
    \node[font=\scriptsize\ttfamily] at (9.2,-0.07) {dfs(2)};
    \node[font=\scriptsize\ttfamily] at (9.2,0.28) {dfs(3)};
    \node[font=\scriptsize\ttfamily] at (9.2,0.63) {dfs(4)};
    \node[font=\scriptsize\ttfamily] at (9.2,1.52) {dfs($10^6$)};
    \node[font=\scriptsize\bfseries\color{recursionred}] at (9.2,2.0) {$10^6$ open frames};
\end{tikzpicture}
\figcap{DFS on a long path: every call is still open when the next one starts.}
\end{diagrambox}

When the depth can exceed about $10^5$, either use BFS (if the order does not matter, as for plain reachability) or replace the recursion by an explicit stack. The iterative version below remembers, for each vertex on the stack, which neighbour to try next, so it visits vertices in \emph{exactly} the same order as the recursive one.

\begin{lstlisting}[style=cpp]
#include <bits/stdc++.h>
using namespace std;

// Iterative DFS with an explicit stack: no recursion, so no stack overflow
// even on a path of 10^6 vertices. It visits exactly the same vertices.
int main() {
    int n, m;
    if (!(cin >> n >> m)) return 0;
    vector<vector<int>> adj(n + 1);
    for (int i = 0; i < m; i++) {
        int u, v;
        cin >> u >> v;
        adj[u].push_back(v);
        adj[v].push_back(u);
    }
    int s;
    cin >> s;
    vector<bool> visited(n + 1, false);
    vector<size_t> it(n + 1, 0);     // it[u] = next neighbour index of u to try
    vector<int> st;                  // the "call stack", made by hand
    visited[s] = true;
    st.push_back(s);
    cout << "DFS order: " << s;
    while (!st.empty()) {
        int u = st.back();
        if (it[u] == adj[u].size()) { st.pop_back(); continue; }   // u finished
        int v = adj[u][it[u]++];
        if (!visited[v]) {
            visited[v] = true;
            cout << " " << v;
            st.push_back(v);                                       // "call" dfs(v)
        }
    }
    cout << "\n";
    return 0;
}
\end{lstlisting}

\begin{lstlisting}[style=trace]
input : the G3 input above
output: DFS order: 1 2 3 5 4 6          (identical to the recursive version)
\end{lstlisting}

Chapters~2 and~3 go much deeper: entry/exit times, edge classification, cycle detection and bridges for DFS; multi-source BFS, 0-1 BFS and BFS on states for BFS. For the rest of this chapter, the two programs above are all you need --- Patterns~2--4 are applications of them.

%------------------------------------------------------------
\subsubsection*{Check your understanding}

\begin{enumerate}[leftmargin=*]
    \item In $G_3$, change the input order so that the edge \texttt{1 4} comes before \texttt{1 2}. What is the new DFS order from 1?
    \item On a star (vertex 1 joined to $2,\dots,n$), what is the maximum recursion depth of DFS from 1? From 2?
    \item A graph has $n=10^5$ vertices and $m=0$ edges. How much work does a BFS from vertex 1 do? How much work does \emph{looping over all vertices and starting a BFS from every unvisited one} do?
    \item Why can BFS compute shortest paths, but DFS cannot? Answer with $G_3$.
    \item You run BFS and forget to set \texttt{dist[s] = 0} before the loop (it stays $-1$). What goes wrong?
\end{enumerate}

\begin{solutionbox}[Answers]
\begin{enumerate}[leftmargin=*]
    \item Then \texttt{adj[1] = 4 2}, and DFS goes $1\to4\to5\to2\to3$, back to 5, then $5\to6$: order \textbf{1, 4, 5, 2, 3, 6}. (Running the program on the reordered input prints exactly this.) A different but equally valid DFS.
    \item From 1: depth 1 (every leaf is a child of 1). From 2: depth 2 ($2\to1\to$ each other leaf).
    \item One BFS from 1 does $\bigO{1}$ work (no edges). Starting a search from every unvisited vertex costs $\bigO{n+m}$ in total, because each vertex is still expanded only once over all the searches --- this is how components are counted (Pattern~2, Chapter~4).
    \item BFS pops vertices in order of distance, so a vertex is first discovered along a shortest path. DFS discovered 4 along 1--2--5--4 (3 edges) although the edge 1--4 gives distance 1.
    \item Everything shifts and $s$ is processed twice. The neighbours of $s$ get $-1+1=0$; $s$ still looks undiscovered, so the first neighbour expanded pushes it again with distance 1. On $G_3$ the modified program prints \texttt{BFS order: 1 2 4 1 3 5 6} and \texttt{dist: 1 0 1 0 1 2 -1}. Always initialise the source before the loop.
\end{enumerate}
\end{solutionbox}

%============================================================
% CH1 · LESSON 5 — MODELLING: SEEING THE GRAPH
%============================================================
\subsection{Modelling: Seeing the Graph}
\label{ssec:ch1-modelling}

Few contest problems say ``here is a graph''. Most describe a board, a number that can be changed, a word that can be edited, a lock, a game --- and the graph is hidden inside. Finding it is the most important step of every solution in this book, which is why each solved problem has a \emph{Graph Model} box. This lesson teaches the habit: \textbf{decide what a vertex is, decide what an edge is, estimate how many there are, and recognise the shape.} Then a standard algorithm --- often the DFS or BFS of Lesson~\ref{ssec:ch1-first-traversal} --- does the rest.

\begin{definitionbox}[Implicit graph]
A graph is \textbf{implicit} when it is not given as a list of edges but by a \emph{rule}: the vertices are all objects of some kind (cells, numbers, strings, configurations), and a function \texttt{neighbours(x)} computes the edges out of $x$ on demand. The graph is never stored; a traversal only needs \texttt{visited} (or \texttt{dist}) for the vertices it actually reaches.
\end{definitionbox}

%------------------------------------------------------------
\subsubsection*{Grids: cells are vertices, shared sides are edges}

A grid of $R\times C$ cells is the most common implicit graph. Each free cell is a vertex; with \textbf{4-neighbourhood} two cells are adjacent if they share a side, with \textbf{8-neighbourhood} also if they share a corner (king moves). The statement decides which: ``up, down, left or right'' means 4, ``including diagonals'' or ``touching at a corner'' means 8.

\begin{diagrambox}
\begin{tikzpicture}[x=0.8cm, y=0.8cm]
    \begin{scope}
        \foreach \r in {0,1,2} \foreach \c in {0,1,2} \node[gcell] at (\c,-\r) {};
        \foreach \c/\r in {1/0, 1/2, 0/1, 2/1} \node[gcell, fill=hitfill] at (\c,-\r) {};
        \node[gcell, fill=currentfill] at (1,-1) {};
        \foreach \dx/\dy in {0/1, 0/-1, -1/0, 1/0} \draw[dep, line width=1.2pt] (1+0.2*\dx,-1+0.2*\dy) -- (1+0.75*\dx,-1+0.75*\dy);
        \node[font=\small\bfseries] at (1,-3) {4 neighbours};
        \node[font=\scriptsize\ttfamily] at (1,-3.55) {dr = \{-1,1,0,0\}};
        \node[font=\scriptsize\ttfamily] at (1,-3.95) {dc = \{0,0,-1,1\}};
    \end{scope}
    \begin{scope}[xshift=3.8cm]
        \foreach \r in {0,1,2} \foreach \c in {0,1,2} \node[gcell, fill=hitfill] at (\c,-\r) {};
        \node[gcell, fill=currentfill] at (1,-1) {};
        \foreach \dx/\dy in {0/1, 0/-1, -1/0, 1/0, 1/1, 1/-1, -1/1, -1/-1} \draw[dep, line width=1.2pt] (1+0.2*\dx,-1+0.2*\dy) -- (1+0.75*\dx,-1+0.75*\dy);
        \node[font=\small\bfseries] at (1,-3) {8 neighbours};
        \node[font=\scriptsize\ttfamily] at (1,-3.55) {dr, dc in -1..1};
        \node[font=\scriptsize\ttfamily] at (1,-3.95) {skip dr = dc = 0};
    \end{scope}
    \begin{scope}[xshift=8cm, yshift=0.2cm]
        \foreach \r in {0,1,2} \foreach \c in {0,1,2} \draw[ge] (\c*0.9,-\r*0.9) -- (\c*0.9+0.9,-\r*0.9);
        \foreach \r in {0,1} \foreach \c in {0,1,2,3} \draw[ge] (\c*0.9,-\r*0.9) -- (\c*0.9,-\r*0.9-0.9);
        \foreach \r in {0,1,2} \foreach \c in {0,1,2,3} \fill[primaryblue] (\c*0.9,-\r*0.9) circle (2.6pt);
        \node[font=\small\bfseries] at (1.35,-2.4) {$3\times4$ grid as a graph};
        \node[font=\scriptsize] at (1.35,-2.95) {$|V|=12$, $|E|=17$ (4 nbrs)};
        \node[font=\scriptsize] at (1.35,-3.35) {$|E|=29$ with 8 neighbours};
    \end{scope}
\end{tikzpicture}
\figcap{The two neighbourhoods and the graph of a $3\times 4$ grid with no walls: $3\cdot3$ horizontal $+\,2\cdot4$ vertical $=17$ edges; adding the $2\cdot(2\cdot3)=12$ diagonals gives 29.}
\end{diagrambox}

In general an $R\times C$ grid has $RC$ vertices and fewer than $2RC$ edges (4-neighbourhood) --- always sparse. A $1000\times1000$ grid has $10^6$ vertices and $1\,998\,000$ edges: one BFS is a few million steps, but a DFS on it can recurse $10^6$ levels deep (Lesson~\ref{ssec:ch1-first-traversal}). Pattern~3 is devoted to grids.

%------------------------------------------------------------
\subsubsection*{Numbers: values are vertices, operations are edges}

``You start with $a$. In one move you may add 1 or multiply by 3. What is the fewest number of moves to reach $b$?'' The vertices are the integers; each operation is a directed edge $x\to x+1$ and $x\to 3x$. Fewest moves $=$ shortest path in an unweighted graph $=$ BFS.

\begin{diagrambox}
\begin{tikzpicture}[x=1.05cm, y=1cm]
    \foreach \i in {1,...,12} {
        \ifnum\i=1 \node[gvsrc, minimum size=6.5mm] (n\i) at (\i,0) {\i};
        \else\ifnum\i=10 \node[gvdst, minimum size=6.5mm] (n\i) at (\i,0) {\i};
        \else \node[gv, minimum size=6.5mm] (n\i) at (\i,0) {\i}; \fi\fi
    }
    \foreach \i [evaluate=\i as \j using int(\i+1)] in {1,...,11} \draw[gde] (n\i) -- (n\j);
    \draw[gde, draw=softpurple] (n2) to[bend left=45] (n6);
    \draw[gde, draw=softpurple] (n4) to[bend right=35] (n12);
    \draw[gpath, -{Stealth[length=2.6mm]}] (n1) to[bend left=50] (n3);
    \draw[gpath, -{Stealth[length=2.6mm]}] (n3) to[bend left=40] (n9);
    \draw[gpath, -{Stealth[length=2.6mm]}] (n9) -- (n10);
    \foreach \i/\d in {1/0, 2/1, 3/1, 9/2, 10/3} \node[gdist, below=1pt of n\i] {\d};
    \node[font=\scriptsize\color{softpurple}] at (8,1.75) {$\times 3$ arcs};
    \node[font=\scriptsize\color{darkgray}] at (6.5,-0.75) {$+1$ arcs};
\end{tikzpicture}
\figcap{The first 12 vertices of the ``$+1$ / $\times3$'' graph ($\times3$ arcs curved; the ones used by the answer are orange). The shortest route from 1 to 10 is $1\to3\to9\to10$ (3 moves; BFS distances in green). Walking with $+1$ only would take 9 moves.}
\end{diagrambox}

The neighbours are computed, not stored. One extra thought is needed: \emph{which values can ever be useful?} Both operations only increase $x$, so values above $b$ can be cut off, and the graph has at most $b$ vertices and $2b$ edges.

\begin{lstlisting}[style=cpp]
#include <bits/stdc++.h>
using namespace std;

// An implicit graph on numbers: from x you may go to x+1 or to 3x.
// Fewest operations from a to b (1 <= a <= b <= 10^6). The graph is never stored:
// the neighbours of x are computed on the fly.
int main() {
    int a, b;
    if (!(cin >> a >> b)) return 0;
    const int LIM = b;                       // larger values can never come back down
    vector<int> dist(LIM + 1, -1), par(LIM + 1, -1);
    queue<int> q;
    dist[a] = 0;
    q.push(a);
    while (!q.empty()) {
        int x = q.front();
        q.pop();
        long long nb[2] = {x + 1LL, 3LL * x};           // the two edges out of x
        for (int k = 0; k < 2; k++) {
            if (nb[k] > LIM) continue;                  // outside the useful range
            int y = (int)nb[k];
            if (dist[y] == -1) {
                dist[y] = dist[x] + 1;
                par[y] = x;
                q.push(y);
            }
        }
    }
    cout << "operations: " << dist[b] << "\npath:";
    vector<int> path;
    for (int x = b; x != -1; x = par[x]) path.push_back(x);
    reverse(path.begin(), path.end());
    for (int x : path) cout << " " << x;
    cout << "\n";
    return 0;
}
\end{lstlisting}

\begin{lstlisting}[style=trace]
input : 1 31      output: operations: 5   path: 1 3 9 10 30 31
input : 2 100     output: operations: 7   path: 2 3 9 10 11 33 99 100
input : 5 5       output: operations: 0   path: 5
\end{lstlisting}

The array \texttt{par} remembers through which edge each vertex was discovered; walking it back from $b$ rebuilds a shortest path. This ``parent array'' trick works for every BFS.

\begin{warningbox}[Bound the state space before you code]
If an operation can \emph{decrease} the value or if values may temporarily exceed the target, the cut-off is less obvious. Prove a bound (for example ``it never helps to go above $2b$'') or the BFS either misses the answer or runs out of memory. Pattern~4 works through these arguments.
\end{warningbox}

%------------------------------------------------------------
\subsubsection*{Strings: words are vertices, small edits are edges}

Words of equal length are adjacent when they differ in exactly one letter. Changing one letter at a time, \textsf{cold} becomes \textsf{warm} in four steps:

\begin{diagrambox}
\begin{tikzpicture}[x=2.1cm, y=1cm,
    w/.style={rectangle, rounded corners=3pt, draw=darkgray!85, fill=white, line width=0.9pt, font=\small\sffamily\bfseries, minimum width=13mm, minimum height=7mm}]
    \node[w, draw=successgreen, fill=basefill] (a) at (0,0) {cold};
    \node[w] (b) at (1,0) {cord};
    \node[w] (c) at (2,0) {card};
    \node[w] (d) at (3,0) {ward};
    \node[w, draw=accentorange, fill=answerfill] (e) at (4,0) {warm};
    \draw[gpath] (a) -- node[above=3pt, font=\scriptsize, text=darkgray] {l$\to$r} (b);
    \draw[gpath] (b) -- node[above=3pt, font=\scriptsize, text=darkgray] {o$\to$a} (c);
    \draw[gpath] (c) -- node[above=3pt, font=\scriptsize, text=darkgray] {c$\to$w} (d);
    \draw[gpath] (d) -- node[above=3pt, font=\scriptsize, text=darkgray] {d$\to$m} (e);
    \node[w, draw=midgray, text=midgray] (f) at (0.5,-1.3) {colt};
    \node[w, draw=midgray, text=midgray] (g) at (2,-1.3) {cart};
    \node[w, draw=midgray, text=midgray] (h) at (3.5,-1.3) {word};
    \draw[gfaded] (a) -- (f); \draw[gfaded] (c) -- (g); \draw[gfaded] (d) -- (h); \draw[gfaded] (b) -- (h);
\end{tikzpicture}
\figcap{Part of the word graph. A shortest path from \textsf{cold} to \textsf{warm} has 4 edges; grey words are other neighbours that a BFS would also explore (\textsf{word} is adjacent to both \textsf{cord} and \textsf{ward}).}
\end{diagrambox}

How to generate the neighbours of a word of length $L$? Either compare it with every word of the dictionary ($\bigO{N\cdot L}$ per word, $\bigO{N^2L}$ in total), or try all $25L$ one-letter changes (100 candidates for a 4-letter word) and look each one up in a hash set. Which is faster depends on $N$ and $L$ --- exactly the kind of estimate this lesson is about.

%------------------------------------------------------------
\subsubsection*{Game and puzzle states: tuples encoded as integers}

When the ``position'' of a puzzle consists of several numbers --- two tokens on a board, the amounts of water in some jugs, a set of switched-on lamps --- the vertex is the whole \emph{tuple}. To use plain arrays for \texttt{visited} and \texttt{dist}, encode the tuple as one integer, like the digits of a number in a mixed base.

\begin{diagrambox}
\begin{tikzpicture}[x=0.55cm, y=0.55cm]
    \foreach \r in {0,...,7} \foreach \c in {0,...,7} {
        \pgfmathtruncatemacro{\parity}{mod(\r+\c,2)}
        \ifnum\parity=0 \fill[lightgray] (\c,-\r) rectangle (\c+1,-\r-1); \fi
    }
    \draw[darkgray] (0,0) rectangle (8,-8);
    \node[gvA, minimum size=4.5mm, font=\scriptsize\bfseries] at (5.5,-2.5) {A};
    \node[gvB, minimum size=4.5mm, font=\scriptsize\bfseries] at (0.5,-7.5) {B};
    \foreach \c in {0,...,7} \node[font=\tiny\color{darkgray}] at (\c+0.5,0.45) {\c};
    \foreach \r in {0,...,7} \node[font=\tiny\color{darkgray}] at (-0.45,-\r-0.5) {\r};
    \begin{scope}[xshift=5.6cm, yshift=-0.5cm]
        \node[anchor=west, font=\small] at (0,0) {state $=(r_A, c_A, r_B, c_B) = (2,5,7,0)$};
        \node[anchor=west, font=\small] at (0,-1.4) {$\mathit{id} = ((r_A\cdot8 + c_A)\cdot 8 + r_B)\cdot 8 + c_B$};
        \node[anchor=west, font=\small] at (0,-2.6) {$\phantom{\mathit{id}} = ((2\cdot8+5)\cdot8+7)\cdot8+0 = \mathbf{1400}$};
        \node[anchor=west, font=\small] at (0,-3.9) {decode: $c_B=\mathit{id}\bmod 8$, $r_B=\lfloor \mathit{id}/8\rfloor \bmod 8$, \ldots};
        \node[anchor=west, font=\small] at (0,-5.2) {number of states: $8^4 = 4096$};
        \node[anchor=west, font=\scriptsize\color{darkgray}, text width=7cm] at (0,-6.7) {one move of A or B changes the state: an edge to another id, generated on the fly by decoding, moving, re-encoding};
    \end{scope}
\end{tikzpicture}
\figcap{Two tokens on an $8\times8$ board: the vertex is the pair of positions, encoded in base 8. Decoding 1400 gives back $(2,5,7,0)$.}
\end{diagrambox}

For a set of $k$ switches use a bitmask ($2^k$ states); for small permutations use a \texttt{map} or a rank; for two jugs with capacities $A$ and $B$ use $\mathit{id}=x\cdot(B+1)+y$. The rule is the same: \textbf{the id must be unique per state and small enough to index an array.}

%------------------------------------------------------------
\subsubsection*{The modelling checklist}

\begin{checklistbox}[What is a vertex? What is an edge?]
\begin{enumerate}[leftmargin=*]
    \item \textbf{Vertex}: what does a ``situation'' consist of? Everything that influences the future must be in it (position \emph{and} keys collected \emph{and} direction faced, \ldots), and nothing else.
    \item \textbf{Edge}: what single action changes the situation? Is it reversible (undirected) or not (directed)? Does it have a cost (weighted) or does every action count as 1?
    \item \textbf{Question}: is it reachability (DFS/BFS), fewest actions (BFS), cheapest actions (Dijkstra, Ch.~7), number of groups (components, Ch.~4), a cycle, an order (Ch.~5)?
    \item \textbf{Size}: $|V|$ and $|E|$ --- do \texttt{visited}/\texttt{dist} arrays of size $|V|$ fit in memory, and does $\bigO{|V|+|E|}$ fit in the time limit?
    \item \textbf{Explicit or implicit}: store edges only when they are given; otherwise generate neighbours on the fly.
\end{enumerate}
\end{checklistbox}

%------------------------------------------------------------
\subsubsection*{Estimating $|V|$ and $|E|$ from the constraints}

A C++ program does roughly $10^8$ simple operations per second. Before coding a traversal, multiply out the size of the state space:

\begin{table}[H]
\centering\small
\begin{tabularx}{\linewidth}{@{}X l r r l@{}}
\toprule
\textbf{Model} & \textbf{vertex} & $|V|$ & $|E|$ (arcs) & \textbf{verdict} \\
\midrule
$1000\times1000$ grid, 4 neighbours & cell & $10^6$ & $\approx 4\cdot10^6$ & BFS: fine \\
values $1..10^6$, two operations & number & $10^6$ & $2\cdot10^6$ & BFS: fine \\
4-digit lock, turn one wheel $\pm1$ & 4 digits & $10^4$ & $8\cdot10^4$ & trivial \\
two tokens on an $8\times8$ board & 4-tuple & $4096$ & small & trivial \\
permutations of 8, adjacent swaps & permutation & $40\,320$ & $282\,240$ & fine (map or rank) \\
20 lamps, toggle one & bitmask & $1\,048\,576$ & $20\,971\,520$ & fine, $\approx 0.2$ s \\
all words of length 5 & string & $11\,881\,376$ & $\times 125$ & dictionary words only \\
values up to $10^{18}$ & number & $10^{18}$ & --- & impossible \\
\bottomrule
\end{tabularx}
\end{table}

If the estimate is far above $10^8$, the plain traversal is not the solution: look for a smaller state, a bound on the useful values, or mathematics (for the last row: a greedy argument, or Chapter~11's binary lifting when the graph is functional).

%------------------------------------------------------------
\subsubsection*{Recognising the shape from the statement}

The constraints and the wording often reveal \emph{which kind} of graph you have, and with it the algorithm family:

\begin{diagrambox}
\begin{tikzpicture}[
    clue/.style={rectangle, rounded corners=2pt, draw=accentorange, fill=memobg, font=\scriptsize, align=left, text width=5.6cm, minimum height=7mm},
    shp/.style={rectangle, rounded corners=2pt, draw=primaryblue, fill=theorybg, font=\scriptsize\bfseries, align=center, text width=2.4cm, minimum height=7mm},
    algo/.style={font=\scriptsize\color{darkgray}, align=left, text width=3.4cm, anchor=west}]
    \foreach \c/\s/\a [count=\i] in {
        {``$n$ vertices, $n-1$ edges, connected''}/tree/{unique paths; Ch.~13},
        {``$a_i$ is the next / the target of $i$''}/functional graph/{follow arrows; Pattern~5},
        {``prerequisites'', ``must come before''}/DAG/{topological order; Ch.~5},
        {``two teams'', ``enemies in different groups''}/bipartite/{2-colouring; Ch.~4},
        {a board of \texttt{.} and \texttt{\#}, ``islands''}/grid graph/{flood fill, BFS; Pattern~3},
        {``minimum number of moves / operations''}/state graph/{BFS on states; Pattern~4},
        {``how many friends'', ``trusted by everyone''}/degrees only/{counting; Pattern~1},
        {``groups'', ``connected?'', ``can reach''}/reachability/{DFS/BFS, DSU; Pattern~2}}{
        \node[clue] (c\i) at (0,-\i*0.95) {\c};
        \node[shp] (s\i) at (5.0,-\i*0.95) {\s};
        \node[algo] at (6.4,-\i*0.95) {\a};
        \draw[flow] (c\i) -- (s\i);
    }
    \node[font=\scriptsize\bfseries] at (0,-0.25) {clue in the statement};
    \node[font=\scriptsize\bfseries] at (5.0,-0.25) {shape};
    \node[font=\scriptsize\bfseries, anchor=west] at (6.4,-0.25) {first idea};
\end{tikzpicture}
\figcap{From wording to shape to algorithm. These are first guesses, not laws --- but they are right surprisingly often.}
\end{diagrambox}

\begin{diagrambox}
\begin{tikzpicture}[x=1cm, y=1cm]
    \begin{scope}
        \node[gv, minimum size=6mm] (a) at (0.9,0.9) {};
        \node[gv, minimum size=6mm] (b) at (0.2,0) {};
        \node[gv, minimum size=6mm] (c) at (1.6,0) {};
        \node[gv, minimum size=6mm] (d) at (1.6,-0.9) {};
        \draw[gtree] (a) -- (b); \draw[gtree] (a) -- (c); \draw[gtree] (c) -- (d);
        \node[font=\scriptsize\bfseries] at (0.9,-1.6) {tree};
    \end{scope}
    \begin{scope}[xshift=2.9cm]
        \node[gvsel, minimum size=6mm] (a) at (0.4,0.7) {};
        \node[gvsel, minimum size=6mm] (b) at (1.4,0.7) {};
        \node[gvsel, minimum size=6mm] (c) at (0.9,-0.1) {};
        \node[gv, minimum size=6mm] (d) at (0.9,-1.0) {};
        \draw[gde] (a) -- (b); \draw[gde] (b) -- (c); \draw[gde] (c) -- (a); \draw[gde] (d) -- (c);
        \node[font=\scriptsize\bfseries] at (0.9,-1.6) {functional};
    \end{scope}
    \begin{scope}[xshift=5.6cm]
        \node[gv, minimum size=6mm] (a) at (0,0.5) {};
        \node[gv, minimum size=6mm] (b) at (0,-0.6) {};
        \node[gv, minimum size=6mm] (c) at (1,0) {};
        \node[gv, minimum size=6mm] (d) at (2,0) {};
        \draw[gde] (a) -- (c); \draw[gde] (b) -- (c); \draw[gde] (c) -- (d); \draw[gde] (a) to[bend left=30] (d);
        \node[font=\scriptsize\bfseries] at (1,-1.6) {DAG};
    \end{scope}
    \begin{scope}[xshift=8.6cm]
        \node[gvA, minimum size=6mm] (a) at (0,0.6) {};
        \node[gvA, minimum size=6mm] (b) at (0,-0.6) {};
        \node[gvB, minimum size=6mm] (c) at (1.2,0.6) {};
        \node[gvB, minimum size=6mm] (d) at (1.2,-0.6) {};
        \draw[ge] (a) -- (c); \draw[ge] (a) -- (d); \draw[ge] (b) -- (c);
        \node[font=\scriptsize\bfseries] at (0.6,-1.6) {bipartite};
    \end{scope}
    \begin{scope}[xshift=11.2cm, yshift=0.5cm]
        \foreach \r in {0,1,2} \foreach \c in {0,1,2} \node[gcell, minimum size=5mm] at (\c*0.5,-\r*0.5) {};
        \node[gwall, minimum size=5mm] at (0.5,-0.5) {};
        \node[font=\scriptsize\bfseries] at (0.5,-2.1) {grid};
    \end{scope}
\end{tikzpicture}
\figcap{The shapes of this chapter's toolbox, in miniature (the purple vertices of the functional graph form its cycle). Learn to see them in a statement before you see them in a drawing.}
\end{diagrambox}

%------------------------------------------------------------
\subsubsection*{Worked example: modelling a statement}

\begin{examplebox}[A kingdom of towns]
\emph{``The kingdom has $n\le 2\cdot10^5$ towns and $n-1$ roads; every town can be reached from the capital (town 1). For each town, how many roads separate it from the capital?''}

\step{Vertex:} a town. \step{Edge:} a road, undirected, unweighted. \step{Shape:} $n-1$ edges and connected $\Rightarrow$ a \emph{tree}. \step{Question:} the number of roads from vertex 1 to every vertex $=$ BFS distances (or DFS depths: in a tree the path is unique, so DFS gives the shortest distance too). \step{Size:} $|V|=2\cdot10^5$, $|E|=2\cdot10^5-1$, so $\bigO{n}$ is instant; adjacency list, not a matrix; BFS or iterative DFS, because a path-shaped tree is $2\cdot10^5$ levels deep.
\end{examplebox}

%------------------------------------------------------------
\subsubsection*{Check your understanding}

For each situation give the vertex, the edge and an estimate of $|V|$.
\begin{enumerate}[leftmargin=*]
    \item A knight on a $100\times100$ board must reach a target square in as few moves as possible.
    \item A robot on an $R\times C$ grid ($R,C\le 50$) collects up to 5 different keys that open doors.
    \item ``$n\le 10^5$ students; student $i$ copies the homework from student $p_i$.''
    \item From $x$ you may go to $x+3$ or $x-5$, but must always stay inside $[0,10^6]$; can you get from $a$ to $b$?
    \item A statement has $n\le 10^5$ cities, and every pair of cities is joined by a road \emph{except} for $k$ listed pairs. Why can you not build the adjacency list?
\end{enumerate}

\begin{solutionbox}[Answers]
\begin{enumerate}[leftmargin=*]
    \item Vertex $=$ square, edge $=$ knight move (up to 8 per square, undirected), $|V|=10^4$; BFS.
    \item Vertex $=$ (row, column, set of keys held), the keys as a 5-bit mask: $|V|=50\cdot50\cdot32=80\,000$. The position alone is not enough: the same cell with different keys has a different future.
    \item Vertex $=$ student, arc $i\to p_i$: every vertex has out-degree 1 --- a functional graph (Pattern~5), $|V|=10^5$.
    \item Vertex $=$ integer in $[0,10^6]$, arcs $x\to x+3$ and $x\to x-5$ when they stay in range: $|V|=10^6+1$, at most $2\cdot10^6$ arcs. A BFS or DFS from $a$ answers it (directed reachability).
    \item The graph is almost complete: about $n^2/2 = 5\cdot10^9$ edges for $n=10^5$, far beyond memory. Work with the $k$ \emph{missing} pairs (the complement) and a cleverer traversal instead.
\end{enumerate}
\end{solutionbox}


%============================================================
\clearpage
\section{Pattern 1: Degrees and Counting}
\label{sec:ch1-p1}
%============================================================
%============================================================
% CH1 · PATTERN 1 — DEGREES AND COUNTING: THEORY
%============================================================
\subsection{Pattern Theory}
\label{ssec:ch1-p1-theory}

\begin{patternbox}{Pattern 1 at a Glance}
Many graph questions are \textbf{local}: the answer for a vertex depends only on the edges that touch it. Then we never walk along the graph --- we read the edge list \textbf{once} and keep a counter per vertex.

\begin{center}
\begin{tikzpicture}[x=1.15cm, y=1cm]
    \node[gv] (a) at (0,0) {0};
    \node[gv] (b) at (1.5,0.8) {1};
    \node[gv] (c) at (1.5,-0.8) {2};
    \node[gv] (d) at (3,0.8) {3};
    \node[gv] (e) at (4.5,0.8) {4};
    \node[gv] (f) at (4.5,-0.8) {5};
    \draw[ge] (a) -- (b); \draw[ge] (a) -- (c); \draw[ge] (b) -- (c);
    \draw[ge] (b) -- (d); \draw[ge] (d) -- (e);
    \node[gdist, left=1mm of a] {2};
    \node[gdist, above=1mm of b] {3};
    \node[gdist, below=1mm of c] {2};
    \node[gdist, above=1mm of d] {2};
    \node[gdist, above=1mm of e] {1};
    \node[gdist, right=1mm of f] {0};
    \node[note, anchor=west, text width=4.4cm, align=left] at (5.5,0)
        {green numbers = degrees, found in one pass over the 5 edges: $2+3+2+2+1+0 = 10 = 2\cdot 5$};
\end{tikzpicture}
\end{center}

\begin{recurrencebox}
\centering
for every edge $(u,v)$:\quad $\deg(u){+}{+},\ \deg(v){+}{+}$ \quad (directed: $\outdeg{u}{+}{+},\ \indeg{v}{+}{+}$)\\[3pt]
then one loop over the vertices answers the question\\[2pt] $\bigO{|V|+|E|}$ time, \ $\bigO{|V|}$ memory
\end{recurrencebox}
Typical answers: ``the vertex with degree $n-1$'', ``in-degree $n-1$ and out-degree $0$'', ``$\deg(u)+\deg(v)$ minus a correction'', ``$\sum_v w(v)\cdot\deg(v)$''. No adjacency list, no DFS, no BFS.
\end{patternbox}

%------------------------------------------------------------
\subsubsection*{From Part A to Pattern 1}

Lesson~\ref{ssec:ch1-vocabulary} defined the \textbf{degree} $\deg(v)$ of a vertex (the number of edges touching it) and, for directed graphs, the \textbf{in-degree} $\indeg{v}$ and \textbf{out-degree} $\outdeg{v}$. Lesson~\ref{ssec:ch1-storing} showed how to store a whole graph. This pattern is the observation that for a surprising number of problems \emph{you do not need to store the graph at all}: an array \texttt{deg[]} of $|V|$ integers is enough. It is the simplest graph algorithm there is, and it is the first thing to try before writing any traversal.

\begin{diagrambox}
\begin{tikzpicture}[x=1cm, y=1cm]
    \node[font=\small\bfseries, text=deepblue] at (0,2.25) {edge list};
    \foreach \u/\v [count=\i] in {0/1, 0/2, 1/2, 1/3, 3/4}{
        \node[cell, minimum width=13mm] (E\i) at (0,2.1-0.5*\i) {(\u, \v)};
    }
    \draw[flow] (1,0.8) -- node[lbl, above]{one pass} (3.2,0.8);
    \node[font=\small\bfseries, text=deepblue] at (6.1,2.25) {\texttt{deg[v]}};
    \foreach \d [count=\i from 0] in {2,3,2,2,1,0}{
        \node[idx] at (3.85+\i*0.9,1.65) {\i};
        \node[compcell] (D\i) at (3.85+\i*0.9,1.05) {\d};
    }
    \node[note, text width=5.6cm] at (6.1,0.05) {every edge adds 1 to \emph{each} of its two ends --- the edge $(1,3)$ is one of the three units in \texttt{deg[1]} and one of the two in \texttt{deg[3]}};
\end{tikzpicture}
\figcap{The whole ``algorithm'': an edge list goes in, a degree array comes out.}
\end{diagrambox}

%------------------------------------------------------------
\subsubsection*{What a degree can tell you}

\begin{table}[H]
\centering\small
\renewcommand{\arraystretch}{1.2}
\begin{tabularx}{\linewidth}{@{}>{\raggedright\arraybackslash}p{3.7cm} >{\raggedright\arraybackslash}p{3.5cm} X@{}}
\toprule
\textbf{Question about vertex $v$} & \textbf{Degree test} & \textbf{Typical wording} \\
\midrule
isolated vertex & $\deg(v) = 0$ & ``a city with no roads'' \\
leaf (end of a path, tree leaf) & $\deg(v) = 1$ & ``dead end'', ``leaf'' \\
connected to everybody & $\deg(v) = n-1$ (simple graph) & ``center of a star'', ``universal vertex'' \\
source of a directed graph & $\indeg{v} = 0$ & ``nobody points to it'', ``no prerequisites'' \\
sink of a directed graph & $\outdeg{v} = 0$ & ``trusts nobody'', ``terminal state'' \\
pointed to by all, points to none & $\indeg{v}=n-1,\ \outdeg{v}=0$ & ``judge'', ``celebrity'' (Problem~1.1.1) \\
every vertex looks the same & all $\deg(v)$ equal ($k$-regular) & ``everyone has exactly $k$ friends'' \\
\bottomrule
\end{tabularx}
\end{table}

For directed graphs we keep \emph{two} counters per vertex. An arrow $u \to v$ leaves $u$ and enters $v$:

\begin{diagrambox}
\begin{tikzpicture}[x=1.25cm, y=1cm]
    \node[gvsrc] (a) at (0,0.9) {0};
    \node[gv] (b) at (1.8,1.6) {1};
    \node[gv] (c) at (1.8,0) {2};
    \node[gvsrc] (d) at (0,-0.7) {3};
    \node[gvdst] (e) at (3.6,0) {4};
    \draw[gde] (a) -- (b); \draw[gde] (a) -- (c); \draw[gde] (b) -- (c);
    \draw[gde] (d) -- (c); \draw[gde] (c) -- (e);
    \foreach \n/\i/\o/\p in {a/0/2/left, b/1/1/above, c/3/1/below, d/0/1/left, e/1/0/below}{
        \node[gdist, \p=1.5mm of \n] {in~\i\ / out~\o};
    }
    \node[note, anchor=west, text width=5.6cm, align=left] at (4.9,0.5)
        {\textcolor{successgreen}{\textbf{green}} vertices 0 and 3: in-degree 0 (\emph{sources})\\
         \textcolor{accentorange}{\textbf{orange}} vertex 4: out-degree 0 (a \emph{sink})\\[3pt]
         $\sum \indeg{v} = \sum \outdeg{v} = |E| = 5$:\\ every arrow has exactly one tail and one head.};
\end{tikzpicture}
\figcap{In-degree and out-degree of every vertex, again from one pass over the edges.}
\end{diagrambox}

%------------------------------------------------------------
\subsubsection*{The one law of degrees: the handshake lemma}

Cut every edge in the middle. Each edge becomes two \emph{half-edges} (``stubs''), one attached to each end. A vertex of degree $d$ holds exactly $d$ stubs, so counting stubs vertex by vertex or edge by edge gives
\[
\sum_{v \in V} \deg(v) \;=\; 2\,|E| .
\]

\begin{diagrambox}
\begin{tikzpicture}[x=1.15cm, y=1cm]
    \node[gv] (a) at (0,0) {0};
    \node[gv] (b) at (1.5,0.8) {1};
    \node[gv] (c) at (1.5,-0.8) {2};
    \node[gv] (d) at (3,0.8) {3};
    \node[gv] (e) at (4.5,0.8) {4};
    \foreach \x/\y in {a/b, a/c, b/c, b/d, d/e}{
        \draw[draw=primaryblue, line width=1.6pt] (\x) -- ($(\x)!0.5!(\y)$);
        \draw[draw=accentorange, line width=1.6pt] ($(\x)!0.5!(\y)$) -- (\y);
        \fill[white] ($(\x)!0.5!(\y)$) circle (1.8pt);
    }
    \node[note, anchor=west, text width=6cm, align=left] at (5.3,0)
        {5 edges $\Rightarrow$ 10 stubs (blue + orange).\\ Vertex 1 holds 3 of them, vertex 4 holds 1, \ldots\\[3pt]
        \textbf{Consequences:} the degree sum is even; the number of odd-degree vertices is even; $n$ vertices of degree $\ge k$ force at least $nk/2$ edges.};
\end{tikzpicture}
\figcap{Counting the same stubs twice proves the handshake lemma.}
\end{diagrambox}

The lemma is a free consistency check (``can these numbers be the degrees of a graph?'') and the reason why the counting loop runs in $\bigO{|V|+|E|}$: it does exactly $2|E|$ increments. For a directed graph the analogue is $\sum_v \indeg{v} = \sum_v \outdeg{v} = |E|$.

%------------------------------------------------------------
\subsubsection*{When to recognise it --- and when not}

Use Pattern 1 when
\begin{itemize}[leftmargin=*]
    \item the answer for a vertex depends only on \textbf{the edges touching it} (``trusts nobody'', ``connected to all'', ``number of roads at a city'');
    \item the answer for a \textbf{pair} of vertices depends only on their two degrees and on whether they are adjacent (Problem~1.1.2);
    \item the quantity is a \textbf{sum over edges} of something that depends on the endpoints: regroup it as a sum over vertices weighted by degree (Exercise~E1.3);
    \item the input \textbf{promises} a shape (``the graph is a star'', ``at most one person is trusted by all'') and you only need to locate the special vertex.
\end{itemize}

\textbf{Do NOT use it} for anything that depends on \emph{paths}: reachability, connectivity, distances, cycles. Degrees are blind to how the edges are wired together:

\begin{diagrambox}
\begin{tikzpicture}[x=1cm, y=1cm]
    \begin{scope}
        \foreach \i in {0,...,5}{ \node[gvA] (h\i) at ({90-60*\i}:1.1) {\i}; }
        \foreach \i [evaluate=\i as \j using {int(mod(\i+1,6))}] in {0,...,5}{ \draw[ge] (h\i) -- (h\j); }
        \node[font=\small\bfseries] at (0,-1.75) {one cycle of length 6: connected};
    \end{scope}
    \begin{scope}[xshift=6.2cm]
        \node[gvB] (t0) at (-1.5,0.7) {0}; \node[gvB] (t1) at (-0.75,-0.5) {1}; \node[gvB] (t2) at (-2.25,-0.5) {2};
        \node[gvB] (s0) at (1.5,0.7) {3}; \node[gvB] (s1) at (2.25,-0.5) {4}; \node[gvB] (s2) at (0.75,-0.5) {5};
        \draw[ge] (t0) -- (t1) -- (t2) -- (t0);
        \draw[ge] (s0) -- (s1) -- (s2) -- (s0);
        \node[font=\small\bfseries] at (0,-1.75) {two triangles: disconnected};
    \end{scope}
    \node[note] at (3.1,1.55) {both graphs: 6 vertices, 6 edges, every degree $= 2$};
\end{tikzpicture}
\figcap{Identical degree sequences, different answers to ``is the graph connected?''. Such questions need a traversal (Pattern~2).}
\end{diagrambox}

%------------------------------------------------------------
\subsubsection*{The algorithm and why it is correct}

\begin{enumerate}[leftmargin=*]
    \item \step{Allocate} \texttt{deg[0..n]} (one extra slot so that 1-indexed labels fit), all zero. Directed: \texttt{in[]} and \texttt{out[]}, or a single \texttt{score[] = in - out}.
    \item \step{Scan} the edge list once; for $(u,v)$ do \texttt{deg[u]++, deg[v]++} (directed: \texttt{out[u]++, in[v]++}).
    \item \step{Answer} with a loop over the vertices (or over the pairs, or over the edges again).
\end{enumerate}

\begin{invariantbox}[Invariant of the counting loop]
After the first $i$ edges $e_1,\dots,e_i$ have been processed, \texttt{deg[v]} equals the number of edges among $e_1,\dots,e_i$ that have $v$ as an endpoint (a self-loop $(v,v)$ counts twice).

\textbf{Proof.} True for $i = 0$ (all zero, no edges). Processing $e_{i+1} = (u,w)$ adds 1 to exactly the counters of its endpoints $u$ and $w$ --- the only vertices whose count of incident edges grows --- and leaves all other counters unchanged. After the last edge, \texttt{deg[v]} $=\deg(v)$. \hfill$\blacksquare$

\textbf{Cost.} $2|E|$ increments plus $\bigO{|V|}$ to initialise and scan: $\bigO{|V|+|E|}$ time, $\bigO{|V|}$ extra memory --- optimal, since we must at least read the input.
\end{invariantbox}

%------------------------------------------------------------
\subsubsection*{Four counting tools}

\textbf{Tool 1 --- one counter per vertex.} The table above. Find the vertices whose counter has a special value.

\textbf{Tool 2 --- a net score.} When the condition is ``in-degree exactly $A$ \emph{and} out-degree exactly $0$'', one array \texttt{score[v] = in[v] - out[v]} often suffices: if in-degrees can never exceed $A$, then $\texttt{score}[v] = A \iff \indeg{v} = A$ and $\outdeg{v} = 0$. This is the whole solution of Problem~1.1.1.

\textbf{Tool 3 --- pairs: inclusion--exclusion.} The number of edges touching $u$ \emph{or} $v$ is
\[
|\,E(u) \cup E(v)\,| \;=\; \deg(u) + \deg(v) - [\,u \text{ and } v \text{ are adjacent}\,],
\]
because the edge $(u,v)$, if present, was counted once at $u$ and once at $v$. Degrees plus an adjacency test answer every such pair question (Problem~1.1.2).

\begin{diagrambox}
\begin{tikzpicture}[x=1.2cm, y=1cm]
    \node[gvsel] (u) at (0,0) {$u$};
    \node[gvsel] (v) at (2.4,0) {$v$};
    \node[gv] (a) at (-1.3,0.8) {}; \node[gv] (b) at (-1.3,-0.8) {};
    \node[gv] (c) at (3.7,0.8) {}; \node[gv] (d) at (3.7,-0.8) {}; \node[gv] (e) at (3.7,0) {};
    \draw[draw=primaryblue, line width=1.3pt] (u) -- (a); \draw[draw=primaryblue, line width=1.3pt] (u) -- (b);
    \draw[draw=accentorange, line width=1.3pt] (v) -- (c); \draw[draw=accentorange, line width=1.3pt] (v) -- (d);
    \draw[draw=accentorange, line width=1.3pt] (v) -- (e);
    \draw[gpath, draw=softpurple] (u) -- node[gw, above=1mm]{counted twice} (v);
    \node[note, anchor=west, text width=4.6cm, align=left] at (4.3,0)
        {$\deg(u) = 3,\ \deg(v) = 4$\\ edges touching $u$ or $v$: $3 + 4 - 1 = 6$\\ (without the purple edge: $2 + 3 - 0 = 5$)};
\end{tikzpicture}
\figcap{Inclusion--exclusion for two vertices: the shared edge is subtracted once.}
\end{diagrambox}

\textbf{Tool 4 --- turn a sum over edges into a sum over vertices.} If every edge contributes $f(u) + f(v)$, then vertex $v$ appears once per incident edge, so
\[
\sum_{(u,v) \in E} \big(f(u) + f(v)\big) \;=\; \sum_{v \in V} f(v)\cdot \deg(v).
\]
In the same spirit the number of \emph{paths with two edges} $u - v - w$ (``cherries'') is $\sum_v \binom{\deg(v)}{2}$: choose the middle vertex, then two of its edges.

\begin{diagrambox}
\begin{tikzpicture}[x=1.15cm, y=1cm]
    \node[gv] (a) at (0,0) {0};
    \node[gvcur] (b) at (1.5,0.8) {1};
    \node[gv] (c) at (1.5,-0.8) {2};
    \node[gv] (d) at (3,0.8) {3};
    \node[gv] (e) at (4.5,0.8) {4};
    \draw[gpath] (a) -- (b); \draw[gpath] (b) -- (c); \draw[gpath] (b) -- (d);
    \draw[ge] (a) -- (c); \draw[ge] (d) -- (e);
    \node[note, anchor=west, text width=6.2cm, align=left] at (5.3,0)
        {the 3 orange edges at vertex 1 give $\binom{3}{2} = 3$ cherries with middle 1: $0{-}1{-}2,\ 0{-}1{-}3,\ 2{-}1{-}3$.\\[2pt]
        All cherries: $\binom22+\binom32+\binom22+\binom22+\binom12 = 1+3+1+1+0 = 6$.};
\end{tikzpicture}
\figcap{Counting small structures by their centre: a cherry is fixed by its middle vertex and two of its edges.}
\end{diagrambox}

%------------------------------------------------------------
\subsubsection*{Variations map}

\begin{diagrambox}
\begin{tikzpicture}[
    var/.style={rectangle, rounded corners=2pt, draw=patternborder, fill=patternbg,
                font=\scriptsize, align=center, text width=3.25cm, minimum height=1.25cm, inner sep=3pt}]
    \node[var, fill=answerfill, draw=accentorange, text width=5cm, font=\scriptsize\bfseries] (root) at (0,0)
        {Degree counting\\one pass: \texttt{deg[u]++, deg[v]++}};
    \node[var] (v1) at (-5.4,-2) {\textbf{special vertex}\\star center, universal vertex, leaf, isolated (E1.1)};
    \node[var] (v2) at (-1.8,-2) {\textbf{in/out-degree}\\judge, sources, sinks (Problem~1.1.1)};
    \node[var] (v3) at (1.8,-2)  {\textbf{pairs}\\$\deg u + \deg v - [uv \in E]$ (Problem~1.1.2)};
    \node[var] (v4) at (5.4,-2)  {\textbf{sum over edges}\\$= \sum_v f(v)\deg(v)$; sort degrees (E1.3)};
    \node[var] (v5) at (-3.6,-3.6) {\textbf{best neighbours}\\top-$k$ neighbour values per vertex (E1.4)};
    \node[var] (v6) at (0,-3.6)    {\textbf{parity \& existence}\\handshake: is it a degree sequence? (E1.5)};
    \node[var] (v7) at (3.6,-3.6)  {\textbf{build from degrees}\\a tree with given degrees (E1.7)};
    \node[draw=patternborder, dashed, rounded corners=4pt, fit=(v1)(v4)(v5)(v7), inner sep=5pt] (grp) {};
    \draw[flow] (root.south) -- (grp.north);
\end{tikzpicture}
\figcap{Every box is solved by degrees alone. Degree tricks on harder problems (triangles, ordering by degree) return in Chapter~22.}
\end{diagrambox}

\textbf{Twists and how to handle them}
\begin{itemize}[leftmargin=*]
    \item \textbf{Labels $1..n$:} allocate $n+1$ slots and ignore slot 0, instead of shifting every label.
    \item \textbf{Self-loops} $(v,v)$ add 2 to $\deg(v)$ (two stubs); \textbf{multi-edges} add 1 each time. If the problem says ``number of \emph{neighbours}'', deduplicate first (a set per vertex, or sort the edge list).
    \item \textbf{Duplicates in a directed list} break the net-score trick: $\indeg{v}$ could exceed $n-1$. Read the constraints (``all pairs are unique'').
    \item \textbf{Large sums:} $\sum_v f(v)\deg(v)$ easily exceeds $2^{31}$ --- use \texttt{long long}.
    \item \textbf{Pairs with $n$ up to $10^5$:} the $\bigO{n^2}$ pair loop is too slow; only vertices of the largest one or two degrees can win (Approach 3 of Problem~1.1.2).
\end{itemize}

%------------------------------------------------------------
\subsubsection*{Code templates}

\begin{lstlisting}[style=cpp]
#include <bits/stdc++.h>
using namespace std;
typedef long long ll;

// Input: n m, then m edges "u v" (1-indexed). Prints a few degree facts.
int main() {
    int n, m;
    if (!(cin >> n >> m)) return 0;
    vector<int> deg(n + 1, 0);                  // undirected view
    vector<int> in(n + 1, 0), out(n + 1, 0);    // directed view (u -> v)
    for (int i = 0; i < m; i++) {
        int u, v;
        cin >> u >> v;
        deg[u]++; deg[v]++;                     // (1) both ends
        out[u]++; in[v]++;                      // (2) tail and head
    }
    // (3) special vertices
    for (int v = 1; v <= n; v++) {
        if (deg[v] == 0) cout << v << " is isolated\n";
        if (deg[v] == 1) cout << v << " is a leaf\n";
        if (deg[v] == n - 1) cout << v << " touches every other vertex\n";
        if (in[v] == 0) cout << v << " is a source\n";
        if (out[v] == 0) cout << v << " is a sink\n";
    }
    // (4) handshake check and cherries (paths with two edges)
    ll sum = 0, cherries = 0;
    for (int v = 1; v <= n; v++) {
        sum += deg[v];
        cherries += (ll)deg[v] * (deg[v] - 1) / 2;
    }
    cout << "degree sum " << sum << " = 2*" << m << ", cherries " << cherries << "\n";
}
\end{lstlisting}

The second template answers \emph{pair} questions. For $n$ up to a few thousand an adjacency matrix gives an $\bigO{1}$ adjacency test; for large $n$ use a hash set of normalised pairs $(\min, \max)$, as here.

\begin{lstlisting}[style=cpp]
#include <bits/stdc++.h>
using namespace std;
typedef long long ll;

struct PairCounter {
    int n;
    vector<int> deg;
    unordered_set<ll> edgeSet;                   // key = min * n + max (0-indexed)
    PairCounter(int n_) : n(n_), deg(n_, 0) {}
    ll key(int u, int v) const { if (u > v) swap(u, v); return (ll)u * n + v; }
    void addEdge(int u, int v) {
        deg[u]++; deg[v]++;
        edgeSet.insert(key(u, v));
    }
    bool adjacent(int u, int v) const { return edgeSet.count(key(u, v)) > 0; }
    // number of edges touching u or v (Tool 3)
    int touching(int u, int v) const { return deg[u] + deg[v] - (adjacent(u, v) ? 1 : 0); }
};

int main() {
    PairCounter g(6);
    int E[5][2] = {{0, 1}, {0, 2}, {1, 2}, {1, 3}, {3, 4}};
    for (auto& e : E) g.addEdge(e[0], e[1]);
    cout << g.touching(0, 1) << " " << g.touching(1, 4) << " " << g.touching(4, 5) << "\n";  // 4 4 1
}
\end{lstlisting}

\begin{insightbox}[Check the templates by hand]
On the running example ($6$ vertices, edges $01, 02, 12, 13, 34$; degrees $2,3,2,2,1,0$): \texttt{touching(0,1)} $= 2+3-1 = 4$ (edges $01,02,12,13$); \texttt{touching(1,4)} $= 3+1-0 = 4$ (edges $01,12,13,34$); \texttt{touching(4,5)} $= 1+0-0 = 1$. Fed with \texttt{6 5} and the same edges shifted to labels $1..6$, the first template reports vertex 6 as isolated, vertex 5 as a leaf, vertices 1 and 6 as sources, vertices 3, 5, 6 as sinks, the degree sum $10 = 2\cdot5$ and $6$ cherries.
\end{insightbox}

%------------------------------------------------------------
\subsubsection*{Common mistakes}

\begin{warningbox}[Mistake 1 --- counting only one endpoint]
\begin{lstlisting}[style=snippet]
for (auto& e : edges) deg[e[0]]++;                  // WRONG for undirected graphs
for (auto& e : edges) { deg[e[0]]++; deg[e[1]]++; } // RIGHT
\end{lstlisting}
Sanity check: the sum of all degrees must be $2|E|$.
\end{warningbox}

\begin{warningbox}[Mistake 2 --- off-by-one in the labels]
Labels $1..n$ with \texttt{vector<int> deg(n)} write past the end at label $n$ (undefined behaviour, often a silent wrong answer). Allocate $n+1$, or subtract 1 when reading --- consistently.
\end{warningbox}

\begin{warningbox}[Mistake 3 --- using degrees for a path question]
``Every vertex has degree $\ge 1$, so the graph is connected'' is false (two separate edges). ``$|E| = n-1$, so it is a tree'' is false too (a triangle plus an isolated vertex has $n=4$, $|E|=3$). Connectivity needs a traversal (Pattern~2) or DSU (Chapter~6).
\end{warningbox}

\begin{warningbox}[Mistake 4 --- double-counting the shared edge]
For a pair $(u,v)$, $\deg(u)+\deg(v)$ counts the edge $(u,v)$ twice. Subtract it --- but only if it exists.
\end{warningbox}

\begin{warningbox}[Mistake 5 --- forgetting the empty and the tiny cases]
With $n = 1$ and no edges, the only vertex has degree $0 = n-1$: it \emph{is} connected to all (zero) others. With no edges at all, every counter is 0 --- make sure the code still returns a valid answer instead of reading \texttt{edges[0]}.
\end{warningbox}

\begin{checklistbox}[``I'm stuck'' checklist for Pattern 1]
\begin{itemize}[leftmargin=*]
    \item[$\square$] Is the question \emph{local}? Can I state the condition using only the edges at one vertex (or at two)?
    \item[$\square$] Directed or undirected? Do I need in-degree, out-degree, both, or their difference?
    \item[$\square$] Is it a sum over edges? Rewrite it as $\sum_v f(v)\deg(v)$.
    \item[$\square$] For pairs: did I subtract the shared edge? Is $\bigO{n^2}$ over pairs affordable, or must I restrict to the top degrees?
    \item[$\square$] Self-loops, multi-edges, duplicates: what do the constraints promise?
    \item[$\square$] Does the answer depend on paths after all? Then go to Pattern~2.
\end{itemize}
\end{checklistbox}

%============================================================
% CH1 · P1 · PROBLEM 1 — FIND THE TOWN JUDGE (LeetCode 997)
%============================================================
\subsection{Problem 1.1.1: Find the Town Judge}
\problemlabel{1.1.1}{ssec:ch1-town-judge}

\problemcard{1.1.1}{Find the Town Judge}{LeetCode 997}{https://leetcode.com/problems/find-the-town-judge/}{\exeasy\ \lc{Easy}}{In-degree and out-degree counting (one net score per vertex)}

\subsubsection*{1. Problem statement}

A town has $n$ people labelled $1..n$. Rumour says one of them is secretly the \textbf{town judge}. If the judge exists, then
\begin{enumerate}[leftmargin=*, label=(\arabic*)]
    \item the judge trusts nobody;
    \item everybody except the judge trusts the judge;
    \item exactly one person satisfies (1) and (2).
\end{enumerate}
You are given the list \texttt{trust}, where \texttt{trust[i]} $= [a_i, b_i]$ means that person $a_i$ trusts person $b_i$ (and no other trust relations exist). Return the label of the judge if the judge exists and can be identified, otherwise $-1$.

\textbf{Constraints:} $1 \le n \le 1000$, \; $0 \le |\texttt{trust}| \le 10^4$, \; all pairs are unique, \; $a_i \ne b_i$, \; $1 \le a_i, b_i \le n$.

\begin{table}[H]
\centering\small
\begin{tabular}{@{}clcl@{}}
\toprule
$n$ & \textbf{trust} & \textbf{Output} & \textbf{Why} \\
\midrule
2 & $[[1,2]]$ & 2 & 1 trusts 2, 2 trusts nobody \\
3 & $[[1,3],[2,3]]$ & 3 & both others trust 3 \\
3 & $[[1,3],[2,3],[3,1]]$ & $-1$ & 3 is trusted by all, but trusts 1 \\
\bottomrule
\end{tabular}
\end{table}

\begin{diagrambox}
\begin{tikzpicture}[x=1cm, y=1cm]
    \begin{scope}
        \node[gv] (a1) at (0,0) {1}; \node[gvdst] (a2) at (1.6,0) {2};
        \draw[gde] (a1) -- (a2);
        \node[font=\scriptsize\bfseries] at (0.8,-1.45) {Example 1: judge 2};
    \end{scope}
    \begin{scope}[xshift=4.2cm]
        \node[gv] (b1) at (0,0.6) {1}; \node[gv] (b2) at (0,-0.6) {2}; \node[gvdst] (b3) at (1.6,0) {3};
        \draw[gde] (b1) -- (b3); \draw[gde] (b2) -- (b3);
        \node[font=\scriptsize\bfseries] at (0.8,-1.45) {Example 2: judge 3};
    \end{scope}
    \begin{scope}[xshift=8.6cm]
        \node[gv] (c1) at (0,0.6) {1}; \node[gv] (c2) at (0,-0.6) {2}; \node[gvdead] (c3) at (1.6,0) {3};
        \draw[gde] (c1) to[bend left=15] (c3); \draw[gde] (c2) -- (c3);
        \draw[gback] (c3) to[bend left=15] (c1);
        \node[font=\scriptsize\bfseries] at (0.8,-1.45) {Example 3: no judge ($-1$)};
    \end{scope}
\end{tikzpicture}
\figcap{The three official examples. An arrow $a \to b$ means ``$a$ trusts $b$''. In Example 3 the dashed arrow $3 \to 1$ disqualifies 3.}
\end{diagrambox}

\subsubsection*{2. The graph model}
\begin{modelbox}[Graph Model --- Find the Town Judge]
\renewcommand{\arraystretch}{1.35}
\begin{tabularx}{\linewidth}{@{}>{\bfseries}l X@{}}
Vertices & the $n$ people, labelled $1..n$ \\
Edges & a \emph{directed} edge $a \to b$ for every pair $[a,b]$ in \texttt{trust}; stored only as the given edge list \\
Weights & none \\
Question & find the vertex $v$ with $\indeg{v} = n-1$ and $\outdeg{v} = 0$ (a sink that every other vertex points to), or report that none exists \\
Size & $|V| \le 1000$, $|E| \le 10^4$: anything up to $\approx 10^7$ operations is fine, $\bigO{|V|+|E|}$ is trivial \\
\end{tabularx}
\end{modelbox}

\begin{diagrambox}
\begin{tikzpicture}[x=1cm, y=1cm]
    \begin{scope}
        \node[font=\small\bfseries, text=deepblue] at (0,1.6) {the rules};
        \node[tag, text width=4.6cm, align=left] at (0,0.2)
            {(1) the judge trusts nobody\\(2) everybody else trusts the judge};
        \draw[flow] (2.6,0.2) -- node[lbl, above]{translate} (4.4,0.2);
    \end{scope}
    \begin{scope}[xshift=7.4cm]
        \node[gvdst] (J) at (0,0) {$J$};
        \foreach \a [count=\i] in {90,150,210,270,330,30}{
            \node[gv, minimum size=5.5mm] (p\i) at (\a:1.45) {};
            \draw[gde] (p\i) -- (J);
        }
        \node[gdist, anchor=west] at (1.9,0.35) {$\indeg{J} = n-1$};
        \node[gdist, anchor=west] at (1.9,-0.25) {$\outdeg{J} = 0$};
    \end{scope}
\end{tikzpicture}
\figcap{The judge is a vertex with every other vertex pointing to it and no arrow leaving it. What the other people trust among themselves does not matter.}
\end{diagrambox}

\subsubsection*{3. Key observations}

\begin{itemize}[leftmargin=*]
    \item \textbf{Both rules are degree conditions.} ``Trusts nobody'' is $\outdeg{v} = 0$; ``trusted by everybody else'' is $\indeg{v} = n-1$, because the pairs are unique and nobody trusts himself, so the $n-1$ arrows into $v$ must come from $n-1$ different people. No path, no traversal --- Pattern 1.
    \item \textbf{Rule (3) is automatic.} Suppose $J_1 \ne J_2$ both satisfy (1) and (2). By (2) for $J_1$, $J_2$ trusts $J_1$ --- but by (1) for $J_2$, $J_2$ trusts nobody. Contradiction. So at most one vertex can pass the degree test, and the first one we find is the answer.
    \item \textbf{In-degree alone is not enough.} In Example 3 vertex 3 has $\indeg{3} = 2 = n-1$ but trusts 1. Both conditions must be checked.
    \item \textbf{$n = 1$, no edges:} the single person trusts nobody and is trusted by all $0$ others --- he \emph{is} the judge; answer 1.
\end{itemize}

\begin{diagrambox}
\begin{tikzpicture}[x=1cm, y=1cm]
    \node[gvdst] (J1) at (0,0) {$J_1$};
    \node[gvdst] (J2) at (3,0) {$J_2$};
    \draw[gde, draw=warningred, line width=1.3pt] (J2) to[bend right=25] node[lbl, above]{required by (2) for $J_1$} (J1);
    \node[note, text width=4.2cm, anchor=west] at (4,0) {but (1) for $J_2$ forbids any arrow leaving $J_2$ \;\Large\contradiction};
\end{tikzpicture}
\figcap{Why there can never be two judges: the arrow $J_2 \to J_1$ would have to exist and not exist at the same time.}
\end{diagrambox}

\subsubsection*{4. Approaches}

\begin{bruteforcebox}
\textbf{Idea.} Try every person $c$ as the candidate; scan the whole trust list to check that $c$ trusts nobody and count how many people trust $c$.
\begin{lstlisting}[style=cpp]
#include <bits/stdc++.h>
using namespace std;

class Solution {
public:
    int findJudge(int n, vector<vector<int>>& trust) {
        for (int c = 1; c <= n; c++) {                 // candidate judge
            bool trustsSomeone = false;
            int trustedBy = 0;
            for (auto& t : trust) {
                if (t[0] == c) trustsSomeone = true;   // rule (1) broken
                if (t[1] == c) trustedBy++;
            }
            if (!trustsSomeone && trustedBy == n - 1) return c;
        }
        return -1;
    }
};
\end{lstlisting}
\begin{complexitybox}
\textbf{Time:} $\bigO{n \cdot |E|}$ --- up to $1000 \times 10^4 = 10^7$ steps: still \ac{} here, but it rescans the same list $n$ times to compute numbers that one scan already contains. \textbf{Space:} $\bigO{1}$.
\end{complexitybox}
\end{bruteforcebox}

\begin{algorithmbox}
\textbf{Idea.} Both counts come from one pass. Even better, combine them into one \textbf{net score} $\texttt{score}[v] = \indeg{v} - \outdeg{v}$: every pair $[a,b]$ does \texttt{score[a]--} and \texttt{score[b]++}. The judge is the vertex with score exactly $n-1$.
\begin{lstlisting}[style=cpp]
#include <bits/stdc++.h>
using namespace std;

class Solution {
public:
    int findJudge(int n, vector<vector<int>>& trust) {
        vector<int> score(n + 1, 0);       // score[v] = in(v) - out(v), labels 1..n
        for (auto& t : trust) {
            score[t[0]]--;                 // t[0] trusts someone: out-degree + 1
            score[t[1]]++;                 // t[1] is trusted:     in-degree  + 1
        }
        for (int v = 1; v <= n; v++)
            if (score[v] == n - 1) return v;
        return -1;
    }
};
\end{lstlisting}
\begin{complexitybox}
\textbf{Time:} $\bigO{n + |E|}$ --- one pass over the pairs, one over the people. \textbf{Space:} $\bigO{n}$ for \texttt{score}. \textbf{Verdict:} \ac{}.
\end{complexitybox}
\end{algorithmbox}

\begin{invariantbox}[Why score $= n-1$ identifies the judge]
After the loop, $\texttt{score}[v] = \indeg{v} - \outdeg{v}$ (each pair changes exactly the two counters of its endpoints, see the invariant of Pattern 1). Since the pairs are unique and $a_i \ne b_i$, at most $n-1$ different people can trust $v$, so $\indeg{v} \le n-1$; and $\outdeg{v} \ge 0$. Hence
\[
\texttt{score}[v] = n-1 \iff \indeg{v} = n-1 \text{ and } \outdeg{v} = 0,
\]
which is exactly rules (1) and (2). By the observation ``at most one judge'', the first such $v$ is the unique answer. \hfill$\blacksquare$
\end{invariantbox}

\begin{alternativebox}
\textbf{Two separate counters} (\texttt{in[]}, \texttt{out[]}, test \texttt{in[v] == n-1 \&\& out[v] == 0}) is the same algorithm, slightly longer but robust when duplicates are allowed (then \texttt{score} could reach $n-1$ by accident).

\textbf{Candidate elimination (the ``celebrity'' algorithm).} Store trust as an $n \times n$ boolean matrix. Walk once over the people, keeping one candidate. Comparing the candidate $c$ with the next person $v$ always eliminates one of them:
\begin{itemize}[leftmargin=*]
    \item if $c$ trusts $v$, then $c$ is not the judge (the judge trusts nobody) $\Rightarrow$ the new candidate is $v$;
    \item otherwise $v$ is not the judge (the judge is trusted by everybody, including $c$) $\Rightarrow$ keep $c$.
\end{itemize}
After the walk only one person can still be the judge; check him explicitly.
\begin{lstlisting}[style=cpp]
#include <bits/stdc++.h>
using namespace std;

class Solution {
public:
    int findJudge(int n, vector<vector<int>>& trust) {
        vector<vector<char>> T(n + 1, vector<char>(n + 1, 0));   // T[a][b] = a trusts b
        for (auto& t : trust) T[t[0]][t[1]] = 1;
        int cand = 1;
        for (int v = 2; v <= n; v++)
            if (T[cand][v]) cand = v;              // cand trusts v: cand eliminated
        for (int v = 1; v <= n; v++) {             // verify the survivor
            if (v == cand) continue;
            if (T[cand][v] || !T[v][cand]) return -1;
        }
        return cand;
    }
};
\end{lstlisting}
This uses $\bigO{n^2}$ memory and time here, so it is worse for this problem. Its point: it asks only $3(n-1)$ questions ``does $a$ trust $b$?'' --- the right tool when trust is hidden behind an expensive query function instead of being listed (variant V3).

\begin{diagrambox}
\begin{tikzpicture}[x=1.55cm, y=1cm]
    \node[gvdead] (p1) at (0,0) {1}; \node[gvdead] (p2) at (1.6,0) {2};
    \node[gvdst] (p3) at (3.2,0) {3}; \node[gvdead] (p4) at (4.8,0) {4};
    \node[lbl, text=darkgray] at (0.8,0.75) {1 vs 2};
    \node[lbl, text=darkgray] at (2.4,0.75) {1 vs 3};
    \node[lbl, text=darkgray] at (4.0,0.75) {3 vs 4};
    \node[note, text width=2.1cm] at (0.8,-0.95) {1 does not trust 2: drop 2};
    \node[note, text width=2.1cm] at (2.4,-0.95) {1 trusts 3: drop 1};
    \node[note, text width=2.1cm] at (4.0,-0.95) {3 does not trust 4: drop 4};
    \node[note, anchor=west, text width=2.6cm, align=left] at (5.3,0) {survivor 3: verified, it is the judge};
\end{tikzpicture}
\figcap{Elimination on the trace example below ($n=4$, pairs $1{\to}3, 1{\to}4, 2{\to}3, 2{\to}4, 4{\to}3$): three comparisons leave one candidate.}
\end{diagrambox}
\end{alternativebox}

\begin{tricksbox}
\begin{itemize}[leftmargin=*]
    \item \textbf{Early exit:} if $|\texttt{trust}| < n-1$, nobody can be trusted by $n-1$ people --- return $-1$ immediately (careful: $n = 1$ with no pairs must still return 1, and indeed $0 < 0$ is false).
    \item \textbf{One array instead of two} halves the memory and makes the final test a single comparison --- valid only because pairs are unique.
    \item The \texttt{vector<vector<int>>\&} input is read-only here; iterate with \texttt{auto\&} to avoid copying each pair.
\end{itemize}
\end{tricksbox}

\subsubsection*{5. Trace}

\begin{tracebox}[{Net scores for $n = 4$, trust $= [1{\to}3,\ 1{\to}4,\ 2{\to}3,\ 2{\to}4,\ 4{\to}3]$}]
\centering
\begin{tikzpicture}[x=1cm, y=1cm]
    % snapshot A: after 2 edges
    \begin{scope}
        \node[gvseen] (a1) at (0,1) {1}; \node[gv] (a2) at (0,-0.4) {2};
        \node[gvcur] (a3) at (1.7,1) {3}; \node[gvcur] (a4) at (1.7,-0.4) {4};
        \draw[gde] (a1) -- (a3); \draw[gde] (a1) -- (a4);
        \node[gdist, left=1mm of a1] {$-2$}; \node[gdist, left=1mm of a2] {$0$};
        \node[gdist, right=1mm of a3] {$1$}; \node[gdist, right=1mm of a4] {$1$};
        \node[font=\scriptsize\bfseries] at (0.85,-1.25) {after $1{\to}3,\ 1{\to}4$};
    \end{scope}
    % snapshot B: after 4 edges
    \begin{scope}[xshift=4.3cm]
        \node[gvseen] (b1) at (0,1) {1}; \node[gvseen] (b2) at (0,-0.4) {2};
        \node[gvcur] (b3) at (1.7,1) {3}; \node[gvcur] (b4) at (1.7,-0.4) {4};
        \draw[gde] (b1) -- (b3); \draw[gde] (b1) -- (b4); \draw[gde] (b2) -- (b3); \draw[gde] (b2) -- (b4);
        \node[gdist, left=1mm of b1] {$-2$}; \node[gdist, left=1mm of b2] {$-2$};
        \node[gdist, right=1mm of b3] {$2$}; \node[gdist, right=1mm of b4] {$2$};
        \node[font=\scriptsize\bfseries] at (0.85,-1.25) {after $2{\to}3,\ 2{\to}4$: a tie};
    \end{scope}
    % snapshot C: final
    \begin{scope}[xshift=8.6cm]
        \node[gvdead] (c1) at (0,1) {1}; \node[gvdead] (c2) at (0,-0.4) {2};
        \node[gvdst] (c3) at (1.7,1) {3}; \node[gvdead] (c4) at (1.7,-0.4) {4};
        \draw[gde] (c1) -- (c3); \draw[gde] (c1) -- (c4); \draw[gde] (c2) -- (c3); \draw[gde] (c2) -- (c4);
        \draw[gpath, -{Stealth[length=2.6mm]}] (c4) -- (c3);
        \node[gdist, left=1mm of c1] {$-2$}; \node[gdist, left=1mm of c2] {$-2$};
        \node[gdist, right=1mm of c3] {$\mathbf{3}$}; \node[gdist, right=1mm of c4] {$1$};
        \node[font=\scriptsize\bfseries] at (0.85,-1.25) {after $4{\to}3$: judge 3};
    \end{scope}
\end{tikzpicture}
\figcap{Green numbers are the scores. The last pair breaks the tie: 4 now trusts somebody (score $2 \to 1$) and 3 reaches $n-1 = 3$.}

\medskip
\small
\begin{tabular}{@{}lcccc@{}}
\toprule
\textbf{pair processed} & \texttt{score[1]} & \texttt{score[2]} & \texttt{score[3]} & \texttt{score[4]} \\
\midrule
(start) & 0 & 0 & 0 & 0 \\
$1 \to 3$ & $-1$ & 0 & 1 & 0 \\
$1 \to 4$ & $-2$ & 0 & 1 & 1 \\
$2 \to 3$ & $-2$ & $-1$ & 2 & 1 \\
$2 \to 4$ & $-2$ & $-2$ & 2 & 2 \\
$4 \to 3$ & $-2$ & $-2$ & \textbf{3} & 1 \\
\bottomrule
\end{tabular}
\end{tracebox}

\begin{tracebox}[{Example 3: $n = 3$, trust $= [1{\to}3,\ 2{\to}3,\ 3{\to}1]$}]
\centering
\begin{tikzpicture}[x=1cm, y=1cm]
    \begin{scope}
        \node[gvseen] (a1) at (0,0.7) {1}; \node[gvseen] (a2) at (0,-0.7) {2}; \node[gvcur] (a3) at (1.8,0) {3};
        \draw[gde] (a1) -- (a3); \draw[gde] (a2) -- (a3);
        \node[gdist, left=1mm of a1] {$-1$}; \node[gdist, left=1mm of a2] {$-1$}; \node[gdist, right=1mm of a3] {$2$};
        \node[font=\scriptsize\bfseries] at (0.9,-1.45) {after two pairs: 3 has score $2 = n-1$};
    \end{scope}
    \draw[flow] (3.3,0) -- node[lbl, above]{$3 \to 1$} (5,0);
    \begin{scope}[xshift=6.6cm]
        \node[gv] (b1) at (0,0.7) {1}; \node[gvdead] (b2) at (0,-0.7) {2}; \node[gvdead] (b3) at (1.8,0) {3};
        \draw[gde] (b1) to[bend left=12] (b3); \draw[gde] (b2) -- (b3);
        \draw[gback] (b3) to[bend left=12] (b1);
        \node[gdist, left=1mm of b1] {$0$}; \node[gdist, left=1mm of b2] {$-1$}; \node[gdist, right=1mm of b3] {$1$};
        \node[font=\scriptsize\bfseries] at (0.9,-1.45) {final: no score equals 2 $\Rightarrow -1$};
    \end{scope}
\end{tikzpicture}
\figcap{One outgoing arrow costs the candidate a point: the net score automatically enforces rule (1).}
\end{tracebox}

\subsubsection*{6. Test suite}

\begin{testbox}
\centering\small
\begin{tabular}{@{}cllc@{}}
\toprule
$n$ & \textbf{trust} & \textbf{Why it is interesting} & \textbf{Expected} \\
\midrule
1 & $[\,]$ & single person, no pairs: he is the judge & 1 \\
2 & $[\,]$ & nobody trusts anybody & $-1$ \\
2 & $[[1,2]]$ & official Example 1 & 2 \\
3 & $[[1,3],[2,3],[3,1]]$ & in-degree is $n-1$, but out-degree is 1 & $-1$ \\
3 & $[[1,2],[2,3]]$ & a chain: 3 is trusted by one person only & $-1$ \\
4 & $[[1,2],[2,3],[3,4],[4,1]]$ & a cycle: every score is 0 & $-1$ \\
4 & $[[1,3],[2,3],[4,3],[1,2]]$ & extra trust between non-judges is fine & 3 \\
1000 & $i \to 1000$ for all $i < 1000$ & largest $n$, judge has the last label & 1000 \\
\bottomrule
\end{tabular}
\end{testbox}

\subsubsection*{7. Common mistakes}

\begin{warningbox}[Mistake 1 --- checking only the in-degree]
\begin{lstlisting}[style=snippet]
if (in[v] == n - 1) return v;                    // WRONG: Example 3 returns 3
if (in[v] == n - 1 && out[v] == 0) return v;     // RIGHT
\end{lstlisting}
\end{warningbox}

\begin{warningbox}[Mistake 2 --- special-casing the empty list]
\texttt{if (trust.empty()) return -1;} fails for $n = 1$, where the answer is 1. The general loop handles it: $\texttt{score}[1] = 0 = n - 1$.
\end{warningbox}

\begin{warningbox}[Mistake 3 --- an array of size $n$]
Labels are $1..n$. \texttt{vector<int> score(n)} makes \texttt{score[n]} an out-of-bounds write. Use \texttt{n + 1} slots.
\end{warningbox}

\begin{warningbox}[Mistake 4 --- trusting the net score when duplicates are possible]
If the same pair could appear twice, a person trusted twice by one friend could reach score $n-1$ without being trusted by everybody. LeetCode promises unique pairs; in other judges deduplicate first or count distinct trusters.
\end{warningbox}

\subsubsection*{8. Takeaways}

\begin{itemize}[leftmargin=*]
    \item Translate every rule into a statement about \emph{one vertex and its edges}: then degrees answer it in $\bigO{|V|+|E|}$.
    \item ``Points to nobody'' $=$ out-degree 0, ``everybody points to it'' $=$ in-degree $n-1$ (in a simple directed graph).
    \item Two counters can often be merged into one net score when one of them is bounded.
    \item Uniqueness claims (``exactly one person'') are often consequences of the rules --- prove them, then stop at the first hit.
\end{itemize}

\begin{variantsbox}
\variant{Find Center of Star Graph}{\varin{https://leetcode.com/problems/find-center-of-star-graph/}}
  {undirected: the edges form a star, find the vertex adjacent to all others.}
  {model: the center is the unique vertex of degree $n-1$; it appears in \emph{every} edge.}
  {algorithm: degree counting works; even simpler, return the common endpoint of the first two edges ($\bigO{1}$).}
\variant{Vertices that reach everything}{\varin{https://leetcode.com/problems/minimum-number-of-vertices-to-reach-all-nodes/}}
  {in a DAG, find the smallest set of vertices from which all vertices are reachable.}
  {model: a vertex with in-degree 0 can be reached from nobody, so it must be chosen; every other vertex is reached from some source.}
  {algorithm: count in-degrees, output all $v$ with $\indeg{v} = 0$ --- the same counting loop as the judge.}
\variant{Celebrity with a query oracle}{\varnone}
  {trust is hidden; you may only call \texttt{knows(a,b)} and must use $\bigO{n}$ calls.}
  {model: the same directed graph, but edges are revealed one query at a time.}
  {algorithm: candidate elimination ($n-1$ queries) plus verification ($2(n-1)$ queries) --- Approach 3.}
\variant{Duplicated trust records}{\varnone}
  {the same pair $[a,b]$ may be listed several times.}
  {model: a directed multigraph; the judge condition is about \emph{distinct} trusters.}
  {algorithm: sort and deduplicate the pairs (or use a set) first, then count; the net score is unsafe on the raw list.}
\end{variantsbox}

%============================================================
% CH1 · P1 · PROBLEM 2 — MAXIMAL NETWORK RANK (LeetCode 1615)
%============================================================
\subsection{Problem 1.1.2: Maximal Network Rank}
\problemlabel{1.1.2}{ssec:ch1-network-rank}

\problemcard{1.1.2}{Maximal Network Rank}{LeetCode 1615}{https://leetcode.com/problems/maximal-network-rank/}{\exmedium\ \lc{Medium}}{Degrees of a pair + inclusion--exclusion on the shared edge}

\subsubsection*{1. Problem statement}

There are $n$ cities and a list of bidirectional roads; \texttt{roads[i]} $= [a_i, b_i]$ connects cities $a_i$ and $b_i$. The \textbf{network rank} of two different cities is the number of roads directly connected to \emph{either} city; a road connected to both cities is counted only once. The \textbf{maximal network rank} is the largest network rank over all pairs of different cities. Return it.

\textbf{Constraints:} $2 \le n \le 100$, \; $0 \le |\texttt{roads}| \le n(n-1)/2$, \; $0 \le a_i, b_i \le n-1$, \; $a_i \ne b_i$, \; each pair of cities has at most one road.

\begin{table}[H]
\centering\small
\begin{tabular}{@{}clcl@{}}
\toprule
$n$ & \textbf{roads} & \textbf{Output} & \textbf{A best pair} \\
\midrule
4 & $[[0,1],[0,3],[1,2],[1,3]]$ & 4 & cities 0 and 1 (road $0{-}1$ counted once) \\
5 & $[[0,1],[0,3],[1,2],[1,3],[2,3],[2,4]]$ & 5 & cities 1 and 2 \\
8 & $[[0,1],[1,2],[2,3],[2,4],[5,6],[5,7]]$ & 5 & cities 2 and 5 (in different parts!) \\
\bottomrule
\end{tabular}
\end{table}

\begin{diagrambox}
\begin{tikzpicture}[x=1cm, y=1cm]
    \begin{scope}
        \node[gvsel] (a0) at (0,0.7) {0}; \node[gvsel] (a1) at (1.6,0.7) {1};
        \node[gv] (a2) at (1.6,-0.7) {2}; \node[gv] (a3) at (0,-0.7) {3};
        \draw[gpath] (a0) -- (a1); \draw[gpath] (a0) -- (a3); \draw[gpath] (a1) -- (a2); \draw[gpath] (a1) -- (a3);
        \node[font=\scriptsize\bfseries] at (0.8,-1.5) {Example 1: rank$(0,1) = 4$};
    \end{scope}
    \begin{scope}[xshift=3.6cm]
        \node[gv] (b0) at (0,0) {0}; \node[gvsel] (b1) at (1.3,0.75) {1}; \node[gv] (b3) at (1.3,-0.75) {3};
        \node[gvsel] (b2) at (2.6,0) {2}; \node[gv] (b4) at (3.9,0) {4};
        \draw[gpath] (b0) -- (b1); \draw[ge] (b0) -- (b3); \draw[gpath] (b1) -- (b3);
        \draw[gpath] (b1) -- (b2); \draw[gpath] (b2) -- (b3); \draw[gpath] (b2) -- (b4);
        \node[font=\scriptsize\bfseries] at (1.95,-1.5) {Example 2: rank$(1,2) = 5$};
    \end{scope}
    \begin{scope}[xshift=9.2cm]
        \node[gv] (c0) at (-1.1,0.75) {0}; \node[gv] (c1) at (0,0.75) {1}; \node[gvsel] (c2) at (1.1,0.75) {2};
        \node[gv] (c3) at (0.45,-0.3) {3}; \node[gv] (c4) at (1.75,-0.3) {4};
        \node[gvsel] (c5) at (3.2,0.75) {5}; \node[gv] (c6) at (2.7,-0.3) {6}; \node[gv] (c7) at (3.7,-0.3) {7};
        \draw[ge] (c0) -- (c1); \draw[gpath] (c1) -- (c2); \draw[gpath] (c2) -- (c3); \draw[gpath] (c2) -- (c4);
        \draw[gpath] (c5) -- (c6); \draw[gpath] (c5) -- (c7);
        \node[font=\scriptsize\bfseries] at (1.25,-1.5) {Example 3: rank$(2,5) = 5$};
    \end{scope}
\end{tikzpicture}
\figcap{The three official examples. Purple: a best pair; orange: the roads it touches. In Example 2 the road $1{-}2$ touches both cities but counts once.}
\end{diagrambox}

\subsubsection*{2. The graph model}
\begin{modelbox}[Graph Model --- Maximal Network Rank]
\renewcommand{\arraystretch}{1.35}
\begin{tabularx}{\linewidth}{@{}>{\bfseries}l X@{}}
Vertices & the $n$ cities, $0..n-1$ \\
Edges & one undirected edge per road; a simple graph (no loops, no parallel roads); need not be connected \\
Weights & none \\
Question & $\displaystyle\max_{u \ne v}\ \big|E(u) \cup E(v)\big| = \max_{u \ne v}\ \deg(u) + \deg(v) - [\,uv \in E\,]$, where $E(x)$ is the set of edges at $x$ \\
Size & $|V| \le 100$, $|E| \le 4950$; $\binom{100}{2} = 4950$ pairs: even $\bigO{|V|^2}$ is tiny \\
\end{tabularx}
\end{modelbox}

\begin{diagrambox}
\begin{tikzpicture}[x=1.1cm, y=1cm]
    \begin{scope}
        \node[gvsel] (u) at (0,0) {$u$}; \node[gvsel] (v) at (2,0) {$v$};
        \node[gv, minimum size=5mm] (x1) at (-1,0.8) {}; \node[gv, minimum size=5mm] (x2) at (-1,-0.8) {};
        \node[gv, minimum size=5mm] (y1) at (3,0.8) {}; \node[gv, minimum size=5mm] (y2) at (3,-0.8) {};
        \draw[draw=primaryblue, line width=1.3pt] (u) -- (x1); \draw[draw=primaryblue, line width=1.3pt] (u) -- (x2);
        \draw[draw=accentorange, line width=1.3pt] (v) -- (y1); \draw[draw=accentorange, line width=1.3pt] (v) -- (y2);
        \node[font=\scriptsize\bfseries, align=center] at (1,-1.45) {not adjacent:\\ rank $= 2 + 2 = 4$};
    \end{scope}
    \begin{scope}[xshift=6.2cm]
        \node[gvsel] (u) at (0,0) {$u$}; \node[gvsel] (v) at (2,0) {$v$};
        \node[gv, minimum size=5mm] (x1) at (-1,0.8) {}; \node[gv, minimum size=5mm] (x2) at (-1,-0.8) {};
        \node[gv, minimum size=5mm] (y1) at (3,0.8) {}; \node[gv, minimum size=5mm] (y2) at (3,-0.8) {};
        \draw[draw=primaryblue, line width=1.3pt] (u) -- (x1); \draw[draw=primaryblue, line width=1.3pt] (u) -- (x2);
        \draw[draw=accentorange, line width=1.3pt] (v) -- (y1); \draw[draw=accentorange, line width=1.3pt] (v) -- (y2);
        \draw[gpath, draw=softpurple] (u) -- (v);
        \node[font=\scriptsize\bfseries, align=center] at (1,-1.45) {adjacent: $\deg = 3, 3$\\ rank $= 3 + 3 - 1 = 5$};
    \end{scope}
\end{tikzpicture}
\figcap{The rank of a pair is a pure degree formula. The only non-local ingredient is one bit: is $uv$ an edge?}
\end{diagrambox}

\subsubsection*{3. Key observations}

\begin{itemize}[leftmargin=*]
    \item \textbf{Inclusion--exclusion.} The roads touching $u$ or $v$ are the $\deg(u)$ roads at $u$ plus the $\deg(v)$ roads at $v$, minus the ones counted twice. Only the road $uv$ itself can touch both (the graph is simple), so $\text{rank}(u,v) = \deg(u) + \deg(v) - [uv \in E]$.
    \item \textbf{Connectivity is irrelevant.} The two cities need not be connected at all (Example 3). Nothing here needs a traversal --- Pattern 1.
    \item \textbf{``Take the two largest degrees'' is not enough.} The two largest degrees might belong to adjacent cities while another pair with the same degree sum is not adjacent. On a 4-cycle every degree is 2; the first two vertices $0,1$ are adjacent (rank 3), but $0,2$ are not (rank 4).
    \item \textbf{But only top degrees matter.} The correction is at most 1, so an optimal pair always uses the largest degree and, if that is unique, the second largest. This gives an $\bigO{|V|+|E|}$ algorithm (Approach 3), needed when $n$ is large.
\end{itemize}

\begin{diagrambox}
\begin{tikzpicture}[x=1cm, y=1cm]
    \begin{scope}
        \node[gvsel] (p0) at (0,1.2) {0}; \node[gvsel] (p1) at (1.4,1.2) {1};
        \node[gv] (p2) at (1.4,-0.2) {2}; \node[gv] (p3) at (0,-0.2) {3};
        \draw[gpath, draw=softpurple] (p0) -- (p1); \draw[ge] (p1) -- (p2); \draw[ge] (p2) -- (p3); \draw[ge] (p3) -- (p0);
        \node[font=\scriptsize\bfseries, align=center] at (0.7,-1.1) {greedy: first two of degree 2\\ $2 + 2 - 1 = 3$ \no};
    \end{scope}
    \begin{scope}[xshift=5cm]
        \node[gvsel] (q0) at (0,1.2) {0}; \node[gv] (q1) at (1.4,1.2) {1};
        \node[gvsel] (q2) at (1.4,-0.2) {2}; \node[gv] (q3) at (0,-0.2) {3};
        \draw[gpath] (q0) -- (q1); \draw[gpath] (q1) -- (q2); \draw[gpath] (q2) -- (q3); \draw[gpath] (q3) -- (q0);
        \node[font=\scriptsize\bfseries, align=center] at (0.7,-1.1) {opposite corners, not adjacent\\ $2 + 2 - 0 = 4$ \yes};
    \end{scope}
    \node[note, anchor=west, text width=3.6cm, align=left] at (8.1,0.5) {Ties among the top degrees must be examined: one non-adjacent pair is enough.};
\end{tikzpicture}
\figcap{The 4-cycle: picking ``the two largest degrees'' blindly gives 3; the answer is 4.}
\end{diagrambox}

\subsubsection*{4. Approaches}

\begin{bruteforcebox}
\textbf{Idea.} Follow the definition literally: for every pair $(u,v)$, scan all roads and count those touching $u$ or $v$.
\begin{lstlisting}[style=cpp]
#include <bits/stdc++.h>
using namespace std;

class Solution {
public:
    int maximalNetworkRank(int n, vector<vector<int>>& roads) {
        int best = 0;
        for (int u = 0; u < n; u++)
            for (int v = u + 1; v < n; v++) {
                int cnt = 0;                               // roads touching u or v
                for (auto& r : roads)
                    if (r[0] == u || r[1] == u || r[0] == v || r[1] == v) cnt++;
                best = max(best, cnt);
            }
        return best;
    }
};
\end{lstlisting}
A road touching both cities satisfies the condition once and is counted once --- correct by definition, no inclusion--exclusion needed.
\begin{complexitybox}
\textbf{Time:} $\bigO{n^2 \cdot |E|}$ --- for $n = 100$, $|E| = 4950$ about $2.5 \cdot 10^7$ checks: \ac{} here, hopeless for $n = 10^5$. \textbf{Space:} $\bigO{1}$.
\end{complexitybox}
\end{bruteforcebox}

\begin{algorithmbox}
\textbf{Idea.} One pass computes all degrees and fills an adjacency matrix; then every pair costs $\bigO{1}$ with the formula $\deg(u) + \deg(v) - [uv \in E]$.
\begin{lstlisting}[style=cpp]
#include <bits/stdc++.h>
using namespace std;

class Solution {
public:
    int maximalNetworkRank(int n, vector<vector<int>>& roads) {
        vector<int> deg(n, 0);
        vector<vector<bool>> adj(n, vector<bool>(n, false));
        for (auto& r : roads) {
            deg[r[0]]++; deg[r[1]]++;                  // Pattern 1 counting
            adj[r[0]][r[1]] = adj[r[1]][r[0]] = true;  // remember the edge (both directions)
        }
        int best = 0;
        for (int u = 0; u < n; u++)
            for (int v = u + 1; v < n; v++)
                best = max(best, deg[u] + deg[v] - (adj[u][v] ? 1 : 0));
        return best;
    }
};
\end{lstlisting}
\begin{complexitybox}
\textbf{Time:} $\bigO{n^2 + |E|}$ --- about $5\,000$ pair evaluations. \textbf{Space:} $\bigO{n^2}$ bits for the matrix ($10^4$ here). \textbf{Verdict:} \ac{}.
\end{complexitybox}
\end{algorithmbox}

\begin{invariantbox}[Why the formula is the network rank]
Let $E(x)$ be the set of roads at city $x$, so $|E(x)| = \deg(x)$. The network rank is $|E(u) \cup E(v)| = |E(u)| + |E(v)| - |E(u) \cap E(v)|$. A road in the intersection has both $u$ and $v$ as endpoints, so it is the road $uv$; since parallel roads are forbidden there is at most one. Hence $|E(u) \cap E(v)| = [uv \in E]$ and the loop evaluates the exact rank of every pair; the maximum over all pairs is the answer. The degrees are correct by the counting invariant of Pattern 1. \hfill$\blacksquare$
\end{invariantbox}

\begin{alternativebox}[{Approach 3: Only the Top Degrees --- $\bigO{|V|+|E|}$}]
Let $d_1$ be the largest degree and $c_1$ the number of cities having it; let $d_2 < d_1$ be the largest smaller degree and $c_2$ its count.
\begin{itemize}[leftmargin=*]
    \item \textbf{$c_1 \ge 2$.} Every pair has rank $\le 2d_1$, reached exactly by a \emph{non-adjacent} pair of top cities. If all $\binom{c_1}{2}$ top pairs are adjacent, two top cities still give $2d_1 - 1$, and any other pair gives at most $d_1 + d_2 \le 2d_1 - 1$. So: count the roads with both ends of degree $d_1$; answer $2d_1$ if this count is $< \binom{c_1}{2}$, else $2d_1 - 1$.
    \item \textbf{$c_1 = 1$} (one top city $x$). Pairs without $x$ give $\le 2d_2 \le d_1 + d_2 - 1$; pairs $(x, w)$ with $\deg w < d_2$ give $\le d_1 + d_2 - 1$. So the answer is $d_1 + d_2$ if $x$ misses at least one city of degree $d_2$, else $d_1 + d_2 - 1$: count the roads joining $x$ to a degree-$d_2$ city and compare with $c_2$.
\end{itemize}
\begin{lstlisting}[style=cpp]
#include <bits/stdc++.h>
using namespace std;

class Solution {
public:
    int maximalNetworkRank(int n, vector<vector<int>>& roads) {
        vector<int> deg(n, 0);
        for (auto& r : roads) { deg[r[0]]++; deg[r[1]]++; }
        int d1 = -1, d2 = -1;                          // top two distinct degrees
        for (int v = 0; v < n; v++) {
            if (deg[v] > d1) { d2 = d1; d1 = deg[v]; }
            else if (deg[v] < d1 && deg[v] > d2) d2 = deg[v];
        }
        long long c1 = 0, c2 = 0;
        for (int v = 0; v < n; v++) {
            if (deg[v] == d1) c1++;
            else if (deg[v] == d2) c2++;
        }
        if (c1 >= 2) {                                 // several top cities
            long long inside = 0;                      // roads between two top cities
            for (auto& r : roads) if (deg[r[0]] == d1 && deg[r[1]] == d1) inside++;
            return inside < c1 * (c1 - 1) / 2 ? 2 * d1 : 2 * d1 - 1;
        }
        long long cross = 0;                           // roads: top city -- d2-city
        for (auto& r : roads)
            if ((deg[r[0]] == d1 && deg[r[1]] == d2) ||
                (deg[r[0]] == d2 && deg[r[1]] == d1)) cross++;
        return cross < c2 ? d1 + d2 : d1 + d2 - 1;
    }
};
\end{lstlisting}
When $c_1 = 1$, $d_2$ always exists because $n \ge 2$. The counting trick replaces ``check every pair'' by ``count edges inside a group and compare with the number of pairs in it''.

\begin{diagrambox}
\begin{tikzpicture}[x=1cm, y=1cm]
    \begin{scope}
        \node[gv] (b0) at (0,0) {0}; \node[gvsel] (b1) at (1.3,0.75) {1}; \node[gvsel] (b3) at (1.3,-0.75) {3};
        \node[gvsel] (b2) at (2.6,0) {2}; \node[gv] (b4) at (3.9,0) {4};
        \draw[ge] (b0) -- (b1); \draw[ge] (b0) -- (b3); \draw[gpath, draw=softpurple] (b1) -- (b3);
        \draw[gpath, draw=softpurple] (b1) -- (b2); \draw[gpath, draw=softpurple] (b2) -- (b3); \draw[ge] (b2) -- (b4);
        \node[font=\scriptsize, align=center] at (1.95,-1.55) {Example 2: $d_1 = 3$, $c_1 = 3$; 3 roads inside, $\binom32 = 3$\\ all top pairs adjacent $\Rightarrow 2\cdot3-1 = \mathbf{5}$};
    \end{scope}
    \begin{scope}[xshift=7.2cm]
        \node[gv] (c0) at (-1.1,0.75) {0}; \node[gvseen] (c1) at (0,0.75) {1}; \node[gvsel] (c2) at (1.1,0.75) {2};
        \node[gv] (c3) at (0.45,-0.3) {3}; \node[gv] (c4) at (1.75,-0.3) {4};
        \node[gvseen] (c5) at (3.2,0.75) {5}; \node[gv] (c6) at (2.7,-0.3) {6}; \node[gv] (c7) at (3.7,-0.3) {7};
        \draw[ge] (c0) -- (c1); \draw[gpath, draw=softpurple] (c1) -- (c2); \draw[ge] (c2) -- (c3); \draw[ge] (c2) -- (c4);
        \draw[ge] (c5) -- (c6); \draw[ge] (c5) -- (c7);
        \node[font=\scriptsize, align=center] at (1.25,-1.55) {Example 3: top city 2 ($d_1=3$); $d_2 = 2$: cities 1, 5\\ 1 cross road $< c_2 = 2 \Rightarrow 3+2 = \mathbf{5}$};
    \end{scope}
\end{tikzpicture}
\figcap{Approach 3 on Examples 2 and 3. Purple vertices: degree $d_1$; yellow: degree $d_2$; purple edges: the roads that are counted.}
\end{diagrambox}
\end{alternativebox}

\begin{tricksbox}
\begin{itemize}[leftmargin=*]
    \item \textbf{Adjacency test without a matrix:} for large $n$, store each road as the key $\min \cdot n + \max$ in a hash set (Pattern 1 template 2), or sort each adjacency list and binary-search.
    \item \textbf{\texttt{vector<vector<bool>>}} packs bits; for speed in tight loops a flat \texttt{vector<char>} of size $n^2$ with index \texttt{u*n+v} is faster.
    \item \textbf{Upper bound pruning:} sort cities by degree; once $\deg(u) + \deg(v) \le \text{best}$ the inner loop can stop. With the correction at most 1, the search ends almost immediately --- that is Approach 3 in disguise.
\end{itemize}
\end{tricksbox}

\subsubsection*{5. Trace}

\begin{tracebox}[{Example 2: $n = 5$, roads $01, 03, 12, 13, 23, 24$}]
\centering
\begin{tikzpicture}[x=1cm, y=1cm]
    \begin{scope}
        \node[gv] (b0) at (0,0) {0}; \node[gv] (b1) at (1.3,0.8) {1}; \node[gv] (b3) at (1.3,-0.8) {3};
        \node[gv] (b2) at (2.6,0) {2}; \node[gv] (b4) at (3.9,0) {4};
        \draw[ge] (b0) -- (b1); \draw[ge] (b0) -- (b3); \draw[ge] (b1) -- (b3);
        \draw[ge] (b1) -- (b2); \draw[ge] (b2) -- (b3); \draw[ge] (b2) -- (b4);
        \node[gdist, left=1mm of b0] {2}; \node[gdist, above=1mm of b1] {3}; \node[gdist, below=1mm of b3] {3};
        \node[gdist, below right=0.3mm of b2] {3}; \node[gdist, above=1mm of b4] {1};
        \node[font=\scriptsize\bfseries] at (1.95,-1.75) {step 1: degrees $2, 3, 3, 3, 1$ (sum $12 = 2\cdot 6$)};
    \end{scope}
    \begin{scope}[xshift=6.6cm, yshift=1.3cm]
        \foreach \i in {0,...,4}{
            \node[idx] at (0.45+\i*0.8,0.55) {\i};
            \node[idx] at (-0.2,-\i*0.6) {\i};
        }
        \foreach \r/\row in {0/{-,4,5,4,3}, 1/{4,-,5,5,4}, 2/{5,5,-,5,3}, 3/{4,5,5,-,4}, 4/{3,4,3,4,-}}{
            \foreach \x [count=\c from 0] in \row {
                \ifnum\c=\r
                    \node[emptycell, minimum width=8mm, minimum height=6mm] at (0.45+\c*0.8,-\r*0.6) {};
                \else
                    \node[cell, minimum width=8mm, minimum height=6mm] at (0.45+\c*0.8,-\r*0.6) {\x};
                \fi
            }
        }
        \foreach \r/\c in {0/2, 1/2, 1/3, 2/3}{
            \node[anscell, minimum width=8mm, minimum height=6mm] at (0.45+\c*0.8,-\r*0.6) {5};
        }
        \node[font=\scriptsize\bfseries] at (2.05,-3.05) {step 2: rank$(u,v)$ for every pair};
    \end{scope}
\end{tikzpicture}
\figcap{Left: degrees from one pass. Right: the rank table (upper triangle highlighted where the maximum 5 is reached: pairs $(0,2), (1,2), (1,3), (2,3)$).}

\medskip
\small
\begin{tabular}{@{}lccccl@{}}
\toprule
\textbf{pair} & $\deg u$ & $\deg v$ & adjacent? & \textbf{rank} & \textbf{best so far} \\
\midrule
$(0,1)$ & 2 & 3 & yes & 4 & 4 \\
$(0,2)$ & 2 & 3 & no  & 5 & 5 \\
$(0,3)$ & 2 & 3 & yes & 4 & 5 \\
$(0,4)$ & 2 & 1 & no  & 3 & 5 \\
$(1,2)$ & 3 & 3 & yes & 5 & 5 \\
$(1,3)$ & 3 & 3 & yes & 5 & 5 \\
$(1,4)$ & 3 & 1 & no  & 4 & 5 \\
$(2,3)$ & 3 & 3 & yes & 5 & 5 \\
$(2,4)$ & 3 & 1 & yes & 3 & 5 \\
$(3,4)$ & 3 & 1 & no  & 4 & 5 \\
\bottomrule
\end{tabular}
\end{tracebox}

\begin{tracebox}[{Example 2: the winning pairs on the map}]
\centering
\begin{tikzpicture}[x=1cm, y=1cm]
    \begin{scope}
        \node[gvsel] (b0) at (0,0) {0}; \node[gv] (b1) at (1.3,0.8) {1}; \node[gv] (b3) at (1.3,-0.8) {3};
        \node[gvsel] (b2) at (2.6,0) {2}; \node[gv] (b4) at (3.9,0) {4};
        \draw[gpath] (b0) -- (b1); \draw[gpath] (b0) -- (b3); \draw[ge] (b1) -- (b3);
        \draw[gpath] (b1) -- (b2); \draw[gpath] (b2) -- (b3); \draw[gpath] (b2) -- (b4);
        \node[font=\scriptsize\bfseries, align=center] at (1.95,-1.6) {$(0,2)$: $2 + 3 - 0 = 5$\\ not adjacent, nothing subtracted};
    \end{scope}
    \begin{scope}[xshift=6.6cm]
        \node[gv] (b0) at (0,0) {0}; \node[gvsel] (b1) at (1.3,0.8) {1}; \node[gv] (b3) at (1.3,-0.8) {3};
        \node[gvsel] (b2) at (2.6,0) {2}; \node[gv] (b4) at (3.9,0) {4};
        \draw[gpath] (b0) -- (b1); \draw[ge] (b0) -- (b3); \draw[gpath] (b1) -- (b3);
        \draw[gpath, draw=softpurple] (b1) -- (b2); \draw[gpath] (b2) -- (b3); \draw[gpath] (b2) -- (b4);
        \node[font=\scriptsize\bfseries, align=center] at (1.95,-1.6) {$(1,2)$: $3 + 3 - 1 = 5$\\ road $1{-}2$ (purple) counted once};
    \end{scope}
\end{tikzpicture}
\figcap{Two different reasons for the same maximum: a smaller degree sum with no overlap, or a larger sum minus the shared road.}
\end{tracebox}

\subsubsection*{6. Test suite}

\begin{testbox}
\centering\small
\begin{tabular}{@{}cllc@{}}
\toprule
$n$ & \textbf{roads} & \textbf{Why it is interesting} & \textbf{Expected} \\
\midrule
2 & $[\,]$ & no roads at all & 0 \\
2 & $[[0,1]]$ & one road touching both cities: counted once & 1 \\
4 & $[[0,1],[1,2],[2,3],[3,0]]$ & 4-cycle: the tie trap of the greedy & 4 \\
4 & all 6 roads ($K_4$) & every top pair adjacent: $2\cdot3-1$ & 5 \\
5 & $[[0,1],[0,2],[0,3],[0,4]]$ & star: center + any leaf, $4 + 1 - 1$ & 4 \\
8 & official Example 3 & best pair in different components & 5 \\
100 & all $4950$ roads ($K_{100}$) & largest input: $99 + 99 - 1$ & 197 \\
\bottomrule
\end{tabular}
\end{testbox}

\subsubsection*{7. Common mistakes}

\begin{warningbox}[Mistake 1 --- forgetting the shared road]
\begin{lstlisting}[style=snippet]
best = max(best, deg[u] + deg[v]);                          // WRONG: Example 1 gives 5
best = max(best, deg[u] + deg[v] - (adj[u][v] ? 1 : 0));    // RIGHT: 4
\end{lstlisting}
\end{warningbox}

\begin{warningbox}[Mistake 2 --- taking the two largest degrees without looking at ties]
Sorting degrees and returning $d_{(1)} + d_{(2)} - [\text{those two cities adjacent}]$ fails on the 4-cycle (3 instead of 4): another pair with the same degree sum may be non-adjacent. Either check all pairs, or count edges inside the top group (Approach 3).
\end{warningbox}

\begin{warningbox}[Mistake 3 --- assuming the graph is connected]
Running a BFS from city 0 and only considering reachable cities misses Example 3. The rank is defined for \emph{every} pair.
\end{warningbox}

\begin{warningbox}[Mistake 4 --- filling the matrix in one direction only]
\texttt{adj[a][b] = true} alone makes \texttt{adj[u][v]} with $u < v$ false whenever the road was given as $[v, u]$. Set both entries (or normalise each pair to $(\min, \max)$).
\end{warningbox}

\subsubsection*{8. Takeaways}

\begin{itemize}[leftmargin=*]
    \item A pair question about ``edges touching $u$ or $v$'' is inclusion--exclusion: $\deg u + \deg v$ minus the shared edges.
    \item Degrees answer the local part; a single adjacency bit fixes the overlap.
    \item When the correction is bounded (here by 1), only the top one or two degree classes can win --- the key to linear time.
    \item Draw the definition on a tiny example (the 4-cycle) before trusting a greedy.
\end{itemize}

\begin{variantsbox}
\variant{Maximum Star Sum of a Graph}{\varin{https://leetcode.com/problems/maximum-star-sum-of-a-graph/}}
  {vertices carry values; choose a center and at most $k$ of its neighbours to maximise the sum.}
  {model: again a purely local question --- each center only looks at its own incident edges.}
  {algorithm: for every vertex keep the positive neighbour values, take the $k$ largest; $\bigO{|V|+|E|\log|E|}$.}
\variant{Maximum Total Importance of Roads}{\varin{https://leetcode.com/problems/maximum-total-importance-of-roads/}}
  {assign the values $1..n$ to the cities to maximise $\sum_{\text{roads}} (\text{val}(a)+\text{val}(b))$.}
  {model: the sum over edges equals $\sum_v \text{val}(v)\deg(v)$ (Tool 4 of Pattern 1).}
  {algorithm: sort degrees; the largest value goes to the largest degree (rearrangement inequality).}
\variant{Network rank for $n, |E| \le 2\cdot10^5$}{\varnone}
  {the same question on a big sparse graph.}
  {model: unchanged; the $n^2$ pair loop and the matrix no longer fit.}
  {algorithm: Approach 3 --- only the top one or two degree classes matter; count edges inside them; $\bigO{|V|+|E|}$.}
\variant{Remove two vertices from a tree}{\varnone}
  {delete two vertices $u \ne v$ (with their edges) from a tree; maximise the number of remaining components.}
  {model: removing $u$ splits the tree into $\deg(u)$ parts; removing $v$ afterwards adds $\deg(v) - 1 - [uv \in E]$.}
  {algorithm: maximise $\deg u + \deg v - 1 - [uv \in E]$ --- exactly this problem's formula shifted by $-1$, solved by Approach 3.}
\end{variantsbox}

%============================================================
% CH1 · PATTERN 1 — EXERCISES, HINTS, ANSWERS
%============================================================
\subsection{Pattern 1: Exercises and Going Further}
\label{ssec:ch1-p1-ex}

\begin{patternbox}{Pattern 1 Summary Card}
\begin{tabularx}{\linewidth}{@{}>{\bfseries}l X@{}}
Recognise & the answer for a vertex (or a pair) depends only on the edges touching it \\
Core loop & for every edge $(u,v)$: \texttt{deg[u]++, deg[v]++} \;(directed: \texttt{out[u]++, in[v]++}) \\
Tools & special values ($0$, $1$, $n-1$); net score $\text{in}-\text{out}$; pairs $\deg u + \deg v - [uv\in E]$; $\sum_{E}(f(u)+f(v)) = \sum_V f(v)\deg(v)$ \\
Law & handshake: $\sum_v \deg(v) = 2|E|$, so the number of odd-degree vertices is even \\
Cost & $\bigO{|V|+|E|}$ time, $\bigO{|V|}$ memory, no adjacency list \\
Top traps & one endpoint only; labels $1..n$ in an array of size $n$; shared edge counted twice; degrees used for connectivity \\
\end{tabularx}
\end{patternbox}

\exlegend
\begin{enumerate}[leftmargin=*, label=\textbf{E1.\arabic*}, itemsep=6pt]
    \item \exeasy\ \pconcept \exprob{Find Center of Star Graph}{LeetCode 1791}{https://leetcode.com/problems/find-center-of-star-graph/}\ \lc{Easy}. An undirected star with $n \ge 3$ vertices (one center joined to all others, $n-1$ edges) is given as its edge list. Return the center.
    \item \exeasy\ \precent \cf{800}\ \exprob{Destroying Bridges}{Codeforces 1944A}{https://codeforces.com/problemset/problem/1944/A}. Islands $1..n$ form a complete graph. An adversary destroys at most $k$ bridges to minimise the number of islands reachable from island 1 (including island 1). Print that minimum, for $t \le 1000$ test cases with $n \le 100$.
    \item \exmedium\ \pconcept \exprob{Maximum Total Importance of Roads}{LeetCode 2285}{https://leetcode.com/problems/maximum-total-importance-of-roads/}\ \lc{Medium}. Assign the values $1..n$ (each once) to the $n$ cities; a road's importance is the sum of its two endpoint values. Maximise the total importance of all roads. $n, |E| \le 5\cdot10^4$.
    \item \exmedium\ \ptricky \exprob{Maximum Star Sum of a Graph}{LeetCode 2497}{https://leetcode.com/problems/maximum-star-sum-of-a-graph/}\ \lc{Medium}. Vertices have values in $[-10^4, 10^4]$. A star is a center plus some of its neighbours; its sum is the sum of the values of its vertices. Return the maximum sum of a star with at most $k$ edges. $n, |E| \le 10^5$.
    \item \exeasy\ \ptricky \exown{Is it a degree sequence?} Which of the following can be the list of degrees of a simple undirected graph? (a) $(3,2,2,2)$, \ (b) $(3,3,1,1)$, \ (c) $(3,3,2,2)$, \ (d) $(4,1,1,1,1)$.
    \item \exmedium\ \pconcept \exown{Counting short paths.} (a) In an undirected simple graph, count the paths with exactly two edges $u - v - w$ ($u \ne w$, a path and its reverse counted once). (b) In a directed graph without opposite arrow pairs, count the directed paths $a \to b \to c$. Both in $\bigO{|V|+|E|}$.
    \item \exmedium\ \pconcept \exown{A tree with prescribed degrees.} Given $n \ge 2$ and $d_1,\dots,d_n \ge 1$ with $\sum d_i = 2n-2$, output the edges of a tree in which vertex $i$ has degree $d_i$ (one always exists).
\end{enumerate}

\subsubsection*{Hints}
\begin{enumerate}[leftmargin=*, label=\textbf{E1.\arabic*}]
    \item The center appears in \emph{every} edge. How many edges do you need to look at?
    \item Island 1 has degree $n-1$. How many bridges must be destroyed to separate it from even one island?
    \item Rewrite the total as a sum over cities. Which city should get the value $n$?
    \item Each center only looks at its own neighbours. Which neighbours can only hurt?
    \item Handshake parity first; then think about what a vertex of degree $n-1$ forces.
    \item Fix the middle vertex $v$ / $b$. How many ways are there to choose the two edges at it?
    \item A vertex of degree 1 is a leaf. Attach leaves to the vertices that still need edges; when such a vertex is satisfied except for one edge, it becomes a leaf itself.
\end{enumerate}

\subsubsection*{Answers}

\begin{solutionbox}[E1.1 Find Center of Star Graph]
The degree solution (count, return the vertex of degree $n-1$) is $\bigO{n}$. But every edge contains the center, so the center is the common endpoint of the first two edges --- $\bigO{1}$:
\begin{lstlisting}[style=cpp]
#include <bits/stdc++.h>
using namespace std;

class Solution {
public:
    int findCenter(vector<vector<int>>& edges) {
        int a = edges[0][0], b = edges[0][1];
        return (a == edges[1][0] || a == edges[1][1]) ? a : b;   // the shared endpoint
    }
};
\end{lstlisting}
Example 1, edges $[1,2],[2,3],[4,2]$: the first two share vertex $2$ $\Rightarrow$ 2. Example 2, $[1,2],[5,1],\ldots$ $\Rightarrow$ 1. This only works because the input is promised to be a star: with a general graph, use the degrees.
\end{solutionbox}

\begin{solutionbox}[E1.2 Destroying Bridges]
\textbf{Answer:} $1$ if $k \ge n-1$, otherwise $n$.

If $k \ge n-1$, destroy all $n-1$ bridges at island 1 (its whole degree): island 1 is alone. If $k < n-1$, every island $v \ne 1$ is still reachable: the direct bridge $1{-}v$ and the $n-2$ two-step routes $1{-}w{-}v$ form $n-1$ routes that share no bridge, and $k < n-1$ destructions cannot hit all of them. So there is nothing in between: either the full degree of island 1 is destroyed, or nothing changes.

\begin{center}
\begin{tikzpicture}[x=1cm, y=1cm]
    \node[gvsrc] (o) at (0,0) {1};
    \foreach \a/\l in {90/2, 18/3, -54/4, -126/5, 162/6}{ \node[gv] (n\l) at (\a:1.4) {\l}; }
    \foreach \l in {2,...,6}{ \draw[gcut] (o) -- (n\l); }
    \foreach \x/\y in {2/3, 3/4, 4/5, 5/6, 6/2, 2/4, 3/5, 4/6, 5/2, 6/3}{ \draw[gfaded] (n\x) -- (n\y); }
    \node[note, anchor=west, text width=5.2cm, align=left] at (2.2,0) {$n = 6$: the $n-1 = 5$ red dashed bridges are exactly $\deg(1)$. With $k = 5$ the answer is 1; with $k = 4$ some bridge at island 1 survives and every island stays reachable: 6.};
\end{tikzpicture}
\end{center}

\begin{lstlisting}[style=cpp]
#include <bits/stdc++.h>
using namespace std;

int main() {
    int t;
    cin >> t;
    while (t--) {
        long long n, k;
        cin >> n >> k;
        cout << (k >= n - 1 ? 1 : n) << "\n";
    }
}
\end{lstlisting}
Official sample (pairs $(n,k)$): $(2,0),(2,1),(4,1),(5,10),(5,3),(4,4) \to 2, 1, 4, 1, 5, 1$. Verified also by trying every set of at most $k$ destroyed bridges for all $n \le 5$.
\end{solutionbox}

\begin{solutionbox}[E1.3 Maximum Total Importance of Roads]
Each road adds the values of its two endpoints, so city $v$ contributes $\text{val}(v)$ once per road: the total is $\sum_v \text{val}(v)\cdot\deg(v)$ (Tool 4). By the rearrangement inequality this is maximal when the values sorted increasingly meet the degrees sorted increasingly: the city with the $i$-th smallest degree gets value $i$. (Exchange argument: if $\deg(a) < \deg(b)$ but $\text{val}(a) > \text{val}(b)$, swapping the two values changes the total by $(\text{val}(a)-\text{val}(b))(\deg(b)-\deg(a)) > 0$.)
\begin{lstlisting}[style=cpp]
#include <bits/stdc++.h>
using namespace std;

class Solution {
public:
    long long maximumImportance(int n, vector<vector<int>>& roads) {
        vector<long long> deg(n, 0);
        for (auto& r : roads) { deg[r[0]]++; deg[r[1]]++; }
        sort(deg.begin(), deg.end());                   // smallest degree first
        long long total = 0;
        for (int i = 0; i < n; i++) total += deg[i] * (i + 1);   // value i+1
        return total;
    }
};
\end{lstlisting}
Example 1: degrees $2,3,4,2,1$; sorted $1,2,2,3,4$; total $1\cdot1 + 2\cdot2 + 2\cdot3 + 3\cdot4 + 4\cdot5 = 43$. Example 2: $\to 20$. The answer reaches about $5\cdot10^4 \cdot 2 \cdot 5\cdot10^4 = 5\cdot10^9$: \texttt{long long} is required.
\end{solutionbox}

\begin{solutionbox}[E1.4 Maximum Star Sum of a Graph]
A star is local to its center, so treat every vertex as the center. Adding a neighbour with value $\le 0$ never helps, so collect only \emph{positive} neighbour values, sort them decreasingly and take at most $k$. The star with zero edges (the center alone) is always allowed --- that is why the answer of Example 2 is $-5$.
\begin{lstlisting}[style=cpp]
#include <bits/stdc++.h>
using namespace std;

class Solution {
public:
    int maxStarSum(vector<int>& vals, vector<vector<int>>& edges, int k) {
        int n = vals.size();
        vector<vector<int>> good(n);                 // positive neighbour values per vertex
        for (auto& e : edges) {
            if (vals[e[1]] > 0) good[e[0]].push_back(vals[e[1]]);
            if (vals[e[0]] > 0) good[e[1]].push_back(vals[e[0]]);
        }
        int best = INT_MIN;
        for (int c = 0; c < n; c++) {
            sort(good[c].rbegin(), good[c].rend());  // largest first
            int sum = vals[c];
            for (int i = 0; i < (int)good[c].size() && i < k; i++) sum += good[c][i];
            best = max(best, sum);
        }
        return best;
    }
};
\end{lstlisting}
Example 1: center 3 (value 4) has neighbours with values $2, 10, -10, -20$; the two best positive ones give $4 + 10 + 2 = 16$. Total work $\bigO{|V| + |E|\log|E|}$; \texttt{nth\_element} instead of sorting gives $\bigO{|V|+|E|}$. Checked against trying every subset of neighbours on random small graphs.
\end{solutionbox}

\begin{solutionbox}[E1.5 Is it a degree sequence?]
\begin{itemize}[leftmargin=*]
    \item[(a)] $(3,2,2,2)$: the sum $9$ is odd --- impossible by the handshake lemma.
    \item[(b)] $(3,3,1,1)$: the sum $8$ is even, yet impossible. With 4 vertices, degree 3 means adjacent to \emph{all} others; two such vertices give each remaining vertex degree $\ge 2$.
    \item[(c)] $(3,3,2,2)$: yes --- $K_4$ minus one edge.
    \item[(d)] $(4,1,1,1,1)$: yes --- a star with 4 leaves.
\end{itemize}
\begin{center}
\begin{tikzpicture}[x=1cm, y=1cm]
    \begin{scope}
        \node[gv] (a) at (0,0) {3}; \node[gv] (b) at (1.5,0) {3}; \node[gv] (c) at (1.5,1.3) {2}; \node[gv] (d) at (0,1.3) {2};
        \draw[ge] (a) -- (b); \draw[ge] (a) -- (c); \draw[ge] (a) -- (d); \draw[ge] (b) -- (c); \draw[ge] (b) -- (d);
        \draw[gfaded, dashed] (c) -- (d);
        \node[font=\scriptsize\bfseries] at (0.75,-0.75) {(c) $K_4$ minus edge};
    \end{scope}
    \begin{scope}[xshift=4.5cm, yshift=0.65cm]
        \node[gv] (o) at (0,0) {4};
        \foreach \a in {45,135,225,315}{ \node[gv] (l\a) at (\a:1.1) {1}; \draw[ge] (o) -- (l\a); }
        \node[font=\scriptsize\bfseries] at (0,-1.4) {(d) star};
    \end{scope}
\end{tikzpicture}
\end{center}
(Labels are degrees.) An even sum is necessary but not sufficient; the exact test for simple graphs is the Erd\H{o}s--Gallai theorem, which we do not need here. All four answers were confirmed by enumerating every simple graph on 4 and 5 vertices.
\end{solutionbox}

\begin{solutionbox}[E1.6 Counting short paths]
(a) Fix the middle vertex $v$; the path is determined by an unordered pair of distinct edges at $v$: $\sum_v \binom{\deg(v)}{2}$. On the running example of the theory (degrees $2,3,2,2,1,0$) this is $1+3+1+1+0+0 = 6$.

(b) Fix the middle vertex $b$; choose an arrow into $b$ and an arrow out of $b$: $\sum_b \indeg{b}\cdot\outdeg{b}$. The two ends differ because no opposite pair $a \to b \to a$ exists. On the directed example of the theory (in/out: $0/2,\ 1/1,\ 3/1,\ 0/1,\ 1/0$): $0 + 1 + 3 + 0 + 0 = 4$, namely $0{\to}1{\to}2$, $0{\to}2{\to}4$, $1{\to}2{\to}4$, $3{\to}2{\to}4$.

Counting \emph{triangles} is a different story: the three edges are not all at one vertex, so degrees alone cannot do it (Chapter~22).
\end{solutionbox}

\begin{solutionbox}[E1.7 A tree with prescribed degrees]
Split the vertices into leaves ($d = 1$) and inner vertices ($d \ge 2$). For every inner vertex $v$: attach $d_v - 1$ leaves to it (each attached leaf is finished and never used again), after which $v$ needs exactly one more edge --- it has become a leaf itself. At the end exactly two leaves remain; join them.
\begin{lstlisting}[style=cpp]
#include <bits/stdc++.h>
using namespace std;

vector<pair<int,int>> treeFromDegrees(vector<int> d) {    // vertices 0..n-1
    int n = d.size();
    vector<pair<int,int>> edges;
    vector<int> leaves, inner;
    for (int v = 0; v < n; v++) (d[v] == 1 ? leaves : inner).push_back(v);
    for (int v : inner) {
        while (d[v] > 1) {                                  // attach a finished leaf to v
            int l = leaves.back(); leaves.pop_back();
            edges.push_back({l, v});
            d[v]--;
        }
        leaves.push_back(v);                                // v needs exactly one more edge
    }
    edges.push_back({leaves[0], leaves[1]});                // the last two leaves
    return edges;
}

int main() {
    vector<int> d = {1, 3, 1, 2, 1, 1, 3};                  // n = 7, sum = 12 = 2*7 - 2
    // prints (5,1) (4,1) (1,3) (3,6) (2,6) (0,6)
    for (auto& e : treeFromDegrees(d)) cout << "(" << e.first << "," << e.second << ") ";
    cout << "\n";
}
\end{lstlisting}
\begin{center}
\begin{tikzpicture}[x=1cm, y=1cm]
    \node[gv] (v1) at (0,0) {1}; \node[gv] (v3) at (1.6,0) {3}; \node[gv] (v6) at (3.2,0) {6};
    \node[gv] (v5) at (-1.1,0.8) {5}; \node[gv] (v4) at (-1.1,-0.8) {4};
    \node[gv] (v2) at (4.3,0.8) {2}; \node[gv] (v0) at (4.3,-0.8) {0};
    \draw[ge] (v5) -- (v1); \draw[ge] (v4) -- (v1); \draw[ge] (v1) -- (v3); \draw[ge] (v3) -- (v6);
    \draw[ge] (v2) -- (v6); \draw[gpath] (v0) -- (v6);
    \node[gdist, below=1mm of v1] {3}; \node[gdist, below=1mm of v3] {2}; \node[gdist, below=1mm of v6] {3};
    \node[note, anchor=west, text width=4.5cm, align=left] at (5.2,0) {degrees $1,3,1,2,1,1,3$ for vertices $0..6$; the orange edge joins the last two leaves 0 and 6.};
\end{tikzpicture}
\end{center}
\textbf{Why it never runs out of leaves.} $\sum_v (d_v - 1) = n - 2$, so the inner vertices need $\sum_{\text{inner}} (d_v - 1) = n - 2$ attachments in total, while the supply is the $L = n - |\text{inner}|$ original leaves plus one new leaf per finished inner vertex. Before inner vertex number $j$ the supply exceeds the demand because every later inner vertex still needs at least one attachment; at the end $L + |\text{inner}| - (n - 2) = 2$ leaves remain. Every edge removes a finished vertex for good, so no cycle can close, and $n - 1$ acyclic edges on $n$ vertices form a tree. Verified on every valid degree sequence with $n \le 7$.
\end{solutionbox}


%============================================================
\clearpage
\section{Pattern 2: Reachability}
\label{sec:ch1-p2}
%============================================================
%============================================================
% CH1 · PATTERN 2 — REACHABILITY: THEORY
%============================================================
\subsection{Pattern Theory}
\label{ssec:ch1-p2-theory}

% ---- one reusable drawing of the running example G (8 vertices, 10 directed edges) ----
% every vertex style g1p2vK and every edge style g1p2eXY can be overridden per picture
\tikzset{
    g1p2v0/.style={gv}, g1p2v1/.style={gv}, g1p2v2/.style={gv}, g1p2v3/.style={gv},
    g1p2v4/.style={gv}, g1p2v5/.style={gv}, g1p2v6/.style={gv}, g1p2v7/.style={gv},
    g1p2e/.style={gde},
    g1p2e01/.style={g1p2e}, g1p2e02/.style={g1p2e}, g1p2e13/.style={g1p2e}, g1p2e23/.style={g1p2e},
    g1p2e34/.style={g1p2e}, g1p2e41/.style={g1p2e}, g1p2e50/.style={g1p2e}, g1p2e56/.style={g1p2e},
    g1p2e67/.style={g1p2e}, g1p2e76/.style={g1p2e},
    g1p2tree/.style={gtree, -{Stealth[length=2.6mm]}},
    pics/g1p2G/.style={code={
        \node[g1p2v0] (-0) at (0,0)        {0};
        \node[g1p2v1] (-1) at (1.5,0.85)   {1};
        \node[g1p2v2] (-2) at (1.5,-0.85)  {2};
        \node[g1p2v3] (-3) at (3,0)        {3};
        \node[g1p2v4] (-4) at (3,1.7)      {4};
        \node[g1p2v5] (-5) at (-1.5,0)     {5};
        \node[g1p2v6] (-6) at (-1.5,-1.4)  {6};
        \node[g1p2v7] (-7) at (-3,-1.4)    {7};
        \draw[g1p2e01] (-0) -- (-1);
        \draw[g1p2e02] (-0) -- (-2);
        \draw[g1p2e13] (-1) -- (-3);
        \draw[g1p2e23] (-2) -- (-3);
        \draw[g1p2e34] (-3) -- (-4);
        \draw[g1p2e41] (-4) -- (-1);
        \draw[g1p2e50] (-5) -- (-0);
        \draw[g1p2e56] (-5) -- (-6);
        \draw[g1p2e67] ([yshift=1.5mm]-6.west) -- ([yshift=1.5mm]-7.east);
        \draw[g1p2e76] ([yshift=-1.5mm]-7.east) -- ([yshift=-1.5mm]-6.west);
    }}
}

\begin{patternbox}{Pattern 2 at a Glance}
\textbf{Question:} \emph{starting from $s$ and following edges, which vertices can I get to?} The answer is the \textbf{reach set} $R(s)$. One search --- DFS or BFS, it does not matter which --- marks exactly the vertices of $R(s)$, touching every vertex and every edge at most once.

\begin{center}
\begin{tikzpicture}[g1p2v0/.style={gvsrc}, g1p2v1/.style={gvdone}, g1p2v2/.style={gvdone},
    g1p2v3/.style={gvdone}, g1p2v4/.style={gvdone}, g1p2v5/.style={gvdead}, g1p2v6/.style={gvdead},
    g1p2v7/.style={gvdead}, g1p2e01/.style={g1p2tree}, g1p2e02/.style={g1p2tree},
    g1p2e23/.style={g1p2tree}, g1p2e34/.style={g1p2tree},
    g1p2e56/.style={gfaded, -{Stealth[length=2.4mm]}}, g1p2e67/.style={gfaded, -{Stealth[length=2.4mm]}},
    g1p2e76/.style={gfaded, -{Stealth[length=2.4mm]}}]
    \pic (G) {g1p2G};
    \node[draw=successgreen, dashed, rounded corners=6pt, line width=1pt, fit=(G-0)(G-1)(G-2)(G-3)(G-4),
          inner sep=6pt, label={[font=\small\bfseries, text=successgreen]right:$R(0)=\{0,1,2,3,4\}$}] {};
    \node[note, anchor=north] at (-2.25,-1.9) {5, 6, 7 are not reachable\\(5 reaches 0, not vice versa)};
\end{tikzpicture}
\end{center}

\begin{lstlisting}[style=snippet]
seen[s] = true;  todo = {s};               // todo = a stack (DFS) or a queue (BFS)
while (todo not empty) {
    u = take one vertex out of todo;
    for (v in adj[u])
        if (!seen[v]) { seen[v] = true; add v to todo; }   // mark when ADDED
}
// now seen[v] == true  <=>  v is reachable from s.      Time O(|V| + |E|).
\end{lstlisting}
The thick blue edges discovered a new vertex: they form a tree rooted at $s$, the \emph{search tree}.
\end{patternbox}

%------------------------------------------------------------
\subsubsection*{From Pattern 1 to Pattern 2: from local to global}

Pattern~1 answered questions by looking at \textbf{one vertex and the edges touching it}: a degree, an in-degree, a pair of degrees. Such questions are \emph{local} --- one pass over the edge list suffices, and no vertex ever needs to know about a vertex two steps away.

Reachability is the first \emph{global} question. Whether $t$ can be reached from $s$ depends on whole chains of edges, possibly very long ones, and a single missing link anywhere in a chain changes the answer. Degree counting cannot see this. The two graphs below have \emph{identical} degrees --- every vertex has degree~2 --- yet in the first one $0$ reaches $4$ and in the second one it does not.

\begin{diagrambox}
\begin{tikzpicture}
    \begin{scope}
        \foreach \i in {0,...,5} \node[gv] (a\i) at ({90-60*\i}:1.25) {\i};
        \node[gvsrc] at (a0) {0}; \node[gvdst] at (a4) {4};
        \foreach \i/\j in {0/1,1/2,2/3,3/4,4/5,5/0} \draw[ge] (a\i) -- (a\j);
        \draw[gpath] (a0) -- (a5) -- (a4);
        \node[gv] at (a5) {5}; \node[gvsrc] at (a0) {0}; \node[gvdst] at (a4) {4};
        \node[note] at (0,-2) {one hexagon: $0 \leadsto 4$ \yes};
    \end{scope}
    \begin{scope}[xshift=5.2cm]
        \foreach \i/\x/\y in {0/0/1.1, 1/0.95/-0.45, 2/-0.95/-0.45} \node[gv] (b\i) at (\x,\y) {\i};
        \foreach \i/\x/\y in {3/3.6/1.1, 4/4.55/-0.45, 5/2.65/-0.45} \node[gv] (b\i) at (\x,\y) {\i};
        \node[gvsrc] at (b0) {0}; \node[gvdst] at (b4) {4};
        \foreach \i/\j in {0/1,1/2,2/0,3/4,4/5,5/3} \draw[ge] (b\i) -- (b\j);
        \draw[gcut] (1.8,1.5) -- (1.8,-1.0);
        \node[note] at (1.8,-1.55) {two triangles: $0 \not\leadsto 4$ \no};
    \end{scope}
\end{tikzpicture}
\figcap{Same number of vertices, same number of edges, same degree at every vertex --- different answers. A \emph{local} statistic cannot tell whether a chain is broken somewhere; a search can.}
\end{diagrambox}

So we need a procedure that \emph{walks}. Lesson~\ref{ssec:ch1-first-traversal} introduced the two classical walkers, DFS and BFS. In this pattern we use them for their most basic job: \textbf{deciding who can reach whom}. Their deeper properties (DFS edge types and timestamps, BFS distances) wait for Chapter~2 (DFS) and Chapter~3 (BFS).

%------------------------------------------------------------
\subsubsection*{Direction matters}

In an \textbf{undirected} graph reachability is symmetric: if $s$ reaches $t$, then $t$ reaches $s$ along the same path walked backwards. The vertices then split into groups that all reach each other --- the \emph{connected components} (Chapter~4). In a \textbf{directed} graph this symmetry is gone:

\begin{diagrambox}
\begin{tikzpicture}[x=1.4cm]
    \begin{scope}
        \node[gvsrc] (u) at (0,0) {u}; \node[gv] (v) at (1,0) {v}; \node[gv] (w) at (2,0) {w};
        \draw[ge] (u) -- (v); \draw[ge] (v) -- (w);
        \node[note] at (1,-0.8) {undirected: $u - v - w$,\\everybody reaches everybody};
    \end{scope}
    \begin{scope}[xshift=4.2cm]
        \node[gvsrc] (u) at (0,0) {u}; \node[gv] (v) at (1,0) {v}; \node[gv] (w) at (2,0) {w};
        \draw[gde] (u) -- (v); \draw[gde] (v) -- (w);
        \node[note] at (1,-0.8) {directed: $u$ reaches $w$,\\but $w$ reaches nobody else};
    \end{scope}
    \begin{scope}[xshift=8.4cm]
        \node[gv] (u) at (0,0) {u}; \node[gv] (v) at (1,0) {v}; \node[gv] (w) at (2,0) {w};
        \draw[gde] (u) -- (v); \draw[gde] (v) -- (w);
        \draw[gde] (w) to[bend right=50] (u);
        \node[note] at (1,-0.8) {directed cycle: symmetric again\\(strongly connected, Chapter 19)};
    \end{scope}
\end{tikzpicture}
\figcap{The same three vertices, three different reachability relations.}
\end{diagrambox}

Two consequences you will use constantly:
\begin{itemize}[leftmargin=*]
    \item An undirected edge $\{u,v\}$ is stored as \textbf{two} directed edges $u\to v$ and $v \to u$; then the same search code serves both kinds of graphs.
    \item In a directed graph, ``\emph{can $s$ reach $t$?}'' and ``\emph{can $t$ reach $s$?}'' are different questions; the second one is answered by searching from $t$ in the \textbf{reverse graph} (every edge flipped) --- see below.
\end{itemize}

%------------------------------------------------------------
\subsubsection*{When to recognise it}

\textbf{Use this pattern when the statement asks}
\begin{itemize}[leftmargin=*]
    \item ``is there a path / a way / a sequence of moves from $A$ to $B$?'' (yes / no);
    \item ``can we visit / open / infect / activate \emph{all} of them?'' --- is the reach set everything?;
    \item ``how many things are affected if this one starts?'' --- the size of a reach set (chain reactions, spreading rumours, detonations);
    \item ``starting from \emph{any} of these, what is reached?'' --- one multi-source search;
    \item ``who can reach $t$?'' --- one search in the reverse graph;
    \item ``for every pair \ldots'' with $|V|$ at most a few thousand --- one search per start vertex.
\end{itemize}

\textbf{Do NOT use a plain reachability search when}
\begin{itemize}[leftmargin=*]
    \item the question is \emph{how far} (fewest moves) --- that is BFS distances (Pattern~4, Chapter~3) or Dijkstra with weights (Chapter~7);
    \item the question only involves the edges of one vertex --- degrees suffice (Pattern~1) and are much simpler;
    \item there are many queries ``are $a$ and $b$ connected?'' in an undirected graph --- label the components once (Chapter~4) or use DSU, especially if edges keep arriving (Chapter~6);
    \item there are many queries ``does $a$ reach $b$?'' in a \emph{large} directed graph --- no simple trick exists in general; on DAGs use bitsets (Chapter~5), in general first compress the cycles (Chapter~19);
    \item every vertex has exactly one outgoing edge --- you are not searching, you are \emph{following arrows} (Pattern~5).
\end{itemize}

%------------------------------------------------------------
\subsubsection*{The algorithm and why it works}

Write $s \leadsto t$ for ``there is a walk from $s$ to $t$'' (a walk of zero edges is allowed, so $s \leadsto s$ always). The reach set is $R(s) = \{ t : s \leadsto t\}$.

The generic search keeps several kinds of vertices, shown in every trace of this book with the same colours:
\begin{center}\small
\legendchip{gvsrc}{start} \quad \legendchip{gv}{unseen} \quad \legendchip{gvseen}{marked, still in the to-do list} \quad \legendchip{gvcur}{being processed} \quad \legendchip{gvdone}{finished} \quad \legendchip{gvdead}{never reached}
\end{center}

Here is the iterative DFS (a stack) on the running example $G$, starting at $0$, with the adjacency lists in input order: $\Adj{0}=(1,2)$, $\Adj{1}=(3)$, $\Adj{2}=(3)$, $\Adj{3}=(4)$, $\Adj{4}=(1)$.

\begin{table}[H]
\centering\small
\begin{tabular}{@{}clll@{}}
\toprule
\textbf{step} & \textbf{pop} & \textbf{newly marked (pushed)} & \textbf{stack after the step (top on the right)} \\
\midrule
1 & 0 & 1, 2 & 1 2 \\
2 & 2 & 3    & 1 3 \\
3 & 3 & 4    & 1 4 \\
4 & 4 & --- \ (1 is already marked) & 1 \\
5 & 1 & --- \ (3 is already marked) & (empty) \\
\bottomrule
\end{tabular}
\end{table}

\begin{diagrambox}
\begin{tikzpicture}[scale=0.7, transform shape]
    \begin{scope}[g1p2v0/.style={gvdone}, g1p2v1/.style={gvseen}, g1p2v2/.style={gvdone},
        g1p2v3/.style={gvcur}, g1p2v4/.style={gvseen},
        g1p2e01/.style={g1p2tree}, g1p2e02/.style={g1p2tree},
        g1p2e23/.style={g1p2tree}, g1p2e34/.style={g1p2tree}]
        \pic (A) {g1p2G};
        \node[font=\large\bfseries] at (0,2.6) {after step 3: stack $= [1, 4]$};
    \end{scope}
    \begin{scope}[xshift=9cm, g1p2v0/.style={gvsrc}, g1p2v1/.style={gvdone}, g1p2v2/.style={gvdone},
        g1p2v3/.style={gvdone}, g1p2v4/.style={gvdone},
        g1p2v5/.style={gvdead}, g1p2v6/.style={gvdead}, g1p2v7/.style={gvdead},
        g1p2e01/.style={g1p2tree}, g1p2e02/.style={g1p2tree},
        g1p2e23/.style={g1p2tree}, g1p2e34/.style={g1p2tree}]
        \pic (B) {g1p2G};
        \node[font=\large\bfseries] at (0,2.6) {after step 5: stack empty};
    \end{scope}
\end{tikzpicture}
\figcap{Left: 0 and 2 are finished, 3 was just processed (it pushed 4), 1 and 4 wait on the stack. Right: the final state; the thick edges $0{\to}1$, $0{\to}2$, $2{\to}3$, $3{\to}4$ form the search tree.}
\end{diagrambox}

\begin{invariantbox}[Why the marked set is exactly the reach set]
\textbf{Invariant (soundness).} Every marked vertex is reachable from $s$.\\
\emph{Proof.} $s$ is marked first and $s \leadsto s$. Any other vertex $v$ is marked while scanning an edge $u \to v$ of an already marked vertex $u$; by induction $s \leadsto u$, hence $s \leadsto v$.

\smallskip
\textbf{Closure at the end (completeness).} When the to-do list is empty, \emph{no edge leads from a marked vertex to an unmarked one}.\\
\emph{Proof.} Every marked vertex was added to the to-do list exactly once, so it was taken out and scanned exactly once; during that scan every out-neighbour that was not yet marked got marked.

\smallskip
\textbf{Conclusion.} Let $s = v_0 \to v_1 \to \dots \to v_k = t$ be any walk. If $t$ were unmarked, take the first unmarked vertex $v_j$ on the walk ($j \ge 1$ because $v_0 = s$ is marked); then $v_{j-1} \to v_j$ is an edge from a marked to an unmarked vertex, contradicting closure. So every reachable vertex is marked, and by soundness only reachable vertices are. \hfill$\blacksquare$

\smallskip
\textbf{Running time.} A vertex is marked once, so it enters the to-do list once and its adjacency list is scanned once: $\sum_u (1 + \outdeg{u}) = |V| + |E|$ steps for a directed graph, $|V| + 2|E|$ for an undirected one --- $\VE$ either way. Memory: the graph, one \texttt{seen} flag per vertex and a to-do list of at most $|V|$ entries.
\end{invariantbox}

Notice what the proof did \emph{not} use: the order in which vertices leave the to-do list. Stack, queue, recursion, even a random bag --- all give the same marked set. The order matters only for the \emph{extra} information a search provides: DFS timestamps (Chapter~2) and BFS distances (Chapter~3). For pure reachability, \textbf{pick whatever you can write without bugs} --- in this book that is usually the iterative stack version.

The closure argument has a picture that is worth remembering:

\begin{diagrambox}
\begin{tikzpicture}[g1p2v0/.style={gvsrc}, g1p2v1/.style={gvdone}, g1p2v2/.style={gvdone},
    g1p2v3/.style={gvdone}, g1p2v4/.style={gvdone}, g1p2v5/.style={gvdead}, g1p2v6/.style={gvdead},
    g1p2v7/.style={gvdead}, g1p2e50/.style={gpath, -{Stealth[length=3mm]}}]
    \pic (G) {g1p2G};
    \draw[gcut] (-0.75,2.2) -- (-0.75,-2.0);
    \node[note, anchor=west, text width=4.6cm, align=left] at (3.9,0.6) {\textbf{inside the red line:} $R(0)$\\\textbf{outside:} 5, 6, 7\\[3pt]The only edge crossing the line is $5 \to 0$, and it points \textbf{inward}. No edge leaves $R(0)$ --- that is exactly why the search stopped there.};
\end{tikzpicture}
\figcap{Any set containing $s$ with no edge leaving it contains all of $R(s)$, and the search returns the \emph{smallest} such set. ``Here is a set with no exit'' is how you \emph{prove} that something is impossible in a reachability problem.}
\end{diagrambox}

%------------------------------------------------------------
\subsubsection*{Two tools that come for free: many sources and the reverse graph}

\textbf{Many sources.} ``What can be reached from \emph{at least one} of $s_1, \dots, s_k$?'' Do \emph{not} run $k$ searches. Mark all sources at the start, put them all into the to-do list and run the same loop once: $\VE$ in total. (Imagine a super-source $\star$ with an edge to every $s_i$; the multi-source search is a single search from $\star$.)

\textbf{The reverse graph.} ``Which vertices can reach $t$?'' Flip every edge: $u \to v$ becomes $v \to u$. A walk $u \leadsto t$ in $G$ is exactly a walk $t \leadsto u$ read backwards in the reverse graph $G^{R}$, so one search from $t$ in $G^{R}$ answers the question.

\begin{diagrambox}
\begin{tikzpicture}[scale=0.7, transform shape]
    \begin{scope}[g1p2v6/.style={gvsrc}, g1p2v2/.style={gvsrc}, g1p2v3/.style={gvdone},
        g1p2v4/.style={gvdone}, g1p2v1/.style={gvdone}, g1p2v7/.style={gvdone}]
        \pic (A) {g1p2G};
        \node[font=\large\bfseries, align=center] at (0,2.8) {multi-source search from $\{6, 2\}$\\reaches $\{1,2,3,4,6,7\}$};
    \end{scope}
    \begin{scope}[xshift=9cm, g1p2e/.style={ge, {Stealth[length=2.4mm]}-},
        g1p2v3/.style={gvsrc}, g1p2v0/.style={gvdone}, g1p2v1/.style={gvdone}, g1p2v2/.style={gvdone},
        g1p2v4/.style={gvdone}, g1p2v5/.style={gvdone}, g1p2v6/.style={gvdead}, g1p2v7/.style={gvdead}]
        \pic (B) {g1p2G};
        \node[font=\large\bfseries, align=center] at (0,2.8) {reverse graph $G^R$, search from $3$:\\who can reach $3$? $\{0,1,2,3,4,5\}$};
    \end{scope}
\end{tikzpicture}
\figcap{Left: both sources start marked; 0 and 5 stay white because no edge enters them from the reached part. Right: every arrow of $G$ is flipped; 6 and 7 cannot reach 3 (they only reach each other).}
\end{diagrambox}

Combining the two tools answers ``\emph{which vertices lie on some walk from $s$ to $t$?}'': exactly those in $R(s) \cap R^{R}(t)$ --- forward from $s$ intersected with backward from $t$ (Exercise~E2.6).

%------------------------------------------------------------
\subsubsection*{Variations map}

\begin{diagrambox}
\begin{tikzpicture}[
    var/.style={rectangle, rounded corners=2pt, draw=patternborder, fill=patternbg,
                font=\scriptsize, align=center, text width=3.25cm, minimum height=1.25cm, inner sep=3pt}]
    \node[var, fill=answerfill, draw=accentorange, text width=4.6cm, font=\scriptsize\bfseries] (root) at (0,0)
        {Reachability from $s$\\one DFS/BFS, mark on push, $\VE$};
    \node[var] (v1) at (-5.4,-2) {\textbf{yes/no: $s \leadsto t$?}\\stop as soon as $t$ appears\\(Problem 1.2.1)};
    \node[var] (v2) at (-1.8,-2) {\textbf{``reach everything?''}\\number marked $= n$?\\(Problem 1.2.2)};
    \node[var] (v3) at (1.8,-2)  {\textbf{size of $R(s)$}\\chain reactions: count marks\\(E2.5)};
    \node[var] (v4) at (5.4,-2)  {\textbf{many sources}\\push them all first,\\one search};
    \node[var] (v5) at (-5.4,-3.6) {\textbf{who reaches $t$?}\\search the reverse graph\\(E2.4, E2.6)};
    \node[var] (v6) at (-1.8,-3.6) {\textbf{all pairs}, small $n$\\one search per vertex\\$\bigO{|V|(|V|+|E|)}$ (E2.3)};
    \node[var] (v7) at (1.8,-3.6)  {\textbf{edges given by a rule}\\keys, radii, portals:\\build adj.\ first (E2.1)};
    \node[var] (v8) at (5.4,-3.6)  {\textbf{many $a$--$b$ queries}\\undirected: components\\or DSU (Ch.~4, Ch.~6)};
    \node[draw=patternborder, dashed, rounded corners=4pt, fit=(v1)(v4)(v5)(v8), inner sep=5pt] (grp) {};
    \draw[flow] (root.south) -- (grp.north);
\end{tikzpicture}
\end{diagrambox}

\textbf{Twists and how to handle them}
\begin{itemize}[leftmargin=*]
    \item \textbf{Forbidden vertices or edges} (walls, closed roads, ``you may not pass through $x$''): never mark them, never follow them. Setting \texttt{seen[x] = 1} for every forbidden $x$ before the search starts does it in one line.
    \item \textbf{Edges defined by a rule} (a key lying in a room, a bomb's radius, a portal $i \to i + a_i$): either build the adjacency lists first or generate the neighbours inside the loop. For $|V| \le 100$ with an $\bigO{|V|^2}$ rule, building the full lists is simplest (E2.5).
    \item \textbf{Every vertex has one out-edge:} the ``search'' degenerates into a loop \texttt{cur = next[cur]} (E2.1, and all of Pattern~5).
    \item \textbf{Repeated queries from different sources:} each query costs a fresh search unless the graph is undirected (then label the components once). If $q \cdot n$ is too big to clear a \texttt{seen} array per query, a \emph{timestamp} array avoids the clearing (Mistake~5).
    \item \textbf{``Fewest starting points that reach everything'':} in a DAG these are exactly the vertices with in-degree~0 (E2.2) --- a reachability question answered by Pattern~1.
\end{itemize}

%------------------------------------------------------------
\subsubsection*{Code templates}

All four routines below are $\VE$; the all-pairs loop at the end of \texttt{main()} is $\bigO{|V|(|V|+|E|)}$. The \texttt{main()} uses the running example $G$; the values in the comments were printed by the program.

\begin{lstlisting}[style=cpp]
#include <bits/stdc++.h>
using namespace std;

// (1) iterative DFS: mark every vertex reachable from s (no recursion: safe for 10^6 vertices)
vector<char> reachFrom(const vector<vector<int>>& adj, int s) {
    vector<char> seen(adj.size(), 0);
    vector<int> st(1, s);
    seen[s] = 1;                                   // mark when PUSHED, never twice
    while (!st.empty()) {
        int u = st.back(); st.pop_back();
        for (int v : adj[u])
            if (!seen[v]) { seen[v] = 1; st.push_back(v); }
    }
    return seen;
}

// (2) BFS from a whole SET of sources at once (one search, not one per source)
vector<char> reachFromSet(const vector<vector<int>>& adj, const vector<int>& sources) {
    vector<char> seen(adj.size(), 0);
    queue<int> q;
    for (int s : sources)
        if (!seen[s]) { seen[s] = 1; q.push(s); }
    while (!q.empty()) {
        int u = q.front(); q.pop();
        for (int v : adj[u])
            if (!seen[v]) { seen[v] = 1; q.push(v); }
    }
    return seen;
}

// (3) the reverse graph: u -> v becomes v -> u.  "Who can reach t?" = reachFrom(rev, t)
vector<vector<int>> reversedGraph(const vector<vector<int>>& adj) {
    vector<vector<int>> rev(adj.size());
    for (int u = 0; u < (int)adj.size(); u++)
        for (int v : adj[u]) rev[v].push_back(u);
    return rev;
}

// (4) recursive DFS: shortest to write, but the depth can reach |V| (see Mistake 3)
void dfsRec(int u, const vector<vector<int>>& adj, vector<char>& seen) {
    seen[u] = 1;
    for (int v : adj[u])
        if (!seen[v]) dfsRec(v, adj, seen);
}

void show(const char* name, const vector<char>& seen) {
    cout << name << ":";
    for (int v = 0; v < (int)seen.size(); v++)
        if (seen[v]) cout << " " << v;
    cout << "\n";
}

int main() {
    int n = 8;
    int E[][2] = {{0,1},{0,2},{1,3},{2,3},{3,4},{4,1},{5,0},{5,6},{6,7},{7,6}};
    vector<vector<int>> adj(n);
    for (auto& e : E) adj[e[0]].push_back(e[1]);         // directed edges

    show("reach(0)", reachFrom(adj, 0));                  // reach(0): 0 1 2 3 4
    vector<char> seen(n, 0);
    dfsRec(0, adj, seen);
    show("dfsRec(0)", seen);                              // dfsRec(0): 0 1 2 3 4
    show("reach(6)", reachFrom(adj, 6));                  // reach(6): 6 7
    show("canReach(3)", reachFrom(reversedGraph(adj), 3)); // canReach(3): 0 1 2 3 4 5
    vector<int> src = {6, 2};
    show("reach({6,2})", reachFromSet(adj, src));         // reach({6,2}): 1 2 3 4 6 7

    long long pairs = 0;                                  // all pairs (s, t) with s ->* t
    for (int s = 0; s < n; s++) {
        vector<char> r = reachFrom(adj, s);
        pairs += count(r.begin(), r.end(), 1);
    }
    cout << "reachable pairs = " << pairs << "\n";        // reachable pairs = 30
}
\end{lstlisting}

\begin{insightbox}[Check the templates by hand]
\begin{itemize}[leftmargin=*]
    \item \textbf{reach(0)} $=\{0,1,2,3,4\}$ (the trace above); \textbf{reach(6)} $= \{6, 7\}$: the 2-cycle $6 \rightleftarrows 7$ has no exit.
    \item \textbf{canReach(3)}: in $G^R$ the search from 3 follows $3 \to 1$ and $3 \to 2$ (from $1{\to}3$, $2{\to}3$), then $1 \to 0$, $1 \to 4$, $2 \to 0$, $0 \to 5$, and stops: $\{0,1,2,3,4,5\}$.
    \item \textbf{All pairs:} $|R(v)|$ for $v = 0, \dots, 7$ is $5, 3, 4, 3, 3, 8, 2, 2$, and the sum is $30$ (pairs $(v,v)$ included).
\end{itemize}
\end{insightbox}

\begin{tricksbox}[Which version should I type?]
\begin{itemize}[leftmargin=*]
    \item \textbf{Yes/no or ``how many are reached'':} the iterative stack version (1). It is as short as the recursive one and never overflows.
    \item \textbf{Recursive DFS (4)} is fine when the depth is provably small: $|V|$ up to about $10^4$ with a default 1--8\,MB stack, or a tree of known small height. A path graph with $2\cdot 10^5$ vertices makes the recursion $2\cdot10^5$ frames deep --- a likely \stackoverflow{} on many judges.
    \item \textbf{BFS (2)} costs the same and gives distances for free --- use it whenever the next question might be ``in how many steps?''.
    \item \textbf{\texttt{vector<char>}} instead of \texttt{vector<bool>} for \texttt{seen}: one byte per vertex, no bit fiddling, usually a little faster.
\end{itemize}
\end{tricksbox}

%------------------------------------------------------------
\subsubsection*{Common mistakes}

\begin{warningbox}[Mistake 1 --- marking when popped instead of when pushed]
\begin{lstlisting}[style=snippet]
// WRONG: v may be pushed once per incoming edge -> the stack grows to O(|E|),
//        and every copy re-scans adj[v] unless you also check on pop
for (int v : adj[u]) if (!seen[v]) st.push_back(v);      // seen[v] set only when popped
// RIGHT: mark at the moment of discovery; each vertex enters the stack once
for (int v : adj[u]) if (!seen[v]) { seen[v] = 1; st.push_back(v); }
\end{lstlisting}
``Mark on pop'' is correct only with an extra \texttt{if (seen[u]) continue;} right after popping, and it still needs $\bigO{|E|}$ memory. Marking on push is simpler and never worse.
\end{warningbox}

\begin{warningbox}[Mistake 2 --- storing an undirected edge once]
\begin{lstlisting}[style=snippet]
adj[a].push_back(b);                          // WRONG if undirected: b cannot go back
adj[a].push_back(b); adj[b].push_back(a);     // RIGHT
\end{lstlisting}
Reread the statement: ``bi-directional'', ``roads can be used in both directions'' $\Rightarrow$ two entries; ``one-way'', ``from $a$ to $b$'', keys, portals $\Rightarrow$ one entry.
\end{warningbox}

\begin{warningbox}[Mistake 3 --- deep recursion]
Recursive DFS on a path of $n = 2\cdot10^5$ vertices nests $2\cdot 10^5$ calls. Each frame costs tens of bytes or more, and the default stack is often 1\,MB (Windows) or 8\,MB (Linux): the result is a crash without a useful message. Use the stack version whenever $|V|$ can exceed about $10^4$.
\end{warningbox}

\begin{warningbox}[{Mistake 4 --- searching in the wrong direction}]
``Can every city reach the capital?'' is \emph{not} answered by a search \emph{from} the capital in a directed graph. Search from the capital in the \textbf{reverse} graph (or flip each edge while reading the input).
\end{warningbox}

\begin{warningbox}[Mistake 5 --- resetting \texttt{seen} the expensive way]
With many test cases or many queries, either re-create \texttt{seen} for each test (\texttt{seen.assign(n, 0)} costs $\bigO{n}$ of the \emph{current} test --- fine) or use a timestamp: \texttt{stamp[v] == query} means ``seen in this query'', and starting a new query is just \texttt{query++}. Never \texttt{memset} a $10^6$-sized array $10^5$ times.
\end{warningbox}

\begin{warningbox}[Mistake 6 --- 1-indexed input in a 0-indexed array]
Vertices numbered $1..n$ in the input: either allocate $n+1$ lists or subtract 1 when reading. Mixing the two silently drops vertex $n$ or creates a phantom vertex $0$ that is ``never reached''.
\end{warningbox}

\begin{checklistbox}[``I'm stuck'' checklist for Pattern 2]
\begin{itemize}[leftmargin=*]
    \item[$\square$] What is a vertex? What is an edge? \textbf{Directed or not?} (Write the Graph Model box.)
    \item[$\square$] Is the question really ``reach'' --- or ``how far'' (BFS distances), ``how many groups'' (components), ``one arrow per vertex'' (Pattern~5)?
    \item[$\square$] From where do I search? One source, many sources (push them all), or from the target in the reverse graph?
    \item[$\square$] Is the answer a yes/no, the number of marked vertices, or ``all $n$ marked''?
    \item[$\square$] Can I afford one search per vertex ($|V| \cdot (|V|+|E|)$ operations)?
    \item[$\square$] Is the recursion depth safe? Are both directions of undirected edges stored? Is \texttt{seen} reset between tests?
    \item[$\square$] To prove impossibility: can I exhibit a set containing $s$, not containing $t$, with no edge leaving it?
\end{itemize}
\end{checklistbox}

%============================================================
% CH1 · P2 · PROBLEM 1 — FIND IF PATH EXISTS IN GRAPH (LeetCode 1971)
%============================================================
\subsection{Problem 1.2.1: Find if Path Exists in Graph}
\problemlabel{1.2.1}{ssec:ch1-path-exists}

\problemcard{1.2.1}{Find if Path Exists in Graph}{LeetCode 1971}{https://leetcode.com/problems/find-if-path-exists-in-graph/}{\exeasy\ \lc{Easy}\ (the canonical reachability question)}{One BFS/DFS from the source with early exit; DSU as an alternative}

\subsubsection*{1. Problem statement}

An undirected graph has $n$ vertices labelled $0..n-1$ and an edge list \texttt{edges}, where $\texttt{edges}[i] = [u_i, v_i]$ is a bi-directional edge. Return \texttt{true} if there is a valid path from vertex \texttt{source} to vertex \texttt{destination}, and \texttt{false} otherwise.

\textbf{Constraints:} $1 \le n \le 2\cdot 10^5$, \; $0 \le |\texttt{edges}| \le 2\cdot 10^5$, \; $0 \le u_i, v_i \le n-1$, \; $u_i \ne v_i$ (no self-loops), no duplicate edges, $0 \le \texttt{source}, \texttt{destination} \le n-1$.

\begin{table}[H]
\centering\small
\begin{tabular}{@{}clccl@{}}
\toprule
$n$ & \textbf{edges} & \textbf{source} & \textbf{destination} & \textbf{Output} \\
\midrule
3 & $[[0,1],[1,2],[2,0]]$ & 0 & 2 & \texttt{true} \quad ($0 \to 2$ directly, or $0 \to 1 \to 2$) \\
6 & $[[0,1],[0,2],[3,5],[5,4],[4,3]]$ & 0 & 5 & \texttt{false} \quad (5 is in another piece) \\
\bottomrule
\end{tabular}
\end{table}

\begin{diagrambox}
\begin{tikzpicture}
    \begin{scope}
        \node[gvsrc] (a0) at (0,0) {0};
        \node[gv] (a1) at (1.4,0.8) {1};
        \node[gvdst] (a2) at (1.4,-0.8) {2};
        \draw[ge] (a0) -- (a1); \draw[ge] (a1) -- (a2); \draw[ge] (a2) -- (a0);
        \node[note] at (0.7,-1.6) {Sample 1: source 0, destination 2\\answer \texttt{true}};
    \end{scope}
    \begin{scope}[xshift=5cm]
        \node[gvsrc] (b0) at (0,0) {0};
        \node[gv] (b1) at (1.4,0.8) {1};
        \node[gv] (b2) at (1.4,-0.8) {2};
        \node[gv] (b3) at (3.6,0.8) {3};
        \node[gvdst] (b5) at (5,0) {5};
        \node[gv] (b4) at (3.6,-0.8) {4};
        \draw[ge] (b0) -- (b1); \draw[ge] (b0) -- (b2);
        \draw[ge] (b3) -- (b5); \draw[ge] (b5) -- (b4); \draw[ge] (b4) -- (b3);
        \draw[gcut] (2.5,1.3) -- (2.5,-1.3);
        \node[note] at (2.5,-1.6) {Sample 2: source 0, destination 5\\no edge crosses the red line $\Rightarrow$ \texttt{false}};
    \end{scope}
\end{tikzpicture}
\figcap{The two official samples. In Sample~2 the set $\{0,1,2\}$ has no edge leaving it, which is a complete proof that 5 is unreachable.}
\end{diagrambox}

\subsubsection*{2. The graph model}

\begin{modelbox}[Graph Model --- Find if Path Exists in Graph]
\renewcommand{\arraystretch}{1.35}
\begin{tabularx}{\linewidth}{@{}>{\bfseries}l X@{}}
Vertices & the $n$ labelled vertices $0..n-1$ --- the graph is given explicitly \\
Edges & every $[u, v]$ is undirected: store it as $u \to v$ \emph{and} $v \to u$ (adjacency lists) \\
Weights & none \\
Question & is \texttt{destination} $\in R(\texttt{source})$? Equivalently: are the two vertices in the same connected component? \\
Size & $|V| \le 2\cdot10^5$, $|E| \le 2\cdot 10^5$ ($4\cdot10^5$ adjacency entries): any $\VE$ algorithm runs in a few milliseconds; $\bigO{|V|\cdot|E|} = 4\cdot10^{10}$ is far too slow \\
\end{tabularx}
\end{modelbox}

The input is an \emph{edge list}; a search needs \emph{adjacency lists} (``who are the neighbours of $u$?''). Converting one into the other is the first five lines of every solution:

\begin{diagrambox}
\begin{tikzpicture}[lst/.style={font=\small\ttfamily, anchor=west}]
    \node[font=\small\bfseries, anchor=west] at (0,1.9) {edge list (input)};
    \foreach \e [count=\i] in {{[0,1]},{[0,2]},{[3,5]},{[5,4]},{[4,3]}}
        \node[lst] at (0,1.9-0.45*\i) {\e};
    \draw[flow] (1.6,0.6) -- node[lbl, above] {each edge twice} (4.0,0.6);
    \node[font=\small\bfseries, anchor=west] at (4.3,1.9) {adjacency lists (built)};
    \foreach \v/\l [count=\i] in {0/{1, 2}, 1/{0}, 2/{0}, 3/{5, 4}, 4/{5, 3}, 5/{3, 4}}
        \node[lst] at (4.3,1.9-0.4*\i) {adj[\v] = \{\l\}};
\end{tikzpicture}
\figcap{Sample 2. Edge $[5,4]$ produces two entries: \texttt{adj[5]} gets 4 and \texttt{adj[4]} gets 5. The order inside each list follows the input order.}
\end{diagrambox}

\subsubsection*{3. Key observations}

\begin{itemize}[leftmargin=*]
    \item \textbf{Only yes/no is asked.} We do not need the path, its length or the number of paths --- exactly the reach-set question of the pattern.
    \item \textbf{Local information is not enough.} ``Both endpoints have at least one edge'' is true in Sample~2, yet the answer is \texttt{false}. Reachability depends on whole chains of edges.
    \item \textbf{$\texttt{source} = \texttt{destination}$ is \texttt{true}}, even with no edges at all (a walk of zero edges). With $n = 1$ this is the only possible query.
    \item \textbf{Early exit.} The search may stop the moment \texttt{destination} is reached; on a \texttt{false} instance it must explore the whole component of \texttt{source}.
    \item \textbf{Depth.} A path graph $0 - 1 - \dots - (n{-}1)$ with $n = 2\cdot10^5$ is allowed, so a recursive DFS may go $2\cdot10^5$ calls deep. Prefer BFS or the iterative stack.
\end{itemize}

\subsubsection*{4. Approaches}

\begin{bruteforcebox}
\textbf{Idea: grow the reached set by sweeping the edge list until nothing changes.} Start with $\{\texttt{source}\}$. In one \emph{pass}, look at every edge $[u,v]$: if exactly one endpoint is reached, mark the other one. Repeat passes until a whole pass marks nothing.

\begin{lstlisting}[style=cpp]
#include <bits/stdc++.h>
using namespace std;

class Solution {
public:
    bool validPath(int n, vector<vector<int>>& edges, int source, int destination) {
        vector<char> reached(n, 0);
        reached[source] = 1;
        bool changed = true;
        while (changed) {                          // one pass over all edges
            changed = false;
            for (auto& e : edges) {
                int u = e[0], v = e[1];
                if (reached[u] && !reached[v]) { reached[v] = 1; changed = true; }
                else if (reached[v] && !reached[u]) { reached[u] = 1; changed = true; }
            }
        }
        return reached[destination];
    }
};
\end{lstlisting}

It is correct (when a pass changes nothing, no edge connects a reached and an unreached vertex --- the closure property of the pattern), but a pass may add only \emph{one} new vertex. On the path $0-1-2-3-4$ with the edges listed backwards, $[3,4], [2,3], [1,2], [0,1]$:

\begin{center}
\begin{tikzpicture}[x=1.05cm]
    \foreach \p/\k in {1/1, 2/2, 3/3, 4/4, 5/4}{
        \begin{scope}[yshift=-\p*0.62cm]
            \node[font=\scriptsize\bfseries, anchor=east] at (-0.5,0) {pass \p};
            \foreach \i in {0,...,4}{
                \ifnum\i>\k \node[gv, minimum size=5.5mm, font=\scriptsize\bfseries] (p\p\i) at (\i,0) {\i};
                \else \ifnum\i=0 \node[gvsrc, minimum size=5.5mm, font=\scriptsize\bfseries] (p\p\i) at (\i,0) {\i};
                \else \node[gvdone, minimum size=5.5mm, font=\scriptsize\bfseries] (p\p\i) at (\i,0) {\i}; \fi\fi
            }
            \foreach \i [evaluate=\i as \j using int(\i+1)] in {0,...,3} \draw[ge] (p\p\i) -- (p\p\j);
        \end{scope}
    }
    \node[note, anchor=west] at (4.6,-0.62) {edge $[0,1]$ is last in the list:\\only vertex 1 is added};
    \node[note, anchor=west] at (4.6,-3.1) {pass 5 changes nothing: stop};
\end{tikzpicture}
\end{center}

With $n$ vertices on such a path we need $n$ passes of $|E|$ edges each (the program counted exactly $1000$ passes for $n = 1000$).

\begin{complexitybox}
\textbf{Time:} $\bigO{|V| \cdot |E|}$ --- up to $2\cdot10^5 \times 2\cdot10^5 = 4\cdot10^{10}$ edge checks. \textbf{Space:} $\bigO{|V|}$. \textbf{Verdict:} \tle{}. (Enumerating all simple paths would be even worse: exponential.)
\end{complexitybox}
\end{bruteforcebox}

\begin{algorithmbox}[Approach 2: BFS from the source (early exit)]
\textbf{Idea.} Build adjacency lists once; then each vertex is processed once and each edge is looked at twice (once from each end). The brute force wasted time re-reading edges whose endpoints had not changed; a search only looks at the edges of a vertex at the moment that vertex becomes reached.

\begin{lstlisting}[style=cpp]
#include <bits/stdc++.h>
using namespace std;

class Solution {
public:
    bool validPath(int n, vector<vector<int>>& edges, int source, int destination) {
        vector<vector<int>> adj(n);
        for (auto& e : edges) {                    // undirected: store both directions
            adj[e[0]].push_back(e[1]);
            adj[e[1]].push_back(e[0]);
        }
        vector<char> seen(n, 0);
        queue<int> q;
        q.push(source);
        seen[source] = 1;
        while (!q.empty()) {
            int u = q.front(); q.pop();
            if (u == destination) return true;     // early exit (covers source == dest)
            for (int v : adj[u])
                if (!seen[v]) { seen[v] = 1; q.push(v); }
        }
        return false;                              // component of source fully explored
    }
};
\end{lstlisting}

\begin{complexitybox}
\textbf{Time:} $\bigO{|V| + |E|}$ --- building the lists is $\bigO{|V|+|E|}$, the search touches every vertex once and every adjacency entry at most once. \textbf{Space:} $\bigO{|V|+|E|}$ for the lists, $\bigO{|V|}$ for \texttt{seen} and the queue. On a path of $2\cdot 10^5$ vertices the program answered \texttt{true} in about $0.06$\,s, including building the graph. \textbf{Verdict:} \ac{}.
\end{complexitybox}
\end{algorithmbox}

\begin{invariantbox}[Why it is correct]
This is the pattern's search with one extra line, so its invariant applies: every vertex in the queue or already popped is reachable from \texttt{source} (soundness), and when the queue runs empty no edge leaves the marked set (closure), so the marked set is exactly the component of \texttt{source}.
\begin{itemize}[leftmargin=*]
    \item If the function returns \texttt{true}, \texttt{destination} was popped, hence marked, hence reachable.
    \item If it returns \texttt{false}, the queue emptied without popping \texttt{destination}. Every marked vertex is eventually popped, so \texttt{destination} was never marked, and by closure it is not reachable.
\end{itemize}
The early exit only stops the loop sooner; it never changes which answer is given.
\end{invariantbox}

\begin{alternativebox}[Approach 3: Disjoint Set Union (preview of Chapter 6)]
\textbf{Idea.} Forget searching: process the edges one by one and maintain the \emph{groups of vertices known to be connected}. Every vertex points to a \emph{parent}; following parents leads to the group's \emph{root} (its representative). An edge $[u,v]$ merges two groups by hanging one root under the other. At the end, \texttt{source} and \texttt{destination} are connected iff they have the same root.

\begin{lstlisting}[style=cpp]
#include <bits/stdc++.h>
using namespace std;

class Solution {
    vector<int> par;
    int findRoot(int x) {
        while (par[x] != x) {
            par[x] = par[par[x]];                  // path halving: shortcut every other step
            x = par[x];
        }
        return x;
    }
public:
    bool validPath(int n, vector<vector<int>>& edges, int source, int destination) {
        par.resize(n);
        for (int i = 0; i < n; i++) par[i] = i;   // every vertex is its own group
        for (auto& e : edges) {
            int a = findRoot(e[0]), b = findRoot(e[1]);
            if (a != b) par[a] = b;                 // merge the two groups
        }
        return findRoot(source) == findRoot(destination);
    }
};
\end{lstlisting}

On Sample~2 the parent array evolves as follows (output of the program):
\begin{center}\small
\begin{tabular}{@{}lcccccc@{}}
\toprule
\textbf{after edge} & \texttt{par[0]} & \texttt{par[1]} & \texttt{par[2]} & \texttt{par[3]} & \texttt{par[4]} & \texttt{par[5]} \\
\midrule
start     & 0 & 1 & 2 & 3 & 4 & 5 \\
$[0,1]$   & 1 & 1 & 2 & 3 & 4 & 5 \\
$[0,2]$   & 1 & 2 & 2 & 3 & 4 & 5 \\
$[3,5]$   & 1 & 2 & 2 & 5 & 4 & 5 \\
$[5,4]$   & 1 & 2 & 2 & 5 & 4 & 4 \\
$[4,3]$   & 1 & 2 & 2 & 4 & 4 & 4 \\
\bottomrule
\end{tabular}
\end{center}
At the last edge, \texttt{findRoot(3)} walks $3 \to 5 \to 4$ and shortcuts \texttt{par[3]} to $4$; both ends already share root $4$, so nothing is merged.

\begin{center}
\begin{tikzpicture}[x=1.2cm, y=0.95cm]
    \node[gvA] (r2) at (0,0) {2}; \node[gvA] (n1) at (0,-1) {1}; \node[gvsrc] (n0) at (0,-2) {0};
    \draw[gde] (n1) -- (r2); \draw[gde] (n0) -- (n1);
    \node[gvB] (r4) at (3,0) {4}; \node[gvB] (n3) at (2.4,-1) {3}; \node[gvdst] (n5) at (3.6,-1) {5};
    \draw[gde] (n3) -- (r4); \draw[gde] (n5) -- (r4);
    \node[note, anchor=west] at (4.3,-0.6) {root$(0) = 2 \ne 4 = $ root$(5)$\\$\Rightarrow$ \texttt{false}};
\end{tikzpicture}
\end{center}

\begin{complexitybox}
\textbf{Time:} $\bigO{(|V|+|E|)\,\log|V|}$ with path halving alone, practically linear; with union by size as well it is $\bigO{(|V|+|E|)\,\alpha(|V|)}$ (Chapter~6). \textbf{Space:} $\bigO{|V|}$ --- no adjacency lists at all. \textbf{Verdict:} \ac{}.
\end{complexitybox}

\textbf{When is DSU the better tool?} When edges arrive \emph{over time} mixed with ``are $a$ and $b$ connected now?'' queries, a search would have to be repeated after every new edge, while DSU just keeps merging. For a single query on a fixed graph, BFS and DSU are equally good.
\end{alternativebox}

\begin{tricksbox}[Approach 4: Implementation Refinements]
\textbf{Iterative DFS.} Replacing the queue by a \texttt{vector} used as a stack gives the same answer (the pattern's proof did not depend on the order):
\begin{lstlisting}[style=snippet]
vector<int> st(1, source);  seen[source] = 1;
while (!st.empty()) {
    int u = st.back(); st.pop_back();
    if (u == destination) return true;
    for (int v : adj[u]) if (!seen[v]) { seen[v] = 1; st.push_back(v); }
}
\end{lstlisting}
\textbf{Why not recursion here?} \texttt{dfs(u)} calling \texttt{dfs(v)} along the path $0 - 1 - \dots - 199999$ is $2 \cdot 10^5$ frames deep. Some judges survive it, others crash; the stack version costs nothing extra.

\textbf{Early exit at discovery time.} Testing \texttt{v == destination} when $v$ is \emph{marked} (instead of when popped) finishes one layer earlier. It then needs a separate check for $\texttt{source} = \texttt{destination}$; checking on pop, as in Approach~2, needs none.
\end{tricksbox}

\subsubsection*{5. Trace}

\begin{tracebox}[{Sample 2: $n = 6$, source $0$, destination $5$}]
Adjacency lists (input order): \texttt{adj[0]=\{1,2\}}, \texttt{adj[1]=\{0\}}, \texttt{adj[2]=\{0\}}, \texttt{adj[3]=\{5,4\}}, \texttt{adj[4]=\{5,3\}}, \texttt{adj[5]=\{3,4\}}.

\begin{center}
\begin{tikzpicture}[scale=0.62, transform shape]
    \foreach \k/\t/\a/\b/\c/\d/\e/\f in {
        0/{start: queue $[0]$}/gvsrc/gv/gv/gv/gv/gvdst,
        1/{pop 0, push 1, 2: queue $[1,2]$}/gvdone/gvseen/gvseen/gv/gv/gvdst,
        2/{pop 1, pop 2: queue empty}/gvdone/gvdone/gvdone/gvdead/gvdead/gvdead}{
        \begin{scope}[xshift=\k*8.3cm]
            \node[\a] (x0) at (0,0) {0};
            \node[\b] (x1) at (1.4,0.8) {1};
            \node[\c] (x2) at (1.4,-0.8) {2};
            \node[\d] (x3) at (3.6,0.8) {3};
            \node[\f] (x5) at (5,0) {5};
            \node[\e] (x4) at (3.6,-0.8) {4};
            \draw[ge] (x3) -- (x5); \draw[ge] (x5) -- (x4); \draw[ge] (x4) -- (x3);
            \ifnum\k>0 \draw[gtree] (x0) -- (x1); \draw[gtree] (x0) -- (x2);
            \else \draw[ge] (x0) -- (x1); \draw[ge] (x0) -- (x2); \fi
            \node[font=\large\bfseries] at (2.5,1.7) {\t};
        \end{scope}
    }
\end{tikzpicture}
\end{center}

\begin{center}\small
\begin{tabular}{@{}clll@{}}
\toprule
\textbf{step} & \textbf{pop} & \textbf{push (newly marked)} & \textbf{queue after (front on the left)} \\
\midrule
1 & 0 & 1, 2 & 1 2 \\
2 & 1 & --- \ (0 is marked) & 2 \\
3 & 2 & --- \ (0 is marked) & (empty) $\Rightarrow$ return \texttt{false} \\
\bottomrule
\end{tabular}
\end{center}
The search never looks at the edges $[3,5], [5,4], [4,3]$: it only pays for the component it explores. The finished set $\{0,1,2\}$ is the ``no edge leaves'' certificate drawn in the statement.
\end{tracebox}

\begin{tracebox}[{Sample 1: $n = 3$, source $0$, destination $2$}]
Adjacency lists: \texttt{adj[0]=\{1,2\}}, \texttt{adj[1]=\{0,2\}}, \texttt{adj[2]=\{1,0\}}.
\begin{center}
\begin{tikzpicture}[scale=0.8, transform shape]
    \foreach \k/\t/\a/\b/\c in {
        0/{start: queue $[0]$}/gvsrc/gv/gvdst,
        1/{pop 0: queue $[1,2]$}/gvdone/gvseen/gvseen,
        2/{pop 1: queue $[2]$}/gvdone/gvdone/gvseen,
        3/{pop 2 $=$ destination}/gvdone/gvdone/gvcur}{
        \begin{scope}[xshift=\k*4.1cm]
            \node[\a] (y0) at (0,0) {0};
            \node[\b] (y1) at (1.4,0.8) {1};
            \node[\c] (y2) at (1.4,-0.8) {2};
            \draw[ge] (y1) -- (y2);
            \ifnum\k>0 \draw[gtree] (y0) -- (y1); \draw[gtree] (y0) -- (y2);
            \else \draw[ge] (y0) -- (y1); \draw[ge] (y0) -- (y2); \fi
            \node[font=\small\bfseries] at (0.7,1.55) {\t};
        \end{scope}
    }
\end{tikzpicture}
\end{center}
Popping 1 finds both neighbours already marked. Popping 2 hits the early exit: \texttt{true}. The search tree ($0{-}1$, $0{-}2$) shows the path BFS found: $0 \to 2$, one edge.
\end{tracebox}

\subsubsection*{6. Test suite}

\begin{testbox}
\centering\small
\begin{tabular}{@{}p{6.6cm}cp{5.0cm}@{}}
\toprule
\textbf{Input} & \textbf{Expected} & \textbf{Why} \\
\midrule
$n=1$, no edges, $0 \to 0$ & \texttt{true} & zero-edge walk; smallest input \\
$n=2$, no edges, $0 \to 1$ & \texttt{false} & no edges at all \\
Sample 2 with source $5$, destination $3$ & \texttt{true} & the other component works \\
path $0-1-\dots-199999$, $0 \to 199999$ & \texttt{true} & depth $2\cdot10^5$: kills naive recursion \\
star with centre 0, $n = 2\cdot10^5$, $1 \to 2$ & \texttt{true} & one vertex with $2\cdot10^5-1$ neighbours \\
two paths of $10^5$ vertices each, $0 \to 199999$ & \texttt{false} & full exploration of a large component \\
\bottomrule
\end{tabular}
\end{testbox}
All six expected values were produced by both the BFS and the DSU program; in addition, BFS, iterative DFS, DSU and the brute force agreed with a Floyd--Warshall-style closure on 20\,000 random graphs with up to 8 vertices.

\subsubsection*{7. Common mistakes}

\begin{warningbox}[Mistake 1 --- one direction only]
\begin{lstlisting}[style=snippet]
for (auto& e : edges) adj[e[0]].push_back(e[1]);                          // WRONG
for (auto& e : edges) { adj[e[0]].push_back(e[1]); adj[e[1]].push_back(e[0]); }  // RIGHT
\end{lstlisting}
Both official samples still pass with the wrong version: Sample~1 keeps the cycle $0\to1\to2\to0$, and Sample~2 is \texttt{false} anyway. Sample~2's graph with source $1$ and destination $2$ catches it: the right answer is \texttt{true} ($1 - 0 - 2$), but the one-way lists give vertex 1 no neighbours at all, so the wrong version says \texttt{false} (checked by running both).
\end{warningbox}

\begin{warningbox}[Mistake 2 --- forgetting the trivial case]
Code that only checks \texttt{v == destination} for \emph{newly discovered} neighbours returns \texttt{false} for $\texttt{source} = \texttt{destination}$ (no neighbour of $0$ is $0$ itself; with $n=1$ there are no neighbours at all). Check when popping, or test \texttt{source == destination} first.
\end{warningbox}

\begin{warningbox}[Mistake 3 --- an adjacency matrix]
\texttt{vector<vector<bool>> g(n, vector<bool>(n))} with $n = 2\cdot 10^5$ needs $4\cdot 10^{10}$ bits $= 5$\,GB, and scanning a row costs $\bigO{n}$ per vertex, $\bigO{n^2}$ in total. Adjacency lists are $\bigO{|V|+|E|}$ in memory and time.
\end{warningbox}

\begin{warningbox}[Mistake 4 --- recursion on a long path]
A recursive DFS passes the samples and small random tests, then crashes on the $2\cdot10^5$-vertex path. Use BFS or an explicit stack (Approach~4).
\end{warningbox}

\subsubsection*{8. Takeaways}

\begin{itemize}[leftmargin=*]
    \item ``Is there a path from $a$ to $b$?'' $=$ ``is $b$ in $R(a)$?'' $=$ one search, $\VE$.
    \item Build adjacency lists from the edge list first; undirected $=$ both directions.
    \item Check the target when popping: it handles $\texttt{source} = \texttt{destination}$ for free.
    \item A failed search \emph{proves} impossibility: the explored set has no edge leaving it.
    \item DSU answers the same question without adjacency lists and shines when edges arrive over time (Chapter~6).
\end{itemize}

\begin{variantsbox}
\variant{One-way edges}{\varin{https://codeforces.com/problemset/problem/500/A}}
  {each edge can be used only from $u$ to $v$ (portals, one-way streets).}
  {same vertices; store only $u \to v$. Reachability is no longer symmetric.}
  {the same search; E2.1 is an instance where every vertex has a single out-edge, so the search becomes a loop.}
\variant{Many queries, fixed graph}{\varnone}
  {$q \le 2\cdot10^5$ questions ``is there a path from $a_i$ to $b_i$?'' on one undirected graph.}
  {label every vertex with the id of its connected component.}
  {one search per \emph{unlabelled} vertex, $\VE$ in total (Chapter~4); then each query is \texttt{comp[a] == comp[b]}, $\bigO{1}$.}
\variant{Edges arrive between queries}{\varnone}
  {the edge list is revealed one edge at a time, interleaved with connectivity questions.}
  {the same vertices; the graph only grows.}
  {DSU: \texttt{unite} for each edge, compare roots for each query, near-$\bigO{1}$ per operation (Chapter~6).}
\variant{Forbidden vertices}{\varnone}
  {the path may not pass through a given set $B$ of blocked vertices.}
  {delete $B$ from the graph --- or simply mark every vertex of $B$ as \texttt{seen} before the search.}
  {unchanged search; if \texttt{source} or \texttt{destination} is blocked, answer \texttt{false} (unless the statement says otherwise).}
\variant{Print a path}{\varnone}
  {output one path from \texttt{source} to \texttt{destination}, not just yes/no.}
  {store \texttt{parent[v] = u} at the moment $v$ is marked from $u$.}
  {walk \texttt{parent} back from \texttt{destination} and reverse; with BFS this path even has the fewest edges (Chapter~3).}
\end{variantsbox}

%============================================================
% CH1 · P2 · PROBLEM 2 — KEYS AND ROOMS (LeetCode 841)
%============================================================
\subsection{Problem 1.2.2: Keys and Rooms}
\problemlabel{1.2.2}{ssec:ch1-keys-rooms}

\problemcard{1.2.2}{Keys and Rooms}{LeetCode 841}{https://leetcode.com/problems/keys-and-rooms/}{\exmedium\ \lc{Medium}\ (the graph is hidden in the story)}{Model ``key to room $k$ found in room $r$'' as an edge $r \to k$; one search from room 0; count the reached rooms}

\subsubsection*{1. Problem statement}

There are $n$ rooms labelled $0..n-1$. All rooms are locked except room~$0$. Inside room $i$ lies a set of keys $\texttt{rooms}[i]$; a key with number $k$ opens room $k$, and you may carry every key you find. Return \texttt{true} if you can visit \emph{all} the rooms, \texttt{false} otherwise.

\textbf{Constraints:} $2 \le n \le 1000$, \; $0 \le |\texttt{rooms}[i]| \le 1000$, \; $1 \le \sum_i |\texttt{rooms}[i]| \le 3000$, \; $0 \le \texttt{rooms}[i][j] < n$, the keys inside one room are distinct.

\begin{table}[H]
\centering\small
\begin{tabular}{@{}lcl@{}}
\toprule
\textbf{rooms} & \textbf{Output} & \textbf{Why} \\
\midrule
$[[1],[2],[3],[\,]]$ & \texttt{true} & room 0 gives key 1, room 1 gives key 2, room 2 gives key 3 \\
$[[1,3],[3,0,1],[2],[0]]$ & \texttt{false} & the only key to room 2 lies inside room 2 \\
\bottomrule
\end{tabular}
\end{table}

\begin{diagrambox}
\begin{tikzpicture}[
    room/.style={rectangle, rounded corners=2pt, draw=deepblue, fill=theorybg, line width=0.9pt,
                 minimum width=1.5cm, minimum height=1.0cm, font=\small\bfseries, align=center},
    lockd/.style={room, draw=recursionred, fill=dupfill},
    key/.style={font=\scriptsize, text=accentorange}]
    \node[room] (r0) at (0,0) {room 0\\[-2pt]{\scriptsize\mdseries open}};
    \node[room] (r1) at (2.6,0) {room 1};
    \node[lockd] (r2) at (5.2,0) {room 2};
    \node[room] (r3) at (7.8,0) {room 3};
    \node[key] at (0,-0.85) {keys: 1, 3};
    \node[key] at (2.6,-0.85) {keys: 3, 0, 1};
    \node[key] at (6.3,-0.85) {keys: 2};
    \node[key] at (7.8,-0.85) {keys: 0};
    \draw[flow] (r0.north) to[bend left=25] node[lbl, above]{key 1} (r1.north);
    \draw[flow] (r0.north east) to[bend left=30] node[lbl, above]{key 3} (r3.north west);
    \draw[warningred, thick, -{Stealth[length=2.5mm]}] ([xshift=4mm]r2.south) to[out=-60, in=-120, looseness=4] node[lbl, below, text=warningred]{key 2} ([xshift=-4mm]r2.south);
\end{tikzpicture}
\figcap{Sample 2 as a building. Room 2 is red: its only key is locked inside it, so it can never be opened.}
\end{diagrambox}

\subsubsection*{2. The graph model}

The statement never says ``graph''. The modelling step is to notice that \emph{finding key $k$ in room $r$} means \emph{``once I am in $r$, I can go to $k$''} --- an edge $r \to k$.

\begin{modelbox}[Graph Model --- Keys and Rooms]
\renewcommand{\arraystretch}{1.35}
\begin{tabularx}{\linewidth}{@{}>{\bfseries}l X@{}}
Vertices & the $n$ rooms \\
Edges & directed: $r \to k$ for every key $k \in \texttt{rooms}[r]$. The input \emph{is} the adjacency list --- nothing to build. Self-loops are possible (room 1 of Sample~2 contains key 1) and harmless \\
Weights & none \\
Question & is $R(0) = \{0, 1, \dots, n-1\}$? i.e.\ does one search from room 0 mark all $n$ vertices? \\
Size & $|V| \le 1000$, $|E| = \sum|\texttt{rooms}[i]| \le 3000$: tiny. Even the $\bigO{|V|\cdot(|V|+|E|)}$ brute force below fits; the point is to see the graph \\
\end{tabularx}
\end{modelbox}

\begin{diagrambox}
\begin{tikzpicture}
    \begin{scope}
        \node[gvsrc] (a0) at (0,0) {0}; \node[gv] (a1) at (1.4,0) {1};
        \node[gv] (a2) at (2.8,0) {2}; \node[gv] (a3) at (4.2,0) {3};
        \draw[gde] (a0) -- (a1); \draw[gde] (a1) -- (a2); \draw[gde] (a2) -- (a3);
        \node[note] at (2.1,-0.9) {Sample 1: a chain $0\to1\to2\to3$, \texttt{true}};
    \end{scope}
    \begin{scope}[xshift=6.6cm]
        \node[gvsrc] (b0) at (0,0.9) {0}; \node[gv] (b1) at (2.0,0.9) {1};
        \node[gv] (b3) at (1.0,-0.7) {3}; \node[gvdst] (b2) at (3.8,0) {2};
        \draw[gde] (b0) to[bend left=15] (b1); \draw[gde] (b1) to[bend left=15] (b0);
        \draw[gde] (b0) to[bend left=15] (b3); \draw[gde] (b3) to[bend left=15] (b0);
        \draw[gde] (b1) -- (b3);
        \draw[gde] (b1) to[out=60, in=120, looseness=6] (b1);
        \draw[gde] (b2) to[out=60, in=120, looseness=6] (b2);
        \draw[gcut] (2.9,1.7) -- (2.9,-1.1);
        \node[note] at (1.9,-1.45) {Sample 2: no edge enters 2 from outside, \texttt{false}};
    \end{scope}
\end{tikzpicture}
\figcap{Left: the rooms of Sample 1. Right: Sample 2. Room 2's only incoming edge is its own self-loop; the red line separates $\{0,1,3\}$, which has no edge leaving it.}
\end{diagrambox}

\subsubsection*{3. Key observations}

\begin{itemize}[leftmargin=*]
    \item \textbf{Keys are never used up and doors never close.} Holding more keys never hurts, so the \emph{order} in which we open rooms does not matter. The set of rooms we can eventually enter is fixed --- it is $R(0)$. (Contrast: if a key disappeared when used, the order would matter and this would be a much harder search over states.)
    \item \textbf{``Visit all'' $=$ ``everything reachable''.} Count the rooms marked by the search and compare with $n$.
    \item \textbf{Direction matters.} Key $k$ in room $r$ helps only after reaching $r$: edge $r \to k$, never $k \to r$. Room 3 of Sample~2 holds key 0, which is useless --- room 0 was open anyway.
    \item \textbf{A local check is not enough.} ``Every room except 0 has a key somewhere'' holds in Sample~2 (key 2 is in room 2) and still the answer is \texttt{false}: the key has to lie in a room you can \emph{reach}.
\end{itemize}

\subsubsection*{4. Approaches}

\begin{bruteforcebox}
\textbf{Idea: simulate a visitor who repeatedly walks along the corridor.} Each pass walks past rooms $0, 1, \dots, n-1$ and enters every room for which a key is in hand and that has not been entered yet, picking up its keys. Stop when a whole pass enters nothing.

\begin{lstlisting}[style=cpp]
#include <bits/stdc++.h>
using namespace std;

class Solution {
public:
    bool canVisitAllRooms(vector<vector<int>>& rooms) {
        int n = rooms.size();
        vector<char> haveKey(n, 0), entered(n, 0);
        haveKey[0] = 1;                                // room 0 is unlocked
        bool changed = true;
        while (changed) {                              // one walk along the corridor
            changed = false;
            for (int r = 0; r < n; r++)
                if (haveKey[r] && !entered[r]) {
                    entered[r] = 1; changed = true;
                    for (int k : rooms[r]) haveKey[k] = 1;
                }
        }
        for (int r = 0; r < n; r++) if (!entered[r]) return false;
        return true;
    }
};
\end{lstlisting}

Correct by the same closure argument as the pattern, but a pass can enter a single room. Worst case: room 0 holds key $n{-}1$, and room $i \ge 2$ holds key $i{-}1$. With $n = 6$ the program printed:
\begin{center}\small
\begin{tabular}{@{}lcccccc@{}}
\toprule
pass & 1 & 2 & 3 & 4 & 5 & 6 \\
rooms entered & 0, 5 & 4 & 3 & 2 & 1 & --- \\
\bottomrule
\end{tabular}
\end{center}
\begin{center}
\begin{tikzpicture}[x=1.25cm]
    \foreach \i/\p in {0/1, 1/5, 2/4, 3/3, 4/2, 5/1}{
        \ifnum\i=0 \node[gvsrc] (q\i) at (\i,0) {\i}; \else \node[gvdone] (q\i) at (\i,0) {\i}; \fi
        \node[gdist] at (\i,-0.6) {pass \p};
    }
    \draw[gde] (q0) to[bend left=30] (q5);
    \foreach \i [evaluate=\i as \j using int(\i-1)] in {5,4,3,2} \draw[gde] (q\i) to[bend left=30] (q\j);
    \draw[->, thick, darkgray] (-0.4,-1.05) -- node[note, below] {direction of every corridor walk} (5.4,-1.05);
\end{tikzpicture}
\end{center}
The key chain $0 \to 5 \to 4 \to 3 \to 2 \to 1$ runs \emph{against} the corridor: each key is found just after the walk has passed its door, so for $n = 1000$ the program needed $1000$ passes.

\begin{complexitybox}
\textbf{Time:} $\bigO{n \cdot (n + \sum|\texttt{rooms}[i]|)}$ --- about $4 \cdot 10^6$ here, so it is even accepted, but it is quadratic: with $n = 10^5$ rooms it would be hopeless. \textbf{Space:} $\bigO{n}$. \textbf{Verdict:} \ac{} only thanks to the small limits.
\end{complexitybox}
\end{bruteforcebox}

\begin{algorithmbox}[Approach 2: DFS with a stack of unused keys]
\textbf{Idea.} Instead of walking past every door again and again, keep a \textbf{pile of keys for rooms we have not entered yet}. Take any key from the pile, enter that room, and add its new keys to the pile. This is the pattern's search: the pile is the to-do stack, ``room opened'' is \texttt{seen}.

\begin{lstlisting}[style=cpp]
#include <bits/stdc++.h>
using namespace std;

class Solution {
public:
    bool canVisitAllRooms(vector<vector<int>>& rooms) {
        int n = rooms.size();
        vector<char> opened(n, 0);
        vector<int> pile(1, 0);                       // keys for rooms not yet entered
        opened[0] = 1;
        int count = 1;                                // rooms opened so far
        while (!pile.empty()) {
            int r = pile.back(); pile.pop_back();     // enter room r
            for (int k : rooms[r])                    // edge r -> k
                if (!opened[k]) { opened[k] = 1; count++; pile.push_back(k); }
        }
        return count == n;
    }
};
\end{lstlisting}

\begin{complexitybox}
\textbf{Time:} $\bigO{n + \sum|\texttt{rooms}[i]|}$ --- every room is entered once, every key is looked at once. \textbf{Space:} $\bigO{n}$ (the input already is the adjacency list). \textbf{Verdict:} \ac{}.
\end{complexitybox}
\end{algorithmbox}

\begin{invariantbox}[Why it is correct]
\textbf{Invariant.} At any moment, \texttt{opened} $=$ the rooms entered so far $\cup$ the rooms whose key is on the pile, and every opened room can really be entered.

\emph{Soundness.} Room 0 can be entered. A room $k$ is marked only when its key is found in a room $r$ we actually entered, so $0 \leadsto r \to k$.

\emph{Completeness.} When the pile is empty, every opened room has been entered and all its keys were examined, so no key for an unopened room lies in an opened room --- the opened set has no edge leaving it. Any room reachable from 0 is therefore opened (the pattern's closure argument). Since \texttt{count} is the number of opened rooms, \texttt{count == n} holds exactly when $R(0)$ contains every room. \hfill$\blacksquare$
\end{invariantbox}

\begin{alternativebox}[Approach 3: BFS or recursive DFS --- same set, different order]
Replacing the pile by a queue visits the rooms in order of ``how many doors away'' they are; on Sample~2 the BFS order is $0, 1, 3$ and the DFS order is $0, 3, 1$ --- the same \emph{set} $\{0,1,3\}$, so the same answer.

\begin{center}
\begin{tikzpicture}
    \foreach \x/\name/\o in {0/BFS (queue)/{0, 1, 3}, 5.5/DFS (stack)/{0, 3, 1}}{
        \begin{scope}[xshift=\x cm]
            \node[gvsrc] (h0) at (0,0.9) {0}; \node[gvdone] (h1) at (2.0,0.9) {1};
            \node[gvdone] (h3) at (1.0,-0.7) {3}; \node[gvdead] (h2) at (3.2,0) {2};
            \draw[g1p2tree] (h0) -- (h1); \draw[g1p2tree] (h0) -- (h3);
            \draw[gfaded, -{Stealth[length=2.2mm]}] (h1) -- (h3);
            \node[font=\small\bfseries] at (1.6,1.7) {\name: order \o};
        \end{scope}
    }
\end{tikzpicture}
\end{center}
Here even the search trees coincide (both consist of the two edges out of room 0); only the visiting order differs. The recursive version is the shortest to write and safe here, because the recursion depth is at most $n \le 1000$:

\begin{lstlisting}[style=cpp]
#include <bits/stdc++.h>
using namespace std;

class Solution {
    vector<char> opened;
    int count = 0;
    void enter(int r, vector<vector<int>>& rooms) {
        opened[r] = 1; count++;
        for (int k : rooms[r])
            if (!opened[k]) enter(k, rooms);
    }
public:
    bool canVisitAllRooms(vector<vector<int>>& rooms) {
        opened.assign(rooms.size(), 0);
        count = 0;
        enter(0, rooms);
        return count == (int)rooms.size();
    }
};
\end{lstlisting}
\begin{complexitybox}
\textbf{Time:} $\bigO{n + \sum|\texttt{rooms}[i]|}$. \textbf{Space:} $\bigO{n}$ plus a recursion depth $\le n$. \textbf{Verdict:} \ac{} --- but with $n = 10^6$ only the iterative versions would be safe.
\end{complexitybox}
\end{alternativebox}

\begin{tricksbox}[Approach 4: Implementation Refinements]
\begin{itemize}[leftmargin=*]
    \item \textbf{No graph building.} When the input is already ``for each vertex, its out-neighbours'', use it directly as \texttt{adj}. Recognising that saves both code and memory.
    \item \textbf{Early success.} One may stop as soon as \texttt{count == n}; it does not change the complexity but helps when most keys are in the first rooms.
    \item \textbf{Report the locked rooms.} After the search, every $r$ with \texttt{opened[r] == 0} is a room that can never be entered --- useful for debugging and for variants that ask ``which rooms?''.
\end{itemize}
\end{tricksbox}

\subsubsection*{5. Trace}

\begin{tracebox}[{Sample 2: rooms $= [[1,3],[3,0,1],[2],[0]]$}]
\begin{center}
\begin{tikzpicture}[scale=0.72, transform shape]
    \foreach \k/\t/\a/\b/\c/\d in {
        0/{start: pile $[0]$}/gvsrc/gv/gv/gv,
        1/{enter 0: pile $[1,3]$}/gvdone/gvseen/gv/gvseen,
        2/{enter 3: pile $[1]$}/gvdone/gvseen/gv/gvdone,
        3/{enter 1: pile empty}/gvdone/gvdone/gvdead/gvdone}{
        \begin{scope}[xshift=\k*5.4cm]
            \node[\a] (z0) at (0,0.9) {0}; \node[\b] (z1) at (2.0,0.9) {1};
            \node[\d] (z3) at (1.0,-0.7) {3}; \node[\c] (z2) at (3.3,0) {2};
            \draw[gde] (z1) to[bend left=15] (z0);
            \draw[gde] (z3) to[bend left=15] (z0);
            \draw[gde] (z1) -- (z3);
            \draw[gde] (z1) to[out=60, in=120, looseness=6] (z1);
            \draw[gde] (z2) to[out=60, in=120, looseness=6] (z2);
            \ifnum\k>0 \draw[g1p2tree] (z0) to[bend left=15] (z1); \draw[g1p2tree] (z0) to[bend left=15] (z3);
            \else \draw[gde] (z0) to[bend left=15] (z1); \draw[gde] (z0) to[bend left=15] (z3); \fi
            \node[font=\large\bfseries] at (1.6,2.2) {\t};
        \end{scope}
    }
\end{tikzpicture}
\end{center}
\begin{center}\small
\begin{tabular}{@{}cllcl@{}}
\toprule
\textbf{step} & \textbf{enter} & \textbf{keys found $\to$ new rooms} & \textbf{count} & \textbf{pile after (top on the right)} \\
\midrule
1 & 0 & $1, 3 \to$ 1, 3 & 3 & 1 3 \\
2 & 3 & $0 \to$ --- & 3 & 1 \\
3 & 1 & $3, 0, 1 \to$ --- & 3 & (empty) \\
\bottomrule
\end{tabular}
\end{center}
$\texttt{count} = 3 \ne 4 \Rightarrow$ \texttt{false}. Room 1's self-key (1) and room 3's key 0 are simply skipped by the \texttt{opened} test. Room 2 is never entered, so its self-loop is never even looked at.
\end{tracebox}

\begin{tracebox}[{Sample 1: rooms $= [[1],[2],[3],[\,]]$}]
\begin{center}
\begin{tikzpicture}[scale=0.8, transform shape]
    \foreach \k/\t/\s in {0/{enter 0, pile $[1]$}/1, 1/{enter 1, pile $[2]$}/2, 2/{enter 2, pile $[3]$}/3, 3/{enter 3, pile empty}/4}{
        \begin{scope}[xshift=\k*4.3cm]
            \foreach \i in {0,...,3}{
                \pgfmathtruncatemacro{\cur}{\s-1}
                \ifnum\i<\cur \node[gvdone] (w\i) at (\i*0.95,0) {\i};
                \else\ifnum\i=\cur \node[gvcur] (w\i) at (\i*0.95,0) {\i};
                \else\ifnum\i=\s \node[gvseen] (w\i) at (\i*0.95,0) {\i};
                \else \node[gv] (w\i) at (\i*0.95,0) {\i};\fi\fi\fi
            }
            \foreach \i [evaluate=\i as \j using int(\i+1)] in {0,...,2}{
                \ifnum\i<\s \draw[g1p2tree] (w\i) -- (w\j); \else \draw[gde] (w\i) -- (w\j); \fi}
            \node[font=\small\bfseries] at (1.45,0.75) {\t};
        \end{scope}
    }
\end{tikzpicture}
\end{center}
Every room yields exactly one new key; the pile never holds more than one key. After room 3, $\texttt{count} = 4 = n \Rightarrow$ \texttt{true}. The search tree is the whole chain.
\end{tracebox}

\subsubsection*{6. Test suite}

\begin{testbox}
\centering\small
\begin{tabular}{@{}p{6.4cm}cp{5.2cm}@{}}
\toprule
\textbf{rooms} & \textbf{Expected} & \textbf{Why} \\
\midrule
$[[1],[\,]]$ & \texttt{true} & smallest \texttt{true} \\
$[[\,],[0]]$ & \texttt{false} & key exists, but in the wrong room (direction) \\
$[[0],[0]]$ & \texttt{false} & self-key in room 0; room 1 unreachable \\
chain $i \to i+1$, $n = 1000$ & \texttt{true} & longest possible search path \\
room 0 holds keys $1..999$ & \texttt{true} & one room with 999 out-edges \\
chain $0 \to \dots \to 998 \to 0$, room 999 holds key 0 & \texttt{false} & the last room's only edge goes \emph{out} of it \\
\bottomrule
\end{tabular}
\end{testbox}
The expected values come from running Approach~2; in addition, Approaches 1--3 agreed with a Floyd--Warshall closure on 30\,000 random instances with $2 \le n \le 8$.

\subsubsection*{7. Common mistakes}

\begin{warningbox}[Mistake 1 --- edges in the wrong direction]
Building \texttt{adj[k].push\_back(r)} (``room $k$ is opened \emph{by} room $r$'') answers the reverse question ``which rooms lead to room 0?''. On the last test above it would mark every room and wrongly return \texttt{true}.
\end{warningbox}

\begin{warningbox}[Mistake 2 --- counting keys instead of rooms]
\begin{lstlisting}[style=snippet]
for (int k : rooms[r]) { count++; ... }         // WRONG: every key counted, even old ones
for (int k : rooms[r]) if (!opened[k]) { opened[k] = 1; count++; ... }   // RIGHT
\end{lstlisting}
On $[[1],[0]]$ the wrong count is $1 + 1 + 1 = 3 \ne 2$, so it answers \texttt{false} although both rooms open; on Sample~2 it reaches 7 (keys 3, 0 and 1 are counted twice). Only \emph{newly opened} rooms may be counted.
\end{warningbox}

\begin{warningbox}[Mistake 3 --- forgetting that room 0 is already open]
Starting with \texttt{count = 0} and \texttt{opened[0] = 0} makes room 0 count only if some room contains key 0. Room 0 must be marked and counted before the loop.
\end{warningbox}

\begin{warningbox}[Mistake 4 --- ``every room has a key somewhere'']
Checking that each room $1..n{-}1$ appears in some key list is a Pattern~1 (in-degree $>0$) test. Sample~2 passes it and is still \texttt{false}: in-degree from \emph{reachable} rooms is what matters, and that needs a search.
\end{warningbox}

\subsubsection*{8. Takeaways}

\begin{itemize}[leftmargin=*]
    \item ``Finding $X$ in $A$ lets you use $B$'' is an edge $A \to B$; ``can we get all of them?'' is ``$|R(\text{start})| = n$?''.
    \item Monotone processes (keys never vanish, infections never heal) are pure reachability: order does not matter.
    \item Sometimes the input already is the adjacency list --- recognise it and skip the building step.
    \item Count \emph{newly marked} vertices, not edges.
\end{itemize}

\begin{variantsbox}
\variant{Fewest extra unlocked rooms}{\varin{https://leetcode.com/problems/minimum-number-of-vertices-to-reach-all-nodes/}}
  {choose as few starting rooms as possible so that every room becomes reachable.}
  {same graph; we want the smallest set $S$ with $\bigcup_{s\in S} R(s) = V$.}
  {if the graph is a DAG: exactly the vertices of in-degree 0 (E2.2); in general one per ``source'' strongly connected component (Chapter~19).}
\variant{How many rooms from every start?}{\varin{https://atcoder.jp/contests/abc204/tasks/abc204_c}}
  {for every room as the start, count the rooms reachable from it (summed over all starts).}
  {the same directed graph with $|V| \le 2000$, $|E| \le 2000$.}
  {one search per start vertex, $\bigO{|V|(|V|+|E|)}$ (E2.3).}
\variant{Several rooms open initially}{\varnone}
  {rooms in a set $S$ are unlocked at the start, not only room 0.}
  {same graph, many sources.}
  {multi-source search: push every room of $S$ before the loop; still $\VE$.}
\variant{Which rooms stay locked?}{\varnone}
  {output the rooms that can never be entered.}
  {same graph.}
  {the same search; output every $r$ with \texttt{opened[r] == 0} (the complement of $R(0)$).}
\variant{Keys break after one use}{\varnone}
  {each key opens one door and is then lost; rooms may need to be re-entered.}
  {the state is no longer ``which room'' but ``which keys are in hand'' --- an implicit state graph.}
  {search over states (Pattern~4); the order of moves now matters, and the state space can be exponential.}
\end{variantsbox}

%============================================================
% CH1 · PATTERN 2 — EXERCISES, HINTS, ANSWERS
%============================================================
\subsection{Pattern 2: Exercises and Going Further}
\label{ssec:ch1-p2-ex}

\begin{patternbox}{Pattern 2 Summary Card}
\begin{tabularx}{\linewidth}{@{}>{\bfseries}l X@{}}
Question & which vertices are reachable from $s$? --- the reach set $R(s)$ \\
Algorithm & mark $s$; repeatedly take a vertex from a stack/queue and mark + add its unmarked out-neighbours; $\VE$ \\
Proof & marked $\Rightarrow$ reachable (induction); at the end no edge leaves the marked set $\Rightarrow$ all reachable are marked \\
Key variations & yes/no with early exit, ``all reached?'' (count $= n$), size of $R(s)$, many sources, reverse graph, one search per vertex \\
Top traps & undirected edge stored once; wrong edge direction; deep recursion; marking on pop; counting edges instead of vertices \\
\end{tabularx}
\end{patternbox}

\exlegend
\begin{enumerate}[leftmargin=*, label=\textbf{E2.\arabic*}, itemsep=6pt]
    \item \exeasy\ \pconcept \cf{1000}\ \exprob{New Year Transportation}{Codeforces 500A}{https://codeforces.com/problemset/problem/500/A}. Cells $1..n$ ($n \le 3\cdot10^4$); from cell $i < n$ a one-way portal leads to cell $i + a_i$ ($1 \le a_i \le n-i$). Starting in cell 1, can you reach cell $t$? Print \texttt{YES} or \texttt{NO}.
    \item \exeasy\ \ptricky \exprob{Minimum Number of Vertices to Reach All Nodes}{LeetCode 1557}{https://leetcode.com/problems/minimum-number-of-vertices-to-reach-all-nodes/}\ \lc{Medium}. Given a DAG with $n \le 10^5$ vertices and up to $10^5$ edges, return the smallest set of vertices from which every vertex is reachable (it is unique).
    \item \exmedium\ \pconcept \exprob{Tour}{AtCoder ABC204 C}{https://atcoder.jp/contests/abc204/tasks/abc204_c}. A directed graph with $N \le 2000$ cities and $M \le 2000$ one-way roads. Count the ordered pairs (origin, destination) such that the destination can be reached from the origin using zero or more roads.
    \item \exmedium\ \pclassic \exprob{Reorder Routes to Make All Paths Lead to the City Zero}{LeetCode 1466}{https://leetcode.com/problems/reorder-routes-to-make-all-paths-lead-to-the-city-zero/}\ \lc{Medium}. $n \le 5\cdot10^4$ cities joined by $n-1$ one-way roads whose undirected version is a tree. Reverse as few roads as possible so that every city can reach city 0; return that number.
    \item \exmedium\ \ptricky \exprob{Detonate the Maximum Bombs}{LeetCode 2101}{https://leetcode.com/problems/detonate-the-maximum-bombs/}\ \lc{Medium}. $n \le 100$ bombs $[x_i, y_i, r_i]$ ($1 \le x_i, y_i, r_i \le 10^5$). A detonated bomb detonates every bomb within Euclidean distance $r_i$ of its centre, and so on. Detonating exactly one bomb by hand, what is the maximum number of bombs that explode?
    \item \exmedium\ \pconcept \exown{Vertices on some $s$--$t$ walk.} Given a directed graph ($|V|, |E| \le 2\cdot 10^5$) and two vertices $s, t$, count the vertices $v$ such that some walk from $s$ to $t$ passes through $v$. On the running example $G$ of the theory section with $s = 5$, $t = 2$ the answer is $3$.
    \item \exmedium\ \precent \exprob{Properties Graph}{LeetCode 3493}{https://leetcode.com/problems/properties-graph/}\ \lc{Medium}. $n \le 100$ arrays $\texttt{properties}[i]$ of length $m \le 100$ with values in $1..100$, and an integer $k$. Vertices $i \ne j$ are joined by an undirected edge iff the two arrays have at least $k$ \emph{distinct} common values. Return the number of connected components.
\end{enumerate}

\subsubsection*{Hints}
\begin{enumerate}[leftmargin=*, label=\textbf{E2.\arabic*}]
    \item Every cell has exactly one outgoing portal, and portals only go to the right. What does a ``search'' look like when there is never a choice?
    \item In a DAG, which vertices can no other vertex reach? Must they be in the set? Is that already enough?
    \item $N + M$ is small. How expensive is one search? How many do you need?
    \item Search from city 0 over the roads \emph{ignoring} their direction. Every road you walk over is either already pointing towards 0 or not.
    \item The ``detonates'' relation is not symmetric: a big bomb can reach a small one but not the other way round. Build a directed graph, then search from every bomb.
    \item $v$ is on some $s$--$t$ walk iff $s \leadsto v$ and $v \leadsto t$. One search answers the first condition for all $v$; what answers the second?
    \item First build the edges (careful: ``distinct'' common values, so $[1,1]$ and $[1,1]$ share only one). Then: how many times does ``search from every vertex that is still unmarked'' start a new search?
\end{enumerate}

\subsubsection*{Answers}

\begin{solutionbox}[E2.1 New Year Transportation]
\textbf{Model:} vertex $=$ cell, edge $i \to i + a_i$; out-degree exactly 1 for $i < n$, so $R(1)$ is a single chain $1 \to 1 + a_1 \to \cdots$ (a \emph{functional graph}, Pattern~5). Positions strictly increase, so the walk ends after at most $n - 1$ jumps; $t$ is reachable iff the walk lands exactly on it.
\begin{lstlisting}[style=cpp]
#include <bits/stdc++.h>
using namespace std;

int main() {
    int n, t;
    scanf("%d %d", &n, &t);
    vector<int> a(n + 1, 0);
    for (int i = 1; i <= n - 1; i++) scanf("%d", &a[i]);
    int cur = 1;                          // every cell has exactly one portal: just follow it
    while (cur < t) cur += a[cur];        // positions only increase, so at most n-1 jumps
    puts(cur == t ? "YES" : "NO");
}
\end{lstlisting}
Samples (checked): $n=8, t=4$, $a = 1\,2\,1\,2\,1\,2\,1$ visits $1, 2, 4$ $\Rightarrow$ \texttt{YES}; $n=8, t=5$, $a = 1\,2\,1\,2\,1\,1\,1$ visits $1, 2, 4, 6$ and jumps over 5 $\Rightarrow$ \texttt{NO}. $\bigO{n}$.
\end{solutionbox}

\begin{solutionbox}[E2.2 Minimum Number of Vertices to Reach All Nodes]
A vertex with in-degree 0 can be reached only from itself, so it \emph{must} be in the set. Conversely, in a DAG every vertex $v$ with in-degree $> 0$ is reachable from some in-degree-0 vertex: walk backwards along incoming edges; without cycles the backward walk cannot repeat a vertex, so it stops --- at a vertex with in-degree 0. Hence the answer is exactly the set of in-degree-0 vertices: a reachability question solved by \textbf{Pattern~1}, without any search.
\begin{lstlisting}[style=cpp]
#include <bits/stdc++.h>
using namespace std;

class Solution {
public:
    vector<int> findSmallestSetOfVertices(int n, vector<vector<int>>& edges) {
        vector<char> hasIn(n, 0);
        for (auto& e : edges) hasIn[e[1]] = 1;
        vector<int> res;
        for (int v = 0; v < n; v++)
            if (!hasIn[v]) res.push_back(v);
        return res;
    }
};
\end{lstlisting}
Samples: $[0,3]$ and $[0,2,3]$ (checked). A brute force over all subsets agreed on 3000 random DAGs. In a graph \emph{with} cycles the rule breaks (a cycle has no in-degree-0 vertex); there one needs one vertex per source strongly connected component (Chapter~19).
\end{solutionbox}

\begin{solutionbox}[E2.3 Tour]
One search per origin: $N$ searches of $\bigO{N + M}$ each, i.e.\ $2000 \cdot 4000 = 8\cdot 10^6$ steps. The answer is $\sum_s |R(s)|$, and $(s,s)$ is counted because a journey may use zero roads.
\begin{lstlisting}[style=cpp]
#include <bits/stdc++.h>
using namespace std;

int main() {
    int n, m;
    scanf("%d %d", &n, &m);
    vector<vector<int>> adj(n);
    for (int i = 0; i < m; i++) {
        int a, b;
        scanf("%d %d", &a, &b);
        adj[a - 1].push_back(b - 1);                 // one-way road, 1-indexed input
    }
    long long answer = 0;
    vector<char> seen(n);
    vector<int> st;
    for (int s = 0; s < n; s++) {                    // one search per origin
        fill(seen.begin(), seen.end(), 0);
        seen[s] = 1; st.assign(1, s);
        int cnt = 1;                                 // (s, s) counts: zero roads
        while (!st.empty()) {
            int u = st.back(); st.pop_back();
            for (int v : adj[u])
                if (!seen[v]) { seen[v] = 1; cnt++; st.push_back(v); }
        }
        answer += cnt;
    }
    cout << answer << "\n";
}
\end{lstlisting}
Samples (checked): $7$, $3$, $16$. In Sample~1 (roads $1{\to}2$, $2{\to}3$, $3{\to}2$) the reach sets have sizes $3, 2, 2$. The theory section's \texttt{reachable pairs = 30} is the same computation on $G$.
\end{solutionbox}

\begin{solutionbox}[E2.4 Reorder Routes to Make All Paths Lead to the City Zero]
Store each road $a \to b$ twice: as $a \to b$ with cost 1 (it points \emph{away} from where we came from) and as $b \to a$ with cost 0. Search from city 0 over this undirected tree. When the search crosses a road from $u$ to a new city $v$, the road must end up pointing $v \to u$ (towards 0); it costs 1 exactly if it was stored as $u \to v$. Every road is crossed exactly once (the graph is a tree), so the sum of costs is the answer --- and no smaller answer exists, because every crossed road \emph{must} point towards 0 for $v$ to reach 0 (the only route from $v$ to 0 in a tree is the tree path).
\begin{lstlisting}[style=cpp]
#include <bits/stdc++.h>
using namespace std;

class Solution {
public:
    int minReorder(int n, vector<vector<int>>& connections) {
        vector<vector<pair<int,int>>> adj(n);        // (neighbour, 1 if the road points away)
        for (auto& c : connections) {
            adj[c[0]].push_back({c[1], 1});
            adj[c[1]].push_back({c[0], 0});
        }
        vector<char> seen(n, 0);
        vector<int> st(1, 0);
        seen[0] = 1;
        int flips = 0;
        while (!st.empty()) {
            int u = st.back(); st.pop_back();
            for (auto& p : adj[u])
                if (!seen[p.first]) {
                    seen[p.first] = 1;
                    flips += p.second;                // road u -> v must be reversed
                    st.push_back(p.first);
                }
        }
        return flips;
    }
};
\end{lstlisting}
\begin{center}
\begin{tikzpicture}[x=1.3cm]
    \node[gvsrc] (c0) at (0,0) {0}; \node[gv] (c1) at (1,0.7) {1}; \node[gv] (c3) at (2,0.7) {3};
    \node[gv] (c2) at (3,0.7) {2}; \node[gv] (c4) at (1,-0.7) {4}; \node[gv] (c5) at (2,-0.7) {5};
    \draw[gpath, -{Stealth[length=3mm]}] (c0) -- (c1); \draw[gpath, -{Stealth[length=3mm]}] (c1) -- (c3);
    \draw[gde] (c2) -- (c3); \draw[gde] (c4) -- (c0); \draw[gpath, -{Stealth[length=3mm]}] (c4) -- (c5);
    \node[note, anchor=west] at (3.5,0) {orange roads point away from 0:\\$0{\to}1$, $1{\to}3$, $4{\to}5$ $\Rightarrow$ \textbf{3}};
\end{tikzpicture}
\end{center}
Samples (checked): $3$, $2$, $0$; a brute force over all $2^{n-1}$ orientations agreed on 3000 random trees. This is the ``search in the right direction'' idea of the theory section turned into a count.
\end{solutionbox}

\begin{solutionbox}[E2.5 Detonate the Maximum Bombs]
\textbf{Model:} vertex $=$ bomb; edge $i \to j$ iff $(x_i - x_j)^2 + (y_i - y_j)^2 \le r_i^2$ --- directed, because the radius belongs to $i$. The answer is $\max_s |R(s)|$: $n$ searches on a graph with at most $n^2 = 10^4$ edges, $\bigO{n^3} = 10^6$.
\begin{lstlisting}[style=cpp]
#include <bits/stdc++.h>
using namespace std;

class Solution {
public:
    int maximumDetonation(vector<vector<int>>& b) {
        int n = b.size();
        vector<vector<int>> adj(n);
        for (int i = 0; i < n; i++)
            for (int j = 0; j < n; j++) {
                if (i == j) continue;
                long long dx = b[i][0] - b[j][0], dy = b[i][1] - b[j][1], r = b[i][2];
                if (dx * dx + dy * dy <= r * r) adj[i].push_back(j);   // 64-bit: up to 2e10
            }
        int best = 0;
        for (int s = 0; s < n; s++) {
            vector<char> seen(n, 0);
            vector<int> st(1, s);
            seen[s] = 1;
            int cnt = 1;
            while (!st.empty()) {
                int u = st.back(); st.pop_back();
                for (int v : adj[u])
                    if (!seen[v]) { seen[v] = 1; cnt++; st.push_back(v); }
            }
            best = max(best, cnt);
        }
        return best;
    }
};
\end{lstlisting}
Sample~3, $[[1,2,3],[2,3,1],[3,4,2],[4,5,3],[5,6,4]]$, builds the edges $0{\to}1, 0{\to}2$; $2{\to}1, 2{\to}3$; $3{\to}1, 3{\to}2, 3{\to}4$; $4{\to}2, 4{\to}3$ (bomb 1 reaches nobody), and $R(0)$ is everything: answer $5$.
\begin{center}
\begin{tikzpicture}[x=1.25cm]
    \node[gvsrc] (d0) at (0,0) {0}; \node[gvdone] (d1) at (1.2,1) {1}; \node[gvdone] (d2) at (1.2,-0.6) {2};
    \node[gvdone] (d3) at (2.6,0.2) {3}; \node[gvdone] (d4) at (4,0.2) {4};
    \draw[gtree, -{Stealth[length=2.6mm]}] (d0) -- (d1); \draw[gtree, -{Stealth[length=2.6mm]}] (d0) -- (d2);
    \draw[gde] (d2) -- (d1); \draw[gtree, -{Stealth[length=2.6mm]}] (d2) to[bend left=15] (d3);
    \draw[gde] (d3) to[bend left=15] (d2); \draw[gde] (d3) -- (d1);
    \draw[gtree, -{Stealth[length=2.6mm]}] (d3) to[bend left=15] (d4); \draw[gde] (d4) to[bend left=15] (d3);
    \draw[gde] (d4) to[bend left=35] (d2);
\end{tikzpicture}
\end{center}
Samples (checked): $2$, $1$, $5$; a closure-based brute force agreed on 3000 random inputs. \textbf{Traps:} building an \emph{undirected} graph (on $[[1,1,5],[5,1,1],[9,1,5]]$ bombs 0 and 2 both reach the small bomb 1 but not each other: the right answer is $2$, while the undirected graph puts all three into one component and claims $3$) and computing $dx^2 + dy^2$ in \texttt{int} ($2\cdot 10^{10}$ overflows).
\end{solutionbox}

\begin{solutionbox}[E2.6 Vertices on some $s$--$t$ walk]
$v$ lies on some walk $s \leadsto v \leadsto t$ iff $v \in R(s)$ and $v \in R^{R}(t)$ (reachable from $t$ in the reverse graph). Two searches, $\VE$:
\begin{lstlisting}[style=cpp]
#include <bits/stdc++.h>
using namespace std;

vector<char> reach(const vector<vector<int>>& g, int s) {
    vector<char> seen(g.size(), 0);
    vector<int> st(1, s);
    seen[s] = 1;
    while (!st.empty()) {
        int u = st.back(); st.pop_back();
        for (int v : g[u]) if (!seen[v]) { seen[v] = 1; st.push_back(v); }
    }
    return seen;
}

int main() {
    int n = 8, s = 5, t = 2;
    int E[][2] = {{0,1},{0,2},{1,3},{2,3},{3,4},{4,1},{5,0},{5,6},{6,7},{7,6}};
    vector<vector<int>> adj(n), rev(n);
    for (auto& e : E) { adj[e[0]].push_back(e[1]); rev[e[1]].push_back(e[0]); }
    vector<char> F = reach(adj, s), B = reach(rev, t);
    int cnt = 0;
    for (int v = 0; v < n; v++)
        if (F[v] && B[v]) { cout << v << " "; cnt++; }
    cout << "-> " << cnt << "\n";                      // 0 2 5 -> 3
}
\end{lstlisting}
Forward from 5 everything is reached; backward from 2 only $\{2, 0, 5\}$. Vertices 1, 3, 4, 6, 7 are reachable from 5 but cannot get back to 2. With $t = 3$ instead, the program finds $\{0,1,2,3,4,5\}$ (6 vertices). A brute force using the transitive closure agreed on 3000 random graphs.
\end{solutionbox}

\begin{solutionbox}[E2.7 Properties Graph]
\textbf{Model:} vertex $=$ array; undirected edge $i - j$ iff $|\mathrm{set}(i) \cap \mathrm{set}(j)| \ge k$. Store every array as a 101-bit \texttt{bitset} of the values it contains --- duplicates collapse by themselves --- so one intersection is \texttt{(has[i] \& has[j]).count()}. \textbf{Counting components with reachability:} scan the vertices; every time a vertex is still unmarked, a new component starts --- count it and mark its whole reach set with one search. Each vertex is marked once, so all searches together cost $\VE$; building the graph costs $\bigO{n^2 \cdot 100/64}$ plus $\bigO{n\,m}$ to read the arrays.
\begin{lstlisting}[style=cpp]
#include <bits/stdc++.h>
using namespace std;

class Solution {
public:
    int numberOfComponents(vector<vector<int>>& properties, int k) {
        int n = properties.size();
        vector<bitset<101>> has(n);                   // has[i][x] = 1 iff x occurs in row i
        for (int i = 0; i < n; i++)
            for (int x : properties[i]) has[i][x] = 1;    // duplicates collapse automatically
        vector<vector<int>> adj(n);
        for (int i = 0; i < n; i++)
            for (int j = i + 1; j < n; j++)
                if ((int)(has[i] & has[j]).count() >= k) {   // distinct common values
                    adj[i].push_back(j);
                    adj[j].push_back(i);
                }
        vector<char> seen(n, 0);
        int components = 0;
        for (int s = 0; s < n; s++) {                 // a new search per unseen vertex
            if (seen[s]) continue;
            components++;
            vector<int> st(1, s);
            seen[s] = 1;
            while (!st.empty()) {
                int u = st.back(); st.pop_back();
                for (int v : adj[u])
                    if (!seen[v]) { seen[v] = 1; st.push_back(v); }
            }
        }
        return components;
    }
};
\end{lstlisting}
Example~1 ($k = 1$): $[1,2]$--$[1,1]$ share 1; $[3,4]$--$[4,5]$--$[5,6]$ form a chain; $[7,7]$ is alone $\Rightarrow$ components $\{0,1\}, \{2,3,4\}, \{5\}$, answer $3$. Example~3 is the trap: $[1,1]$ and $[1,1]$ have one \emph{distinct} common value, less than $k = 2$, so the answer is $2$ (counting matching positions would wrongly join them). The three official examples give $3, 1, 2$ (checked), and a DSU brute force with \texttt{std::set} intersections agreed on 5000 random inputs. Counting components is studied properly in Chapter~4.
\end{solutionbox}


%============================================================
\clearpage
\section{Pattern 3: Grids as Graphs}
\label{sec:ch1-p3}
%============================================================
%============================================================
% CH1 · PATTERN 3 — GRIDS AS GRAPHS: THEORY
%============================================================
% Local helpers of Pattern 3 (TikZ styles g1p3..., macros \gIpIII...):
%   \gIpIIIgrid{row,row,...}  draws a grid inside a tikzpicture; one character per cell:
%       o open   w wall/water   A..F component colours with the letter, a..f same colours without it
%       s start cell   x current cell   q on the stack / in the queue   k finished
%   cell (r,c) is centred at (\gIpIIIs*c, -\gIpIIIs*r);  \gIpIIIat{r}{c}{text} writes into a cell
\tikzset{
    g1p3o/.style={gcell},
    g1p3w/.style={gwall},
    g1p3a/.style={gcell, fill=primaryblue!22},
    g1p3b/.style={gcell, fill=successgreen!25},
    g1p3c/.style={gcell, fill=accentorange!35},
    g1p3d/.style={gcell, fill=softpurple!25},
    g1p3e/.style={gcell, fill=recursionred!22},
    g1p3f/.style={gcell, fill=insightborder!30},
    g1p3A/.style={g1p3a, label={[font=\scriptsize\bfseries, text=darkgray]center:A}},
    g1p3B/.style={g1p3b, label={[font=\scriptsize\bfseries, text=darkgray]center:B}},
    g1p3C/.style={g1p3c, label={[font=\scriptsize\bfseries, text=darkgray]center:C}},
    g1p3D/.style={g1p3d, label={[font=\scriptsize\bfseries, text=darkgray]center:D}},
    g1p3E/.style={g1p3e, label={[font=\scriptsize\bfseries, text=darkgray]center:E}},
    g1p3F/.style={g1p3f, label={[font=\scriptsize\bfseries, text=darkgray]center:F}},
    g1p3s/.style={gcell, fill=basefill, draw=successgreen, line width=1.2pt},
    g1p3x/.style={gcell, fill=currentfill, draw=accentorange, line width=1.2pt},
    g1p3q/.style={gcell, fill=hitfill},
    g1p3k/.style={gcell, fill=computedfill},
    g1p3lbl/.style={font=\scriptsize\bfseries, text=darkgray, inner sep=0pt},
    g1p3var/.style={rectangle, rounded corners=2pt, draw=patternborder, fill=patternbg,
                    font=\scriptsize, align=center, text width=3.25cm, minimum height=1.25cm, inner sep=3pt}
}
\newcount\gIpIIIc
\def\gIpIIIs{0.7}
\def\gIpIIIscan#1{\ifx#1\relax\else
    \node[g1p3#1, minimum size=\gIpIIIs cm] at (\gIpIIIs*\the\gIpIIIc, -\gIpIIIs*\gIpIIIr) {};%
    \global\advance\gIpIIIc by 1 \expandafter\gIpIIIscan\fi}
\newcommand{\gIpIIIgrid}[1]{%
    \foreach \gIpIIIrowtxt [count=\gIpIIIr from 0] in {#1}{%
        \global\gIpIIIc=0 \expandafter\gIpIIIscan\gIpIIIrowtxt\relax}}
\newcommand{\gIpIIIat}[3]{\node[g1p3lbl] at (\gIpIIIs*#2, -\gIpIIIs*#1) {#3};}

\subsection{Pattern Theory}
\label{ssec:ch1-p3-theory}

\begin{patternbox}{Pattern 3 at a Glance}
A grid \emph{is} a graph: every open \textbf{cell is a vertex}, and two open cells that \textbf{share a side are joined by an edge}. We never store this graph --- the neighbours of $(r,c)$ are computed on the fly. A ``room'', ``island'', ``lake'' or ``region'' is a \textbf{connected component}, found by a \textbf{flood fill} (a DFS or BFS that spreads from one cell to its open neighbours).

\begin{center}
\begin{tikzpicture}
    \begin{scope}
        \def\gIpIIIs{0.75}
        \gIpIIIgrid{ooo,oxo,ooo}
        \draw[flow] (0.75,-0.75) -- (0.75,-0.1);
        \draw[flow] (0.75,-0.75) -- (0.75,-1.4);
        \draw[flow] (0.75,-0.75) -- (0.1,-0.75);
        \draw[flow] (0.75,-0.75) -- (1.4,-0.75);
        \node[font=\scriptsize] at (0.75,0.65) {$(r{-}1,c)$};
        \node[font=\scriptsize] at (0.75,-2.15) {$(r{+}1,c)$};
        \node[font=\scriptsize, anchor=east] at (-0.45,-0.75) {$(r,c{-}1)$};
        \node[font=\scriptsize, anchor=west] at (1.95,-0.75) {$(r,c{+}1)$};
    \end{scope}
    \begin{scope}[xshift=6.3cm, yshift=0.3cm]
        \def\gIpIIIs{0.6}
        \gIpIIIgrid{AAwwB,wAwBB,CwwwB,CCwDw}
        \node[font=\scriptsize, align=center] at (1.2,-2.75) {4 components: sizes 3, 4, 3, 1};
    \end{scope}
\end{tikzpicture}
\end{center}

\begin{recurrencebox}
\small
\textbf{for} every cell $(r,c)$ in row-major order: \textbf{if} it is open and unvisited \textbf{then}\\
\hspace*{1.5em}\texttt{components++}; mark $(r,c)$; push it on a stack (or queue)\\
\hspace*{1.5em}\textbf{while} the stack is not empty: pop $(x,y)$; \texttt{size++};\\
\hspace*{3em}\textbf{for} each of the 4 neighbours $(x{+}dr_k,\,y{+}dc_k)$ that is \emph{inside}, \emph{open} and \emph{unvisited}: mark it, push it.\\[2pt]
\hspace*{\fill}$\bigO{n\cdot m}$ time, $\bigO{n\cdot m}$ memory --- every cell is marked once and looks at 4 neighbours.
\end{recurrencebox}
\end{patternbox}

%------------------------------------------------------------
\subsubsection*{From Pattern 2 to Pattern 3}

In Pattern 2 the graph was \emph{handed to us}: $n$ vertices and a list of edges, stored in an adjacency list, and we asked ``can $s$ reach $t$?''. In this pattern the input is a \textbf{picture} --- a map of $n$ rows and $m$ columns of characters or numbers --- and nobody mentions vertices or edges. The graph is still there; we only have to \emph{see} it:

\begin{center}
\renewcommand{\arraystretch}{1.25}
\begin{tabular}{@{}ll@{}}
\toprule
\textbf{Picture word} & \textbf{Graph word} \\
\midrule
open cell (floor `\texttt{.}', land \texttt{1}, a pixel of colour $x$) & vertex \\
two open cells sharing a side & edge (undirected) \\
``you can walk from $A$ to $B$'' & $A$ and $B$ are in the same connected component \\
room / island / lake / region / blob & connected component \\
area of an island & number of vertices of a component \\
number of rooms & number of components \\
\bottomrule
\end{tabular}
\end{center}

\begin{diagrambox}
\begin{tikzpicture}
    \begin{scope}
        \gIpIIIgrid{oowo,owoo,oooo}
        \foreach \c in {0,1,2,3} \node[font=\scriptsize, text=darkgray] at (0.7*\c, 0.55) {$c{=}\c$};
        \foreach \r in {0,1,2} \node[font=\scriptsize, text=darkgray, anchor=east] at (-0.45, -0.7*\r) {$r{=}\r$};
        \node[font=\small\bfseries] at (1.05,-2.25) {the map ($3\times4$, two walls)};
    \end{scope}
    \draw[-{Stealth[length=3mm]}, line width=1.4pt, primaryblue] (3.0,-0.7) -- node[above, font=\scriptsize]{see it as} (4.4,-0.7);
    \begin{scope}[xshift=5.2cm]
        \foreach \r/\c in {0/0,0/1,0/3,1/0,1/2,1/3,2/0,2/1,2/2,2/3}
            \node[gv, minimum size=6.5mm, font=\scriptsize] (v\r\c) at (1.25*\c, -1.0*\r) {\r,\c};
        \foreach \a/\b in {00/01, 00/10, 03/13, 10/20, 12/13, 12/22, 13/23, 20/21, 21/22, 22/23}
            \draw[ge] (v\a) -- (v\b);
        \node[font=\small\bfseries] at (1.9,-2.75) {the graph: 10 vertices, 10 edges};
    \end{scope}
\end{tikzpicture}
\figcap{Each open cell becomes a vertex named by its coordinates; walls simply have no vertex. An edge joins every pair of open cells that share a side. Cells $(0,1)$ and $(1,0)$ touch only at a corner: no edge.}
\end{diagrambox}

\textbf{How big is this graph?} At most $n\cdot m$ vertices, and every vertex has at most 4 neighbours, so $|E| \le 4nm/2 = 2nm$ (a fully open grid has exactly $n(m-1) + m(n-1)$ edges). For the usual limits $n, m \le 1000$ that is $10^6$ vertices and $2\cdot10^6$ edges --- well within an $\bigO{|V|+|E|}$ algorithm, but \emph{building} an adjacency list for it would waste time and tens of megabytes on information we can recompute instantly. So we keep the grid itself as the graph: it is an \textbf{implicit graph}, exactly the kind described in Lesson~\ref{ssec:ch1-modelling}.

%------------------------------------------------------------
\subsubsection*{Neighbours with two small arrays}

Instead of four copies of the same \texttt{if}, store the four moves as \emph{direction vectors} and loop over them:

\begin{diagrambox}
\begin{tikzpicture}
    \begin{scope}
        \gIpIIIgrid{ooo,oxo,ooo}
        \foreach \x/\y/\t in {0.7/0/{$-1,0$}, 0.7/-1.4/{$+1,0$}, 0/-0.7/{$0,-1$}, 1.4/-0.7/{$0,+1$}}
            \node[font=\tiny\bfseries, text=deepblue] at (\x,\y) {\t};
        \node[font=\scriptsize] at (0.7,-0.7) {$r,c$};
        \node[font=\small\bfseries, align=center] at (0.7,-2.15) {4-neighbourhood};
        \node[font=\scriptsize, align=center] at (0.7,-2.8) {\texttt{dr = \{-1, 1, 0, 0\}}\\\texttt{dc = \{ 0, 0,-1, 1\}}};
    \end{scope}
    \begin{scope}[xshift=4.6cm]
        \gIpIIIgrid{qqq,qxq,qqq}
        \foreach \x/\y/\t in {0/0/{$-1,-1$}, 0.7/0/{$-1,0$}, 1.4/0/{$-1,+1$}, 0/-0.7/{$0,-1$}, 1.4/-0.7/{$0,+1$}, 0/-1.4/{$+1,-1$}, 0.7/-1.4/{$+1,0$}, 1.4/-1.4/{$+1,+1$}}
            \node[font=\tiny\bfseries, text=deepblue] at (\x,\y) {\t};
        \node[font=\scriptsize] at (0.7,-0.7) {$r,c$};
        \node[font=\small\bfseries, align=center] at (0.7,-2.15) {8-neighbourhood};
        \node[font=\scriptsize, align=center] at (0.7,-2.8) {add the 4 diagonals\\(``king moves'')};
    \end{scope}
    \begin{scope}[xshift=9.2cm]
        \gIpIIIgrid{ooooo,ooooo,ooxoo,ooooo,ooooo}
        \foreach \r/\c in {0/1,0/3,1/0,1/4,3/0,3/4,4/1,4/3} \node[g1p3q, minimum size=0.7cm] at (0.7*\c,-0.7*\r) {};
        \node[font=\scriptsize] at (1.4,-1.4) {$r,c$};
        \node[font=\small\bfseries, align=center] at (1.4,-3.55) {knight moves};
        \node[font=\scriptsize, align=center] at (1.4,-4.1) {$(\pm1,\pm2),\ (\pm2,\pm1)$};
    \end{scope}
\end{tikzpicture}
\figcap{The neighbourhood is part of the \emph{model}, not of the algorithm. Read the statement: ``share a side'' / ``up, down, left, right'' = 4 neighbours; ``touch'' / ``including diagonally'' = 8; a chess piece = its own move list. The flood fill code stays identical --- only the arrays change.}
\end{diagrambox}

The neighbour test has three parts, and their \textbf{order matters}: first \emph{inside the grid} (otherwise \texttt{g[nr][nc]} reads outside the array), then \emph{open}, then \emph{not yet visited}.

\subsubsection*{Naming a cell with one integer}

Pairs $(r,c)$ are readable, but sometimes one integer per vertex is handier (a \texttt{vector<int>} stack, a DSU, a \texttt{visited} array of size $nm$). Number the cells row by row:
\[
\text{id} = r\cdot m + c, \qquad r = \text{id} \;/\; m, \qquad c = \text{id} \bmod m \qquad (0 \le \text{id} < nm).
\]
\begin{diagrambox}
\begin{tikzpicture}
    \gIpIIIgrid{oooo,oooo,oooo}
    \foreach \r in {0,1,2} \foreach \c in {0,1,2,3} {
        \pgfmathtruncatemacro{\id}{\r*4+\c}
        \gIpIIIat{\r}{\c}{\id}
    }
    \foreach \c in {0,1,2,3} \node[font=\scriptsize, text=darkgray] at (0.7*\c, 0.55) {$c{=}\c$};
    \foreach \r in {0,1,2} \node[font=\scriptsize, text=darkgray, anchor=east] at (-0.45, -0.7*\r) {$r{=}\r$};
    \node[note, anchor=west, text width=7.2cm, align=left] at (3.0,-0.7) {$n=3$ rows, $m=4$ columns. Cell $(2,1)$ has id $2\cdot4+1 = 9$; id $6$ is $(6/4,\ 6 \bmod 4) = (1,2)$. Multiply by the number of \textbf{columns} $m$, not by $n$ --- with $n \ne m$ the wrong factor silently mixes cells up.};
\end{tikzpicture}
\end{diagrambox}

%------------------------------------------------------------
\subsubsection*{When to recognise it}

\begin{itemize}[leftmargin=*]
    \item the input is a \textbf{2D map}: characters (\texttt{.} / \texttt{\#}), 0/1 values, colours, heights, depths;
    \item movement is between \textbf{adjacent cells} (4 or 8 directions, or a fixed move set);
    \item the question is about \textbf{regions}: count them (rooms, islands), measure them (area, volume, perimeter), recolour one (flood fill), or decide whether two cells are in the same region.
\end{itemize}
\textbf{Not this pattern (yet):}
\begin{itemize}[leftmargin=*]
    \item ``\textbf{minimum number of steps}'' from one cell to another: the order of exploration matters, use BFS with distances --- Chapter 3 (BFS);
    \item cells with \textbf{different costs} to enter: Dijkstra or 0-1 BFS, later chapters;
    \item the grid \textbf{changes} over time (walls added or removed between queries): Chapter 6 (DSU);
    \item the cell alone is not enough to describe where you are (you carry keys, you have used $k$ wall-breaks, you face a direction): the vertex becomes (cell, extra information) --- a state graph, Pattern 4.
\end{itemize}

%------------------------------------------------------------
\subsubsection*{The algorithm: scan and fill}

\begin{enumerate}[leftmargin=*]
    \item Create \texttt{visited} (or reuse the grid: overwrite a visited open cell with a wall character).
    \item Scan the cells in row-major order. An open, unvisited cell is the \textbf{first cell of a new component}: count it and start a flood fill there.
    \item The flood fill marks the start cell and puts it on a stack. While the stack is not empty, pop a cell and, for each of its neighbours that is inside, open and unvisited, \textbf{mark it immediately} and push it.
    \item When the stack is empty, the whole component is marked; the scan continues and skips every marked cell.
\end{enumerate}

\begin{diagrambox}
\begin{tikzpicture}
    \begin{scope}
        \gIpIIIgrid{sowwo,wowoo,owwwo,oowow}
        \draw[flow, dashed] (-0.35,0.5) -- (3.15,0.5);
        \node[font=\scriptsize, text=primaryblue] at (1.4,0.75) {scan row by row};
        \node[font=\scriptsize\bfseries, align=center] at (1.4,-3.0) {scan reaches $(0,0)$:\\open, unvisited $\Rightarrow$ component A};
    \end{scope}
    \begin{scope}[xshift=4.6cm]
        \gIpIIIgrid{AAwws,wAwoo,owwwo,oowow}
        \node[font=\scriptsize\bfseries, align=center] at (1.4,-3.0) {A is filled (3 cells);\\next new start: $(0,4)$};
    \end{scope}
    \begin{scope}[xshift=9.2cm]
        \gIpIIIgrid{AAwwB,wAwBB,CwwwB,CCwDw}
        \foreach \r/\c in {0/0,0/4,2/0,3/3} \draw[successgreen, line width=1.3pt] (0.7*\c-0.35,-0.7*\r-0.35) rectangle (0.7*\c+0.35,-0.7*\r+0.35);
        \node[font=\scriptsize\bfseries, align=center] at (1.4,-3.0) {done: 4 starts (green frames)\\= 4 components, sizes 3, 4, 3, 1};
    \end{scope}
\end{tikzpicture}
\figcap{The scan finds the first cell of each component in row-major order; the fill then swallows the whole component, so the scan never starts twice inside it. The template below prints ``4 components, sizes: 3 4 3 1'' for this grid.}
\end{diagrambox}

\begin{invariantbox}[Why scan-and-fill counts components exactly]
\textbf{Claim 1 (a fill marks exactly one component).} A fill started at $s$ marks precisely the open cells reachable from $s$.
\emph{Only reachable cells:} a cell is marked only when it is a neighbour of a marked cell, so by induction every marked cell has a path from $s$.
\emph{Every reachable cell:} suppose a reachable cell $v$ stays unmarked; on a path $s = u_0, u_1, \dots, u_k = v$ take the first unmarked $u_{i+1}$. Its predecessor $u_i$ was marked, hence pushed, hence popped, and when it was popped all its open unvisited neighbours --- $u_{i+1}$ among them --- were marked. Contradiction.

\textbf{Claim 2 (one start per component).} A fill starts only at an unmarked cell; by Claim~1 that cell's component is still completely unmarked (fills mark whole components), and after the fill it is completely marked, so no later start lies in it. Every component gets a start because the scan visits all cells. Hence \#starts $=$ \#components.

\textbf{Cost.} Each cell is marked at most once and therefore pushed and popped at most once; each pop inspects 4 neighbours. Total $\bigO{nm}$ time and $\bigO{nm}$ memory for \texttt{visited} and the stack. \hfill$\blacksquare$
\end{invariantbox}

\subsubsection*{Stack or queue? Recursion?}

The fill works with \emph{any} container: the set of marked cells at the end is the same, only the \emph{order} differs.

\begin{diagrambox}
\begin{tikzpicture}
    \begin{scope}
        \gIpIIIgrid{kkkk,kskk,kkkk,kkkk}
        \foreach \r/\c/\t in {0/0/14,0/1/16,0/2/12,0/3/10, 1/0/13,1/1/1,1/2/2,1/3/3, 2/0/9,2/1/15,2/2/11,2/3/4, 3/0/8,3/1/7,3/2/6,3/3/5}
            \gIpIIIat{\r}{\c}{\t}
        \node[font=\small\bfseries, align=center] at (1.05,-3.2) {stack (DFS order)\\[-1pt]{\scriptsize\mdseries runs away along a long path}};
    \end{scope}
    \begin{scope}[xshift=4.6cm]
        \gIpIIIgrid{kkkk,kskk,kkkk,kkkk}
        \foreach \r/\c/\t in {0/0/6,0/1/2,0/2/7,0/3/12, 1/0/4,1/1/1,1/2/5,1/3/11, 2/0/9,2/1/3,2/2/10,2/3/15, 3/0/13,3/1/8,3/2/14,3/3/16}
            \gIpIIIat{\r}{\c}{\t}
        \node[font=\small\bfseries, align=center] at (1.05,-3.2) {queue (BFS order)\\[-1pt]{\scriptsize\mdseries grows in rings around the start}};
    \end{scope}
    \begin{scope}[xshift=9.2cm]
        \gIpIIIgrid{skkk,kkkk,kkkk,kkkk}
        \draw[gpath, opacity=0.5, -{Stealth[length=2.5mm]}] (0,0) -- (0,-2.1) -- (0.7,-2.1) -- (0.7,0) -- (1.4,0) -- (1.4,-2.1) -- (2.1,-2.1) -- (2.1,0);
        \foreach \r/\c/\t in {0/0/1,0/1/8,0/2/9,0/3/16, 1/0/2,1/1/7,1/2/10,1/3/15, 2/0/3,2/1/6,2/2/11,2/3/14, 3/0/4,3/1/5,3/2/12,3/3/13}
            \gIpIIIat{\r}{\c}{\t}
        \node[font=\small\bfseries, align=center] at (1.05,-3.2) {recursive DFS from $(0,0)$\\[-1pt]{\scriptsize\mdseries depth 16 = every cell!}};
    \end{scope}
\end{tikzpicture}
\figcap{Pop order on an open $4\times4$ grid (neighbour order up, down, left, right; green = start). Left and middle start at $(1,1)$. Right: the recursive DFS from $(0,0)$ snakes through all 16 cells before returning even once, so the recursion depth equals the number of cells.}
\end{diagrambox}

\begin{warningbox}[Recursion depth on grids]
The picture on the right is not a worst case, it is the \emph{normal} case: on an open $1000\times1000$ room the recursive fill reaches depth $10^6$ (we measured it: the program crashes with the default stack, while depth $10^4$ on a $100\times100$ grid is fine). Rule of thumb: recursion is safe for $nm$ up to about $10^5$ (LeetCode limits like $50\times50$ or $300\times300$ are fine); for $nm$ up to $10^6$ (CSES, Codeforces) use the \textbf{explicit stack or the queue}. The short recursive version is still worth knowing: it is the fastest to type when it is safe.
\end{warningbox}

%------------------------------------------------------------
\subsubsection*{Variations map}

\begin{diagrambox}
\begin{tikzpicture}
    \node[g1p3var, fill=answerfill, draw=accentorange, text width=4.6cm, font=\scriptsize\bfseries] (root) at (0,0)
        {Scan + flood fill\\cells = vertices, sides = edges};
    \node[g1p3var] (v1) at (-5.4,-2) {\textbf{count components}\\rooms, islands\\(Problem 1.3.1)};
    \node[g1p3var] (v2) at (-1.8,-2) {\textbf{size / sum per component}\\largest island, lake volume\\(Problem 1.3.2, E3.3)};
    \node[g1p3var] (v3) at (1.8,-2)  {\textbf{fill from one cell}\\recolour a region\\(Flood Fill, E3.2)};
    \node[g1p3var] (v4) at (5.4,-2)  {\textbf{boundary of a region}\\perimeter = sides facing\\water (E3.1)};
    \node[g1p3var] (v5) at (-5.4,-4.1) {\textbf{a property of a whole component}\\``every cell is also land in grid 1''\\(E3.4)};
    \node[g1p3var] (v6) at (-1.8,-4.1) {\textbf{label + sizes, then combine}\\flip one cell to join islands\\(E3.5)};
    \node[g1p3var] (v7) at (1.8,-4.1)  {\textbf{fill from the border}\\regions that cannot escape\\(Chapter 4)};
    \node[g1p3var] (v8) at (5.4,-4.1)  {\textbf{shortest path in a maze}\\BFS with distances\\(Chapter 3)};
    \node[draw=patternborder, dashed, rounded corners=4pt, fit=(v1)(v4)(v5)(v8), inner sep=5pt] (grp) {};
    \draw[flow] (root.south) -- (grp.north);
\end{tikzpicture}
\end{diagrambox}

\textbf{Twists and how to handle them}
\begin{itemize}[leftmargin=*]
    \item \textbf{8-connectivity} (``cells touching at a corner are connected''): use the 8-entry direction arrays; nothing else changes.
    \item \textbf{Same colour instead of open/blocked} (images, height maps): ``open'' becomes ``has the same value as the start cell''; store the start value before you overwrite anything.
    \item \textbf{Wrap-around} (a torus: leaving on the right enters on the left): replace the inside test by \texttt{nr = (r + dr[d] + n) \% n} and \texttt{nc = (c + dc[d] + m) \% m}.
    \item \textbf{Two grids at once} (sub-islands, comparing pictures): fill in one grid while checking the other.
    \item \textbf{Per-component information used later} (size of the island next to a water cell): first label every cell with its component id and store the sizes, then answer each question in $\bigO{1}$.
    \item \textbf{Many test cases with a total size limit}: allocate \texttt{visited} per test with the real $n\times m$, never a global $1000\times1000$ array cleared every time.
\end{itemize}

%------------------------------------------------------------
\subsubsection*{Code templates}

The toolkit below contains everything this pattern needs: the inside test, the direction arrays, the recursive fill (small grids only), the iterative fill with an explicit stack, the BFS fill, and the scan that labels all components.

\begin{lstlisting}[style=cpp]
#include <bits/stdc++.h>
using namespace std;

int n, m;                                    // rows, columns
vector<string> g;                            // '.' = open cell, '#' = blocked
const int dr[4] = {-1, 1, 0, 0};             // up, down, left, right
const int dc[4] = {0, 0, -1, 1};

bool inside(int r, int c) { return 0 <= r && r < n && 0 <= c && c < m; }
bool open(int r, int c) { return inside(r, c) && g[r][c] == '.'; }

// (1) recursive DFS: fine when the grid is small (depth can reach n*m!)
void dfs(int r, int c, int id, vector<vector<int>>& comp) {
    comp[r][c] = id;
    for (int d = 0; d < 4; d++) {
        int nr = r + dr[d], nc = c + dc[d];
        if (open(nr, nc) && comp[nr][nc] == -1) dfs(nr, nc, id, comp);
    }
}

// (2) iterative DFS with an explicit stack: safe for 10^6 cells; returns the size
int fillStack(int sr, int sc, int id, vector<vector<int>>& comp) {
    vector<pair<int,int>> st;
    comp[sr][sc] = id;                       // mark when PUSHED
    st.push_back(make_pair(sr, sc));
    int size = 0;
    while (!st.empty()) {
        int r = st.back().first, c = st.back().second;
        st.pop_back();
        size++;
        for (int d = 0; d < 4; d++) {
            int nr = r + dr[d], nc = c + dc[d];
            if (!open(nr, nc) || comp[nr][nc] != -1) continue;
            comp[nr][nc] = id;
            st.push_back(make_pair(nr, nc));
        }
    }
    return size;
}

// (3) BFS: same marking rule, a queue instead of a stack (Chapter 3 needs this one)
int fillQueue(int sr, int sc, int id, vector<vector<int>>& comp) {
    queue<pair<int,int>> q;
    comp[sr][sc] = id;
    q.push(make_pair(sr, sc));
    int size = 0;
    while (!q.empty()) {
        int r = q.front().first, c = q.front().second;
        q.pop();
        size++;
        for (int d = 0; d < 4; d++) {
            int nr = r + dr[d], nc = c + dc[d];
            if (!open(nr, nc) || comp[nr][nc] != -1) continue;
            comp[nr][nc] = id;
            q.push(make_pair(nr, nc));
        }
    }
    return size;
}

// (4) scan + fill: label every component, return their sizes
vector<int> labelAll(vector<vector<int>>& comp) {
    comp.assign(n, vector<int>(m, -1));
    vector<int> sizes;
    for (int r = 0; r < n; r++)
        for (int c = 0; c < m; c++)
            if (g[r][c] == '.' && comp[r][c] == -1)
                sizes.push_back(fillStack(r, c, sizes.size(), comp));
    return sizes;
}

int main() {
    n = 4; m = 5;
    g = {"..##.",
         "#.#..",
         ".###.",
         "..#.#"};
    vector<vector<int>> comp;
    vector<int> sizes = labelAll(comp);
    printf("%d components, sizes:", (int)sizes.size());
    for (int s : sizes) printf(" %d", s);                  // 4 components, sizes: 3 4 3 1
    printf("\n");
    for (int r = 0; r < n; r++, printf("\n"))
        for (int c = 0; c < m; c++)
            printf("%c", comp[r][c] < 0 ? '#' : char('A' + comp[r][c]));

    vector<vector<int>> c2(n, vector<int>(m, -1)), c3(n, vector<int>(m, -1));
    dfs(2, 0, 0, c2);                                       // recursive version
    printf("BFS from (0,4): %d cells\n", fillQueue(0, 4, 0, c3));   // 4 cells
}
\end{lstlisting}

\begin{insightbox}[Check the template by hand]
The grid in \texttt{main()} is the one of the scan picture. The labels printed are
\texttt{AA\#\#B / \#A\#BB / C\#\#\#B / CC\#D\#}: A $=\{(0,0),(0,1),(1,1)\}$, B $=\{(0,4),(1,3),(1,4),(2,4)\}$, C $=\{(2,0),(3,0),(3,1)\}$, D $=\{(3,3)\}$ --- sizes 3, 4, 3, 1. Note that $(1,1)$ and $(2,0)$ touch only at a corner, so A and C stay apart.
\end{insightbox}

\textbf{Reading grids.} Rows without spaces (\texttt{\#..\#}) are read with \texttt{cin >> row} into a \texttt{string}; numbers separated by spaces (\texttt{0 1 1}) need \texttt{cin >> a[r][c]} into an \texttt{int} array. For $10^6$ cells add \texttt{ios::sync\_with\_stdio(false); cin.tie(nullptr);} or use \texttt{scanf}.

%------------------------------------------------------------
\subsubsection*{Common mistakes}

\begin{warningbox}[Mistake 1 --- reading the cell before checking the bounds]
\begin{lstlisting}[style=snippet]
if (g[nr][nc] == '.' && nr >= 0 && nr < n && nc >= 0 && nc < m)   // WRONG: g[-1][..] is read first
if (nr >= 0 && nr < n && nc >= 0 && nc < m && g[nr][nc] == '.')   // RIGHT: && stops early
\end{lstlisting}
Out-of-range reads on a \texttt{vector<string>} may ``work'' locally and crash on the judge.
\end{warningbox}

\begin{warningbox}[Mistake 2 --- marking a cell when it is popped instead of when it is pushed]
If \texttt{visited} is set only on \texttt{pop}, a cell next to several marked cells is pushed several times and \emph{processed} several times. On a $2\times2$ island the buggy fill reports area 5, on a $4\times4$ island 25 instead of 16 (we ran it). Mark on push --- every cell enters the container once. (Marking on pop is correct only together with an ``\texttt{if (visited) continue;}'' right after the pop, at up to 4 times the memory.)
\end{warningbox}

\begin{warningbox}[Mistake 3 --- 4 versus 8 neighbours]
\begin{center}
\begin{tikzpicture}
    \begin{scope}
        \gIpIIIgrid{Awww,wBBw,wBww,Cwww}
        \node[font=\scriptsize\bfseries] at (1.05,-3.05) {4-connected: 3 islands};
    \end{scope}
    \begin{scope}[xshift=4.6cm]
        \gIpIIIgrid{Awww,wAAw,wAww,Awww}
        \node[font=\scriptsize\bfseries] at (1.05,-3.05) {8-connected: 1 island};
    \end{scope}
\end{tikzpicture}
\end{center}
Corner contacts join regions only under 8-connectivity. Using the wrong arrays gives a plausible but wrong count --- the samples often contain no diagonal contact, so test one yourself.
\end{warningbox}

\begin{warningbox}[{Mistake 4 --- \texttt{r * n + c} instead of \texttt{r * m + c}}]
The id multiplies by the row \emph{length} (the number of columns). On square grids both are equal, so square samples pass and rectangular tests fail.
\end{warningbox}

\begin{warningbox}[{Mistake 5 --- recursion on a $1000\times1000$ grid}]
See the warning above: the depth can reach $nm$. Use the explicit stack (template 2) or BFS (template 3) as soon as $nm$ exceeds about $10^5$.
\end{warningbox}

\begin{warningbox}[Mistake 6 --- recolouring with the same colour]
In a ``recolour the region'' fill, the visited test is usually ``still has the old colour''. If the new colour equals the old one, nothing ever changes and the recursion never stops. Return immediately when the two colours are equal (Exercise E3.2).
\end{warningbox}

\begin{checklistbox}[``I'm stuck'' checklist for Pattern 3]
\begin{itemize}[leftmargin=*]
    \item[$\square$] What is a vertex: every cell, or only the open / land / same-colour cells?
    \item[$\square$] Which cells are neighbours: 4 sides, 8 directions, knight moves, wrap-around?
    \item[$\square$] Is the question about \emph{regions} (count, size, sum, shape)? Then scan + fill. About \emph{steps}? Then BFS with distances (Chapter 3).
    \item[$\square$] Do I check \emph{inside} before reading the cell, and mark on \emph{push}?
    \item[$\square$] May I overwrite the input as my \texttt{visited} array, or do I need it afterwards?
    \item[$\square$] $nm$ up to $10^6$? Then no recursion, fast input, and \texttt{long long} if I sum values.
    \item[$\square$] Do I need information about a component \emph{later}? Then label every cell with a component id first and store per-id data.
\end{itemize}
\end{checklistbox}

%============================================================
% CH1 · P3 · PROBLEM 1 — COUNTING ROOMS (CSES 1192)
%============================================================
\subsection{Problem 1.3.1: Counting Rooms}
\problemlabel{1.3.1}{ssec:ch1-counting-rooms}

\problemcard{1.3.1}{Counting Rooms}{CSES 1192}{https://cses.fi/problemset/task/1192}{\exeasy\ (the canonical flood-fill problem)}{Grid as an implicit graph; count connected components with scan + iterative flood fill}

\subsubsection*{1. Problem statement}

You are given the map of a building: $n$ rows and $m$ columns of squares, each either floor (`\texttt{.}') or wall (`\texttt{\#}'). You can walk left, right, up and down through floor squares. Count the \textbf{rooms}: maximal groups of floor squares that are mutually reachable.

\textbf{Constraints:} $1 \le n, m \le 1000$. \quad (The same task appears on LeetCode as ``Number of Islands'', with land \texttt{1} instead of floor.)

\begin{center}
\small
\begin{tabular}{@{}l c@{}}
\toprule
\textbf{Input} & \textbf{Output} \\
\midrule
\texttt{5 8} & \\
\texttt{\#\#\#\#\#\#\#\#} & \\
\texttt{\#..\#...\#} & \\
\texttt{\#\#\#\#.\#.\#} & 3 \\
\texttt{\#..\#...\#} & \\
\texttt{\#\#\#\#\#\#\#\#} & \\
\bottomrule
\end{tabular}
\end{center}

\begin{diagrambox}
\begin{tikzpicture}
    \begin{scope}
        \gIpIIIgrid{wwwwwwww,woowooow,wwwwowow,woowooow,wwwwwwww}
        \foreach \c in {0,...,7} \node[font=\tiny, text=darkgray] at (0.7*\c, 0.5) {\c};
        \foreach \r in {0,...,4} \node[font=\tiny, text=darkgray] at (-0.55, -0.7*\r) {\r};
        \node[font=\small\bfseries] at (2.45,-3.45) {the map};
    \end{scope}
    \begin{scope}[xshift=7.2cm]
        \gIpIIIgrid{wwwwwwww,wAAwBBBw,wwwwBwBw,wCCwBBBw,wwwwwwww}
        \node[font=\small\bfseries] at (2.45,-3.45) {its 3 rooms};
    \end{scope}
\end{tikzpicture}
\figcap{Sample 1. Room A has 2 squares, room B (a ring around the wall at $(2,5)$) has 8, room C has 2. A wall square inside a room does not split it: B is still one room because you can walk around the pillar.}
\end{diagrambox}

\subsubsection*{2. The graph model}

\begin{modelbox}[Graph Model --- Counting Rooms]
\renewcommand{\arraystretch}{1.35}
\begin{tabularx}{\linewidth}{@{}>{\bfseries}l X@{}}
Vertices & the floor squares $(r,c)$; at most $nm = 10^6$ \\
Edges & between two floor squares that share a side (4-neighbourhood); undirected; \textbf{implicit} --- computed from the map, never stored \\
Weights & none \\
Question & the number of \textbf{connected components} \\
Size & $|V| \le 10^6$, $|E| \le 2\cdot10^6$: an $\bigO{|V|+|E|} = \bigO{nm}$ algorithm, about $5\cdot10^6$ simple steps \\
\end{tabularx}
\end{modelbox}

\begin{diagrambox}
\begin{tikzpicture}
    \foreach \r/\c/\s in {1/1/gvA,1/2/gvA, 1/4/gvB,1/5/gvB,1/6/gvB,2/4/gvB,2/6/gvB,3/4/gvB,3/5/gvB,3/6/gvB, 3/1/gvsel,3/2/gvsel}
        \node[\s, minimum size=6.5mm, font=\tiny] (v\r\c) at (1.1*\c, -1.0*\r) {\r,\c};
    \foreach \a/\b in {11/12, 14/15, 15/16, 14/24, 16/26, 24/34, 26/36, 34/35, 35/36, 31/32}
        \draw[ge] (v\a) -- (v\b);
    \node[font=\scriptsize\bfseries, text=primaryblue] at (1.65,-0.45) {A};
    \node[font=\scriptsize\bfseries, text=recursionred] at (5.5,-2.0) {B};
    \node[font=\scriptsize\bfseries, text=softpurple] at (1.65,-3.55) {C};
    \node[note, anchor=west, text width=5.2cm, align=left] at (7.6,-2.0) {12 vertices, 10 edges, 3 components. Component B is a \textbf{cycle} of 8 vertices --- cycles are no problem for a fill, because \texttt{visited} stops it from going round forever.};
\end{tikzpicture}
\figcap{The implicit graph of the sample. Walls have no vertex; the answer is the number of pieces of this graph.}
\end{diagrambox}

\subsubsection*{3. Key observations}

\begin{itemize}[leftmargin=*]
    \item \textbf{A room is a connected component} of the floor graph --- ``you can walk between any two of its squares, and it cannot be extended'' is word for word the definition of a component.
    \item \textbf{Counting components = counting fills.} Scan the map; every floor square that is not yet in a counted room starts a new room, and a fill marks that whole room (Pattern 3 invariant).
    \item \textbf{The grid can be its own \texttt{visited} array}: turning a visited `\texttt{.}' into `\texttt{\#}' never changes the answer, because that room is already counted.
    \item \textbf{Depth trap:} $nm$ reaches $10^6$, and a recursive fill on an open $1000\times1000$ room recurses $10^6$ levels deep. Use an explicit stack.
\end{itemize}

\subsubsection*{4. Approaches}

\begin{bruteforcebox}[Approach 1: Brute Force --- Label Propagation]
\textbf{Idea.} Give every floor square its own id as a label. Repeat sweeps over the map: each square takes the \emph{smallest} label among itself and its floor neighbours. When a sweep changes nothing, every room carries the smallest id of the room; the answer is the number of squares whose label is still their own id. No graph search at all --- just a local rule, repeated.

\begin{lstlisting}[style=cpp]
#include <bits/stdc++.h>
using namespace std;

int main() {
    int n, m;
    cin >> n >> m;
    vector<string> g(n);
    for (auto& row : g) cin >> row;
    vector<int> lab(n * m, -1);
    for (int r = 0; r < n; r++)
        for (int c = 0; c < m; c++)
            if (g[r][c] == '.') lab[r * m + c] = r * m + c;

    const int dr[4] = {-1, 1, 0, 0}, dc[4] = {0, 0, -1, 1};
    bool changed = true;
    while (changed) {                       // one round = one sweep over all cells
        changed = false;
        for (int r = 0; r < n; r++)
            for (int c = 0; c < m; c++) {
                if (g[r][c] != '.') continue;
                for (int d = 0; d < 4; d++) {
                    int nr = r + dr[d], nc = c + dc[d];
                    if (nr < 0 || nr >= n || nc < 0 || nc >= m || g[nr][nc] != '.') continue;
                    if (lab[nr * m + nc] < lab[r * m + c]) {
                        lab[r * m + c] = lab[nr * m + nc];
                        changed = true;
                    }
                }
            }
    }
    int rooms = 0;
    for (int i = 0; i < n * m; i++)
        if (lab[i] == i) rooms++;           // the smallest id of each room survives
    cout << rooms << "\n";
}
\end{lstlisting}

\begin{diagrambox}
\begin{tikzpicture}
    \begin{scope}
        \def\gIpIIIs{0.32}
        \gIpIIIgrid{ooooooooo,wwwwwwwwo,ooooooooo,owwwwwwww,ooooooooo,wwwwwwwwo,ooooooooo,owwwwwwww,ooooooooo}
    \end{scope}
    \node[note, anchor=west, text width=4.6cm, align=left] at (3.3,-1.3) {A ``snake'' room: one long corridor folded back and forth. A small label must travel along the whole corridor, and against the sweep direction it moves only about one square per sweep.};
    \node[font=\scriptsize] at (10.4,-1.3) {
        \begin{tabular}{@{}rr@{}}
        \toprule
        snake $n\times n$ & sweeps \\
        \midrule
        $20\times20$ & 97 \\
        $40\times40$ & 392 \\
        $80\times80$ & 1582 \\
        \bottomrule
        \end{tabular}};
\end{tikzpicture}
\figcap{Measured number of sweeps on snakes: it grows like $nm/4$. For $1000\times1000$ that is about $2.5\cdot10^5$ sweeps of $10^6$ squares each.}
\end{diagrambox}

\begin{complexitybox}
\textbf{Time:} $\bigO{(nm)^2}$ in the worst case (up to $\approx nm/4$ sweeps of $nm$ squares) --- about $2.5\cdot10^{11}$ steps for $1000\times1000$. \textbf{Space:} $\bigO{nm}$. \textbf{Verdict:} \tle{} --- but correct and very easy to trust, so it is our reference on small random maps.
\end{complexitybox}
\end{bruteforcebox}

\begin{algorithmbox}[Approach 2: Scan + Iterative Flood Fill]
\textbf{Idea.} Scan the map; a `\texttt{.}' square starts a new room. Fill that room with an explicit stack, overwriting every reached floor square with `\texttt{\#}' the moment it is pushed. When the scan ends, the counter is the answer.

\begin{lstlisting}[style=cpp]
#include <bits/stdc++.h>
using namespace std;

int main() {
    int n, m;
    scanf("%d %d", &n, &m);
    vector<string> g(n);
    char buf[1005];
    for (int i = 0; i < n; i++) { scanf("%s", buf); g[i] = buf; }

    const int dr[4] = {-1, 1, 0, 0}, dc[4] = {0, 0, -1, 1};
    vector<int> st;                        // explicit stack of cell ids r*m+c
    st.reserve(n * m);
    int rooms = 0;
    for (int r = 0; r < n; r++)
        for (int c = 0; c < m; c++) {
            if (g[r][c] != '.') continue;  // wall, or floor already in a counted room
            rooms++;                       // (r,c) starts a new room
            g[r][c] = '#';                 // mark on PUSH: never pushed twice
            st.push_back(r * m + c);
            while (!st.empty()) {
                int id = st.back(); st.pop_back();
                int cr = id / m, cc = id % m;
                for (int d = 0; d < 4; d++) {
                    int nr = cr + dr[d], nc = cc + dc[d];
                    if (nr < 0 || nr >= n || nc < 0 || nc >= m) continue;  // outside
                    if (g[nr][nc] != '.') continue;                       // wall / seen
                    g[nr][nc] = '#';
                    st.push_back(nr * m + nc);
                }
            }
        }
    printf("%d\n", rooms);
}
\end{lstlisting}

\begin{complexitybox}
\textbf{Time:} $\bigO{nm}$ --- every square is pushed and popped at most once and looks at 4 neighbours. Measured: 0.16\,s for an all-floor $1000\times1000$ map, input included. \textbf{Space:} $\bigO{nm}$ for the map and the stack (the stack never holds more than $nm$ ids). \textbf{Verdict:} \ac{}.
\end{complexitybox}
\end{algorithmbox}

\begin{invariantbox}[Why it is correct]
\textbf{Invariant (between two iterations of the scan):} every floor square of an already counted room has been overwritten with `\texttt{\#}', and no square of an uncounted room has been touched.

\emph{Initially} nothing is counted and nothing is touched. \emph{Step:} when the scan meets a `\texttt{.}' at $(r,c)$, the invariant says its room is uncounted, so \texttt{rooms++} counts it exactly once; the fill then overwrites exactly the squares reachable from $(r,c)$ --- that room and nothing else (Pattern 3, Claim 1). \emph{End:} every floor square was scanned, so every room was met at least once while still untouched, and was counted exactly once. \hfill$\blacksquare$
\end{invariantbox}

\begin{alternativebox}[Approach 3: Alternative Algorithms --- BFS Fill, or DSU]
\textbf{BFS fill.} Replace the stack by a \texttt{queue} and mark on enqueue (template 3 of the theory). Same $\bigO{nm}$, same answer --- checked on 500 random maps against the brute force. BFS is the version you will extend in Chapter 3 to compute distances.

\textbf{DSU preview (Chapter 6).} Start with every floor square as its own room. For each floor square, unite it with its upper and its left floor neighbour; every union of two \emph{different} sets merges two rooms, so \texttt{rooms--}. There is no search order at all --- and it keeps working if floor squares are \emph{added} one at a time, which a fill cannot handle cheaply.
\begin{lstlisting}[style=cpp]
#include <bits/stdc++.h>
using namespace std;

vector<int> par;
int findRoot(int x) {                      // with path halving
    while (par[x] != x) { par[x] = par[par[x]]; x = par[x]; }
    return x;
}
bool unite(int a, int b) {                 // true if two different rooms merged
    a = findRoot(a); b = findRoot(b);
    if (a == b) return false;
    par[a] = b;
    return true;
}

int main() {
    ios::sync_with_stdio(false);
    cin.tie(nullptr);
    int n, m;
    cin >> n >> m;
    vector<string> g(n);
    for (auto& row : g) cin >> row;
    par.resize(n * m);
    for (int i = 0; i < n * m; i++) par[i] = i;
    int rooms = 0;
    for (int r = 0; r < n; r++)
        for (int c = 0; c < m; c++) {
            if (g[r][c] != '.') continue;
            rooms++;                                   // every floor cell starts alone
            if (r > 0 && g[r - 1][c] == '.' && unite(r * m + c, (r - 1) * m + c)) rooms--;
            if (c > 0 && g[r][c - 1] == '.' && unite(r * m + c, r * m + c - 1)) rooms--;
        }
    cout << rooms << "\n";
}
\end{lstlisting}
\end{alternativebox}

\begin{tricksbox}
\begin{itemize}[leftmargin=*]
    \item \textbf{No separate \texttt{visited}:} overwriting the map saves $10^6$ bytes and one comparison per neighbour.
    \item \textbf{Sentinel border:} the sample is surrounded by walls, but that is not promised. If you pad the map yourself with a ring of `\texttt{\#}' (size $(n+2)\times(m+2)$), the four bounds checks disappear from the inner loop.
    \item \textbf{One integer per stack entry} (\texttt{r*m+c}) instead of a \texttt{pair}: half the memory traffic.
    \item \textbf{Fast input:} $10^6$ characters --- read whole rows with \texttt{scanf("\%s")}, or \texttt{cin >> row} after \texttt{sync\_with\_stdio(false)}; never character by character with slow streams.
\end{itemize}
\end{tricksbox}

\subsubsection*{5. Trace}

\begin{tracebox}[{Sample 1 --- the fill of room B (stack, neighbour order up, down, left, right)}]
The scan reaches $(1,1)$ first: room A $=\{(1,1),(1,2)\}$ is filled in two pops. The next `\texttt{.}' is $(1,4)$: \texttt{rooms} becomes 2 and the fill of B starts.

\begin{center}
\begin{tikzpicture}
    \def\gIpIIIs{0.5}
    \begin{scope}
        \gIpIIIgrid{wwwwwwww,wAAwxqow,wwwwqwow,woowooow,wwwwwwww}
        \node[font=\scriptsize\bfseries, align=center] at (1.75,-2.6) {pop 1: $(1,4)$\\push $(2,4)$, $(1,5)$};
    \end{scope}
    \begin{scope}[xshift=4.8cm]
        \gIpIIIgrid{wwwwwwww,wAAwkkxw,wwwwqwqw,woowooow,wwwwwwww}
        \node[font=\scriptsize\bfseries, align=center] at (1.75,-2.6) {pops 2--3: $(1,5)$, $(1,6)$\\push $(1,6)$, then $(2,6)$};
    \end{scope}
    \begin{scope}[xshift=9.6cm]
        \gIpIIIgrid{wwwwwwww,wAAwkkkw,wwwwqwkw,woowqxkw,wwwwwwww}
        \node[font=\scriptsize\bfseries, align=center] at (1.75,-2.6) {pops 4--6: $(2,6)$, $(3,6)$, $(3,5)$\\push $(3,6)$, $(3,5)$, $(3,4)$};
    \end{scope}
    \begin{scope}[yshift=-3.7cm, xshift=2.4cm]
        \gIpIIIgrid{wwwwwwww,wAAwBBBw,wwwwBwBw,woowBBBw,wwwwwwww}
        \node[font=\scriptsize\bfseries, align=center] at (1.75,-2.6) {pops 7--8: $(3,4)$, $(2,4)$\\stack empty: room B has 8 squares};
    \end{scope}
    \begin{scope}[yshift=-3.7cm, xshift=7.4cm]
        \gIpIIIgrid{wwwwwwww,wAAwBBBw,wwwwBwBw,wCCwBBBw,wwwwwwww}
        \node[font=\scriptsize\bfseries, align=center] at (1.75,-2.6) {the scan continues; $(3,1)$ starts C\\\texttt{rooms} $= 3$};
    \end{scope}
\end{tikzpicture}
\end{center}
{\small Orange frame = the square just popped; yellow = on the stack (already marked); light blue = popped earlier.}

\begin{center}
\small
\begin{tabular}{@{}clll@{}}
\toprule
\textbf{pop} & \textbf{square} & \textbf{newly pushed} & \textbf{stack after (bottom $\to$ top)} \\
\midrule
1 & $(1,4)$ & $(2,4)$, $(1,5)$ & $(2,4)\ (1,5)$ \\
2 & $(1,5)$ & $(1,6)$ & $(2,4)\ (1,6)$ \\
3 & $(1,6)$ & $(2,6)$ & $(2,4)\ (2,6)$ \\
4 & $(2,6)$ & $(3,6)$ & $(2,4)\ (3,6)$ \\
5 & $(3,6)$ & $(3,5)$ & $(2,4)\ (3,5)$ \\
6 & $(3,5)$ & $(3,4)$ & $(2,4)\ (3,4)$ \\
7 & $(3,4)$ & --- ($(2,4)$ is already marked) & $(2,4)$ \\
8 & $(2,4)$ & --- & empty \\
\bottomrule
\end{tabular}
\end{center}
Pop 7 is where ``mark on push'' pays off: $(2,4)$ is a floor neighbour of $(3,4)$, but it was marked at pop 1, so it is not pushed a second time. The fill went the long way round the ring and closed the cycle without noticing it. Output: \textbf{3}.
\end{tracebox}

\subsubsection*{6. Test suite}

\begin{testbox}
\centering\small
\begin{tabular}{@{}lll@{}}
\toprule
\textbf{Map} & \textbf{Expected} & \textbf{Why} \\
\midrule
the sample & 3 & a room with a pillar is still one room \\
$1\times1$ `\texttt{.}' & 1 & smallest map, no wall border \\
$1\times1$ `\texttt{\#}' & 0 & no floor at all \\
$1000\times1000$ all `\texttt{\#}' & 0 & scan only, no fill \\
$1000\times1000$ all `\texttt{.}' & 1 & one huge room; recursion depth would be $10^6$ \\
$1000\times1000$ checkerboard & 500000 & one-square rooms; diagonal contacts do not count \\
$1000\times1000$ snake & 1 & long corridor; brute force: $\approx 2.5\cdot10^5$ sweeps \\
500 random maps ($\le 8\times8$) & = brute & all four approaches agree \\
\bottomrule
\end{tabular}
\end{testbox}

\subsubsection*{7. Common mistakes}

\begin{warningbox}[{Mistake 1 --- recursive fill on a $1000\times1000$ map}]
\begin{lstlisting}[style=snippet]
void dfs(int r, int c) {                      // depth up to n*m = 10^6
    if (r < 0 || r >= n || c < 0 || c >= m || g[r][c] != '.') return;
    g[r][c] = '#';
    dfs(r - 1, c); dfs(r + 1, c); dfs(r, c - 1); dfs(r, c + 1);
}
\end{lstlisting}
Correct on the sample (maximum depth 8), but on an all-floor $1000\times1000$ map the depth is $10^6$ and the program crashes with the usual stack (we ran it). Use the explicit stack of Approach 2.
\end{warningbox}

\begin{warningbox}[Mistake 2 --- counting squares instead of fills]
``Number of rooms'' counts \emph{fills started}, not squares filled. Put \texttt{rooms++} next to the start of a fill, outside the \texttt{while} loop.
\end{warningbox}

\begin{warningbox}[Mistake 3 --- slow character-by-character input]
\texttt{cin >> ch} for $10^6$ characters on synchronised streams can cost more than the whole algorithm. Read a row at a time.
\end{warningbox}

\begin{warningbox}[Mistake 4 --- adding diagonal neighbours]
The statement says left, right, up, down. With 8 directions the checkerboard test gives 1 instead of 500000.
\begin{center}
\begin{tikzpicture}
    \def\gIpIIIs{0.45}
    \begin{scope}
        \gIpIIIgrid{awbwc,wdwew,fwawb,wcwdw}
        \node[font=\scriptsize\bfseries] at (0.9,-2.15) {4 directions: 10 rooms};
    \end{scope}
    \begin{scope}[xshift=4.2cm]
        \gIpIIIgrid{awawa,wawaw,awawa,wawaw}
        \node[font=\scriptsize\bfseries] at (0.9,-2.15) {8 directions: 1 room};
    \end{scope}
\end{tikzpicture}
\end{center}
\end{warningbox}

\subsubsection*{8. Takeaways}

\begin{itemize}[leftmargin=*]
    \item Squares = vertices, shared sides = edges; rooms = connected components.
    \item Number of components = number of fills started by the row-major scan.
    \item Mark on push, and let the map be its own \texttt{visited} array.
    \item For $nm$ up to $10^6$: explicit stack or queue, not recursion.
\end{itemize}

\begin{variantsbox}
\variant{Eight directions}{\varnone}
  {squares touching at a corner also belong to the same room.}
  {vertices unchanged; edges also join diagonal neighbours (up to 8 per square).}
  {use the 8-entry direction arrays; everything else is identical, still $\bigO{nm}$.}
\variant{Volume of the biggest lake}{\varin{https://codeforces.com/problemset/problem/1829/E}}
  {cells carry depths; find the component with the largest \emph{sum} of depths.}
  {vertices = cells with depth $> 0$; same 4-neighbour edges.}
  {the fill adds up the values it reaches; answer = the largest sum (Exercise E3.3).}
\variant{Sub-islands}{\varin{https://leetcode.com/problems/count-sub-islands/}}
  {count the islands of grid 2 whose every cell is also land in grid 1.}
  {components of grid 2; a component is ``good'' if all its vertices are land in grid 1.}
  {fill each island of grid 2 completely, AND-ing a flag over its cells; count the good ones (Exercise E3.4).}
\variant{Rooms away from the outside}{\varnone}
  {count only rooms that contain no square of the outer border.}
  {same graph; a component is ``closed'' if none of its vertices lies on row $0$, row $n{-}1$, column $0$ or column $m{-}1$.}
  {first fill from every border floor square and discard those rooms, then count the remaining fills (Chapter 4 studies this ``fill from the border'' idea).}
\variant{Walls broken one by one}{\varnone}
  {after each of $q$ updates ``square $(r,c)$ becomes floor'', print the number of rooms.}
  {the graph only gains vertices and edges.}
  {DSU: a new floor square adds 1 room, every union with a different neighbouring room subtracts 1 (Chapter 6).}
\end{variantsbox}

%============================================================
% CH1 · P3 · PROBLEM 2 — MAX AREA OF ISLAND (LeetCode 695)
%============================================================
\subsection{Problem 1.3.2: Max Area of Island}
\problemlabel{1.3.2}{ssec:ch1-max-area}

\problemcard{1.3.2}{Max Area of Island}{LeetCode 695}{https://leetcode.com/problems/max-area-of-island/}{\exmedium\ \lc{Medium}}{Flood fill that returns the size of the component; maximum over all components}

\subsubsection*{1. Problem statement}

You are given an $R \times C$ binary matrix \texttt{grid}. An \textbf{island} is a group of \texttt{1}s (land) connected 4-directionally (horizontally or vertically); everything outside the grid is water. The \textbf{area} of an island is the number of its cells. Return the maximum area of an island, or $0$ if there is no island.

\textbf{Constraints:} $1 \le R, C \le 50$, \; \texttt{grid[i][j]} is $0$ or $1$.

\begin{center}
\small
\begin{tabular}{@{}lcl@{}}
\toprule
\textbf{grid} & \textbf{Output} & \textbf{Note} \\
\midrule
the $8\times13$ map below & 6 & not 11: two islands touch only at a corner \\
\texttt{[[0,0,0,0,0,0,0,0]]} & 0 & no land at all \\
\bottomrule
\end{tabular}
\end{center}

\begin{diagrambox}
\begin{tikzpicture}
    \def\gIpIIIs{0.55}
    \gIpIIIgrid{wwowwwwowwwww,wwwwwwwooowww,woowowwwwwwww,wowwoowwowoww,wowwoowwoooww,wwwwwwwwwwoww,wwwwwwwooowww,wwwwwwwoowwww}
    \foreach \c in {0,...,12} \node[font=\tiny, text=darkgray] at (0.55*\c, 0.42) {\c};
    \foreach \r in {0,...,7} \node[font=\tiny, text=darkgray] at (-0.45, -0.55*\r) {\r};
    \node[font=\small\bfseries] at (3.3,-4.5) {Example 1 (white = land, grey = water)};
\end{tikzpicture}
\figcap{The first example of the statement. How many islands do you see, and which is the largest?}
\end{diagrambox}

\subsubsection*{2. The graph model}

\begin{modelbox}[Graph Model --- Max Area of Island]
\renewcommand{\arraystretch}{1.35}
\begin{tabularx}{\linewidth}{@{}>{\bfseries}l X@{}}
Vertices & the land cells; at most $RC = 2500$ \\
Edges & between land cells sharing a side; undirected; implicit \\
Weights & none (every vertex counts 1 towards the area) \\
Question & the \textbf{largest number of vertices in a connected component} (0 if there is no vertex) \\
Size & $|V| \le 2500$, $|E| \le 5000$: anything up to $\bigO{(RC)^2} \approx 6\cdot10^6$ fits; $\bigO{RC}$ is the clean target \\
\end{tabularx}
\end{modelbox}

\begin{diagrambox}
\begin{tikzpicture}
    \def\gIpIIIs{0.55}
    \gIpIIIgrid{wwAwwwwBwwwww,wwwwwwwBBBwww,wCCwDwwwwwwww,wCwwDDwwEwEww,wCwwDDwwEEEww,wwwwwwwwwwEww,wwwwwwwFFFwww,wwwwwwwFFwwww}
    \draw[warningred, line width=1.4pt] (5.225,-3.025) circle (0.5);
    \draw[warningred, -{Stealth[length=2mm]}] (6.6,-3.6) -- (5.7,-3.2);
    \node[font=\scriptsize, text=warningred, anchor=west] at (6.6,-3.7) {corner contact: no edge};
    \node[font=\scriptsize, anchor=west] at (7.6,-1.2) {
        \begin{tabular}{@{}clc@{}}
        \toprule
        island & first cell & area \\
        \midrule
        A & $(0,2)$ & 1 \\
        B & $(0,7)$ & 4 \\
        C & $(2,1)$ & 4 \\
        D & $(2,4)$ & 5 \\
        E & $(3,8)$ & \textbf{6} \\
        F & $(6,7)$ & 5 \\
        \bottomrule
        \end{tabular}};
\end{tikzpicture}
\figcap{The components, named in the order the row-major scan meets them. E's cell $(5,10)$ and F's cell $(6,9)$ touch only at a corner, so they are \emph{different} islands: counting them together would give the wrong answer $6+5=11$.}
\end{diagrambox}

\subsubsection*{3. Key observations}

\begin{itemize}[leftmargin=*]
    \item \textbf{Area = size of a component.} Problem 1.3.1 counted components; here we \emph{measure} each one and keep the maximum. The only change: the fill returns how many cells it marked.
    \item \textbf{A recursive fill can return the size directly:} area $= 1 + $ the areas returned by the calls on the four neighbours. Each cell contributes its 1 exactly once because it is marked before the recursion continues.
    \item \textbf{Recursion is safe here:} at most $RC = 2500$ land cells, so the depth is at most 2500 frames --- tiny.
    \item \textbf{Sinking:} setting a visited land cell to $0$ removes it from the graph; no extra \texttt{visited} array is needed (LeetCode lets us modify \texttt{grid}).
\end{itemize}

\subsubsection*{4. Approaches}

\begin{bruteforcebox}[Approach 1: Brute Force --- a Fresh Fill From Every Land Cell]
\textbf{Idea.} For every land cell, start a new fill with a \emph{fresh} \texttt{seen} array and count the cells it reaches; the answer is the largest count. Every island is measured once per cell it contains.

\begin{lstlisting}[style=cpp]
#include <bits/stdc++.h>
using namespace std;

class Solution {
public:
    int maxAreaOfIsland(vector<vector<int>>& grid) {
        int R = grid.size(), C = grid[0].size();
        const int dr[4] = {-1, 1, 0, 0}, dc[4] = {0, 0, -1, 1};
        int best = 0;
        for (int r = 0; r < R; r++)
            for (int c = 0; c < C; c++) {
                if (grid[r][c] != 1) continue;
                vector<vector<char>> seen(R, vector<char>(C, 0));   // fresh every time
                vector<pair<int,int>> st(1, make_pair(r, c));
                seen[r][c] = 1;
                int area = 0;
                while (!st.empty()) {
                    int x = st.back().first, y = st.back().second;
                    st.pop_back();
                    area++;
                    for (int d = 0; d < 4; d++) {
                        int nx = x + dr[d], ny = y + dc[d];
                        if (nx < 0 || nx >= R || ny < 0 || ny >= C) continue;
                        if (grid[nx][ny] != 1 || seen[nx][ny]) continue;
                        seen[nx][ny] = 1;
                        st.push_back(make_pair(nx, ny));
                    }
                }
                best = max(best, area);
            }
        return best;
    }
};
\end{lstlisting}

\begin{diagrambox}
\begin{tikzpicture}
    \def\gIpIIIs{0.5}
    \foreach \x in {0, 2.6, 5.2} {
        \begin{scope}[xshift=\x cm]
            \gIpIIIgrid{aaw,aww,wwb}
        \end{scope}
    }
    \gIpIIIat{0}{0}{s}
    \begin{scope}[xshift=2.6cm]\gIpIIIat{0}{1}{s}\end{scope}
    \begin{scope}[xshift=5.2cm]\gIpIIIat{1}{0}{s}\end{scope}
    \node[note, anchor=west, text width=4.2cm, align=left] at (7.0,-0.5) {3 starts $\times$ 3 cells = 9 steps for the blue island, then one more fill for the green one.};
\end{tikzpicture}
\figcap{The same 3-cell island (blue) is filled again from each of its cells (s = start). An island of $k$ cells costs $k^2$ steps, plus a fresh $R\times C$ array per start.}
\end{diagrambox}

\begin{complexitybox}
\textbf{Time:} $\bigO{(RC)^2}$ --- each of up to $RC$ starts allocates and fills up to $RC$ cells. For $50\times50$ that is about $6\cdot10^6$: accepted on LeetCode, but the same idea on a $1000\times1000$ map would need $10^{12}$ steps. \textbf{Space:} $\bigO{RC}$. \textbf{Verdict:} \ac{} here, \tle{} in general --- and our reference for random tests.
\end{complexitybox}
\end{bruteforcebox}

\begin{algorithmbox}[Approach 2: Scan + Recursive Fill That Returns the Area]
\textbf{Idea.} Scan the grid; each land cell still present starts a fill that \emph{sinks} its island (sets every cell to $0$) and returns the number of cells sunk. Keep the maximum.

\begin{lstlisting}[style=cpp]
#include <bits/stdc++.h>
using namespace std;

class Solution {
    int R, C;
    int sink(vector<vector<int>>& g, int r, int c) {   // returns the area it sinks
        if (r < 0 || r >= R || c < 0 || c >= C || g[r][c] != 1) return 0;
        g[r][c] = 0;                                   // visited = turned into water
        return 1 + sink(g, r - 1, c) + sink(g, r + 1, c)
                 + sink(g, r, c - 1) + sink(g, r, c + 1);
    }
public:
    int maxAreaOfIsland(vector<vector<int>>& grid) {
        R = grid.size(); C = grid[0].size();
        int best = 0;
        for (int r = 0; r < R; r++)
            for (int c = 0; c < C; c++)
                if (grid[r][c] == 1) best = max(best, sink(grid, r, c));
        return best;
    }
};
\end{lstlisting}

Here the checks are made \emph{at the top of the call} (``am I outside, or not land?'') instead of before the call. That costs a few extra calls that return 0 immediately, but the fill stays three lines long. It is the standard style for small grids.

\begin{complexitybox}
\textbf{Time:} $\bigO{RC}$ --- every land cell is sunk once and makes 4 calls; every water cell is only looked at. \textbf{Space:} $\bigO{RC}$ recursion depth in the worst case ($\le 2500$ frames, safe). \textbf{Verdict:} \ac{}.
\end{complexitybox}
\end{algorithmbox}

\begin{invariantbox}[Why it is correct]
\textbf{Lemma.} If $(r,c)$ is land when \texttt{sink(r,c)} is called, the call sinks exactly the land cells that are reachable from $(r,c)$ through land \emph{at the moment of the call}, and returns their number.

\emph{Proof by induction on the number of land cells.} The call sinks $(r,c)$ itself (1 cell) and then calls its four neighbours one after another. Each call on a neighbour that is still land sinks, by induction (the grid now has fewer land cells), the part of the island reachable from that neighbour without passing through cells already sunk, and returns its size; calls on water return 0. The parts are disjoint, because a cell is sunk at most once, and together with $(r,c)$ they cover everything reachable from $(r,c)$: such a cell is reached through some neighbour of $(r,c)$, and is either sunk by that neighbour's call or already sunk by an earlier sibling call. So $1 + \sum$ counts every island cell exactly once.

\textbf{Main loop.} By the lemma, each call from the scan sinks one complete island and returns its area; islands sunk earlier are never measured again, and every island is met by the scan. Hence \texttt{best} is the maximum area. \hfill$\blacksquare$
\end{invariantbox}

\begin{alternativebox}[Approach 3: Alternative Algorithms --- BFS, or DSU With Sizes]
\textbf{BFS, input untouched.} If the grid must not be modified, use a separate \texttt{seen} array and count a cell when it leaves the queue:
\begin{lstlisting}[style=cpp]
#include <bits/stdc++.h>
using namespace std;

class Solution {
public:
    int maxAreaOfIsland(vector<vector<int>>& grid) {
        int R = grid.size(), C = grid[0].size();
        const int dr[4] = {-1, 1, 0, 0}, dc[4] = {0, 0, -1, 1};
        vector<vector<char>> seen(R, vector<char>(C, 0));
        int best = 0;
        for (int r = 0; r < R; r++)
            for (int c = 0; c < C; c++) {
                if (grid[r][c] != 1 || seen[r][c]) continue;
                queue<int> q;
                q.push(r * C + c);
                seen[r][c] = 1;
                int area = 0;
                while (!q.empty()) {
                    int id = q.front(); q.pop();
                    area++;                                    // count when DEQUEUED
                    int x = id / C, y = id % C;
                    for (int d = 0; d < 4; d++) {
                        int nx = x + dr[d], ny = y + dc[d];
                        if (nx < 0 || nx >= R || ny < 0 || ny >= C) continue;
                        if (grid[nx][ny] != 1 || seen[nx][ny]) continue;
                        seen[nx][ny] = 1;
                        q.push(nx * C + ny);
                    }
                }
                best = max(best, area);
            }
        return best;
    }
};
\end{lstlisting}

\textbf{DSU with sizes (preview of Chapter 6).} Unite every land cell with its land neighbours below and to the right; union by size keeps \texttt{sz[root]} = area of the island. The answer is the largest \texttt{sz} of a root.
\begin{lstlisting}[style=cpp]
#include <bits/stdc++.h>
using namespace std;

class Solution {
    vector<int> par, sz;
    int find(int x) { while (par[x] != x) x = par[x] = par[par[x]]; return x; }
    void unite(int a, int b) {
        a = find(a); b = find(b);
        if (a == b) return;
        if (sz[a] < sz[b]) swap(a, b);
        par[b] = a; sz[a] += sz[b];                    // union by size
    }
public:
    int maxAreaOfIsland(vector<vector<int>>& grid) {
        int R = grid.size(), C = grid[0].size();
        par.resize(R * C); sz.assign(R * C, 1);
        for (int i = 0; i < R * C; i++) par[i] = i;
        for (int r = 0; r < R; r++)
            for (int c = 0; c < C; c++) {
                if (grid[r][c] != 1) continue;
                if (r + 1 < R && grid[r + 1][c] == 1) unite(r * C + c, (r + 1) * C + c);
                if (c + 1 < C && grid[r][c + 1] == 1) unite(r * C + c, r * C + c + 1);
            }
        int best = 0;
        for (int i = 0; i < R * C; i++)
            if (grid[i / C][i % C] == 1 && find(i) == i) best = max(best, sz[i]);
        return best;
    }
};
\end{lstlisting}

\begin{diagrambox}
\begin{tikzpicture}
    \def\gIpIIIs{0.6}
    \begin{scope}
        \gIpIIIgrid{owo,ooo}
        \draw[gtree] (0,0) -- (0,-0.6);
        \draw[gtree] (0,-0.6) -- (0.6,-0.6);
        \draw[gtree] (0.6,-0.6) -- (1.2,-0.6);
        \draw[gtree] (1.2,0) -- (1.2,-0.6);
        \node[font=\scriptsize\bfseries, align=center] at (0.6,-1.4) {unite with the land\\below and to the right};
    \end{scope}
    \draw[-{Stealth[length=3mm]}, line width=1.2pt, primaryblue] (2.0,-0.3) -- (3.0,-0.3);
    \begin{scope}[xshift=3.6cm, yshift=0.3cm]
        \node[gvdst, minimum size=7mm, font=\scriptsize] (rt) at (1.2,0) {};
        \node[gv, minimum size=7mm, font=\scriptsize] (a) at (0,-1.2) {};
        \node[gv, minimum size=7mm, font=\scriptsize] (b) at (0.8,-1.2) {};
        \node[gv, minimum size=7mm, font=\scriptsize] (c) at (1.6,-1.2) {};
        \node[gv, minimum size=7mm, font=\scriptsize] (d) at (2.4,-1.2) {};
        \foreach \x in {a,b,c,d} \draw[gde] (\x) -- (rt);
        \node[gdist, anchor=west] at ([xshift=2pt]rt.east) {root: sz $=5$};
    \end{scope}
    \node[note, anchor=west, text width=5.2cm, align=left] at (7.6,-0.4) {One DSU tree per island; its root stores the island's size. The exact shape depends on the union order (path compression flattens it), the size does not.};
\end{tikzpicture}
\figcap{DSU view of a 5-cell island: no traversal at all, only ``merge two sets'' operations, each adding the sizes.}
\end{diagrambox}
All four versions agree with each other on $20\,000$ random grids up to $7\times7$ (we ran them).
\end{alternativebox}

\begin{tricksbox}
\begin{itemize}[leftmargin=*]
    \item \textbf{Bigger grids:} with $RC$ up to $10^6$ replace the recursion by the explicit stack of Problem 1.3.1; the area is simply the number of pops.
    \item \textbf{Keep the input:} if the caller still needs \texttt{grid}, sink with the value $2$ and turn all $2$s back into $1$ at the end, or use \texttt{seen} as in the BFS version.
    \item \textbf{Early stop:} once \texttt{best} is at least the number of land cells not yet sunk, no later island can beat it --- rarely needed, but a cheap bound when the grid is huge and mostly one island.
\end{itemize}
\end{tricksbox}

\subsubsection*{5. Trace}

\begin{tracebox}[{Example 1 --- the scan and the fill of island E}]
\begin{center}
\small
\begin{tabular}{@{}clcc@{}}
\toprule
\textbf{island} & \textbf{scan meets} & \textbf{\texttt{sink} returns} & \textbf{\texttt{best} after} \\
\midrule
A & $(0,2)$ & 1 & 1 \\
B & $(0,7)$ & 4 & 4 \\
C & $(2,1)$ & 4 & 4 \\
D & $(2,4)$ & 5 & 5 \\
E & $(3,8)$ & 6 & \textbf{6} \\
F & $(6,7)$ & 5 & 6 \\
\bottomrule
\end{tabular}
\end{center}

\begin{center}
\begin{tikzpicture}
    \def\gIpIIIs{0.5}
    \begin{scope}
        \gIpIIIgrid{wwkwwwwkwwwww,wwwwwwwkkkwww,wkkwkwwwwwwww,wkwwkkwwxwoww,wkwwkkwwoooww,wwwwwwwwwwoww,wwwwwwwooowww,wwwwwwwoowwww}
        \node[font=\scriptsize\bfseries, align=center] at (3.0,-4.3) {the scan reaches $(3,8)$ (orange); A--D are\\already sunk (light blue = now water)};
    \end{scope}
    \begin{scope}[xshift=8.6cm]
        \node[gvsrc, minimum size=9mm, font=\tiny] (s1) at (0,0) {3,8};
        \node[gv, minimum size=9mm, font=\tiny] (s2) at (0,-1.15) {4,8};
        \node[gv, minimum size=9mm, font=\tiny] (s3) at (0,-2.3) {4,9};
        \node[gv, minimum size=9mm, font=\tiny] (s4) at (0,-3.45) {4,10};
        \node[gv, minimum size=9mm, font=\tiny] (s5) at (-1.0,-4.6) {3,10};
        \node[gv, minimum size=9mm, font=\tiny] (s6) at (1.0,-4.6) {5,10};
        \draw[gtree] (s1) -- node[lbl, left]{down} (s2);
        \draw[gtree] (s2) -- node[lbl, left]{right} (s3);
        \draw[gtree] (s3) -- node[lbl, left]{right} (s4);
        \draw[gtree] (s4) -- node[lbl, left]{up} (s5);
        \draw[gtree] (s4) -- node[lbl, right]{down} (s6);
        \foreach \n/\v in {s1/6, s2/5, s3/4, s4/3, s5/1, s6/1}
            \node[gdist, anchor=west] at ([xshift=1pt]\n.east) {$\to$ \v};
        \node[font=\scriptsize\bfseries, align=center] at (0.2,-5.5) {recursion tree of \texttt{sink(3,8)}};
    \end{scope}
\end{tikzpicture}
\end{center}
The call \texttt{sink(3,8)} goes down to $(4,8)$, right twice to $(4,10)$, and there splits: up to $(3,10)$ and down to $(5,10)$. The leaves return 1; on the way back every call adds 1 for itself: $1+1+1 = 3$, then 4, 5 and finally \textbf{6}. The maximum depth is 5 frames. The diagonal neighbour $(6,9)$ of $(5,10)$ is never called --- only side neighbours are --- so island F is measured separately (5) and the answer stays \textbf{6}.
\end{tracebox}

\subsubsection*{6. Test suite}

\begin{testbox}
\centering\small
\begin{tabular}{@{}lll@{}}
\toprule
\textbf{grid} & \textbf{Expected} & \textbf{Why} \\
\midrule
Example 1 & 6 & a diagonal contact must not merge E and F \\
\texttt{[[0,0,0,0,0,0,0,0]]} & 0 & no island \\
\texttt{[[1]]} & 1 & single cell \\
\texttt{[[1,1],[1,1]]} & 4 & the mark-on-pop bug returns 5 here \\
\texttt{[[1,0,1],[0,1,0],[1,0,1]]} & 1 & five islands touching at corners only \\
\texttt{[[1,1,0],[1,1,0],[0,0,1]]} & 4 & a square and a lone cell touching at a corner \\
$50\times50$ all \texttt{1} & 2500 & largest island, deepest recursion \\
$50\times50$ checkerboard & 1 & 1250 islands of area 1 \\
\bottomrule
\end{tabular}
\end{testbox}

\subsubsection*{7. Common mistakes}

\begin{warningbox}[Mistake 1 --- sinking after the recursive calls]
\begin{lstlisting}[style=snippet]
int a = 1 + sink(g, r - 1, c) + sink(g, r + 1, c) + ...;
g[r][c] = 0;                         // WRONG: too late, the neighbours call back into (r,c)
\end{lstlisting}
The neighbour's call still sees $(r,c)$ as land and recurses back: infinite recursion. Sink the cell \emph{first}.
\end{warningbox}

\begin{warningbox}[Mistake 2 --- forgetting the ``no island'' case]
Initialise \texttt{best = 0}, not with \texttt{INT\_MIN} or with the first island's area: the second example has no land and expects 0.
\end{warningbox}

\begin{warningbox}[Mistake 3 --- counting diagonal neighbours]
With 8 directions Example 1 returns 11 (E and F merge). The statement says ``horizontal or vertical''.
\end{warningbox}

\begin{warningbox}[Mistake 4 --- a global counter that is not reset]
A popular style keeps \texttt{int area} as a class member and increments it inside the fill. Forgetting \texttt{area = 0} before each new island makes every island ``inherit'' the previous ones: the answer becomes the total amount of land. Returning the area from the function, as above, makes this mistake impossible.
\end{warningbox}

\subsubsection*{8. Takeaways}

\begin{itemize}[leftmargin=*]
    \item Size of a component = number of cells a fill marks; a recursive fill can \emph{return} it as $1 + \sum$ children.
    \item Sink a cell \emph{before} recursing; check bounds before reading.
    \item Small grids ($\le 2500$ cells): recursion is fine. Large grids: explicit stack.
    \item Count, size, sum, maximum --- all are ``scan + fill + one number per fill''.
\end{itemize}

\begin{variantsbox}
\variant{Flip one water cell}{\varin{https://leetcode.com/problems/making-a-large-island/}}
  {you may turn at most one $0$ into $1$; maximise the largest island.}
  {label every island with an id and store its area; a water cell joins its distinct neighbouring islands.}
  {for each $0$: $1 + $ the areas of its \emph{distinct} neighbouring ids; answer = max with the largest island (Exercise E3.5).}
\variant{Fish instead of area}{\exmedium\ \lc{Medium}\ \varlink{Maximum Number of Fish in a Grid}{LeetCode 2658}{https://leetcode.com/problems/maximum-number-of-fish-in-a-grid/}}
  {water cells hold fish; collect the most fish in one connected water region.}
  {vertices = cells with a positive value; same 4-neighbour edges.}
  {the fill returns the sum of the values instead of the number of cells.}
\variant{Perimeter instead of area}{\varin{https://leetcode.com/problems/island-perimeter/}}
  {measure the boundary of the (single) island.}
  {count edges between a land cell and water or the outside, instead of vertices.}
  {$4\cdot\#\text{land} - 2\cdot\#\text{land--land edges}$, no traversal needed (Exercise E3.1).}
\variant{Histogram of island sizes}{\varnone}
  {report how many islands have each area $1, 2, \dots$}
  {same components.}
  {\texttt{cnt[sink(r,c)]++} for every start of the scan; $\bigO{RC}$.}
\variant{Largest 8-connected island}{\varnone}
  {cells touching at a corner belong to the same island.}
  {edges also between diagonal neighbours.}
  {8 recursive calls instead of 4; Example 1 then returns 11.}
\end{variantsbox}

%============================================================
% CH1 · PATTERN 3 — EXERCISES, HINTS, ANSWERS
%============================================================
\subsection{Pattern 3: Exercises and Going Further}
\label{ssec:ch1-p3-ex}

\begin{patternbox}{Pattern 3 Summary Card}
\begin{tabularx}{\linewidth}{@{}>{\bfseries}l X@{}}
Model & open cell = vertex; two open cells sharing a side = edge (8 directions if corners count); the graph stays implicit \\
Algorithm & scan in row-major order; an unvisited open cell starts a flood fill (stack, queue or recursion) that marks its whole component \\
Answers & \#fills = \#components; pops per fill = size; sum of values per fill = volume \\
Cost & $\bigO{nm}$ time and memory; $|E| \le 2nm$ \\
Top traps & bounds after access; marking on pop; recursion depth $nm$; 4 vs 8 neighbours; \texttt{r*n+c}; recolouring with the same colour \\
\end{tabularx}
\end{patternbox}

\exlegend
\begin{enumerate}[leftmargin=*, label=\textbf{E3.\arabic*}, itemsep=6pt]
    \item \exeasy\ \pconcept \exprob{Island Perimeter}{LeetCode 463}{https://leetcode.com/problems/island-perimeter/}\ \lc{Easy}. A $0/1$ grid ($\le 100\times100$) contains exactly one island (4-connected, no lakes inside). Return the length of its boundary, every cell being a $1\times1$ square. (\texttt{[[0,1,0,0],[1,1,1,0],[0,1,0,0],[1,1,0,0]]} $\to 16$.)
    \item \exeasy\ \ptricky \exprob{Flood Fill}{LeetCode 733}{https://leetcode.com/problems/flood-fill/}\ \lc{Easy}. Recolour the 4-connected region of pixels that have the same colour as \texttt{image[sr][sc]} and contain it, with a new \texttt{color}; return the image. (\texttt{[[1,1,1],[1,1,0],[1,0,1]]}, $(1,1)$, $2$ $\to$ \texttt{[[2,2,2],[2,2,0],[2,0,1]]}.)
    \item \exeasy\ \pconcept \cf{1100}\\exprob{The Lakes}{Codeforces 1829E}{https://codeforces.com/problemset/problem/1829/E}. $t$ test cases; each is an $n\times m$ grid of depths $0 \le a_{ij} \le 1000$ with $\sum nm \le 10^6$. A lake is a 4-connected component of cells with positive depth, its volume is the sum of the depths. Print the largest volume (0 if there is no lake).
    \item \exmedium\ \ptricky \exprob{Count Sub Islands}{LeetCode 1905}{https://leetcode.com/problems/count-sub-islands/}\ \lc{Medium}. Two $0/1$ grids of the same size ($\le 500\times500$). An island of \texttt{grid2} is a \emph{sub-island} if every one of its cells is land in \texttt{grid1}. Count the sub-islands.
    \item \exhard\ \pclassic \exprob{Making A Large Island}{LeetCode 827}{https://leetcode.com/problems/making-a-large-island/}\ \lc{Hard}. In an $n\times n$ $0/1$ grid ($n \le 500$) you may change at most one $0$ into $1$. Return the largest possible island area. (\texttt{[[1,0],[0,1]]} $\to 3$.)
    \item \exmedium\ \pconcept \exown{Four or eight neighbours?} For the grid below, count the islands under 4-connectivity and under 8-connectivity, by hand, then with code. Which single direction-array change turns one program into the other?
    \begin{center}\ttfamily 1 0 0 1 \quad / \quad 0 1 1 0 \quad / \quad 0 1 0 0 \quad / \quad 1 0 0 1\end{center}
    \item \exmedium\ \precent \exprob{Count Islands With Total Value Divisible by K}{LeetCode 3619}{https://leetcode.com/problems/count-islands-with-total-value-divisible-by-k/}\ \lc{Medium}. An $m\times n$ grid of values $0 \le v \le 10^6$ ($mn \le 10^5$); an island is a 4-connected group of positive cells, and its total value is the sum of its cells. Count the islands whose total value is divisible by $k$ ($1 \le k \le 10^6$).
\end{enumerate}

\subsubsection*{Hints}
\begin{enumerate}[leftmargin=*, label=\textbf{E3.\arabic*}]
    \item Every land cell brings 4 sides; every edge between two land cells hides 2 of them. No traversal needed.
    \item ``Visited'' is ``already has the new colour'' --- but what if the new colour equals the old one?
    \item Problem 1.3.2 with the depth added instead of 1. Watch the total size of all tests and the maximum sum ($10^6\cdot1000$).
    \item Fill each island of \texttt{grid2} \emph{completely}, even after you have found a cell that is water in \texttt{grid1}.
    \item Paint each island with its own id and remember its area. A $0$ cell touches at most 4 islands --- but maybe the same island twice.
    \item Draw the corner contacts. The 4-direction array is a prefix of the 8-direction array.
    \item The Lakes (E3.3) again, but you test each sum instead of taking the maximum. How large can one sum get?
\end{enumerate}

\subsubsection*{Answers}

\begin{solutionbox}[E3.1 Island Perimeter]
Count land cells $L$ and land--land edges $S$ (look only down and right, so each edge is counted once). Perimeter $= 4L - 2S$.
\begin{lstlisting}[style=cpp]
#include <bits/stdc++.h>
using namespace std;

class Solution {
public:
    int islandPerimeter(vector<vector<int>>& grid) {
        int R = grid.size(), C = grid[0].size(), land = 0, shared = 0;
        for (int r = 0; r < R; r++)
            for (int c = 0; c < C; c++) {
                if (grid[r][c] != 1) continue;
                land++;
                if (r + 1 < R && grid[r + 1][c] == 1) shared++;   // edge below
                if (c + 1 < C && grid[r][c + 1] == 1) shared++;   // edge right
            }
        return 4 * land - 2 * shared;
    }
};
\end{lstlisting}
\begin{center}
\begin{tikzpicture}
    \def\gIpIIIs{0.55}
    \gIpIIIgrid{waww,aaaw,waww,aaww}
    \node[note, anchor=west, text width=6.5cm, align=left] at (2.3,-0.8) {Example 1: $L = 7$ land cells, $S = 6$ shared edges ($(0,1){-}(1,1)$, $(1,0){-}(1,1)$, $(1,1){-}(1,2)$, $(1,1){-}(2,1)$, $(2,1){-}(3,1)$, $(3,0){-}(3,1)$): $28 - 12 = 16$. The other examples give 4 and 4.};
\end{tikzpicture}
\end{center}
This is a pure counting argument on the grid graph (Pattern 1 thinking): no fill at all, because the island is the only one. The ``no lakes'' promise matters: with a lake, its shore would also count as perimeter --- and the formula would \emph{still} give the total boundary length.
\end{solutionbox}

\begin{solutionbox}[E3.2 Flood Fill]
The region is the component of $(sr,sc)$ in the graph ``same-coloured neighbours''. Recolour while filling; a pixel that already has the new colour is never entered again. If the old colour equals the new one, return at once --- otherwise the recursion would bounce between two pixels forever.
\begin{lstlisting}[style=cpp]
#include <bits/stdc++.h>
using namespace std;

class Solution {
    void fill(vector<vector<int>>& img, int r, int c, int from, int to) {
        if (r < 0 || r >= (int)img.size() || c < 0 || c >= (int)img[0].size()) return;
        if (img[r][c] != from) return;
        img[r][c] = to;
        fill(img, r - 1, c, from, to); fill(img, r + 1, c, from, to);
        fill(img, r, c - 1, from, to); fill(img, r, c + 1, from, to);
    }
public:
    vector<vector<int>> floodFill(vector<vector<int>>& image, int sr, int sc, int color) {
        int from = image[sr][sc];
        if (from != color) fill(image, sr, sc, from, color);   // equal colours: nothing to do
        return image;
    }
};
\end{lstlisting}
Example 1 gives \texttt{[[2,2,2],[2,2,0],[2,0,1]]} (the bottom-right 1 touches the region only at a corner); Example 2 (\texttt{color} equal to the old colour 0) returns the image unchanged. The grid is at most $50\times50$, so recursion is safe. Checked against a BFS recolouring on 5000 random images.
\end{solutionbox}

\begin{solutionbox}[E3.3 The Lakes]
Scan; a positive cell starts a fill that adds up the depths of the cells it reaches and sets them to 0. With $\sum nm \le 10^6$ use an explicit stack (an open $1000\times1000$ lake would recurse $10^6$ deep) and allocate the grid per test. The largest volume is at most $10^6 \cdot 1000 = 10^9$, which still fits in \texttt{int}; \texttt{long long} costs nothing and removes the doubt.
\begin{lstlisting}[style=cpp]
#include <bits/stdc++.h>
using namespace std;

int main() {
    ios::sync_with_stdio(false);
    cin.tie(nullptr);
    int t;
    cin >> t;
    while (t--) {
        int n, m;
        cin >> n >> m;
        vector<vector<int>> a(n, vector<int>(m));
        for (auto& row : a) for (auto& x : row) cin >> x;
        const int dr[4] = {-1, 1, 0, 0}, dc[4] = {0, 0, -1, 1};
        long long best = 0;                        // no lake at all -> 0
        vector<int> st;
        for (int r = 0; r < n; r++)
            for (int c = 0; c < m; c++) {
                if (a[r][c] == 0) continue;
                long long volume = 0;
                st.push_back(r * m + c);
                volume += a[r][c];
                a[r][c] = 0;                       // visited = depth set to 0
                while (!st.empty()) {
                    int x = st.back() / m, y = st.back() % m;
                    st.pop_back();
                    for (int d = 0; d < 4; d++) {
                        int nx = x + dr[d], ny = y + dc[d];
                        if (nx < 0 || nx >= n || ny < 0 || ny >= m || a[nx][ny] == 0) continue;
                        volume += a[nx][ny];
                        a[nx][ny] = 0;
                        st.push_back(nx * m + ny);
                    }
                }
                best = max(best, volume);
            }
        cout << best << "\n";
    }
}
\end{lstlisting}
On the five official samples the program prints 10, 0, 7, 16, 21. In the fourth sample the ring of 1s around the 5 is one lake of volume 16, and the 5 is a separate lake (its neighbours are 0): $\max(16,5) = 16$. In the fifth the 4 touches the ring, so everything is one lake: $17 + 4 = 21$.
\end{solutionbox}

\begin{solutionbox}[E3.4 Count Sub Islands]
Fill every island of \texttt{grid2}; while filling, AND a flag ``this cell is land in \texttt{grid1}''. Do \textbf{not} stop the fill when the flag becomes false: the rest of the island must still be sunk, otherwise its remaining cells would later be counted as a separate (wrong) island.
\begin{lstlisting}[style=cpp]
#include <bits/stdc++.h>
using namespace std;

class Solution {
public:
    int countSubIslands(vector<vector<int>>& grid1, vector<vector<int>>& grid2) {
        int R = grid2.size(), C = grid2[0].size(), cnt = 0;
        const int dr[4] = {-1, 1, 0, 0}, dc[4] = {0, 0, -1, 1};
        for (int r = 0; r < R; r++)
            for (int c = 0; c < C; c++) {
                if (grid2[r][c] != 1) continue;
                bool inside = true;                    // every cell so far is land in grid1
                vector<pair<int,int>> st(1, make_pair(r, c));
                grid2[r][c] = 0;                       // sink grid2's island as we go
                while (!st.empty()) {
                    int x = st.back().first, y = st.back().second;
                    st.pop_back();
                    if (grid1[x][y] == 0) inside = false;   // do NOT stop: finish the fill
                    for (int d = 0; d < 4; d++) {
                        int nx = x + dr[d], ny = y + dc[d];
                        if (nx < 0 || nx >= R || ny < 0 || ny >= C || grid2[nx][ny] != 1) continue;
                        grid2[nx][ny] = 0;
                        st.push_back(make_pair(nx, ny));
                    }
                }
                if (inside) cnt++;
            }
        return cnt;
    }
};
\end{lstlisting}
Why is ``all cells are land in \texttt{grid1}'' enough --- must they lie in the \emph{same} island of \texttt{grid1}? Yes, automatically: the island of \texttt{grid2} is connected, and if all its cells are land in \texttt{grid1}, the same paths connect them in \texttt{grid1}. Examples: 3 and 2. Checked against a brute force (label both grids, compare label sets) on 5000 random pairs.
\end{solutionbox}

\begin{solutionbox}[E3.5 Making A Large Island]
\textbf{Step 1:} paint every island with its own id $2, 3, \dots$ (ids 0 and 1 are taken) and store \texttt{area[id]}. \textbf{Step 2:} for every $0$ cell, collect the \emph{distinct} ids among its 4 neighbours; flipping it gives $1 + \sum$ their areas. The answer is the maximum of these values and of the largest existing island (the grid may contain no $0$ at all).
\begin{lstlisting}[style=cpp]
#include <bits/stdc++.h>
using namespace std;

class Solution {
public:
    int largestIsland(vector<vector<int>>& grid) {
        int n = grid.size();
        const int dr[4] = {-1, 1, 0, 0}, dc[4] = {0, 0, -1, 1};
        vector<int> area(2, 0);                    // area[id] for island ids 2, 3, ...
        for (int r = 0; r < n; r++)
            for (int c = 0; c < n; c++) {
                if (grid[r][c] != 1) continue;
                int id = area.size(), a = 0;       // paint this island with colour id
                vector<int> st(1, r * n + c);
                grid[r][c] = id;
                while (!st.empty()) {
                    int x = st.back() / n, y = st.back() % n;
                    st.pop_back();
                    a++;
                    for (int d = 0; d < 4; d++) {
                        int nx = x + dr[d], ny = y + dc[d];
                        if (nx < 0 || nx >= n || ny < 0 || ny >= n || grid[nx][ny] != 1) continue;
                        grid[nx][ny] = id;
                        st.push_back(nx * n + ny);
                    }
                }
                area.push_back(a);
            }
        int best = 0;
        for (size_t id = 2; id < area.size(); id++) best = max(best, area[id]);   // no flip needed
        for (int r = 0; r < n; r++)
            for (int c = 0; c < n; c++) {
                if (grid[r][c] != 0) continue;
                int ids[4], k = 0, total = 1;      // the flipped cell itself
                for (int d = 0; d < 4; d++) {
                    int nr = r + dr[d], nc = c + dc[d];
                    if (nr < 0 || nr >= n || nc < 0 || nc >= n || grid[nr][nc] == 0) continue;
                    int id = grid[nr][nc];
                    bool dup = false;              // the same island may touch twice
                    for (int j = 0; j < k; j++) if (ids[j] == id) dup = true;
                    if (!dup) { ids[k++] = id; total += area[id]; }
                }
                best = max(best, total);
            }
        return best;
    }
};
\end{lstlisting}
\begin{center}
\begin{tikzpicture}
    \def\gIpIIIs{0.55}
    \begin{scope}
        \gIpIIIgrid{aaa,axa,aaa}
        \node[font=\scriptsize\bfseries, align=center] at (0.55,-1.65) {ring: one island (8)\\touches the hole 4 times};
    \end{scope}
    \begin{scope}[xshift=4.2cm]
        \gIpIIIgrid{awb,wxw,cwd}
        \node[font=\scriptsize\bfseries, align=center] at (0.55,-1.65) {four 1-cell islands, but the\\centre touches none of them};
    \end{scope}
\end{tikzpicture}
\end{center}
The ring ($3\times3$ of 1s with a 0 in the middle) must give $1 + 8 = 9$, not $1 + 4\cdot8 = 33$ --- that is what the \texttt{dup} test is for. In the second picture the centre has only water neighbours, so the best flip is next to a corner island: $1+1+1 = 3$. Official examples: 3, 4, 4. All verified against a brute force that tries every flip and recomputes the largest island, on 5000 random grids. Total $\bigO{n^2}$.
\end{solutionbox}

\begin{solutionbox}[E3.6 Four or eight neighbours?]
\begin{center}
\begin{tikzpicture}
    \def\gIpIIIs{0.55}
    \begin{scope}
        \gIpIIIgrid{AwwB,wCCw,wCww,Dwwe}
        \gIpIIIat{3}{3}{E}
        \node[font=\scriptsize\bfseries] at (0.82,-2.4) {4-connected: 5 islands};
    \end{scope}
    \begin{scope}[xshift=4.2cm]
        \gIpIIIgrid{Awwa,wAAw,wAww,awwb}
        \gIpIIIat{0}{3}{A}\gIpIIIat{3}{0}{A}\gIpIIIat{3}{3}{B}
        \node[font=\scriptsize\bfseries] at (0.82,-2.4) {8-connected: 2 islands};
    \end{scope}
\end{tikzpicture}
\end{center}
Under 4-connectivity: the three middle cells form one island and the four corners are alone: \textbf{5}. Under 8-connectivity the corners $(0,0)$, $(0,3)$ and $(3,0)$ each touch the middle island diagonally, but $(3,3)$ touches only water: \textbf{2}. (Our program prints 5 and 2.) With the arrays
\texttt{dr = \{-1,1,0,0,-1,-1,1,1\}}, \texttt{dc = \{0,0,-1,1,-1,1,-1,1\}}, the loop \texttt{for (d = 0; d < K; d++)} gives 4-connectivity for $K = 4$ and 8-connectivity for $K = 8$ --- one number is the whole difference.
\end{solutionbox}

\begin{solutionbox}[E3.7 Count Islands With Total Value Divisible by K]
Scan; every positive cell starts a fill that adds up the values and sets the cells to $0$; after the fill, test \texttt{sum \% k == 0}. The trap is the size of the sum: one island can hold $10^5$ cells of value $10^6$, a total of $10^{11}$, which overflows a 32-bit \texttt{int}. Use \texttt{long long} (or keep only \texttt{sum \% k}). The fill uses an explicit stack, because an island can have $10^5$ cells.
\begin{lstlisting}[style=cpp]
#include <bits/stdc++.h>
using namespace std;

class Solution {
public:
    int countIslands(vector<vector<int>>& grid, int k) {
        int R = grid.size(), C = grid[0].size(), cnt = 0;
        const int dr[4] = {-1, 1, 0, 0}, dc[4] = {0, 0, -1, 1};
        vector<int> st;
        for (int r = 0; r < R; r++)
            for (int c = 0; c < C; c++) {
                if (grid[r][c] == 0) continue;
                long long sum = grid[r][c];            // up to 1e5 * 1e6: needs 64 bits
                grid[r][c] = 0;                        // visited = value set to 0
                st.push_back(r * C + c);
                while (!st.empty()) {
                    int x = st.back() / C, y = st.back() % C;
                    st.pop_back();
                    for (int d = 0; d < 4; d++) {
                        int nx = x + dr[d], ny = y + dc[d];
                        if (nx < 0 || nx >= R || ny < 0 || ny >= C || grid[nx][ny] == 0) continue;
                        sum += grid[nx][ny];
                        grid[nx][ny] = 0;
                        st.push_back(nx * C + ny);
                    }
                }
                if (sum % k == 0) cnt++;
            }
        return cnt;
    }
};
\end{lstlisting}
Example 1 ($k = 5$) has four islands with sums $2+1+5 = 8$, $5$, $1+1+4+7+2 = 15$ and $8$; two of them are divisible by 5, so the answer is \textbf{2}. In Example 2 ($k = 3$) there are six single cells of value 3, all divisible: \textbf{6}. Checked against a brute force (DSU labels, then the sum per label) on 5000 random grids, and on one $1\times10^5$ island of $10^6$s (sum $10^{11}$).
\end{solutionbox}


%============================================================
\clearpage
\section{Pattern 4: Implicit State Graphs}
\label{sec:ch1-p4}
%============================================================
%============================================================
% CH1 · PATTERN 4 — IMPLICIT STATE GRAPHS: THEORY
%============================================================
\subsection{Pattern Theory}
\label{ssec:ch1-p4-theory}

\begin{patternbox}{Pattern 4 at a Glance}
The input does \textbf{not} contain a graph. It describes a \textbf{situation} (a number on a display, the digits of a lock, the water in two jugs) and the \textbf{moves} that change it. Every possible situation is a \textbf{vertex} (a \emph{state}), every move is an \textbf{edge}. ``Fewest moves'' $=$ shortest path, and because every move costs one, \textbf{BFS} finds it --- generating the neighbours of a state on the fly instead of storing the graph.

\begin{center}
\begin{tikzpicture}[x=2.3cm, y=0.95cm]
    \node[font=\scriptsize\bfseries, text=darkgray] at (0,1.15) {dist 0};
    \node[font=\scriptsize\bfseries, text=darkgray] at (1,1.15) {dist 1};
    \node[font=\scriptsize\bfseries, text=darkgray] at (2,1.15) {dist 2};
    \node[font=\scriptsize\bfseries, text=darkgray] at (3,1.15) {dist 3};
    \node[gvsrc] (s3) at (0,-1.5) {3};
    \node[gvdone] (s6) at (1,0) {6};
    \node[gvdone] (s2) at (1,-3) {2};
    \node[gvdone] (s12) at (2,0.2) {12};
    \node[gvdone] (s5) at (2,-1.0) {5};
    \node[gvdone] (s4) at (2,-2.2) {4};
    \node[gvdone] (s1) at (2,-3.4) {1};
    \node[gvseen] (s11) at (3,0.2) {11};
    \node[gvdst] (s10) at (3,-1.0) {10};
    \node[gvseen] (s8) at (3,-2.2) {8};
    \draw[gtree, -{Stealth[length=2.4mm]}] (s3) -- node[gw, above, sloped]{$\times2$} (s6);
    \draw[gtree, -{Stealth[length=2.4mm]}] (s3) -- node[gw, below, sloped]{$-1$} (s2);
    \draw[gtree, -{Stealth[length=2.4mm]}] (s6) -- (s12);
    \draw[gtree, -{Stealth[length=2.4mm]}] (s6) -- (s5);
    \draw[gtree, -{Stealth[length=2.4mm]}] (s2) -- (s4);
    \draw[gtree, -{Stealth[length=2.4mm]}] (s2) -- (s1);
    \draw[gtree, -{Stealth[length=2.4mm]}] (s12) -- (s11);
    \draw[gpath, -{Stealth[length=2.6mm]}] (s5) -- node[gw, above]{$\times2$} (s10);
    \draw[gtree, -{Stealth[length=2.4mm]}] (s4) -- (s8);
    \draw[gfaded, dashed, -{Stealth[length=2mm]}] (s5) -- (s4);
    \node[note, text width=3.3cm, anchor=west] at (3.35,-1.9) {moves: $\times2$ and $-1$\\start $3$, goal $10$\\answer $=$ layer of the goal $=3$\\(dashed: $5\to4$ finds an\\already discovered state)};
\end{tikzpicture}
\end{center}

\begin{lstlisting}[style=snippet]
dist[start] = 0;  queue = {start};               // layer 0
while (!queue.empty()) {
    s = pop_front();  if (s == goal) return dist[s];
    for (each move m)  { t = m(s);                // neighbour generated, not stored
        if (valid(t) && dist[t] == -1) { dist[t] = dist[s] + 1; push_back(t); } }
}
return -1;                                       // goal not reachable
\end{lstlisting}
Time $\bigO{\#\text{states} \times \#\text{moves}}$ --- so the first job is always to \textbf{bound the number of states}.
\end{patternbox}

%------------------------------------------------------------
\subsubsection*{From Pattern 3 to Pattern 4}

In Pattern~3 the graph was hidden in a grid, but at least you could \emph{see} it: the cells were in the input and the neighbours were the four cells around. Here even that disappears. In Two Buttons (Problem~1.4.1) the input is just two numbers, $n$ and $m$; in Open the Lock (Problem~1.4.2) it is a target code and a list of forbidden codes. Nobody hands you vertices or edges --- you have to \emph{invent} them.

\begin{diagrambox}
\begin{tikzpicture}
    % left: explicit graph
    \node[font=\small\bfseries, text=deepblue] at (1.5,2.2) {Patterns 1--3: the graph is given};
    \node[gv] (a) at (0,1) {1};
    \node[gv] (b) at (1.5,1.5) {2};
    \node[gv] (c) at (3,1) {3};
    \node[gv] (d) at (1.5,0) {4};
    \draw[ge] (a)--(b) (b)--(c) (a)--(d) (d)--(c);
    \node[draw=darkgray!60, rounded corners=2pt, font=\scriptsize\ttfamily, align=left, fill=lightgray, anchor=north]
        at (1.5,-0.6) {1 2\\2 3\\1 4\\4 3};
    \node[note] at (1.5,-2.3) {edges are read from the input\\and stored in \texttt{adj}};
    % right: implicit graph
    \node[font=\small\bfseries, text=deepblue] at (8.2,2.2) {Pattern 4: the graph is a \emph{rule}};
    \node[gvcur, minimum size=9mm] (s) at (6.4,0.5) {$s$};
    \node[tag=accentorange, text width=2.5cm] (f) at (8.4,0.5) {\scriptsize\texttt{neighbours(s)}\\[-1pt]\scriptsize apply every move};
    \node[gvseen, minimum size=9mm, font=\scriptsize\bfseries] (t1) at (10.6,1.5) {$2s$};
    \node[gvseen, minimum size=9mm, font=\scriptsize\bfseries] (t2) at (10.6,-0.5) {$s{-}1$};
    \draw[flow] (s) -- (f);
    \draw[flow] (f) -- (t1);
    \draw[flow] (f) -- (t2);
    \node[note] at (8.4,-1.6) {nothing is stored: the edges of $s$\\are \emph{computed} when BFS reaches $s$};
\end{tikzpicture}
\figcap{An explicit graph is data; an implicit graph is a function. BFS only ever asks ``what are the neighbours of this vertex?'' --- and a function can answer that as well as an adjacency list.}
\end{diagrambox}

This is the most important modelling skill of the whole book: whenever a problem says ``minimum number of operations / moves / presses / steps to turn $A$ into $B$'', ask yourself \emph{what the situation between two moves looks like}. That snapshot is a vertex.

%------------------------------------------------------------
\subsubsection*{The three questions}

Before writing any code, answer three questions on paper. They are the Graph Model of Pattern~4.

\begin{diagrambox}
\begin{tikzpicture}[node distance=5mm]
    \node[fcproc=3.3cm] (q1) {\textbf{1. State}\\what must I know\\between two moves?};
    \node[fcproc=3.3cm, right=of q1] (q2) {\textbf{2. Moves}\\which states can one\\move reach? (edges)};
    \node[fcproc=3.3cm, right=of q2] (q3) {\textbf{3. Size}\\how many states\\can ever matter?};
    \node[fcout=2.2cm, right=of q3] (bfs) {BFS\\from start};
    \draw[fcarrow] (q1) -- (q2);
    \draw[fcarrow] (q2) -- (q3);
    \draw[fcarrow] (q3) -- (bfs);
    \node[note, below=2mm of q1, text width=3.3cm] {display: one number\\lock: 4 digits\\jugs: a pair $(a,b)$};
    \node[note, below=2mm of q2, text width=3.3cm] {$\times2$, $-1$\\turn one wheel $\pm1$\\fill / empty / pour};
    \node[note, below=2mm of q3, text width=3.3cm] {$\le 2\cdot10^4$\\$10^4$\\$(x{+}1)(y{+}1)$};
\end{tikzpicture}
\figcap{The three questions, with the answers for the two solved problems and for the water jugs of Exercise~E4.3.}
\end{diagrambox}

\begin{enumerate}[leftmargin=*]
    \item \textbf{State} --- the \emph{smallest} description of the situation that decides which moves are possible and whether we have arrived. Too little information and two different situations get merged (wrong answers); too much and the state space explodes (TLE). The \emph{history} of how you got there is \textbf{not} part of the state --- BFS remembers the number of moves for you.
    \item \textbf{Moves} --- every move of the statement becomes one line of code that turns $s$ into a neighbour $t$, plus a validity test (``the display must stay positive'', ``a dead-end code may not be entered'').
    \item \textbf{Size} --- multiply the ranges of the components (4 digits: $10^4$; two jugs: $(x+1)(y+1)$). If a component is unbounded (a number that can grow forever), you must \emph{prove} a window outside which states are useless. Then $\#\text{states} \times \#\text{moves}$ must fit in the time limit (about $10^8$ simple steps per second).
\end{enumerate}

%------------------------------------------------------------
\subsubsection*{Turning a state into an array index}

BFS needs a \texttt{dist} entry for each state. If the state is a small integer, it \emph{is} the index. Otherwise \textbf{encode} it as an integer --- a mixed-radix number --- so that \texttt{dist} can be a plain vector instead of a slow \texttt{map}.

\begin{diagrambox}
\begin{tikzpicture}
    % lock code
    \node[font=\small\bfseries, text=deepblue] at (1.4,1.15) {lock wheels};
    \foreach \d [count=\i from 0] in {0,2,0,2}
        \node[gcell, minimum size=8mm, font=\large\bfseries] (w\i) at (\i*0.85,0) {\d};
    \foreach \p [count=\i from 0] in {1000,100,10,1}
        \node[font=\scriptsize, text=darkgray] at (\i*0.85,-0.65) {$\times\p$};
    \draw[flow] (3.2,0) -- (4.2,0);
    \node[tag=successgreen] at (5.35,0) {\texttt{id} $=202$};
    \node[note] at (2.6,-1.35) {4 digits $\Rightarrow$ ids $0..9999$};
    % jug pair
    \node[font=\small\bfseries, text=deepblue] at (9.3,1.15) {pair $(a,b)$, $0\le b\le y$};
    \node[tag=primaryblue] (p) at (8.3,0) {$(a,b)=(2,4)$};
    \draw[flow] (9.4,0) -- (10.3,0);
    \node[tag=successgreen, align=center] at (11.75,0) {\texttt{id} $= a\cdot(y{+}1)+b$\\$=2\cdot6+4=16$};
    \node[note] at (10.2,-1.35) {$y=5$: ids $0..(x{+}1)(y{+}1){-}1$; decode:\\$a=\texttt{id}/(y{+}1)$, \ $b=\texttt{id}\bmod(y{+}1)$};
\end{tikzpicture}
\figcap{Two encodings used in this pattern. Any tuple $(c_1,\dots,c_k)$ with $0 \le c_i < R_i$ becomes $(\dots(c_1 R_2 + c_2)R_3 + \dots)R_k + c_k$; a tuple of booleans becomes a bitmask.}
\end{diagrambox}

%------------------------------------------------------------
\subsubsection*{Why BFS gives the fewest moves}

BFS explores the states in \textbf{layers}: first the start (0 moves), then everything one move away, then everything two moves away, and so on --- like ripples on water. The first time BFS discovers a state is therefore along a shortest sequence of moves.

\begin{diagrambox}
\begin{tikzpicture}
    \foreach \r/\c in {3.0/primaryblue!6, 2.1/primaryblue!12, 1.2/primaryblue!20}
        \fill[\c] (0,0) circle (\r);
    \foreach \r in {1.2, 2.1, 3.0} \draw[primaryblue!50, dashed] (0,0) circle (\r);
    \node[gvsrc] at (0,0) {$s$};
    \node[font=\scriptsize\bfseries, text=primaryblue] at (0,0.85) {dist 1};
    \node[font=\scriptsize\bfseries, text=primaryblue] at (0,1.75) {dist 2};
    \node[font=\scriptsize\bfseries, text=primaryblue] at (0,2.65) {dist 3};
    \node[gvdst, minimum size=6mm] (g) at (1.9,-1.6) {$g$};
    % queue picture
    \begin{scope}[xshift=4.6cm, yshift=0.3cm]
        \node[font=\small\bfseries, text=deepblue, anchor=west] at (0,1.2) {the queue at any moment};
        \foreach \v/\s [count=\i from 0] in {d/gcell, d/gcell, d/gcell, d{+}1/gcell, d{+}1/gcell}
            \node[\s, minimum width=1.1cm, font=\small] (q\i) at (\i*1.1+0.55,0.3) {$\v$};
        \node[font=\scriptsize, text=darkgray] at (0.55,-0.4) {front};
        \node[font=\scriptsize, text=darkgray] at (4.95,-0.4) {back};
        \node[note, anchor=west, text width=6cm, align=left] at (0,-1.5)
            {Distances in the queue are \emph{sorted} and differ by at most 1: some states at distance $d$, followed by some at $d+1$. Popping from the front always takes a closest unfinished state.};
    \end{scope}
\end{tikzpicture}
\figcap{Left: layers of a BFS. Right: the queue invariant that makes the layers come out in order.}
\end{diagrambox}

\begin{invariantbox}[{Invariant of BFS with unit-cost moves}]
\textbf{(I1)} At every moment the queue holds states with distances $d, \dots, d, d{+}1, \dots, d{+}1$ (non-decreasing, spread at most 1).
\textbf{(I2)} When a state $t$ gets $\texttt{dist}[t] = \texttt{dist}[s]+1$, this is the true minimum number of moves to $t$.

\emph{Proof sketch.} (I1) holds for the queue $\{\text{start}\}$. Popping a state at distance $d$ and pushing its new neighbours at $d{+}1$ to the back keeps the order and the spread. For (I2): suppose $t$ could be reached in $k < \texttt{dist}[s]+1$ moves, and let $p$ be the state just before $t$ on such a shorter route; $p$ is at distance $k-1 < \texttt{dist}[s]$. By (I1) $p$ was popped \emph{before} $s$, and at that moment it would have discovered $t$ first, with the smaller value --- contradiction. \hfill$\blacksquare$

Consequence: mark a state as discovered \textbf{when it is pushed}, not when it is popped; and the first time the goal is popped (or pushed), its distance is the answer.
\end{invariantbox}

%------------------------------------------------------------
\subsubsection*{Bounding an infinite state space}

A number that can be doubled forever gives infinitely many states, and BFS would never stop when the goal is unreachable. The fix is never ``pick a big constant and hope'': it is a short \textbf{proof} that leaving some window $[\mathrm{LO}, \mathrm{HI}]$ cannot help.

\begin{diagrambox}
\begin{tikzpicture}[x=1cm]
    \fill[recursionbg] (-0.4,-0.25) rectangle (1,0.25);
    \fill[solutionbg] (1,-0.25) rectangle (9,0.25);
    \fill[recursionbg] (9,-0.25) rectangle (12.4,0.25);
    \draw[thick, -{Stealth}] (-0.4,0) -- (12.5,0);
    \foreach \x/\l in {1/LO, 9/HI} {
        \draw[thick] (\x,-0.35) -- (\x,0.35);
        \node[font=\small\bfseries, text=deepblue] at (\x,0.65) {\l};
    }
    \node[font=\scriptsize, text=warningred, align=center] at (0.3,-0.75) {invalid\\(e.g.\ $\le 0$)};
    \node[font=\scriptsize, text=successgreen!70!black, align=center] at (5,-0.75) {every state an optimal solution can visit};
    \node[font=\scriptsize, text=warningred, align=center] at (10.7,-0.75) {provably useless\\(``too far to come back cheaply'')};
    \node[gvsrc, minimum size=6mm] at (3,0) {};
    \node[gvdst, minimum size=6mm] at (6.5,0) {};
\end{tikzpicture}
\figcap{BFS only on the green window. Two Buttons: HI $=\max(n, 2m)$ (Lemma in Problem~1.4.1). Exercise E4.5: only values in $[0,1000]$ can be continued, everything else is a leaf.}
\end{diagrambox}

Typical arguments:
\begin{itemize}[leftmargin=*]
    \item \textbf{Only-increasing moves} ($+3$, $\times2$, appending a digit): a value above the goal can never come back, so HI $=$ goal.
    \item \textbf{Overshoot costs more than it saves}: in Two Buttons, a doubling made at a value $\ge m$ adds at least $m$ that must later be removed one by one, so it is never optimal; hence HI $= \max(n, 2m)$.
    \item \textbf{Modular states}: ``$\bmod\ 32768$'' or a wheel with 10 digits makes the space finite automatically.
    \item \textbf{The statement does it for you}: ``operations are possible only while $0\le x\le 1000$''.
\end{itemize}

%------------------------------------------------------------
\subsubsection*{When to recognise it}

\begin{itemize}[leftmargin=*]
    \item ``\textbf{minimum number} of operations / moves / presses / turns / steps'' to go from a start configuration to a goal;
    \item a \textbf{small menu of moves}, each costing the same (one press, one turn, one pour);
    \item the configuration fits in a few small numbers, so the number of states is at most a few million;
    \item ``is it \textbf{possible} at all?'' --- the same BFS (or a DFS) answers reachability.
\end{itemize}
\textbf{Not this pattern (or not only):}
\begin{itemize}[leftmargin=*]
    \item moves have \textbf{different costs} --- BFS is wrong; use 0-1 BFS or Dijkstra (Chapters 7--8);
    \item the number of states is astronomically large (values up to $10^{18}$, permutations of 20 items) --- look for a mathematical shortcut or a greedy in the \emph{reverse} direction (Problem~1.4.1, Approach~3);
    \item you must \textbf{count} the ways or optimise a sum over a path in a graph without cycles --- that is DP on a DAG (Chapter 5 and the DP book);
    \item a single long walk where each state has exactly one successor --- that is a functional graph (Pattern~5).
\end{itemize}

%------------------------------------------------------------
\subsubsection*{Variations map}

\begin{diagrambox}
\begin{tikzpicture}[
    var/.style={rectangle, rounded corners=2pt, draw=patternborder, fill=patternbg,
                font=\scriptsize, align=center, text width=3.25cm, minimum height=1.25cm, inner sep=3pt}]
    \node[var, fill=answerfill, draw=accentorange, text width=4.6cm, font=\scriptsize\bfseries] (root) at (0,0)
        {BFS over states\\fewest unit-cost moves from start to goal};
    \node[var] (v1) at (-5.4,-2) {\textbf{numbers}\\(Two Buttons, E4.1, E4.4, E4.5)\\bound the window first};
    \node[var] (v2) at (-1.8,-2) {\textbf{digit tuples / codes}\\(Open the Lock)\\mixed-radix id};
    \node[var] (v3) at (1.8,-2)  {\textbf{resource pairs}\\(water jugs, E4.3)\\$(a,b)$ as $a(y{+}1)+b$};
    \node[var] (v4) at (5.4,-2)  {\textbf{strings / boards}\\(word ladders, sliding puzzles)\\Chapter 3};
    \node[var] (v5) at (-3.6,-3.6) {\textbf{reverse the moves}\\multi-source / one target\\(E4.4, greedy of 1.4.1)};
    \node[var] (v6) at (0,-3.6)    {\textbf{meet in the middle}\\bidirectional BFS\\(Problem 1.4.2, Approach 3)};
    \node[var] (v7) at (3.6,-3.6)  {\textbf{weighted moves}\\0-1 BFS / Dijkstra\\(Chapters 7--8)};
    \node[draw=patternborder, dashed, rounded corners=4pt, fit=(v1)(v4)(v5)(v7), inner sep=5pt] (grp) {};
    \draw[flow] (root.south) -- (grp.north);
\end{tikzpicture}
\figcap{Top row: what a state looks like. Bottom row: how the search itself changes.}
\end{diagrambox}

\textbf{Twists and how to handle them}
\begin{itemize}[leftmargin=*]
    \item \textbf{Forbidden states} (dead ends, walls, ``the device breaks''): simply never push them. If the \emph{start} is forbidden, answer immediately.
    \item \textbf{Many queries, one goal}: run one BFS \emph{backwards} from the goal over reversed moves; then every start is answered by a lookup (Exercise~E4.4).
    \item \textbf{Print the sequence of moves}: store \texttt{par[t] = s} when $t$ is discovered and walk back from the goal (Template~1).
    \item \textbf{Goal described by a property} (``total water equals $T$''): test the property when a state is popped instead of comparing with one goal state.
    \item \textbf{Out-of-window results that are still allowed} (Exercise~E4.5): check them against the goal, but do not push them.
\end{itemize}

%------------------------------------------------------------
\subsubsection*{Code templates}

\textbf{Template 1 --- integer states in a window, with the path.} Moves of the demo: $+3$ and $\times2$. Both only increase the value, so the window is $[\text{start}, \text{goal}]$.

\begin{lstlisting}[style=cpp]
#include <bits/stdc++.h>
using namespace std;

// fewest moves from start to goal inside the window [LO, HI]; path receives the states
int bfsInt(int start, int goal, int LO, int HI, vector<int>& path) {
    vector<int> dist(HI - LO + 1, -1), par(HI - LO + 1, -1);
    queue<int> q;
    dist[start - LO] = 0;
    q.push(start);
    while (!q.empty()) {
        int s = q.front(); q.pop();
        if (s == goal) break;                          // goal popped: distance final
        int next[2] = {s + 3, s * 2};                  // the moves of this demo
        for (int t : next) {
            if (t < LO || t > HI) continue;            // outside the proven window
            if (dist[t - LO] != -1) continue;          // already discovered
            dist[t - LO] = dist[s - LO] + 1;           // mark when PUSHED
            par[t - LO] = s;
            q.push(t);
        }
    }
    path.clear();
    if (dist[goal - LO] == -1) return -1;
    for (int v = goal; v != -1; v = par[v - LO]) path.push_back(v);
    reverse(path.begin(), path.end());
    return dist[goal - LO];
}

int main() {
    vector<int> path;
    cout << bfsInt(1, 25, 1, 25, path) << ":";
    for (int v : path) cout << " " << v;               // 5: 1 4 8 11 22 25
    cout << "\n";
}
\end{lstlisting}

\textbf{Template 2 --- a tuple of booleans as a bitmask.} Five lamps in a row; pressing button $i$ toggles lamp $i$ and its neighbours. Fewest presses from ``all off'' to a goal pattern. There are $2^5 = 32$ states and 5 moves each.

\begin{lstlisting}[style=cpp]
#include <bits/stdc++.h>
using namespace std;

int main() {
    const int N = 5, S = 1 << N;                       // state = bitmask of lamps that are on
    vector<int> dist(S, -1);
    queue<int> q;
    dist[0] = 0;
    q.push(0);
    while (!q.empty()) {
        int s = q.front(); q.pop();
        for (int i = 0; i < N; i++) {
            int mask = 1 << i;                         // lamp i ...
            if (i > 0) mask |= 1 << (i - 1);           // ... and its neighbours
            if (i + 1 < N) mask |= 1 << (i + 1);
            int t = s ^ mask;
            if (dist[t] == -1) { dist[t] = dist[s] + 1; q.push(t); }
        }
    }
    int reachable = 0;
    for (int d : dist) if (d != -1) reachable++;
    cout << dist[S - 1] << " " << dist[1] << " " << reachable << "\n";   // 2 -1 16
}
\end{lstlisting}
All five lamps on needs 2 presses (buttons 0 and 3); lamp 0 alone (\texttt{dist[1]}) is \emph{impossible}: only 16 of the 32 patterns are reachable. Reachability questions are answered by the same BFS.

\textbf{Template 3 --- string states with a hash map.} Fewest swaps of adjacent letters turning one word into another. Convenient when the state has no obvious small encoding; slower by a constant factor, so prefer an integer id when the state space is large.

\begin{lstlisting}[style=cpp]
#include <bits/stdc++.h>
using namespace std;

int fewestSwaps(const string& from, const string& to) {
    unordered_map<string, int> dist;                   // state -> number of moves
    queue<string> q;
    dist[from] = 0;
    q.push(from);
    while (!q.empty()) {
        string s = q.front(); q.pop();
        if (s == to) return dist[s];
        for (size_t i = 0; i + 1 < s.size(); i++) {
            string t = s;
            swap(t[i], t[i + 1]);                      // one move = one adjacent swap
            if (!dist.count(t)) { dist[t] = dist[s] + 1; q.push(t); }
        }
    }
    return -1;
}

int main() {
    cout << fewestSwaps("aabb", "bbaa") << " " << fewestSwaps("abcd", "dcba") << "\n";  // 4 6
}
\end{lstlisting}

%------------------------------------------------------------
\subsubsection*{Common mistakes}

\begin{warningbox}[Mistake 1 --- no bound on the state space]
\begin{lstlisting}[style=snippet]
q.push(2 * s);                               // WRONG: values grow forever, unreachable goal = infinite loop
if (2 * s <= HI) q.push(2 * s);              // RIGHT: HI comes from a proof, not from a guess
\end{lstlisting}
A window that is too \emph{small} is just as bad: in Two Buttons with HI $= m$ the answer for $n=3$, $m=5$ (path $3\to6\to5$) is not found at all.
\end{warningbox}

\begin{warningbox}[Mistake 2 --- marking visited when popping]
If \texttt{dist[t]} is set only when $t$ is popped, the same state is pushed many times (once per incoming edge). The answer stays correct, but the queue can grow to $\#\text{states}\times\#\text{moves}$ entries --- \mle{} or \tle{} on large spaces. Mark on push.
\end{warningbox}

\begin{warningbox}[Mistake 3 --- DFS for ``fewest moves'']
A DFS with a visited array finds \emph{a} way, usually a long one: on the lock it may wander through thousands of codes. Only BFS (unit costs) guarantees the minimum.
\end{warningbox}

\begin{warningbox}[Mistake 4 --- a state that forgets something]
If the moves or the goal test depend on a quantity, it must be in the state. Example: with a limited number of ``special'' moves, the state is (position, special moves left), not just the position --- otherwise BFS merges situations that behave differently.
\end{warningbox}

\begin{warningbox}[Mistake 5 --- forgetting the start or the goal edge cases]
Start $=$ goal $\Rightarrow$ 0 moves. Start forbidden $\Rightarrow$ $-1$ before the loop (Open the Lock with \texttt{"0000"} among the dead ends). Goal outside the window $\Rightarrow$ check it explicitly.
\end{warningbox}

\begin{checklistbox}[``I'm stuck'' checklist for Pattern 4]
\begin{itemize}[leftmargin=*]
    \item[$\square$] What is the smallest snapshot that decides the next moves and the goal test? That is the state.
    \item[$\square$] Did I write down every move of the statement, including wrap-around and ``until full / until empty''?
    \item[$\square$] How many states are there? Is there a \emph{proof} for my window? Does $\#\text{states}\times\#\text{moves}$ fit?
    \item[$\square$] Can I encode the state as an integer for a flat \texttt{dist} array?
    \item[$\square$] Do all moves cost the same? (Otherwise: not plain BFS.)
    \item[$\square$] Many starts and one goal --- should I run BFS backwards from the goal?
    \item[$\square$] Is the state space too big? Try the reverse direction, a greedy, or meet in the middle.
\end{itemize}
\end{checklistbox}

%============================================================
% CH1 · P4 · PROBLEM 1 — TWO BUTTONS (Codeforces 520B)
%============================================================
\subsection{Problem 1.4.1: Two Buttons}
\problemlabel{1.4.1}{ssec:ch1-two-buttons}

\problemcard{1.4.1}{Two Buttons}{Codeforces 520B}{https://codeforces.com/problemset/problem/520/B}{\exmedium\ \cf{1400}}{Numbers as vertices, button presses as edges; BFS in a proven window; reverse greedy}

\subsubsection*{1. Problem statement}

A device shows a positive integer on its display. The \textbf{red} button multiplies the number by $2$, the \textbf{blue} button subtracts $1$. If the number ever stops being positive, the device breaks. The display can show arbitrarily large numbers. Initially the display shows $n$. What is the \textbf{minimum number of presses} needed to make it show $m$?

\textbf{Constraints:} $1 \le n, m \le 10^4$, \ $n \ne m$. \quad \textbf{Input:} one line with $n$ and $m$. \quad \textbf{Output:} the minimum number of presses.

\begin{table}[H]
\centering\small
\begin{tabular}{@{}cccl@{}}
\toprule
$n$ & $m$ & \textbf{Output} & \textbf{One optimal sequence} \\
\midrule
4 & 6 & 2 & blue, red: \ $4 \to 3 \to 6$ \\
10 & 1 & 9 & blue nine times: \ $10 \to 9 \to \dots \to 1$ \\
\bottomrule
\end{tabular}
\end{table}

\begin{diagrambox}
\begin{tikzpicture}
    % the device
    \draw[rounded corners=4pt, draw=deepblue, fill=theorybg, line width=1pt] (0,0) rectangle (4.2,2.4);
    \draw[rounded corners=2pt, fill=darkgray!85] (0.5,1.3) rectangle (3.7,2.1);
    \node[font=\Large\ttfamily\bfseries, text=successgreen!60] at (2.1,1.7) {4};
    \node[circle, draw=recursionred!80!black, fill=recursionred, minimum size=7mm, text=white, font=\scriptsize\bfseries] at (1.2,0.6) {$\times2$};
    \node[circle, draw=primaryblue!80!black, fill=primaryblue, minimum size=7mm, text=white, font=\scriptsize\bfseries] at (3.0,0.6) {$-1$};
    \node[font=\scriptsize, text=darkgray] at (1.2,-0.25) {red};
    \node[font=\scriptsize, text=darkgray] at (3.0,-0.25) {blue};
    % sample 1 as a path
    \begin{scope}[xshift=6cm, yshift=1.2cm]
        \node[gvsrc] (a) at (0,0) {4};
        \node[gv] (b) at (2,0) {3};
        \node[gvdst] (c) at (4,0) {6};
        \draw[gpath, -{Stealth[length=2.6mm]}] (a) -- node[gw, above=1pt]{blue $-1$} (b);
        \draw[gpath, -{Stealth[length=2.6mm]}] (b) -- node[gw, above=1pt]{red $\times2$} (c);
        \node[gv] (d) at (2,-1.3) {8};
        \node[gv] (e) at (4,-1.3) {7};
        \draw[gde] (a) -- node[gw, below left=-1pt]{$\times2$} (d);
        \draw[gde] (d) -- node[gw, below]{$-1$} (e);
        \draw[gde] (e) -- node[gw, right]{$-1$} (c);
        \node[note, anchor=west] at (4.5,-0.6) {orange: 2 presses\\gray: 3 presses};
    \end{scope}
\end{tikzpicture}
\figcap{Sample 1. Doubling first ($4\to8\to7\to6$) looks natural but costs 3 presses; stepping down first costs 2.}
\end{diagrambox}

\subsubsection*{2. The graph model}

\begin{modelbox}[Graph Model --- Two Buttons]
\renewcommand{\arraystretch}{1.35}
\begin{tabularx}{\linewidth}{@{}>{\bfseries}l X@{}}
Vertices & one vertex per number the display can show: $x \in \{1, 2, 3, \dots\}$; only the window $1 \le x \le \max(n, 2m)$ can matter (Lemma below) \\
Edges & directed: $x \to 2x$ (red) and $x \to x-1$ (blue, only if $x - 1 \ge 1$); generated on the fly, never stored \\
Weights & none --- every press costs 1 \\
Question & length of a shortest path from vertex $n$ to vertex $m$ \\
Size & $|V| \le 2\cdot10^4$, \ $|E| \le 4\cdot10^4$: BFS is instant \\
\end{tabularx}
\end{modelbox}

\begin{diagrambox}
\begin{tikzpicture}[x=1.25cm]
    \foreach \x in {1,...,9}
        \node[gv, minimum size=6.5mm] (n\x) at (\x,0) {\x};
    \node[font=\small, text=darkgray] at (10,0) {$\cdots$};
    \foreach \x in {2,...,9} {
        \pgfmathtruncatemacro{\y}{\x-1}
        \draw[gde, draw=primaryblue] (n\x) to[bend left=35] (n\y);
    }
    \draw[gde, draw=recursionred, line width=1.1pt] (n1) to[bend left=45] (n2);
    \draw[gde, draw=recursionred, line width=1.1pt] (n2) to[bend left=50] (n4);
    \draw[gde, draw=recursionred, line width=1.1pt] (n3) to[bend left=50] (n6);
    \draw[gde, draw=recursionred, line width=1.1pt] (n4) to[bend left=50] (n8);
    \node[font=\scriptsize\bfseries, text=recursionred] at (5.5,1.75) {red: $x \to 2x$};
    \node[font=\scriptsize\bfseries, text=primaryblue] at (5.5,-0.85) {blue: $x \to x-1$};
    \node[gvdead, minimum size=6.5mm] (z) at (0,0) {0};
    \draw[gcut] (0.5,-0.5) -- (0.5,0.5);
    \node[font=\scriptsize, text=warningred] at (0,-0.6) {breaks};
\end{tikzpicture}
\figcap{The implicit graph on the number line (only the first edges are drawn). Every vertex has at most two outgoing edges; the edge $1 \to 0$ is missing because the device would break.}
\end{diagrambox}

\subsubsection*{3. Key observations}

\begin{itemize}[leftmargin=*]
    \item \textbf{It is a shortest path with unit weights.} Each press is one edge, all presses cost the same, so BFS from $n$ gives the answer.
    \item \textbf{Forward greedy fails.} ``Double while below $m$, then step down'' gives $4\to8\to7\to6$ (3 presses) on sample 1; the optimum steps down \emph{first}. Choosing the right moment to subtract is exactly what BFS tries for us.
    \item \textbf{The graph is infinite} --- red can be pressed forever. Without a bound, BFS never stops. The lemma below gives the window $[1, \max(n, 2m)]$.
    \item \textbf{If $n > m$, red is useless:} the answer is $n - m$ (sample 2). This is the lemma applied at the very start.
\end{itemize}

\begin{insightbox}[{Lemma (the window) --- once the display shows $x \ge m$, red is never worth pressing}]
Take any way to go from $x \ge m$ to $m$ that uses $k \ge 1$ red presses, made at values $z_1, \dots, z_k \ge 1$, and $b$ blue presses. Each red press adds $z_i$ to the value, each blue press removes 1, so
\[
x + \sum_i z_i - b = m \quad\Longrightarrow\quad b = x - m + \sum_i z_i \;\ge\; x - m + k,
\]
and the total is $k + b \ge x - m + 2k > x - m$. But pressing blue $x - m$ times reaches $m$ directly (and never goes below $m \ge 1$). So in an optimal sequence, red is pressed only at values $x < m$, which produces at most $2(m-1)$; all other values come from $n$ by subtracting. \textbf{Every value on an optimal path is at most $\max(n, 2m-2)$.} \hfill$\blacksquare$
\end{insightbox}

\begin{diagrambox}
\begin{tikzpicture}[x=0.55cm]
    \fill[solutionbg] (1,-0.25) rectangle (20,0.25);
    \fill[recursionbg] (20,-0.25) rectangle (23.5,0.25);
    \draw[thick, -{Stealth}] (0,0) -- (23.7,0);
    \foreach \x/\l in {1/1, 6/m, 12/2m, 20/{\text{HI}=2\max(n,m)}} {
        \draw[thick] (\x,-0.3) -- (\x,0.3);
        \node[font=\scriptsize\bfseries, text=deepblue] at (\x,0.6) {$\l$};
    }
    \node[gvsrc, minimum size=5mm, font=\tiny\bfseries] (A) at (4,0) {$n$};
    \node[gvdst, minimum size=5mm, font=\tiny\bfseries] (B) at (6,0) {$m$};
    \draw[-{Stealth[length=2mm]}, recursionred, thick] (3,0.3) to[bend left=40] node[font=\tiny, above]{red at $x<m$} (6,0.3);
    \draw[-{Stealth[length=2mm]}, warningred, thick, dashed] (8,-0.3) to[bend right=30] node[font=\tiny, below]{red at $x \ge m$: never optimal} (16,-0.3);
    \node[font=\scriptsize, text=warningred] at (21.8,-0.6) {never needed};
\end{tikzpicture}
\figcap{The window when $n < m$: a red press made below $m$ lands below $2m$; a red press made at or above $m$ is never part of an optimal answer. The code uses $\text{HI} = 2\max(n, m) \ge \max(n, 2m)$.}
\end{diagrambox}

\subsubsection*{4. Approaches}

\begin{bruteforcebox}
\textbf{Idea: try every sequence of presses}, shortest first (iterative deepening). For $d = 0, 1, 2, \dots$ ask ``can $m$ be reached with at most $d$ presses?'' by trying red and blue at every step.

\begin{lstlisting}[style=cpp]
#include <bits/stdc++.h>
using namespace std;

// can the display go from x to m with at most k more presses?
bool canReach(long long x, long long m, int k) {
    if (x == m) return true;
    if (k == 0) return false;
    if (canReach(2 * x, m, k - 1)) return true;              // red
    return x - 1 >= 1 && canReach(x - 1, m, k - 1);          // blue (must stay positive)
}

int main() {
    long long n, m;
    cin >> n >> m;
    int d = 0;
    while (!canReach(n, m, d)) d++;                          // first depth that works
    cout << d << "\n";
}
\end{lstlisting}

\begin{diagrambox}
\begin{forest} dptree, for tree={font=\scriptsize\ttfamily}
[4, ans
  [8, edge label={node[midway, left, font=\tiny]{red}}
    [16] [7]]
  [3, edge label={node[midway, right, font=\tiny]{blue}}
    [6, bc] [2]]]
\end{forest}
\figcap{Depth 2 on sample 1: four sequences, one of them ends at 6 (green).}
\end{diagrambox}

\begin{complexitybox}
\textbf{Time:} $\Theta(2^{d})$ where $d$ is the answer. Correct on both samples, but $n = 10^4$, $m = 1$ needs $d = 9999$ --- about $2^{9999}$ sequences. The same number (say $7$) is re-explored along many different sequences. \textbf{Verdict:} \tle{}.
\end{complexitybox}
\end{bruteforcebox}

\begin{algorithmbox}
\textbf{Idea: BFS over the numbers $1..\text{HI}$.} The brute force re-explores the same number many times; BFS visits each number once, in order of the number of presses. By the Lemma, $\text{HI} = 2\max(n, m)$ is a safe window.

\begin{lstlisting}[style=cpp]
#include <bits/stdc++.h>
using namespace std;

int main() {
    int n, m;
    cin >> n >> m;
    const int HI = 2 * max(n, m);                  // every optimal path stays in [1, HI]
    vector<int> dist(HI + 1, -1);                  // dist[x] = fewest presses n -> x
    queue<int> q;
    dist[n] = 0;
    q.push(n);
    while (!q.empty()) {
        int x = q.front(); q.pop();
        if (x == m) break;                         // goal popped: distance is final
        int next[2] = {2 * x, x - 1};              // red, blue
        for (int y : next) {
            if (y < 1 || y > HI) continue;         // device breaks / outside the window
            if (dist[y] != -1) continue;           // already discovered
            dist[y] = dist[x] + 1;
            q.push(y);
        }
    }
    cout << dist[m] << "\n";
}
\end{lstlisting}

\begin{complexitybox}
\textbf{Time:} $\bigO{\text{HI}} = \bigO{n + m}$ --- at most $2\cdot10^4$ states, 2 edges each. \textbf{Space:} $\bigO{n + m}$. \textbf{Verdict:} \ac{}.
\end{complexitybox}
\end{algorithmbox}

\begin{invariantbox}[Why it is correct]
\textbf{Window:} by the Lemma some optimal sequence stays inside $[1, \text{HI}]$, so restricting BFS to the window removes no optimal path, and the window adds no invalid one (values $\ge 1$, and every edge is a real button press).
\textbf{Distances:} inside the window this is plain BFS with unit costs, so by the BFS invariant of Pattern~4 (\S\ref{ssec:ch1-p4-theory}) every \texttt{dist[y]} is set once, to the true fewest number of presses. $m$ is always reachable (blue from $n > m$; doubling past $m$ and stepping down from $n < m$), so \texttt{dist[m]} $\ne -1$.
\end{invariantbox}

\begin{alternativebox}[Approach 3: Reverse Greedy in $\bigO{\log m}$]
\textbf{Idea: run the film backwards.} Undoing red means \emph{halving} (only possible on an even number), undoing blue means \emph{adding 1}. Going backwards from $m$ to $n$, the choice becomes forced:
\begin{itemize}[leftmargin=*]
    \item $m$ odd $\Rightarrow$ it cannot be the result of a red press: add 1;
    \item $m$ even and $m > n$ $\Rightarrow$ halve (proof below);
    \item $m \le n$ $\Rightarrow$ the rest is $n - m$ blue presses (Lemma).
\end{itemize}
\begin{lstlisting}[style=cpp]
#include <bits/stdc++.h>
using namespace std;

int main() {
    int n, m;
    cin >> n >> m;
    int presses = 0;
    while (m > n) {                       // walk backwards from m
        if (m % 2 == 1) m += 1;           // undo a blue press
        else m /= 2;                      // undo a red press
        presses++;
    }
    cout << presses + (n - m) << "\n";    // the remaining part: blue only
}
\end{lstlisting}
\textbf{Why halving is right when $m > n$ is even.} Suppose an optimal backward sequence starts with $+1$. Plus-ones alone never get down to $n < m$, so some halving follows after $j \ge 1$ plus-ones, and $m + j$ must be even, so $j \ge 2$. Replace ``$j$ times $+1$, then halve'' ($j + 1$ steps) by ``halve, then $j/2$ times $+1$'' ($1 + j/2$ steps): both reach $(m + j)/2$, and the second is strictly shorter --- contradiction.

\begin{diagrambox}
\begin{tikzpicture}[node distance=1.9cm]
    \node[gvdst] (a) {10};
    \node[gv, right=of a] (b) {5};
    \node[gv, right=of b] (c) {6};
    \node[gvsrc, right=of c] (d) {3};
    \draw[back] (a) -- node[lbl, above]{even: $\div2$} (b);
    \draw[back] (b) -- node[lbl, above]{odd: $+1$} (c);
    \draw[back] (c) -- node[lbl, above]{even: $\div2$} (d);
    \node[note, below=3mm of b, xshift=1.9cm] {$n=3$, $m=10$: 3 backward steps, then $m = n$ (0 blue presses). Forward: $3 \xrightarrow{\times2} 6 \xrightarrow{-1} 5 \xrightarrow{\times2} 10$.};
\end{tikzpicture}
\end{diagrambox}

\begin{complexitybox}
\textbf{Time:} $\bigO{\log m}$ --- every two steps at least halve $m$. \textbf{Space:} $\bigO{1}$. \textbf{Verdict:} \ac{}; it also works for $m$ up to $10^{18}$, where BFS is impossible. The same task appears on LeetCode as ``Broken Calculator'' with values up to $10^9$ --- only the greedy survives there.
\end{complexitybox}
\end{alternativebox}

\begin{tricksbox}[Approach 4: Printing the Presses]
Contest problems often ask for the sequence itself. Store the parent when a number is discovered and walk back from $m$:
\begin{lstlisting}[style=cpp]
#include <bits/stdc++.h>
using namespace std;

int main() {
    int n, m;
    cin >> n >> m;
    const int HI = 2 * max(n, m);
    vector<int> dist(HI + 1, -1), par(HI + 1, -1);
    queue<int> q;
    dist[n] = 0;
    q.push(n);
    while (!q.empty()) {
        int x = q.front(); q.pop();
        int next[2] = {2 * x, x - 1};
        for (int y : next)
            if (y >= 1 && y <= HI && dist[y] == -1) {
                dist[y] = dist[x] + 1;
                par[y] = x;
                q.push(y);
            }
    }
    string presses;
    for (int v = m; v != n; v = par[v])
        presses += (par[v] * 2 == v ? 'R' : 'B');           // which button led to v?
    reverse(presses.begin(), presses.end());
    cout << dist[m] << " " << presses << "\n";               // "4 6" -> "2 BR"
}
\end{lstlisting}
The question ``was it red?'' is answered by $\texttt{par}[v] \cdot 2 = v$; red and blue can never lead from the same parent to the same $v$ (that would need $2x = x - 1$).
\end{tricksbox}

\subsubsection*{5. Trace}

\begin{tracebox}[{Sample 1: BFS from $n=4$ to $m=6$, window $[1, 12]$}]
\centering
\begin{tikzpicture}[x=1.15cm, y=0.75cm]
    \foreach \k/\px/\py/\title in {1/0/0/{pop 4}, 2/6.2/0/{pop 8}, 3/0/-4.4/{pop 3}, 4/6.2/-4.4/{pop 7, then pop 6}} {
        \begin{scope}[xshift=\px cm, yshift=\py cm]
            \node[font=\scriptsize\bfseries, text=deepblue, anchor=west] at (-0.4,1.9) {Step \k: \title};
        \end{scope}
    }
    % ---- step 1 ----
    \begin{scope}
        \node[gvcur] (a) at (0,0) {4};
        \node[gvseen] (b) at (1.3,1) {8};
        \node[gvseen] (c) at (1.3,-1) {3};
        \draw[gtree, -{Stealth[length=2.2mm]}] (a) -- (b);
        \draw[gtree, -{Stealth[length=2.2mm]}] (a) -- (c);
        \node[font=\scriptsize, anchor=west] at (2.3,0) {queue: \texttt{[8, 3]}};
    \end{scope}
    % ---- step 2 ----
    \begin{scope}[xshift=6.2cm]
        \node[gvdone] (a) at (0,0) {4};
        \node[gvcur] (b) at (1.3,1) {8};
        \node[gvseen] (c) at (1.3,-1) {3};
        \node[gvdead] (p) at (2.6,1.6) {16};
        \node[gvseen] (e) at (2.6,0.4) {7};
        \draw[gtree] (a) -- (b); \draw[gtree] (a) -- (c);
        \draw[gfaded, dashed] (b) -- (p);
        \draw[gtree, -{Stealth[length=2.2mm]}] (b) -- (e);
        \node[font=\scriptsize, anchor=west, align=left] at (3.2,-0.5) {16 $>$ HI: rejected\\queue: \texttt{[3, 7]}};
    \end{scope}
    % ---- step 3 ----
    \begin{scope}[yshift=-4.4cm]
        \node[gvdone] (a) at (0,0) {4};
        \node[gvdone] (b) at (1.3,1) {8};
        \node[gvcur] (c) at (1.3,-1) {3};
        \node[gvseen] (e) at (2.6,0.4) {7};
        \node[gvseen] (f) at (2.6,-0.6) {6};
        \node[gvseen] (g) at (2.6,-1.6) {2};
        \draw[gtree] (a) -- (b); \draw[gtree] (a) -- (c); \draw[gtree] (b) -- (e);
        \draw[gtree, -{Stealth[length=2.2mm]}] (c) -- (f);
        \draw[gtree, -{Stealth[length=2.2mm]}] (c) -- (g);
        \node[font=\scriptsize, anchor=west, align=left] at (3.2,-0.5) {6 discovered,\\dist 2\\queue: \texttt{[7, 6, 2]}};
    \end{scope}
    % ---- step 4 ----
    \begin{scope}[xshift=6.2cm, yshift=-4.4cm]
        \node[gvsrc] (a) at (0,0) {4};
        \node[gvdone] (b) at (1.3,1) {8};
        \node[gvdone] (c) at (1.3,-1) {3};
        \node[gvdone] (e) at (2.6,0.4) {7};
        \node[gvdst] (f) at (2.6,-0.6) {6};
        \node[gvseen] (g) at (2.6,-1.6) {2};
        \draw[gtree] (a) -- (b); \draw[gtree] (b) -- (e); \draw[gtree] (c) -- (g);
        \draw[gpath, -{Stealth[length=2.4mm]}] (a) -- (c);
        \draw[gpath, -{Stealth[length=2.4mm]}] (c) -- (f);
        \draw[gfaded, dashed, -{Stealth[length=2mm]}] (e) -- (f);
        \node[font=\scriptsize, anchor=west, align=left] at (3.2,-0.5) {$7\to6$: already seen\\6 popped: answer \textbf{2}\\path $4 \to 3 \to 6$};
    \end{scope}
\end{tikzpicture}

\begin{lstlisting}[style=trace]
pop  dist  red (2x)         blue (x-1)        queue after
 4    0    8  new, dist 1   3  new, dist 1    [8, 3]
 8    1    16 > HI, skip    7  new, dist 2    [3, 7]
 3    1    6  new, dist 2   2  new, dist 2    [7, 6, 2]
 7    2    14 > HI, skip    6  already seen   [6, 2]
 6    2    goal popped -> answer 2
\end{lstlisting}
\end{tracebox}

\begin{tracebox}[{Sample 2: $n = 10$, $m = 1$ --- BFS explores more than it needs}]
BFS also tries the red button: it discovers 20 numbers ($1..20$) before popping $1$ at distance 9. The layers are
\begin{center}\small
\begin{tabular}{@{}l*{10}{c}@{}}
\toprule
dist & 0 & 1 & 2 & 3 & 4 & 5 & 6 & 7 & 8 & 9 \\
numbers & 10 & 9, 20 & 8, 18, 19 & 7, 16, 17 & 6, 14, 15 & 5, 12, 13 & 4, 11 & 3 & 2 & 1 \\
\bottomrule
\end{tabular}
\end{center}
The reverse greedy needs no loop at all here ($m \le n$): the answer is $n - m = 9$ immediately.
\end{tracebox}

\subsubsection*{6. Test suite}

\begin{testbox}
\centering\small
\begin{tabular}{@{}cccl@{}}
\toprule
$n$ & $m$ & \textbf{Expected} & \textbf{Why} \\
\midrule
4 & 6 & 2 & sample 1; defeats ``double first'' \\
10 & 1 & 9 & sample 2; $n > m$: only blue \\
3 & 5 & 2 & $3\to6\to5$ goes \emph{above} $m$: kills the window $\text{HI} = m$ \\
7 & 13 & 2 & $7\to14\to13$, same trap \\
3 & 10 & 3 & reverse greedy uses both rules \\
1 & 10000 & 20 & long climb, many halvings backwards \\
10000 & 1 & 9999 & largest answer; brute force is hopeless \\
9999 & 10000 & 5000 & step down to 5000, then one red press \\
2 & 1 & 1 & smallest $n > m$ \\
1 & 2 & 1 & smallest $n < m$ \\
\bottomrule
\end{tabular}
\end{testbox}

All three solutions (iterative deepening on answers up to 14, BFS, reverse greedy) agree on every pair $1 \le n, m \le 200$, and BFS with the window agrees with BFS on $[1, 2\cdot10^5]$ and with the greedy on 3000 random pairs up to $10^4$.

\subsubsection*{7. Common mistakes}

\begin{warningbox}[Mistake 1 --- window too small]
\begin{lstlisting}[style=snippet]
const int HI = max(n, m);          // WRONG: n=3, m=5 needs 6; BFS prints -1
const int HI = 2 * max(n, m);      // RIGHT: Lemma
\end{lstlisting}
\end{warningbox}

\begin{warningbox}[Mistake 2 --- letting the display reach 0]
Without the test $y \ge 1$ the code reads \texttt{dist[0]} (and from there \texttt{dist[-1]}): undefined behaviour. Even if the array had room, 0 is a forbidden state --- the device breaks.
\end{warningbox}

\begin{warningbox}[Mistake 3 --- a greedy in the wrong direction]
Forward, neither ``double while below $m$'' nor ``subtract when it is odd'' is correct (sample 1). Backwards the choices are forced; forwards they are not. When a greedy fails in one direction, try the other.
\end{warningbox}

\begin{warningbox}[Mistake 4 --- forgetting the tail in the reverse greedy]
When the loop stops, $m$ may be \emph{below} $n$ (e.g.\ $n = 4$, $m = 6$: $6 \to 3$). The remaining $n - m$ blue presses must be added: $1 + (4 - 3) = 2$.
\end{warningbox}

\subsubsection*{8. Takeaways}

\begin{itemize}[leftmargin=*]
    \item ``Fewest operations from $n$ to $m$'' with unit-cost operations: numbers are vertices, operations are edges, BFS.
    \item An infinite state space needs a \emph{proven} window. Here: red pressed at $x \ge m$ never pays off, so $\text{HI} = \max(n, 2m)$.
    \item BFS is the safe general method; when the moves are invertible and become forced backwards, a reverse greedy can be exponentially faster.
    \item Store parents during BFS to print the actual sequence of moves.
\end{itemize}

\begin{variantsbox}
\variant{Multiply by 2 or append a 1}{\varin{https://codeforces.com/problemset/problem/727/A}}
  {red becomes $x \to 2x$, the second button becomes $x \to 10x + 1$; print any sequence from $a$ to $b$.}
  {same vertices (numbers up to $b$); both moves only increase, so the window is $[a, b]$.}
  {backwards, every number has at most one predecessor ($b/2$ if even, $(b-1)/10$ if it ends in 1): follow it (Exercise~E4.1).}
\variant{Three buttons}{\varnone}
  {a third button multiplies by 3.}
  {vertices $1..\max(n, 3m)$: the Lemma's argument works for any multiplication.}
  {BFS with three edges per vertex; the reverse greedy is no longer obviously optimal (Exercise~E4.2).}
\variant{Wrap-around display, many queries}{\varin{https://codeforces.com/problemset/problem/1661/B}}
  {operations $+1$ and $\times2$ modulo $32768$; fewest steps to reach 0 from each of many numbers.}
  {all $32768$ residues are vertices; the space is finite, no window needed.}
  {one BFS \emph{backwards} from 0 answers every query by a lookup (Exercise~E4.4).}
\variant{Buttons with different costs}{\varnone}
  {red costs $a$, blue costs $b$.}
  {the same graph, now with weights $a$ and $b$ on the edges.}
  {BFS no longer gives the cheapest sequence: Dijkstra's algorithm (Chapter 7).}
\end{variantsbox}

%============================================================
% CH1 · P4 · PROBLEM 2 — OPEN THE LOCK (LeetCode 752)
%============================================================
\subsection{Problem 1.4.2: Open the Lock}
\problemlabel{1.4.2}{ssec:ch1-open-lock}

\problemcard{1.4.2}{Open the Lock}{LeetCode 752}{https://leetcode.com/problems/open-the-lock/}{\exmedium\ \lc{Medium}}{Codes as vertices of a 4-dimensional torus, forbidden vertices, BFS; bidirectional BFS}

\subsubsection*{1. Problem statement}

A lock has 4 circular wheels, each showing a digit \texttt{0}--\texttt{9}. One move turns \textbf{one} wheel by \textbf{one} slot, up or down, with wrap-around (\texttt{9} $\to$ \texttt{0} and \texttt{0} $\to$ \texttt{9}). The lock starts at \texttt{"0000"}. A list of \textbf{dead ends} is given: if the lock ever displays one of them, the wheels jam and it can never be opened. Return the minimum number of moves to reach the code \texttt{target}, or $-1$ if it is impossible.

\textbf{Constraints:} $1 \le |\texttt{deadends}| \le 500$; all codes have 4 digits; \texttt{target} is not a dead end.

\begin{table}[H]
\centering\small
\begin{tabular}{@{}p{5.6cm}ccp{5.6cm}@{}}
\toprule
\textbf{deadends} & \textbf{target} & \textbf{Output} & \textbf{Comment} \\
\midrule
\texttt{0201 0101 0102 1212 2002} & \texttt{0202} & 6 & direct route (4 moves) is blocked \\
\texttt{8888} & \texttt{0009} & 1 & turn the last wheel down: $0 \to 9$ \\
\texttt{8887 8889 8878 8898 8788 8988 7888 9888} & \texttt{8888} & $-1$ & all 8 neighbours of the target are dead \\
\bottomrule
\end{tabular}
\end{table}

\begin{diagrambox}
\begin{tikzpicture}
    % lock body
    \draw[rounded corners=5pt, draw=deepblue, fill=theorybg, line width=1pt] (-0.4,-0.7) rectangle (4.0,0.9);
    \draw[line width=3pt, draw=midgray, rounded corners=8pt] (0.6,0.9) -- (0.6,1.7) -- (3.0,1.7) -- (3.0,0.9);
    \foreach \d [count=\i from 0] in {0,0,0,0}
        \node[gcell, minimum width=8mm, minimum height=10mm, font=\Large\ttfamily\bfseries] at (\i*0.9+0.4,0.1) {\d};
    \node[font=\scriptsize, text=darkgray] at (1.75,-1.05) {start \texttt{0000}, target \texttt{0202}};
    % one wheel as a ring
    \begin{scope}[xshift=7.3cm, yshift=0.2cm]
        \foreach \k in {0,...,9} {
            \pgfmathsetmacro{\ang}{90-36*\k}
            \node[gv, minimum size=6mm, font=\scriptsize\bfseries] (r\k) at (\ang:1.35) {\k};
        }
        \foreach \k in {0,...,9} {
            \pgfmathtruncatemacro{\nx}{mod(\k+1,10)}
            \draw[ge] (r\k) -- (r\nx);
        }
        \draw[gpath] (r9) -- (r0);
        \node[font=\scriptsize, text=accentorange, anchor=west] at (1.75,1.45) {\texttt{9} $\leftrightarrow$ \texttt{0}: one move};
        \node[note] at (0,-2.0) {each wheel is a cycle of length 10};
    \end{scope}
\end{tikzpicture}
\figcap{The lock and one of its wheels. Turning a wheel moves its digit to a neighbour on the cycle; \texttt{0} and \texttt{9} are neighbours too.}
\end{diagrambox}

\subsubsection*{2. The graph model}

\begin{modelbox}[Graph Model --- Open the Lock]
\renewcommand{\arraystretch}{1.35}
\begin{tabularx}{\linewidth}{@{}>{\bfseries}l X@{}}
Vertices & all $10^4$ codes \texttt{0000}..\texttt{9999}, encoded as the integer $0..9999$; dead ends are \textbf{deleted} vertices \\
Edges & undirected: two codes are adjacent if they differ in exactly one wheel by $\pm1 \pmod{10}$; every code has exactly 8 neighbours, generated on the fly \\
Weights & none --- each turn costs 1 \\
Question & shortest path from \texttt{0000} to \texttt{target} avoiding deleted vertices ($-1$ if disconnected) \\
Size & $|V| = 10^4$, \ $|E| = 4\cdot10^4$ (8 per vertex, each counted twice): tiny \\
\end{tabularx}
\end{modelbox}

\begin{diagrambox}
\begin{tikzpicture}
    \node[gvsrc, minimum size=11mm, font=\small\ttfamily\bfseries] (c) at (0,0) {0000};
    \foreach \lab/\ang/\w in {1000/90/1, 9000/45/1, 0100/0/2, 0900/-45/2, 0010/-90/3, 0090/-135/3, 0001/180/4, 0009/135/4} {
        \node[gvseen, minimum size=11mm, font=\scriptsize\ttfamily\bfseries] (n\lab) at (\ang:2.3) {\lab};
        \draw[ge] (c) -- (n\lab);
    }
    \node[note, anchor=west, text width=6.2cm, align=left] at (3.4,0.4) {
        \textbf{8 neighbours} of every code: wheel $1..4$, turned up or down.\\[3pt]
        Seen as a whole, the graph is a 4-dimensional grid of side 10 whose rows wrap around (a \emph{torus}): Pattern~3's grid, but in four dimensions, with dead ends playing the role of walls.};
\end{tikzpicture}
\figcap{The neighbourhood of \texttt{0000}. All 8 neighbours are at distance 1 (unless they are dead ends).}
\end{diagrambox}

\subsubsection*{3. Key observations}

\begin{itemize}[leftmargin=*]
    \item \textbf{Fewest moves, all moves equal} $\Rightarrow$ BFS from \texttt{0000}.
    \item \textbf{Without dead ends the answer is a formula:} each wheel is independent, so the distance is $\sum_{i} \min(d_i, 10 - d_i)$ over the target digits $d_i$. For \texttt{0202} that is $2 + 2 = 4$. Dead ends force \textbf{detours} --- sample 1 needs 6 --- and only a search sees them.
    \item \textbf{Dead ends are deleted vertices}, not just ``bad answers'': a path may not pass \emph{through} them. Never push them into the queue.
    \item \textbf{If \texttt{0000} itself is a dead end}, the lock is jammed from the start: answer $-1$ (unless the target were \texttt{0000}, which the statement excludes by forbidding a dead target).
    \item \textbf{The state space is fixed and small} ($10^4$ codes): no window proof is needed, unlike Problem~1.4.1.
\end{itemize}

\begin{diagrambox}
\begin{tikzpicture}
    \node[gvdst, minimum size=11mm, font=\small\ttfamily\bfseries] (c) at (0,0) {0202};
    \foreach \lab/\ang/\st/\d in {1202/90/gvseen/5, 9202/45/gvseen/5, 0302/0/gvseen/5, 0102/-45/gvdead/{dead}, 0212/-90/gvseen/5, 0292/-135/gvseen/5, 0201/180/gvdead/{dead}, 0203/135/gvseen/5} {
        \node[\st, minimum size=11mm, font=\scriptsize\ttfamily\bfseries] (n\lab) at (\ang:2.4) {\lab};
        \node[gdist] at ($(n\lab)+(\ang:0.85)$) {\d};
    }
    \foreach \lab in {1202, 9202, 0302, 0212, 0292, 0203} \draw[ge] (c) -- (n\lab);
    \foreach \lab in {0102, 0201} \draw[gfaded, dashed] (c) -- (n\lab);
    \node[note, anchor=west, text width=6cm, align=left] at (3.5,0.2) {
        The two neighbours on a 4-move route, \texttt{0201} and \texttt{0102}, are dead.
        Every live neighbour is at distance 5 (green), so the answer is $5 + 1 = 6$.};
\end{tikzpicture}
\figcap{Sample 1 seen from the target. Distances in green were computed by a full BFS from \texttt{0000}.}
\end{diagrambox}

\subsubsection*{4. Approaches}

\begin{bruteforcebox}
\textbf{Idea: try all move sequences}, shortest first (iterative deepening): for $d = 0, 1, 2, \dots$ enumerate every sequence of $d$ turns that avoids dead ends.

\begin{lstlisting}[style=cpp]
#include <bits/stdc++.h>
using namespace std;

class Solution {
    set<string> dead;
    string target;
    bool reach(string code, int k) {                 // target within k more moves?
        if (code == target) return true;
        if (k == 0) return false;
        for (int w = 0; w < 4; w++)
            for (int step : {1, 9}) {                // +1 or -1 (mod 10)
                string nxt = code;
                nxt[w] = char('0' + (code[w] - '0' + step) % 10);
                if (!dead.count(nxt) && reach(nxt, k - 1)) return true;
            }
        return false;
    }
public:
    int openLock(vector<string>& deadends, string target) {
        dead = set<string>(deadends.begin(), deadends.end());
        this->target = target;
        if (dead.count("0000")) return -1;
        for (int d = 0; d <= 20; d++)                // no answer exceeds 20 (see trace)
            if (reach("0000", d)) return d;
        return -1;                                   // may take forever to get here
    }
};
\end{lstlisting}

\begin{complexitybox}
\textbf{Time:} $\Theta(8^{d})$ sequences for answer $d$: sample 1 ($d=6$) already makes 51135 recursive calls (depths $0..6$ together), an answer of 20 would need $8^{20} \approx 10^{18}$, and an impossible case runs all depths up to 20. The same code is reached by a huge number of different sequences. \textbf{Verdict:} \tle{}.
\end{complexitybox}
\end{bruteforcebox}

\begin{algorithmbox}
\textbf{Idea: BFS over the $10^4$ codes}, each code an integer, \texttt{dist} a flat array, dead ends pre-marked as forbidden. Turning wheel $w$ (place value $p = 1000, 100, 10, 1$) from digit $g$ to digit $g'$ changes the integer by $(g' - g)\cdot p$.

\begin{lstlisting}[style=cpp]
#include <bits/stdc++.h>
using namespace std;

class Solution {
public:
    int openLock(vector<string>& deadends, string target) {
        const int P[4] = {1000, 100, 10, 1};           // place value of each wheel
        vector<char> dead(10000, 0);
        for (const string& s : deadends) dead[stoi(s)] = 1;
        int goal = stoi(target);
        if (dead[0]) return -1;                         // jammed from the start
        vector<int> dist(10000, -1);
        queue<int> q;
        dist[0] = 0;
        q.push(0);
        while (!q.empty()) {
            int u = q.front(); q.pop();
            if (u == goal) return dist[u];
            for (int w = 0; w < 4; w++) {
                int g = u / P[w] % 10;                  // digit on wheel w
                for (int step : {1, 9}) {               // up, down (9 = -1 mod 10)
                    int v = u + ((g + step) % 10 - g) * P[w];
                    if (dead[v] || dist[v] != -1) continue;
                    dist[v] = dist[u] + 1;
                    q.push(v);
                }
            }
        }
        return -1;                                      // target unreachable
    }
};
\end{lstlisting}

\begin{complexitybox}
\textbf{Time:} $\bigO{10^4 \cdot 8}$ plus reading the dead ends --- about $10^5$ steps whatever the input. \textbf{Space:} $\bigO{10^4}$. \textbf{Verdict:} \ac{}.
\end{complexitybox}
\end{algorithmbox}

\begin{invariantbox}[Why it is correct]
The graph has exactly the live codes as vertices and the legal single turns as edges, so a sequence of moves that never shows a dead end is exactly a path from \texttt{0000} in this graph. Dead codes are never pushed, so no path through them is ever produced. BFS with unit costs sets \texttt{dist[v]} to the true minimum number of moves the first time $v$ is discovered (Pattern~4 invariant, \S\ref{ssec:ch1-p4-theory}); when \texttt{goal} is popped its distance is final. If the queue empties first, the goal is not in the component of \texttt{0000}: $-1$.
\end{invariantbox}

\begin{alternativebox}[Approach 3: Bidirectional BFS]
\textbf{Idea: search from both ends and meet in the middle.} The graph is undirected and the goal is a single known code, so we may grow a BFS ball around \texttt{0000} \emph{and} one around \texttt{target}, always expanding one full layer of the \textbf{smaller} frontier. As soon as an edge connects the two balls, the shortest path is found.

\begin{diagrambox}
\begin{tikzpicture}
    \fill[primaryblue!12] (0,0) circle (1.5);
    \fill[successgreen!12] (3.2,0) circle (1.5);
    \draw[primaryblue!60, dashed] (0,0) circle (1.5);
    \draw[successgreen!60, dashed] (3.2,0) circle (1.5);
    \node[gvsrc, minimum size=8mm, font=\scriptsize\ttfamily\bfseries] at (0,0) {0000};
    \node[gvdst, minimum size=8mm, font=\scriptsize\ttfamily\bfseries] at (3.2,0) {0202};
    \node[gv, minimum size=4mm] (x) at (1.45,0.3) {};
    \node[gv, minimum size=4mm] (y) at (1.8,0.3) {};
    \draw[gpath] (x) -- (y);
    \node[font=\scriptsize, text=primaryblue] at (0,-1.8) {radius 3 from start};
    \node[font=\scriptsize, text=successgreen!70!black] at (3.2,-1.8) {radius 2 from target};
    \draw[primaryblue!40, dashed] (9.0,0) circle (2.4);
    \fill[primaryblue!8] (9.0,0) circle (2.4);
    \node[gvsrc, minimum size=8mm, font=\scriptsize\ttfamily\bfseries] at (7.8,0) {0000};
    \node[gvdst, minimum size=8mm, font=\scriptsize\ttfamily\bfseries] at (10.2,0) {0202};
    \node[font=\scriptsize, text=primaryblue] at (9.0,-2.7) {one-sided BFS: radius 6};
\end{tikzpicture}
\figcap{Two balls of radius about $d/2$ instead of one ball of radius $d$. In a graph where the number of states grows quickly with the radius, that is a huge saving.}
\end{diagrambox}

\begin{lstlisting}[style=cpp]
#include <bits/stdc++.h>
using namespace std;

class Solution {
public:
    int openLock(vector<string>& deadends, string target) {
        const int P[4] = {1000, 100, 10, 1};
        vector<char> dead(10000, 0);
        for (const string& s : deadends) dead[stoi(s)] = 1;
        int goal = stoi(target);
        if (dead[0]) return -1;
        if (goal == 0) return 0;
        vector<int> da(10000, -1), db(10000, -1);       // distances from start / from goal
        vector<int> fa = {0}, fb = {goal};              // current frontiers (one layer each)
        da[0] = 0;
        db[goal] = 0;
        while (!fa.empty() && !fb.empty()) {
            bool fromStart = fa.size() <= fb.size();    // expand the smaller side
            vector<int>& front = fromStart ? fa : fb;
            vector<int>& mine = fromStart ? da : db;
            vector<int>& other = fromStart ? db : da;
            vector<int> next;
            int best = INT_MAX;
            for (int u : front)
                for (int w = 0; w < 4; w++) {
                    int g = u / P[w] % 10;
                    for (int step : {1, 9}) {
                        int v = u + ((g + step) % 10 - g) * P[w];
                        if (dead[v]) continue;
                        if (other[v] != -1) best = min(best, mine[u] + 1 + other[v]);
                        if (mine[v] == -1) { mine[v] = mine[u] + 1; next.push_back(v); }
                    }
                }
            if (best != INT_MAX) return best;           // the two balls touched
            front.swap(next);
        }
        return -1;                                      // one side got stuck: no path
    }
};
\end{lstlisting}

\textbf{Correctness.} Every layer is expanded \emph{completely} before checking, and the minimum over all touching edges of that layer is returned, so the first contact gives the shortest total. If one frontier becomes empty, its whole component has been explored without touching the other side: $-1$ (this is how the walled-in target of sample 3 is detected after a single layer).

\begin{complexitybox}
\textbf{Time:} at most $\bigO{10^4 \cdot 8}$ like BFS, often much less. On sample 1 one-sided BFS discovers 1562 codes before popping the target; bidirectional BFS discovers 245 (counting both sides). On \texttt{5555} with no relevant dead ends (distance 20, the maximum) there is no gain: both balls together cover the whole lock. \textbf{Verdict:} \ac{}; the technique pays off in huge state spaces such as sliding puzzles and word ladders (Chapter 3).
\end{complexitybox}
\end{alternativebox}

\begin{tricksbox}[Approach 4: Implementation Refinements]
\begin{itemize}[leftmargin=*]
    \item \textbf{Integers, not strings.} \texttt{unordered\_set<string>} and \texttt{string} copies work (Template~3 of the theory) but are several times slower; the integer id with \texttt{P[w]} arithmetic needs no allocation at all.
    \item \textbf{Down $=$ up by 9.} Writing $(g + 9) \bmod 10$ for ``down'' avoids negative remainders: in C++, $(0 - 1) \bmod 10 = -1$, not 9.
    \item \textbf{Pre-mark dead ends as visited.} Setting \texttt{dist[dead] = -2} merges the two tests into one; the code above keeps them separate for clarity.
    \item \textbf{Quick lower bound.} $\sum_i \min(d_i, 10-d_i)$ is a lower bound on the answer; if BFS reaches it, no detour was needed. Using it to guide the search is the idea of A* search.
\end{itemize}
\end{tricksbox}

\subsubsection*{5. Trace}

The lock graph is 4-dimensional, but sample 1 lives in a 2-dimensional \emph{slice}: the codes \texttt{0a0b} (wheels 1 and 3 stay at 0). In this slice the graph is an ordinary grid and the dead ends are walls --- a Pattern~3 picture.

\begin{tracebox}[{Sample 1 on the slice \texttt{0a0b}: distances from \texttt{0000}}]
\centering
\begin{tikzpicture}[x=0.95cm, y=0.95cm]
    % values row a (0..3), col b (0..4); X = dead
    \foreach \a/\row in {0/{0,1,2,3,4}, 1/{1,-,-,4,5}, 2/{2,-,6,5,6}, 3/{3,4,5,6,7}} {
        \foreach \v [count=\b from 0] in \row
            \node[gcell, minimum size=9mm] at (\b,\a) {\v};
    }
    % dead ends, path cells, start, target
    \foreach \b/\a in {1/1, 2/1, 1/2}
        \node[gwall, minimum size=9mm, font=\scriptsize\bfseries] at (\b,\a) {dead};
    \foreach \b/\a/\v in {0/1/1, 0/2/2, 0/3/3, 1/3/4, 2/3/5}
        \node[gcell, minimum size=9mm, fill=pathfill, font=\small\bfseries] at (\b,\a) {\v};
    \node[gcell, minimum size=9mm, fill=basefill, draw=successgreen, line width=1.2pt] at (0,0) {\textbf{0}};
    \node[gcell, minimum size=9mm, fill=answerfill, draw=accentorange, line width=1.2pt] at (2,2) {\textbf{6}};
    \foreach \b in {0,...,4} \node[font=\scriptsize, text=darkgray] at (\b,-0.75) {$b{=}\b$};
    \foreach \a in {0,...,3} \node[font=\scriptsize, text=darkgray] at (-0.95,\a) {$a{=}\a$};
    \node[note, anchor=west, text width=5.6cm, align=left] at (5.0,1.5) {
        cell $(a,b)$ = code \texttt{0a0b}, number = BFS distance.\\[3pt]
        The direct routes into \texttt{0202} pass \texttt{0201} or \texttt{0102} --- both dead.
        One shortest detour (shaded cells): \texttt{0000}$\to$\texttt{0100}$\to$\texttt{0200}$\to$\texttt{0300}$\to$\texttt{0301}$\to$\texttt{0302}$\to$\texttt{0202}.\\[3pt]
        LeetCode's explanation leaves the slice instead (\texttt{1000}, \texttt{1100}, \ldots, \texttt{1202}): also 6 moves.};
\end{tikzpicture}
\end{tracebox}

\begin{tracebox}[{Sample 1: the BFS layers}]
A full BFS from \texttt{0000} (without stopping at the target) gives these layer sizes; 5 codes (the dead ends) are never reached:
\begin{center}
\begin{tikzpicture}[x=0.6cm, y=0.0028cm]
    \foreach \d/\c in {0/1, 1/8, 2/31, 3/86, 4/190, 5/356, 6/576, 7/824, 8/1056, 9/1224, 10/1286,
                       11/1224, 12/1056, 13/824, 14/576, 15/356, 16/192, 17/88, 18/32, 19/8, 20/1} {
        \ifnum\d=6
            \fill[accentorange] (\d-0.35,0) rectangle (\d+0.35,\c);
        \else
            \fill[primaryblue!55] (\d-0.35,0) rectangle (\d+0.35,\c);
        \fi
        \node[font=\tiny, text=darkgray] at (\d,-45) {\d};
    }
    \draw[darkgray] (-0.6,0) -- (20.6,0);
    \foreach \d/\c in {2/31, 6/576, 10/1286, 18/32} \node[font=\tiny, above] at (\d,\c) {\c};
    \node[font=\scriptsize, text=darkgray] at (10,-140) {distance from \texttt{0000}};
\end{tikzpicture}
\end{center}
Layer 2 has 31 codes instead of 32 because \texttt{0101} is dead. The target sits in layer 6 (orange): the early-stopping BFS pops 865 codes before reaching it. The deepest code, \texttt{5555}, is at distance $5\cdot4 = 20$ --- no answer can exceed 20.
\end{tracebox}

\begin{tracebox}[{Sample 3: target walled in}]
\begin{lstlisting}[style=trace]
neighbours of 8888: 9888 7888 8988 8788 8898 8878 8889 8887  -> all dead
one-sided BFS : explores all 9991 live codes reachable from 0000, never sees 8888 -> -1
bidirectional : the target side's first layer is empty -> -1 immediately
\end{lstlisting}
\end{tracebox}

\subsubsection*{6. Test suite}

\begin{testbox}
\centering\small
\begin{tabular}{@{}llcl@{}}
\toprule
\textbf{deadends} & \textbf{target} & \textbf{Expected} & \textbf{Why} \\
\midrule
sample 1 & \texttt{0202} & 6 & detour around dead ends \\
\texttt{8888} & \texttt{0009} & 1 & wrap-around $0 \to 9$ \\
8 neighbours of \texttt{8888} & \texttt{8888} & $-1$ & target isolated \\
\texttt{0000} & \texttt{8888} & $-1$ & start is a dead end \\
\texttt{1234} & \texttt{0000} & 0 & target $=$ start \\
\texttt{9999} & \texttt{5555} & 20 & farthest code \\
8 neighbours of \texttt{0202} & \texttt{0202} & $-1$ & walled target near the start \\
500 random codes & random & --- & stress: BFS $=$ bidirectional \\
\bottomrule
\end{tabular}
\end{testbox}

BFS and bidirectional BFS agree on 4000 random instances (random dead-end sets of up to 500 codes, and dense sets concentrated around the target), and iterative deepening confirms every answer up to 5.

\subsubsection*{7. Common mistakes}

\begin{warningbox}[Mistake 1 --- negative digit on ``down'']
\begin{lstlisting}[style=snippet]
int nd = (g - 1) % 10;             // WRONG: g = 0 gives -1, not 9
int nd = (g + 9) % 10;             // RIGHT
\end{lstlisting}
\end{warningbox}

\begin{warningbox}[Mistake 2 --- forgetting that the start can be dead]
If \texttt{"0000"} is a dead end and the code pushes it anyway, BFS happily walks away from a jammed lock and returns a positive number. Test \texttt{dead[0]} before the loop.
\end{warningbox}

\begin{warningbox}[Mistake 3 --- checking dead ends only for the answer]
Testing ``is the \emph{target} dead?'' is not enough: a dead code in the \emph{middle} of a route jams the lock. Dead codes must never enter the queue.
\end{warningbox}

\begin{warningbox}[Mistake 4 --- bidirectional BFS that returns too early]
Returning at the very first touching edge \emph{in the middle of a layer} can give a non-shortest answer in general graphs. Finish the layer and take the minimum (as above), or test the meeting condition at whole-layer granularity.
\end{warningbox}

\subsubsection*{8. Takeaways}

\begin{itemize}[leftmargin=*]
    \item A combination lock is a 4-dimensional wrap-around grid: $10^4$ vertices, 8 edges each; forbidden codes are deleted vertices.
    \item Encode tuples of digits as integers; the move is a place-value addition.
    \item Without obstacles the distance has a closed formula; obstacles are what make a search necessary.
    \item Bidirectional BFS explores two balls of radius $d/2$ instead of one of radius $d$ --- a big win in fast-growing state spaces.
\end{itemize}

\begin{variantsbox}
\variant{Water jugs}{\varin{https://leetcode.com/problems/water-and-jug-problem/}}
  {instead of 4 wheels, two jugs; moves are fill, empty and pour.}
  {state $(a, b)$ = water in each jug, $(x{+}1)(y{+}1)$ states, encoded as $a(y{+}1)+b$; 6 moves per state.}
  {BFS (or DFS) for reachability of a state with $a + b = T$; a gcd formula also exists (Exercise~E4.3).}
\variant{Operations on one number with a bounded range}{\varin{https://leetcode.com/problems/minimum-operations-to-convert-number/}}
  {one number $x$; moves $x \pm a_i$, $x \oplus a_i$; only values in $[0, 1000]$ can be continued.}
  {vertices $0..1000$; values outside are leaves that only count if they equal the goal.}
  {BFS; test the goal on every generated value before the range check (Exercise~E4.5).}
\variant{More wheels, more digits}{\varnone}
  {$k$ wheels with $b$ digits each.}
  {$b^k$ codes as mixed-radix integers; $2k$ neighbours each.}
  {the same BFS; it is feasible while $b^k \cdot 2k$ fits (e.g.\ $k = 6$, $b = 10$: $1.2\cdot10^7$).}
\variant{Turning two adjacent wheels together}{\varnone}
  {a move may turn wheels $i$ and $i+1$ simultaneously by the same $\pm1$.}
  {same vertices; 6 more edges per code (3 adjacent pairs, 2 directions).}
  {unchanged BFS --- adding a kind of move only adds lines to the neighbour generator.}
\variant{Turns with different costs}{\varnone}
  {turning wheel $i$ costs $c_i$.}
  {same graph, edge weights $c_i$.}
  {BFS no longer works; Dijkstra's algorithm (Chapter 7).}
\end{variantsbox}

%============================================================
% CH1 · PATTERN 4 — EXERCISES, HINTS, ANSWERS
%============================================================
\subsection{Pattern 4: Exercises and Going Further}
\label{ssec:ch1-p4-ex}

\begin{patternbox}{Pattern 4 Summary Card}
\begin{tabularx}{\linewidth}{@{}>{\bfseries}l X@{}}
Vertex & a \emph{state}: the smallest snapshot that decides the next moves and the goal test \\
Edge & one move of the statement, generated on the fly by a neighbour function \\
Algorithm & BFS from the start (unit costs); \texttt{dist} as a flat array over integer-encoded states \\
Size & $\#\text{states}\times\#\text{moves}$; unbounded values need a \emph{proven} window \\
Key variations & reverse BFS from the goal, multi-query lookups, path printing, bidirectional BFS, reverse greedy \\
Top traps & no / too small window; marking on pop; forbidden start; negative \texttt{\%}; DFS for ``fewest'' \\
\end{tabularx}
\end{patternbox}

\exlegend
\begin{enumerate}[leftmargin=*, label=\textbf{E4.\arabic*}, itemsep=6pt]
    \setlength{\emergencystretch}{2em}
    \item \exeasy\ \pconcept \cf{1000}\ \exprob{Transformation: from A to B}{Codeforces 727A}{https://codeforces.com/problemset/problem/727/A}. Turn $a$ into $b$ ($1 \le a < b \le 10^9$) using $x \to 2x$ and $x \to 10x + 1$. Print \texttt{NO}, or \texttt{YES}, the length $k$ of a sequence $a = x_1, \dots, x_k = b$ and the sequence itself (any valid one).
    \item \exmedium\ \ptricky \exown{Three buttons.} Two Buttons with a third button that multiplies by 3: moves $x \to 2x$, $x \to 3x$, $x \to x-1$ (the display must stay positive), $1 \le n, m \le 10^4$. Fewest presses from $n$ to $m$. ($5 \to 7$: 3; \ $1 \to 26$: 4.)
    \item \exmedium\ \pclassic \exprob{Water and Jug Problem}{LeetCode 365}{https://leetcode.com/problems/water-and-jug-problem/}\ \lc{Medium}. Jugs of capacities $x$ and $y$ ($\le 10^3$), unlimited water. Moves: fill a jug, empty a jug, pour one into the other until the source is empty or the receiver is full. Can the total amount of water in the two jugs become exactly \texttt{target}?
    \item \exmedium\ \pconcept \cf{1300}\ \exprob{Getting Zero}{Codeforces 1661B}{https://codeforces.com/problemset/problem/1661/B}. One step turns $v$ into $(v+1) \bmod 32768$ or $2v \bmod 32768$. For each of $n \le 32768$ given numbers, print the fewest steps to reach 0.
    \item \exmedium\ \ptricky \exprob{Minimum Operations to Convert Number}{LeetCode 2059}{https://leetcode.com/problems/minimum-operations-to-convert-number/}\ \lc{Medium}. From $x = \texttt{start}$ ($0 \le \texttt{start} \le 1000$), one operation replaces $x$ by $x + a_i$, $x - a_i$ or $x \oplus a_i$ for some $a_i$ of the array, but only while $0 \le x \le 1000$; results outside the range are allowed but final. Fewest operations to obtain \texttt{goal}, or $-1$.
    \item \exmedium\ \precent \exprob{Minimum Number of Operations to Make X and Y Equal}{LeetCode 2998}{https://leetcode.com/problems/minimum-number-of-operations-to-make-x-and-y-equal/}\ \lc{Medium}. One operation on $x$: divide by 11 (if divisible), divide by 5 (if divisible), $-1$ or $+1$. Fewest operations to turn $x$ into $y$ ($1 \le x, y \le 10^4$).
\end{enumerate}

\subsubsection*{Hints}
\begin{enumerate}[leftmargin=*, label=\textbf{E4.\arabic*}]
    \item Both moves only increase the number. Look at the \emph{last} digit of $b$: which move could have produced it?
    \item Repeat the Lemma of Problem~1.4.1 for any multiplication: what is the window now? Is the reverse greedy still safe?
    \item State $=$ the pair of amounts $(a, b)$. There are only $(x{+}1)(y{+}1)$ of them. Bonus: which totals are reachable at all (think of $\gcd(x,y)$)?
    \item $n$ separate BFS runs are $32768^2$ steps. All queries share the same goal --- search backwards from 0 once.
    \item The vertices are $0..1000$. A generated value outside the range is not a vertex, but it may still be the goal.
    \item $+1$ can push $x$ \emph{above} both $x$ and $y$ (to reach a multiple of 11). How far up can an optimal path go, if the answer is never more than $|x - y|$?
\end{enumerate}

\subsubsection*{Answers}

\begin{solutionbox}[E4.1 Transformation: from A to B]
\textbf{Model:} vertices are the numbers $a..b$, edges $x \to 2x$ and $x \to 10x+1$; both increase, so nothing above $b$ matters. \textbf{Backwards} every number has at most one predecessor: an even number can only come from $\times2$ (since $10x + 1$ is odd), a number ending in 1 only from appending 1, any other odd number has none. So the reverse walk is forced --- a path in a graph where every vertex has out-degree at most one (Pattern~5 flavour).
\begin{lstlisting}[style=cpp]
#include <bits/stdc++.h>
using namespace std;

int main() {
    long long a, b;
    cin >> a >> b;
    vector<long long> path = {b};
    while (b > a) {
        if (b % 2 == 0) b /= 2;                 // undo x -> 2x
        else if (b % 10 == 1) b /= 10;          // undo x -> 10x + 1
        else break;                             // no predecessor
        path.push_back(b);
    }
    if (b != a) { cout << "NO\n"; return 0; }
    reverse(path.begin(), path.end());
    cout << "YES\n" << path.size() << "\n";
    for (long long v : path) cout << v << " ";
    cout << "\n";
}
\end{lstlisting}
\begin{center}
\begin{tikzpicture}[node distance=1.35cm]
    \node[gvdst, minimum size=10mm, font=\scriptsize\bfseries] (a) {162};
    \node[gv, minimum size=10mm, font=\scriptsize\bfseries, right=of a] (b) {81};
    \node[gv, minimum size=10mm, font=\scriptsize\bfseries, right=of b] (c) {8};
    \node[gv, minimum size=10mm, font=\scriptsize\bfseries, right=of c] (d) {4};
    \node[gvsrc, minimum size=10mm, font=\scriptsize\bfseries, right=of d] (e) {2};
    \draw[back] (a) -- node[lbl, above]{$\div2$} (b);
    \draw[back] (b) -- node[lbl, above]{drop 1} (c);
    \draw[back] (c) -- node[lbl, above]{$\div2$} (d);
    \draw[back] (d) -- node[lbl, above]{$\div2$} (e);
\end{tikzpicture}
\end{center}
Samples: $(2, 162)$ $\to$ \texttt{YES}, 5, \texttt{2 4 8 81 162}; \ $(4, 42)$ $\to$ \texttt{NO} ($42 \to 21 \to 2 < 4$); \ $(100, 40021)$ $\to$ \texttt{YES}, 5, \texttt{100 200 2001 4002 40021}. Time $\bigO{\log b}$. A forward BFS over $a..b$ would also be correct but has up to $10^9$ vertices.
\end{solutionbox}

\begin{solutionbox}[E4.2 Three buttons]
\textbf{Window.} The Lemma's proof never used the factor 2: a multiplication made at $x \ge m$ adds at least $x \ge 1$ per press that must be removed by blue presses, so it is never optimal. Multiplications happen only at $x < m$, giving at most $3(m-1)$: $\text{HI} = 3\max(n, m)$.
\begin{lstlisting}[style=cpp]
#include <bits/stdc++.h>
using namespace std;

int main() {
    int n, m;
    cin >> n >> m;
    const int HI = 3 * max(n, m);
    vector<int> dist(HI + 1, -1);
    queue<int> q;
    dist[n] = 0;
    q.push(n);
    while (!q.empty()) {
        int x = q.front(); q.pop();
        if (x == m) break;
        int next[3] = {2 * x, 3 * x, x - 1};
        for (int y : next)
            if (y >= 1 && y <= HI && dist[y] == -1) { dist[y] = dist[x] + 1; q.push(y); }
    }
    cout << dist[m] << "\n";
}
\end{lstlisting}
$5 \to 7$: 3 ($5\to4\to8\to7$); \ $4 \to 11$: 2 ($4\to12\to11$); \ $1 \to 26$: 4 ($1\to3\to9\to27\to26$). The tempting reverse greedy ``divide by 3 if possible, else by 2 if even, else add 1'' gives 7 on $1 \to 26$ ($26\to13\to14\to7\to8\to4\to2\to1$): with two multiplications the backward choice is no longer forced, and BFS is the safe method. (Checked: BFS with $\text{HI} = 3\max(n,m)$ equals BFS on $[1, 10^5]$ for all $n, m \le 150$.)
\end{solutionbox}

\begin{solutionbox}[E4.3 Water and Jug Problem]
\textbf{Model:} state $(a, b)$, $0 \le a \le x$, $0 \le b \le y$, id $a(y{+}1) + b$; six moves; goal test ``$a + b = \texttt{target}$'' on every popped state.
\begin{lstlisting}[style=cpp]
#include <bits/stdc++.h>
using namespace std;

class Solution {
public:
    bool canMeasureWater(int x, int y, int target) {
        if (target > x + y) return false;
        vector<char> seen((x + 1) * (y + 1), 0);
        queue<pair<int, int>> q;
        q.push(make_pair(0, 0));
        seen[0] = 1;
        while (!q.empty()) {
            int a = q.front().first, b = q.front().second;
            q.pop();
            if (a + b == target) return true;
            int ab = min(a, y - b), ba = min(b, x - a);       // amounts poured
            int nxt[6][2] = {{x, b}, {a, y}, {0, b}, {a, 0},  // fill / empty
                             {a - ab, b + ab}, {a + ba, b - ba}};  // pour a->b, b->a
            for (auto& s : nxt) {
                int id = s[0] * (y + 1) + s[1];
                if (!seen[id]) { seen[id] = 1; q.push(make_pair(s[0], s[1])); }
            }
        }
        return false;
    }
};
\end{lstlisting}
At most $1001^2 \approx 10^6$ states $\times$ 6 moves: fast. Samples: $(3,5,4) \to$ \texttt{true}, $(2,6,5) \to$ \texttt{false}, $(1,2,3) \to$ \texttt{true}.

\begin{center}
\begin{tikzpicture}[node distance=0.9cm]
    \foreach \s/\lab [count=\i from 0] in {gvsrc/{0,0}, gv/{0,5}, gv/{3,2}, gv/{0,2}, gv/{2,0}, gv/{2,5}, gv/{3,4}, gvdst/{0,4}}
        \node[\s, minimum size=10mm, font=\scriptsize\bfseries] (j\i) at (\i*1.65,0) {\lab};
    \foreach \i [evaluate=\i as \j using int(\i+1)] in {0,...,6} \draw[gpath, -{Stealth[length=2.2mm]}] (j\i) -- (j\j);
\end{tikzpicture}
\end{center}
For $x=3$, $y=5$ only 16 of the 24 states are reachable, and the first state with total 4 is at distance 7, e.g.\ the path above (also $(3,1)$ at distance 7).

\textbf{Mathematical shortcut.} Every reachable amount is a combination $px + qy$, hence a multiple of $g = \gcd(x, y)$, and conversely every multiple of $g$ up to $x + y$ is reachable (B\'ezout). So the answer is $\texttt{target} \le x + y$ and $g \mid \texttt{target}$ --- checked against the BFS for all $x, y \le 40$, $\texttt{target} \le 90$.
\end{solutionbox}

\begin{solutionbox}[E4.4 Getting Zero]
\textbf{Model:} $32768$ residues, edges $v \to v+1$ and $v \to 2v$ (mod $32768$). Every query has the same goal 0, so run \textbf{one BFS from 0 on the reversed edges}: the predecessors of $v$ are $v - 1$ and, if $v$ is even, the two solutions of $2u \equiv v$: $u = v/2$ and $u = v/2 + 16384$.

\begin{center}
\begin{tikzpicture}[node distance=1.4cm]
    \node[gvsrc, minimum size=11mm, font=\scriptsize\bfseries] (z) {0};
    \node[gvseen, minimum size=11mm, font=\scriptsize\bfseries, right=3.2cm of z, yshift=1.6cm] (a) {32767};
    \node[gvseen, minimum size=11mm, font=\scriptsize\bfseries, right=3.2cm of z] (b) {16384};
    \node[gvdead, minimum size=11mm, font=\scriptsize\bfseries, right=3.2cm of z, yshift=-1.6cm] (c) {0};
    \draw[gde] (z) -- node[gw, above, sloped]{undo $+1$} (a);
    \draw[gde] (z) -- node[gw, above]{undo $\times2$} (b);
    \draw[gfaded, dashed, -{Stealth[length=2mm]}] (z) -- node[gw, below, sloped]{$0/2$} (c);
    \node[note, anchor=west, text width=5cm, align=left] at (6.0,0) {reverse BFS from 0: layer 1 is $\{32767, 16384\}$; afterwards every query is \texttt{dist[a\_i]}.};
\end{tikzpicture}
\end{center}
\begin{lstlisting}[style=cpp]
#include <bits/stdc++.h>
using namespace std;

int main() {
    const int M = 32768;
    vector<int> dist(M, -1);
    queue<int> q;
    dist[0] = 0;
    q.push(0);
    while (!q.empty()) {
        int v = q.front(); q.pop();
        vector<int> pre = {(v - 1 + M) % M};          // u + 1 = v
        if (v % 2 == 0) { pre.push_back(v / 2); pre.push_back(v / 2 + M / 2); }  // 2u = v
        for (int u : pre)
            if (dist[u] == -1) { dist[u] = dist[v] + 1; q.push(u); }
    }
    int n;
    cin >> n;
    for (int i = 0; i < n; i++) {
        int a;
        cin >> a;
        cout << dist[a] << (i + 1 < n ? ' ' : '\n');
    }
}
\end{lstlisting}
Sample: $19, 32764, 10240, 49 \to 14\ 4\ 4\ 15$. Time $\bigO{M + n}$. \textbf{Alternative:} since $15$ doublings always reach 0, the answer is at most 15, so it is $\min$ over $\text{add} = 0..15$ of $\text{add} + $ (doublings needed for $v + \text{add}$, which is $15 - $ the number of trailing zero bits); both methods agree on all $32768$ inputs, the maximum being 15.
\end{solutionbox}

\begin{solutionbox}[E4.5 Minimum Operations to Convert Number]
\textbf{Model:} vertices $0..1000$; from $x$ there are $3|\texttt{nums}|$ edges. Values outside $[0, 1000]$ are leaves: compare them with \texttt{goal} \emph{before} the range check (the goal may be $-4$ or $10^9$).
\begin{lstlisting}[style=cpp]
#include <bits/stdc++.h>
using namespace std;

class Solution {
public:
    int minimumOperations(vector<int>& nums, int start, int goal) {
        vector<int> dist(1001, -1);
        queue<int> q;
        dist[start] = 0;
        q.push(start);
        while (!q.empty()) {
            int x = q.front(); q.pop();
            for (int a : nums) {
                long long cand[3] = {(long long)x + a, (long long)x - a, (long long)(x ^ a)};
                for (long long y : cand) {
                    if (y == goal) return dist[x] + 1;         // goal may be out of range
                    if (y < 0 || y > 1000 || dist[y] != -1) continue;
                    dist[y] = dist[x] + 1;
                    q.push((int)y);
                }
            }
        }
        return -1;
    }
};
\end{lstlisting}
Returning when the goal is \emph{generated} is safe here: BFS generates vertices in non-decreasing distance, so the first generation of the goal uses the fewest operations (the goal is never pushed, so the ``popped'' test would miss out-of-range goals). Samples: $([2,4,12], 2, 12) \to 2$, \ $([3,5,7], 0, -4) \to 2$, \ $([2,8,16], 0, 1) \to -1$. Time $\bigO{1001 \cdot 3|\texttt{nums}|} \approx 3\cdot10^6$.
\end{solutionbox}

\begin{solutionbox}[E4.6 Minimum Number of Operations to Make X and Y Equal]
\textbf{Model:} vertices are numbers, edges $v \to v-1$, $v \to v+1$, $v \to v/11$ and $v \to v/5$ (when divisible); BFS from $x$ to $y$. \textbf{Window.} Stepping by $\pm1$ alone costs $|x - y|$, so the answer is at most $|x - y| < M = \max(x, y)$. Only $+1$ increases the value, so an optimal path never climbs above $x + |x - y| \le 2M$; and 0 is the smallest value one can reach. BFS on $[0, 2M]$ is exact. The window \emph{must} go above $\max(x,y)$: in sample 2, $54 \to 55$ is the first move.

\begin{center}
\begin{tikzpicture}[node distance=1.5cm]
    \node[gvsrc] (a) {54};
    \node[gv, right=of a] (b) {55};
    \node[gv, right=of b] (c) {5};
    \node[gv, right=of c] (d) {1};
    \node[gvdst, right=of d] (e) {2};
    \draw[gpath, -{Stealth[length=2.4mm]}] (a) -- node[gw, above=1pt]{$+1$} (b);
    \draw[gpath, -{Stealth[length=2.4mm]}] (b) -- node[gw, above=1pt]{$\div11$} (c);
    \draw[gpath, -{Stealth[length=2.4mm]}] (c) -- node[gw, above=1pt]{$\div5$} (d);
    \draw[gpath, -{Stealth[length=2.4mm]}] (d) -- node[gw, above=1pt]{$+1$} (e);
\end{tikzpicture}
\end{center}
\begin{lstlisting}[style=cpp]
#include <bits/stdc++.h>
using namespace std;

class Solution {
public:
    int minimumOperationsToMakeEqual(int x, int y) {
        const int HI = 2 * max(x, y);                   // proven window [0, HI]
        vector<int> dist(HI + 1, -1);
        queue<int> q;
        dist[x] = 0;
        q.push(x);
        while (!q.empty()) {
            int v = q.front(); q.pop();
            if (v == y) return dist[v];
            vector<int> nxt = {v - 1, v + 1};
            if (v % 11 == 0) nxt.push_back(v / 11);
            if (v % 5 == 0) nxt.push_back(v / 5);
            for (int w : nxt)
                if (w >= 0 && w <= HI && dist[w] == -1) { dist[w] = dist[v] + 1; q.push(w); }
        }
        return -1;                                      // unreachable: +1/-1 connect all
    }
};
\end{lstlisting}
Samples: $(26, 1) \to 3$, \ $(54, 2) \to 4$, \ $(25, 30) \to 5$. With $x \le y$ no division helps (it only lowers the value), and the answer is $y - x$. Checked: the window $[0, 2M]$ gives the same answers as $[0, 10^5]$ for all $x, y \le 300$ and 300 random pairs up to $10^4$, whereas the window $[0, M]$ fails on 30382 pairs with $x \le 2000$, $y \le 50$. Time $\bigO{M}$.
\end{solutionbox}


%============================================================
\clearpage
\section{Pattern 5: Functional Graphs --- One Outgoing Edge}
\label{sec:ch1-p5}
%============================================================
%============================================================
% CH1 · PATTERN 5 — FUNCTIONAL GRAPHS (ONE OUTGOING EDGE): THEORY
%============================================================
\subsection{Pattern Theory}
\label{ssec:ch1-p5-theory}

\begin{patternbox}{Pattern 5 at a Glance}
Every vertex has \textbf{exactly one outgoing edge} $v \to f(v)$, so the whole graph is just an array \texttt{f[]}. There is nothing to choose: from a start vertex the walk $v, f(v), f(f(v)), \dots$ is forced. Because there are only $n$ vertices the walk must repeat, and from then on it goes round and round. Every walk looks like the Greek letter $\rho$ (rho): a \textbf{tail} of $\mu$ vertices followed by a \textbf{cycle} of $\lambda$ vertices.

\begin{center}
\begin{tikzpicture}[x=1.1cm, y=1.1cm]
    \node[gvsrc] (s) at (0,0) {$s$};
    \node[gv] (t1) at (1.2,0) {};
    \node[gv] (t2) at (2.4,0) {};
    \node[gvdst] (e) at (3.6,0) {$e$};
    \node[gvseen] (c1) at (4.5,0.75) {};
    \node[gvseen] (c2) at (5.6,0.45) {};
    \node[gvseen] (c3) at (5.6,-0.45) {};
    \node[gvseen] (c4) at (4.5,-0.75) {};
    \draw[gde] (s) -- (t1); \draw[gde] (t1) -- (t2); \draw[gde] (t2) -- (e);
    \draw[gde] (e) -- (c1); \draw[gde] (c1) -- (c2); \draw[gde] (c2) -- (c3);
    \draw[gde] (c3) -- (c4); \draw[gde] (c4) -- (e);
    \draw[decorate, decoration={brace, amplitude=5pt, mirror}, darkgray] (-0.3,-0.55) -- (3.2,-0.55)
        node[midway, below=5pt, font=\scriptsize]{tail: $\mu = 3$ vertices};
    \node[font=\scriptsize] at (5.05,-1.25) {cycle: $\lambda = 5$};
    \node[note, text width=4.3cm, align=left, anchor=west] at (6.4,0)
        {$e$ = \textbf{entry point}: the first vertex the walk visits twice. The walk from $s$ visits $\mu+\lambda \le n$ different vertices.};
\end{tikzpicture}
\end{center}
\begin{recurrencebox}
\centering\small
walk from $s$ and remember the step $\texttt{first}[v]$ at which each vertex is reached; at the first vertex $v$ seen twice:\\[3pt]
$\mu = \texttt{first}[v]$, \qquad $\lambda = \text{step} - \texttt{first}[v]$,\\[2pt]
$f^{k}(s) = $ the vertex at step $\mu + (k-\mu) \bmod \lambda$ \;(for $k \ge \mu$).\\[3pt]
All vertices at once: walk from every new vertex with three colours \emph{new / on this walk / finished} --- \bigO{n} in total.
\end{recurrencebox}
\end{patternbox}

%------------------------------------------------------------
\subsubsection*{From Pattern 4 to Pattern 5}

In Pattern~4 a state had \emph{several} possible next states (press button A or button B, turn any of the four wheels), and we needed BFS to explore all of them. Now look at processes with \textbf{no choice at all}: a teacher who always goes to the student just accused, a teleporter that always sends you to the same town, a pseudo-random generator $x \mapsto (x^2+1) \bmod m$, a linked list whose \texttt{next} pointers may loop. Each state has exactly one successor --- the state graph has out-degree~1. Such a graph is called a \textbf{functional graph}, because its edges are the graph of a function $f : V \to V$ (Lesson~\ref{ssec:ch1-vocabulary}).

With no choices there is no search: we simply \emph{follow the arrows}. The interesting questions are about the \emph{shape} the arrows form. Where does the walk start repeating? How long is the loop? Which vertices lie on loops at all? Where will I be after $10^{18}$ steps?

\subsubsection*{From an array to a picture}

The input of a functional-graph problem is almost always a single array: \texttt{p[i]}, \texttt{next[i]}, \texttt{a[i]}, \texttt{favorite[i]}, \texttt{edges[i]}. Index $i$ is a vertex, the value is where its arrow points. No adjacency list is needed --- the array \emph{is} the adjacency list.

\begin{diagrambox}
\begin{tikzpicture}
    \dpindex{ix}{1.5}{$v$}{0}{9}
    \dparray{fa}{0.8}{$f[v]$}{1/cell, 2/cell, 3/cell, 4/cell, 2/cell, 4/cell, 5/cell, 8/cell, 9/cell, 8/cell}
    \begin{scope}[yshift=-1.9cm, xshift=3.6cm]
        \foreach \v/\x/\y in {0/-2.8/0.8, 1/-1.4/0.8, 2/0/0.8, 3/1.3/0, 4/0/-0.8, 5/-1.4/-0.8, 6/-2.8/-0.8, 7/2.9/0, 8/4.2/0, 9/5.5/0}
            \node[gv] (n\v) at (\x,\y) {\v};
        \foreach \u/\w in {0/1, 1/2, 2/3, 3/4, 4/2, 5/4, 6/5, 7/8}
            \draw[gde] (n\u) -- (n\w);
        \draw[gde] (n8) to[bend left=35] (n9);
        \draw[gde] (n9) to[bend left=35] (n8);
    \end{scope}
\end{tikzpicture}
\figcap{The array $f = [1,2,3,4,2,4,5,8,9,8]$ and its functional graph: 10 vertices, 10 arrows, two pieces. Each piece is one cycle ($2\to3\to4\to2$ and $8 \leftrightarrows 9$) with trees hanging into it.}
\end{diagrambox}

We call this example graph $G$ and use it throughout the section.

%------------------------------------------------------------
\subsubsection*{The shape theorem}

\begin{theorembox}[Shape of a functional graph]
\begin{enumerate}[leftmargin=*]
    \item \textbf{Every walk is a $\rho$.} From any start $s$ the walk $s, f(s), f^2(s), \dots$ visits $\mu$ tail vertices and then cycles forever through $\lambda$ cycle vertices, with $\mu \ge 0$, $\lambda \ge 1$ and $\mu + \lambda \le n$.
    \item \textbf{One cycle per piece.} Every connected piece (ignoring arrow directions) contains \emph{exactly one} cycle; all its other vertices form trees whose arrows point towards the cycle.
\end{enumerate}
\end{theorembox}

\textbf{Proof of 1.} Among the $n+1$ vertices $f^0(s), f^1(s), \dots, f^n(s)$ two must be equal (pigeonhole). Let $j$ be the first step that reaches an already visited vertex, say $f^j(s) = f^i(s)$ with $i < j$. The walk is deterministic, so from step $i$ on it repeats the steps $i, \dots, j-1$ forever. Set $\mu = i$ and $\lambda = j - i$; the vertices at steps $0..j-1$ are all different, so $\mu + \lambda = j \le n$.

\textbf{Proof of 2.} A piece with $k$ vertices has exactly $k$ edges (one leaves each vertex). A connected graph with $k$ vertices and $k-1$ edges is a tree; one more edge closes exactly one cycle (Lesson~\ref{ssec:ch1-vocabulary}). By part 1 the walk from any vertex of the piece ends in a \emph{directed} cycle, so that unique cycle is a directed one, and every vertex of the piece flows into it. \hfill$\blacksquare$

\begin{diagrambox}
\begin{tikzpicture}[x=1cm, y=1cm]
    % (a) permutation
    \begin{scope}
        \node[font=\small\bfseries] at (1.2,1.75) {(a) permutation};
        \foreach \v/\x/\y in {0/0/1, 1/1.1/1, 2/0.55/0.1, 3/1.9/0.1, 4/2.4/1}
            \node[gv] (a\v) at (\x,\y) {\v};
        \draw[gde] (a0) -- (a1); \draw[gde] (a1) -- (a2); \draw[gde] (a2) -- (a0);
        \draw[gde] (a3) to[bend left=30] (a4); \draw[gde] (a4) to[bend left=30] (a3);
        \node[note, text width=3.8cm] at (1.2,-0.85) {in-degree also 1:\\ \textbf{only cycles}, no tails};
    \end{scope}
    % (b) functional
    \begin{scope}[xshift=4.9cm]
        \node[font=\small\bfseries] at (1.2,1.75) {(b) functional graph};
        \foreach \v/\x/\y in {0/0/1, 1/1.1/1, 2/2.2/1, 3/2.2/0, 4/1.1/0}
            \node[gv] (b\v) at (\x,\y) {\v};
        \draw[gde] (b0) -- (b1); \draw[gde] (b1) -- (b2); \draw[gde] (b2) -- (b3);
        \draw[gde] (b3) -- (b4); \draw[gde] (b4) -- (b1);
        \node[note, text width=3.8cm] at (1.2,-0.85) {out-degree exactly 1:\\ \textbf{every piece is a $\rho$-shape} (one cycle + trees)};
    \end{scope}
    % (c) out-degree at most 1
    \begin{scope}[xshift=9.8cm]
        \node[font=\small\bfseries] at (1.2,1.75) {(c) out-degree $\le 1$};
        \foreach \v/\x/\y in {0/0/1, 1/1.1/1, 3/0/0, 4/1.1/0}
            \node[gv] (c\v) at (\x,\y) {\v};
        \node[gvdst] (c2) at (2.2,0.5) {2};
        \draw[gde] (c0) -- (c1); \draw[gde] (c1) -- (c2);
        \draw[gde] (c3) -- (c4); \draw[gde] (c4) -- (c2);
        \draw[-{Stealth[length=2mm]}, warningred, dashed] (c2) -- (3.0,0.5);
        \node[font=\scriptsize\bfseries, text=warningred, anchor=west] at (3.0,0.5) {$-1$};
        \node[note, text width=3.8cm] at (1.2,-0.85) {some \texttt{f[v] = -1}: a piece may be a \textbf{tree ending in a sink} (no cycle)};
    \end{scope}
\end{tikzpicture}
\figcap{Three relatives. Permutations are the special case where every vertex also has in-degree 1. In the ``at most one'' version (LeetCode writes \texttt{edges[i] = -1}) each piece has one cycle or none.}
\end{diagrambox}

%------------------------------------------------------------
\subsubsection*{When to recognise it (and when not)}

\begin{itemize}[leftmargin=*]
    \item the input is \textbf{one array} \texttt{a[1..n]} with ``$i$ points to / reports / likes / teleports to / is followed by $a_i$'';
    \item a \textbf{deterministic process}: the next state depends only on the current one ($x \mapsto g(x)$), and the question is about repeating, looping, or the state after $k$ steps;
    \item words like ``eventually'', ``forever'', ``the first student visited twice'', ``after $K = 10^{18}$ moves'', ``longest cycle'', ``everyone has exactly one favourite''.
\end{itemize}
\textbf{Not this pattern:} if a vertex can have \emph{two or more} outgoing edges, the walk is not forced; cycle detection then needs a real DFS with colours (Chapter~2 (DFS)) and the cycle structure is studied with strongly connected components later in the book. If the graph is undirected, ``one cycle per piece'' is a different object (a pseudotree, Chapter~13 (trees)). And if there are \emph{many} queries ``where am I after $k$ steps from $v$?'', walking once per query is too slow --- that is binary lifting, Chapter~11 (functional graphs and binary lifting); we only preview the idea below.

%------------------------------------------------------------
\subsubsection*{Algorithm 1 --- one start: measure the $\rho$}

Walk from $s$ and write down at which step each vertex is reached. The first time we arrive at a vertex that already has a step number, we have found the entry point $e$, the tail length $\mu = \texttt{first}[e]$ and the cycle length $\lambda = \text{step} - \texttt{first}[e]$.

\begin{diagrambox}
\begin{tikzpicture}[x=1.45cm]
    \foreach \v/\i/\s in {0/0/gvsrc, 1/1/gv, 2/2/gvdst, 3/3/gvseen, 4/4/gvseen}{
        \node[\s] (w\i) at (\i,0) {\v};
        \node[gdist, below=1mm of w\i] {step \i};
    }
    \node[gvdst, dashed] (w5) at (5,0) {2};
    \node[gdist, below=1mm of w5, text=warningred] {step 5};
    \foreach \i [evaluate=\i as \j using int(\i+1)] in {0,...,4} \draw[gde] (w\i) -- (w\j);
    \draw[gback] (w5.north) to[bend right=35] node[lbl, above]{seen before, at step 2} (w2.north);
    \draw[decorate, decoration={brace, amplitude=4pt, mirror}, darkgray] (-0.25,-0.85) -- (1.25,-0.85)
        node[midway, below=4pt, font=\scriptsize]{$\mu = \texttt{first}[2] = 2$};
    \draw[decorate, decoration={brace, amplitude=4pt, mirror}, darkgray] (1.75,-0.85) -- (4.25,-0.85)
        node[midway, below=4pt, font=\scriptsize]{$\lambda = 5 - 2 = 3$};
\end{tikzpicture}
\figcap{The walk from $0$ in $G$: $0,1,2,3,4$ and then $2$ again. Entry point $e = 2$, tail $\mu = 2$, cycle $\lambda = 3$. From $6$ the walk is $6,5,4,2,3$ and then $4$: entry point $4$, again $\mu = 2$, $\lambda = 3$ --- same cycle, different entry point.}
\end{diagrambox}

\textbf{Huge $k$.} Once we know $\mu$ and $\lambda$, the position after $k$ steps needs no long simulation: after the tail the walk is periodic with period $\lambda$, so for $k \ge \mu$ we may replace $k$ by $\mu + (k - \mu) \bmod \lambda$, which is smaller than $\mu + \lambda \le n$. In $G$ from $0$: $k = 10$ becomes $2 + 8 \bmod 3 = 4$, and step $4$ of the walk is vertex $4$; even $k = 10^{18}$ becomes $2 + (10^{18}-2) \bmod 3 = 4$, so $f^{10^{18}}(0) = 4$.

%------------------------------------------------------------
\subsubsection*{Algorithm 2 --- all vertices: walk with three colours}

Running Algorithm~1 from every vertex costs up to $\bigO{n}$ per start, $\bigO{n^2}$ in total --- fine for $n = 1000$, hopeless for $n = 2\cdot 10^5$. The fix is the idea behind \texttt{visited} in Lesson~\ref{ssec:ch1-first-traversal}: \textbf{never walk over a vertex twice}. Give each vertex one of three states:
\begin{center}
\legendchip{gv}{0 = new} \qquad \legendchip{gvseen}{1 = on the current walk} \qquad \legendchip{gvdone}{2 = finished}
\end{center}
For each start $s$ that is still new, walk while the vertices are new, colouring them 1 and storing them in \texttt{path}. The walk stops at the first vertex $v$ that is not new (or at $-1$):
\begin{itemize}[leftmargin=*]
    \item $v$ has \textbf{state 1}: we bit our own tail --- the part of \texttt{path} from $v$ to the end is a \textbf{new cycle};
    \item $v$ has \textbf{state 2}: we ran into territory explored by an earlier walk; its cycle was already found, so there is \textbf{no new cycle};
\end{itemize}
then colour the whole \texttt{path} 2.

\begin{diagrambox}
\begin{tikzpicture}
    % snapshot 1
    \begin{scope}
        \foreach \v/\x/\y/\s in {0/-2.8/0.8/gvseen, 1/-1.4/0.8/gvseen, 2/0/0.8/gvseen, 3/1.3/0/gvseen, 4/0/-0.8/gvseen,
                                 5/-1.4/-0.8/gv, 6/-2.8/-0.8/gv, 7/2.9/0/gv, 8/4.2/0/gv, 9/5.5/0/gv}
            \node[\s] (a\v) at (\x,\y) {\v};
        \foreach \u/\w in {0/1, 1/2, 2/3, 3/4, 5/4, 6/5, 7/8} \draw[gde] (a\u) -- (a\w);
        \draw[gback] (a4) -- (a2);
        \draw[gde] (a8) to[bend left=35] (a9); \draw[gde] (a9) to[bend left=35] (a8);
        \node[note, text width=5.2cm, align=left, anchor=west] at (6.3,0) {\textbf{$s=0$:} \texttt{path} $= 0,1,2,3,4$; the walk stops at $2$, which has state 1 $\Rightarrow$ new cycle $2,3,4$ of length $3$ (red arrow closes it).};
    \end{scope}
    % snapshot 2
    \begin{scope}[yshift=-2.75cm]
        \foreach \v/\x/\y/\s in {0/-2.8/0.8/gvdone, 1/-1.4/0.8/gvdone, 2/0/0.8/gvdone, 3/1.3/0/gvdone, 4/0/-0.8/gvdone,
                                 5/-1.4/-0.8/gvseen, 6/-2.8/-0.8/gv, 7/2.9/0/gv, 8/4.2/0/gv, 9/5.5/0/gv}
            \node[\s] (b\v) at (\x,\y) {\v};
        \foreach \u/\w in {0/1, 1/2, 2/3, 3/4, 4/2, 6/5, 7/8} \draw[gde] (b\u) -- (b\w);
        \draw[gcross] (b5) -- (b4);
        \draw[gde] (b8) to[bend left=35] (b9); \draw[gde] (b9) to[bend left=35] (b8);
        \node[note, text width=5.2cm, align=left, anchor=west] at (6.3,0) {\textbf{$s=5$:} \texttt{path} $= 5$; the walk stops at $4$, which has state 2 $\Rightarrow$ old territory, \textbf{no} new cycle. The same happens for $s = 6$ (stops at $5$).};
    \end{scope}
    % snapshot 3
    \begin{scope}[yshift=-5.5cm]
        \foreach \v/\x/\y/\s in {0/-2.8/0.8/gvdone, 1/-1.4/0.8/gvdone, 2/0/0.8/gvdone, 3/1.3/0/gvdone, 4/0/-0.8/gvdone,
                                 5/-1.4/-0.8/gvdone, 6/-2.8/-0.8/gvdone, 7/2.9/0/gvseen, 8/4.2/0/gvseen, 9/5.5/0/gvseen}
            \node[\s] (c\v) at (\x,\y) {\v};
        \foreach \u/\w in {0/1, 1/2, 2/3, 3/4, 4/2, 5/4, 6/5, 7/8} \draw[gde] (c\u) -- (c\w);
        \draw[gde] (c8) to[bend left=35] (c9);
        \draw[gback] (c9) to[bend left=35] (c8);
        \node[note, text width=5.2cm, align=left, anchor=west] at (6.3,0) {\textbf{$s=7$:} \texttt{path} $= 7,8,9$; the walk stops at $8$, state 1 $\Rightarrow$ new cycle $8,9$ of length $2$. Starts $8$ and $9$ are skipped later (already finished).};
    \end{scope}
\end{tikzpicture}
\figcap{Algorithm 2 on $G$. Orange = state 1 (current walk), blue = state 2 (finished), white = new. Each vertex is put on a \texttt{path} exactly once; result \texttt{cyc} $= [0,0,3,3,3,0,0,0,2,2]$.}
\end{diagrambox}

\begin{invariantbox}[Why Algorithm 2 is correct and linear]
\textbf{Invariant:} at every moment the vertices with state 1 are exactly the vertices of \texttt{path}, in walk order, and every cycle made of finished (state 2) vertices has already been reported.

\textbf{New cycle $\Leftrightarrow$ state 1.} If the walk stops at a state-1 vertex $v$, then $v$ is in \texttt{path}, and following arrows from $v$ along \texttt{path} leads back to $v$: that is a cycle, and it is new because its vertices were new when this walk started. If the walk stops at a state-2 vertex, every vertex reachable from it is finished too (an earlier walk followed exactly these arrows), so the $\rho$ of the current start ends in an old cycle; by the shape theorem there is no other cycle it could reach.

\textbf{Linear time:} a vertex is appended to \texttt{path} only while it is new, and then it leaves state 0 forever --- at most $n$ appends in total. Finding where the cycle starts scans the current \texttt{path} once. Total \bigO{n} time and memory.
\end{invariantbox}

%------------------------------------------------------------
\subsubsection*{Algorithm 3 --- peeling: who is on a cycle?}

A vertex with in-degree $0$ cannot lie on a cycle (a cycle would enter it). Delete it; this may lower the in-degree of its successor to $0$; delete that one too; and so on --- this is Kahn's queue from Chapter~5 (DAGs). What survives are precisely the cycle vertices.

\begin{diagrambox}
\begin{tikzpicture}
    \begin{scope}
        \foreach \v/\x/\y/\s in {0/-2.8/0.8/gvseen, 1/-1.4/0.8/gv, 2/0/0.8/gv, 3/1.3/0/gv, 4/0/-0.8/gv,
                                 5/-1.4/-0.8/gv, 6/-2.8/-0.8/gvseen, 7/2.9/0/gvseen, 8/4.2/0/gv, 9/5.5/0/gv}
            \node[\s] (a\v) at (\x,\y) {\v};
        \foreach \u/\w in {0/1, 1/2, 2/3, 3/4, 4/2, 5/4, 6/5, 7/8} \draw[gde] (a\u) -- (a\w);
        \draw[gde] (a8) to[bend left=35] (a9); \draw[gde] (a9) to[bend left=35] (a8);
        \node[note, text width=5.2cm, align=left, anchor=west] at (6.3,0) {in-degrees $[0,1,2,1,2,1,0,0,2,1]$;\\ the queue starts with $0, 6, 7$ (orange).};
    \end{scope}
    \begin{scope}[yshift=-2.75cm]
        \foreach \v/\x/\y/\s in {0/-2.8/0.8/gvdead, 1/-1.4/0.8/gvseen, 2/0/0.8/gv, 3/1.3/0/gv, 4/0/-0.8/gv,
                                 5/-1.4/-0.8/gvseen, 6/-2.8/-0.8/gvdead, 7/2.9/0/gvdead, 8/4.2/0/gv, 9/5.5/0/gv}
            \node[\s] (b\v) at (\x,\y) {\v};
        \foreach \u/\w in {1/2, 2/3, 3/4, 4/2, 5/4} \draw[gde] (b\u) -- (b\w);
        \foreach \u/\w in {0/1, 6/5, 7/8} \draw[gfaded, dashed] (b\u) -- (b\w);
        \draw[gde] (b8) to[bend left=35] (b9); \draw[gde] (b9) to[bend left=35] (b8);
        \node[note, text width=5.2cm, align=left, anchor=west] at (6.3,0) {peel $0, 6, 7$: the in-degrees of $1$ and $5$ drop to $0$, they join the queue; $8$ drops to $1$ only.};
    \end{scope}
    \begin{scope}[yshift=-5.5cm]
        \foreach \v/\x/\y/\s in {0/-2.8/0.8/gvdead, 1/-1.4/0.8/gvdead, 2/0/0.8/gvsel, 3/1.3/0/gvsel, 4/0/-0.8/gvsel,
                                 5/-1.4/-0.8/gvdead, 6/-2.8/-0.8/gvdead, 7/2.9/0/gvdead, 8/4.2/0/gvsel, 9/5.5/0/gvsel}
            \node[\s] (c\v) at (\x,\y) {\v};
        \foreach \u/\w in {2/3, 3/4, 4/2} \draw[gde] (c\u) -- (c\w);
        \foreach \u/\w in {0/1, 1/2, 5/4, 6/5, 7/8} \draw[gfaded, dashed] (c\u) -- (c\w);
        \draw[gde] (c8) to[bend left=35] (c9); \draw[gde] (c9) to[bend left=35] (c8);
        \node[note, text width=5.2cm, align=left, anchor=west] at (6.3,0) {peel $1, 5$: $2$ and $4$ keep in-degree $1$; the queue is empty. Survivors (purple) $2,3,4,8,9$ = the cycle vertices.};
    \end{scope}
\end{tikzpicture}
\figcap{Peeling $G$. Grey vertices are deleted, in the order $0, 6, 7, 1, 5$.}
\end{diagrambox}

\textbf{Why it works.} A cycle vertex always keeps the arrow from its cycle predecessor, which is never deleted (induction on the deletion order), so its in-degree never reaches 0. A tree vertex hangs in a tree whose arrows point to the cycle; the vertices farthest from the cycle have in-degree 0, and peeling them exposes the next layer --- every tree vertex is eventually deleted. Peeling is handy when information must \textbf{flow along the trees}, e.g.\ the longest chain that ends at each cycle vertex (Exercise E5.7), because a vertex is always deleted \emph{before} its successor.

%------------------------------------------------------------
\subsubsection*{Algorithm 4 --- Floyd's tortoise and hare (\bigO{1} memory)}

Two pointers start at $s$; per round the tortoise moves 1 step and the hare 2 steps. Inside the cycle the hare gains one step per round, so they meet. Then put the tortoise back at $s$ and move both one step at a time: they meet exactly at the entry point $e$, after $\mu$ steps. One more lap around the cycle gives $\lambda$.

\begin{diagrambox}
\begin{tikzpicture}
    \node[font=\small] at (0,0) {
    \begin{tabular}{@{}lcccc@{}}
    \toprule
    round $t$ & 0 & 1 & 2 & 3 \\
    \midrule
    tortoise $f^{t}(0)$ & 0 & 1 & 2 & \textbf{3} \\
    hare $f^{2t}(0)$ & 0 & 2 & 4 & \textbf{3} \\
    \bottomrule
    \end{tabular}};
    \node[note, text width=6.6cm, align=left, anchor=west] at (3.0,0)
        {In $G$ from $s = 0$ they meet at vertex $3$ after 3 rounds. Restart the tortoise at $0$, both move 1 step: $(0,3) \to (1,4) \to (2,2)$ --- they meet at the entry point $2$ after $\mu = 2$ steps; one lap from $2$ gives $\lambda = 3$.};
\end{tikzpicture}
\end{diagrambox}

Why the second phase works: when they first meet after $t$ rounds, the hare has walked $t$ more steps than the tortoise, and that surplus is a whole number of laps, a multiple of $\lambda$. So a pointer that walks $\mu$ more steps from the meeting point stands at step $t + \mu$ of the walk, which is on the cycle and equivalent to step $\mu$ modulo $\lambda$: the entry point --- exactly where the restarted tortoise arrives after $\mu$ steps. Floyd's method is the classic answer to ``detect a loop in a linked list with \bigO{1} memory''; in contests memory is rarely that tight, and the \texttt{first[]} array of Algorithm~1 is simpler. When the input is read-only and only \bigO{1} extra memory is allowed, Floyd is the tool; \hyperref[ssec:ch1-find-duplicate]{Problem~1.5.3} (Find the Duplicate Number) is such a case and proves the second phase in detail.

%------------------------------------------------------------
\subsubsection*{A preview: many $k$-th successor queries}

Algorithm~1 answers \emph{one} question ``where is $f^k(s)$?'' in \bigO{n}. With $q = 2\cdot10^5$ different starts that is too slow. The idea of Chapter~11 (functional graphs and binary lifting) is \textbf{doubling}: precompute $f^{1}, f^{2}, f^{4}, f^{8}, \dots$ for every vertex (using $f^{2^{j+1}}(v) = f^{2^j}(f^{2^j}(v))$), then write $k$ in binary and jump: $f^{13} = f^{8} \circ f^{4} \circ f^{1}$ --- only about $\log_2 k$ jumps per query.

\begin{diagrambox}
\begin{tikzpicture}[x=0.95cm]
    \foreach \i in {0,...,13}{
        \ifnum\i=0 \node[gvsrc, minimum size=6mm, font=\scriptsize\bfseries] (q\i) at (\i,0) {\i};
        \else\ifnum\i=13 \node[gvdst, minimum size=6mm, font=\scriptsize\bfseries] (q\i) at (\i,0) {\i};
        \else \node[gv, minimum size=6mm, font=\scriptsize\bfseries] (q\i) at (\i,0) {\i}; \fi\fi
    }
    \foreach \i [evaluate=\i as \j using int(\i+1)] in {0,...,12} \draw[gfaded, -{Stealth[length=1.6mm]}] (q\i) -- (q\j);
    \draw[gpath, -{Stealth[length=3mm]}] (q0.north) to[bend left=30] node[lbl, above]{jump $2^3 = 8$} (q8.north);
    \draw[gpath, -{Stealth[length=3mm]}] (q8.north) to[bend left=40] node[lbl, above]{$2^2 = 4$} (q12.north);
    \draw[gpath, -{Stealth[length=3mm]}] (q12.south) to[bend right=70] node[lbl, below]{$2^0 = 1$} (q13.south);
\end{tikzpicture}
\figcap{Doubling: 13 single steps (grey) become three precomputed jumps, $13 = 8 + 4 + 1$. The walk is drawn as a line of step numbers; in a functional graph the jumps simply wrap around the cycle. Details and proofs: Chapter~11.}
\end{diagrambox}

%------------------------------------------------------------
\subsubsection*{Variations map}

\begin{diagrambox}
\begin{tikzpicture}[
    var/.style={rectangle, rounded corners=2pt, draw=patternborder, fill=patternbg,
                font=\scriptsize, align=center, text width=3.25cm, minimum height=1.3cm, inner sep=3pt}]
    \node[var, fill=answerfill, draw=accentorange, text width=5cm, font=\scriptsize\bfseries] (root) at (0,0)
        {Follow the arrows $v \to f(v)$\\every piece = one cycle + trees};
    \node[var] (v1) at (-5.4,-2) {\textbf{first repeated vertex}\\from every start\\(Badge, 1.5.1)};
    \node[var] (v2) at (-1.8,-2) {\textbf{longest cycle}\\out-degree $\le 1$\\(Longest Cycle, 1.5.2)};
    \node[var] (v3) at (1.8,-2)  {\textbf{permutation}\\only cycles: lengths, lcm\\(E5.1, E5.5)};
    \node[var] (v4) at (5.4,-2)  {\textbf{position after $K$ steps}\\$K$ huge: skip laps\\(E5.4)};
    \node[var] (v5) at (-5.4,-3.6) {\textbf{implicit $f$}\\$x \mapsto g(x)$: $\rho$ of a generator\\(E5.2)};
    \node[var] (v6) at (-1.8,-3.6) {\textbf{distances along walks}\\where two walkers meet\\(E5.3)};
    \node[var] (v7) at (1.8,-3.6)  {\textbf{one choice per cycle}\\cheapest vertex of each cycle\\(E5.6)};
    \node[var] (v8) at (5.4,-3.6)  {\textbf{cycles + longest tails}\\peeling, chain lengths, layers\\(E5.7, E5.8)};
    \node[var, text width=5cm] (v9) at (0,-5.2) {\textbf{entry point with \bigO{1} memory}\\read-only array: Floyd's two pointers\\(Find the Duplicate, 1.5.3)};
    \node[draw=patternborder, dashed, rounded corners=4pt, fit=(v1)(v4)(v5)(v8)(v9), inner sep=5pt] (grp) {};
    \draw[flow] (root.south) -- (grp.north);
    \node[note, below=2mm of grp] {many $k$-th successor queries (``after $k$ steps'' for many starts) $\Rightarrow$ binary lifting, Chapter~11};
\end{tikzpicture}
\end{diagrambox}

\textbf{Twists and how to handle them}
\begin{itemize}[leftmargin=*]
    \item \textbf{Missing arrows} (\texttt{f[v] = -1}): stop the walk at $-1$; such a walk contains no cycle.
    \item \textbf{Self-loops} (\texttt{f[v] = v}) are cycles of length 1: a student who blames himself gets his second hole immediately.
    \item \textbf{1-indexed input}: allocate $n+1$ entries, or subtract 1 while reading --- but do it consistently.
    \item \textbf{``Something per cycle''} (cheapest vertex, cycle length, number of cycles): Algorithm 2 reports each cycle exactly once.
    \item \textbf{``Something per vertex''} (its entry point, its distance to the cycle): first find the cycles, then fill the tails from the cycle outwards (Problem~1.5.1).
    \item \textbf{Read-only input, \bigO{1} extra memory}: no \texttt{first[]} or \texttt{state[]} array is allowed; walk with two pointers (Algorithm 4, \hyperref[ssec:ch1-find-duplicate]{Problem~1.5.3}).
\end{itemize}

%------------------------------------------------------------
\subsubsection*{Code templates}

\begin{lstlisting}[style=cpp]
#include <bits/stdc++.h>
using namespace std;
typedef long long ll;

// (1) Walk from ONE start s: tail length mu and cycle length lam.
//     Needs out-degree exactly 1 (f[v] is always a vertex).
void rho(const vector<int>& f, int s, int& mu, int& lam) {
    vector<int> first(f.size(), -1);      // first[v] = step at which v was first reached
    int v = s, step = 0;
    while (first[v] == -1) {              // keep walking until a vertex repeats
        first[v] = step++;
        v = f[v];
    }
    mu = first[v];                        // v is the entry point of the cycle
    lam = step - first[v];                // steps between the two visits of v
}

// (2) k-th successor f^k(s) for huge k: skip whole laps of the cycle
int kth(const vector<int>& f, int s, ll k) {
    int mu, lam;
    rho(f, s, mu, lam);
    if (k > mu) k = mu + (k - mu) % lam;  // after the tail the walk has period lam
    int v = s;
    for (ll i = 0; i < k; i++) v = f[v];  // now k < mu + lam <= n
    return v;
}

// (3) ALL vertices: walk with three states. f[v] == -1 means "no outgoing edge".
//     cyc[v] = length of the cycle through v, or 0 if v is not on a cycle.
vector<int> cycleLengths(const vector<int>& f) {
    int n = f.size();
    vector<int> state(n, 0), cyc(n, 0);   // 0 = new, 1 = on the current walk, 2 = finished
    for (int s = 0; s < n; s++) {
        if (state[s] != 0) continue;
        vector<int> path;                 // the current walk, in order
        int v = s;
        while (v != -1 && state[v] == 0) {
            state[v] = 1;
            path.push_back(v);
            v = f[v];
        }
        if (v != -1 && state[v] == 1) {   // we bit our own tail: a NEW cycle
            int len = path.end() - find(path.begin(), path.end(), v);
            for (int i = (int)path.size() - len; i < (int)path.size(); i++) cyc[path[i]] = len;
        }
        for (int u : path) state[u] = 2;  // the whole walk is finished
    }
    return cyc;
}

// (4) Peeling: delete vertices of in-degree 0 again and again (Kahn).
//     The survivors are exactly the vertices that lie on cycles.
vector<bool> onCycle(const vector<int>& f) {
    int n = f.size();
    vector<int> indeg(n, 0);
    for (int v = 0; v < n; v++)
        if (f[v] != -1) indeg[f[v]]++;
    queue<int> q;
    for (int v = 0; v < n; v++)
        if (indeg[v] == 0) q.push(v);
    vector<bool> alive(n, true);
    while (!q.empty()) {
        int v = q.front(); q.pop();
        alive[v] = false;                 // nobody points to v: v is not on a cycle
        if (f[v] != -1 && --indeg[f[v]] == 0) q.push(f[v]);
    }
    return alive;
}

// (5) Floyd's tortoise and hare: mu and lam with O(1) extra memory
void floyd(const vector<int>& f, int s, int& mu, int& lam) {
    int t = f[s], h = f[f[s]];            // tortoise 1 step, hare 2 steps
    while (t != h) { t = f[t]; h = f[f[h]]; }
    mu = 0; t = s;                        // restart the tortoise; both now move 1 step
    while (t != h) { t = f[t]; h = f[h]; mu++; }
    lam = 1; h = f[t];                    // walk once around the cycle
    while (t != h) { h = f[h]; lam++; }
}

int main() {
    vector<int> f = {1, 2, 3, 4, 2, 4, 5, 8, 9, 8};
    int mu, lam;
    rho(f, 0, mu, lam);   cout << mu << " " << lam << "\n";            // 2 3
    rho(f, 6, mu, lam);   cout << mu << " " << lam << "\n";            // 2 3
    floyd(f, 7, mu, lam); cout << mu << " " << lam << "\n";            // 1 2
    cout << kth(f, 0, 10) << " " << kth(f, 0, 1000000000000000000LL) << "\n";  // 4 4
    vector<int> cyc = cycleLengths(f);
    for (int x : cyc) cout << x << " ";                                // 0 0 3 3 3 0 0 0 2 2
    cout << "\n";
    vector<bool> alive = onCycle(f);
    for (int v = 0; v < 10; v++) if (alive[v]) cout << v << " ";       // 2 3 4 8 9
    cout << "\n";
}
\end{lstlisting}

\begin{insightbox}[Which template when?]
\renewcommand{\arraystretch}{1.25}
\begin{tabularx}{\linewidth}{@{}>{\bfseries}l X@{}}
(1) \texttt{rho} & one start: where does it start repeating, how long is the loop? \\
(2) \texttt{kth} & one start: position after a huge number of steps \\
(3) \texttt{cycleLengths} & every cycle exactly once, or a value per vertex; works with \texttt{-1} \\
(4) \texttt{onCycle} & which vertices are on cycles, plus DP along the trees towards the cycles \\
(5) \texttt{floyd} & only \bigO{1} memory allowed (linked lists, generators you cannot store, read-only arrays as in \hyperref[ssec:ch1-find-duplicate]{Problem~1.5.3}) \\
\end{tabularx}
\end{insightbox}

%------------------------------------------------------------
\subsubsection*{Common mistakes}

\begin{warningbox}[Mistake 1 --- a fresh \texttt{visited} array for every start]
\begin{lstlisting}[style=snippet]
for (int s = 0; s < n; s++) {
    vector<bool> seen(n, false);          // O(n) to create, O(n) to walk: O(n^2) overall
    ...
}
\end{lstlisting}
Fine for $n \le 1000$ (Badge), \tle{} for $n = 10^5$. Keep \emph{one} state array for all walks (Algorithm 2).
\end{warningbox}

\begin{warningbox}[Mistake 2 --- ``I met a visited vertex, so I found a new cycle'']
With a two-state \texttt{visited} array, the walk from $5$ in $G$ stops at $4$ (visited by the walk from $0$) and would wrongly report a ``cycle'' through $5$. Only a vertex visited \emph{during the current walk} closes a new cycle: that is why we need three states (or visit times compared with the start time of the walk, Problem~1.5.2).
\end{warningbox}

\begin{warningbox}[{Mistake 3 --- forgetting \texttt{f[v] = -1}}]
In out-degree-$\le 1$ problems, \texttt{state[f[v]]} with \texttt{f[v] = -1} reads outside the array. Test \texttt{v != -1} \emph{before} every access.
\end{warningbox}

\begin{warningbox}[Mistake 4 --- simulating $K$ steps, or overflowing $K$]
$K = 10^{18}$ steps cannot be simulated. Reduce $K$ with $\mu$ and $\lambda$ first, and keep $K$ in \texttt{long long} --- including the expression $(K - \mu) \bmod \lambda$.
\end{warningbox}

\begin{warningbox}[Mistake 5 --- recursion along the arrows]
A recursive \texttt{solve(v)} that calls \texttt{solve(f[v])} recurses as deep as the longest tail plus cycle --- up to $n = 2\cdot10^5$ levels on one long chain, a classic \stackoverflow. Use a loop with an explicit \texttt{path} vector, as all templates above do.
\end{warningbox}

\begin{checklistbox}[``I'm stuck'' checklist for Pattern 5]
\begin{itemize}[leftmargin=*]
    \item[$\square$] Does every vertex have exactly one outgoing arrow (or at most one)? Then draw a few $\rho$'s.
    \item[$\square$] Is the question about \emph{one} start (Algorithm 1) or about \emph{all} vertices / all cycles (Algorithm 2 or 3)?
    \item[$\square$] Am I visiting each vertex once overall, not once per start?
    \item[$\square$] When the walk stops: at a vertex of \emph{this} walk (new cycle), of an older walk (old cycle) or at $-1$ (no cycle)?
    \item[$\square$] Is the answer a property of the cycle (length, minimum) or of the tail (entry point, distance)? Fill tails from the cycle outwards.
    \item[$\square$] Huge $K$? Reduce it modulo the cycle length after the tail. Many queries? Chapter~11.
\end{itemize}
\end{checklistbox}

%============================================================
% CH1 · P5 · PROBLEM 1 — BADGE (Codeforces 1020B)
%============================================================
\subsection{Problem 1.5.1: Badge}
\problemlabel{1.5.1}{ssec:ch1-badge}

\problemcard{1.5.1}{Badge}{Codeforces 1020B}{https://codeforces.com/problemset/problem/1020/B}{\exeasy\ \cf{1000}}{Follow the arrows until a vertex repeats: the answer is the entry point of the $\rho$; \bigO{n} by filling tails from the cycles}

\subsubsection*{1. Problem statement}

A teacher caught $n$ students (numbered $1..n$). She goes to some student $a$ and makes a hole in his badge. That student blames student $p_a$; the teacher goes there and makes a hole in $p_a$'s badge; he blames $p_{p_a}$, and so on. Eventually she comes to a student who \emph{already} has a hole; she makes a second hole in his badge and stops. For \textbf{every} possible first student $a = 1..n$, print who gets the second hole.

\textbf{Constraints:} $1 \le n \le 1000$, \; $1 \le p_i \le n$ (a student may blame himself).

\begin{table}[H]
\centering\small
\begin{tabular}{@{}lll@{}}
\toprule
\textbf{Input} & \textbf{Output} & \textbf{Explanation} \\
\midrule
\texttt{3} / \texttt{2 3 2} & \texttt{2 2 3} & from 1: $1,2,3,\mathbf{2}$; from 2: $2,3,\mathbf{2}$; from 3: $3,2,\mathbf{3}$ \\
\texttt{3} / \texttt{1 2 3} & \texttt{1 2 3} & everybody blames himself: the first student gets both holes \\
\bottomrule
\end{tabular}
\end{table}

\begin{diagrambox}
\begin{tikzpicture}
    \begin{scope}
        \node[gv] (s1) at (0,0) {1};
        \node[gv] (s2) at (1.6,0) {2};
        \node[gv] (s3) at (3.2,0) {3};
        \draw[gde] (s1) -- (s2);
        \draw[gde] (s2) to[bend left=35] (s3);
        \draw[gde] (s3) to[bend left=35] (s2);
        \node[font=\small\bfseries] at (1.6,1.2) {Sample 1: $p = [2,3,2]$};
    \end{scope}
    \begin{scope}[xshift=6.6cm]
        \foreach \v/\x in {1/0, 2/1.6, 3/3.2} {
            \node[gv] (t\v) at (\x,0) {\v};
            \draw[gde] (t\v) to[out=60, in=120, looseness=6] (t\v);
        }
        \node[font=\small\bfseries] at (1.6,1.2) {Sample 2: $p = [1,2,3]$};
    \end{scope}
\end{tikzpicture}
\figcap{Each student is a vertex with one arrow ``blames''. Sample 1 is one $\rho$: tail $1$, cycle $2 \leftrightarrows 3$. Sample 2 has three self-loops (cycles of length 1).}
\end{diagrambox}

\subsubsection*{2. The graph model}

\begin{modelbox}[Graph Model --- Badge]
\renewcommand{\arraystretch}{1.35}
\begin{tabularx}{\linewidth}{@{}>{\bfseries}l X@{}}
Vertices & the $n$ students \\
Edges & $i \to p_i$ (``$i$ blames $p_i$''); exactly one outgoing edge per vertex: a \textbf{functional graph}, stored as the array \texttt{p} itself \\
Weights & none \\
Question & for every start $a$: the first vertex visited twice by the walk $a, p_a, p_{p_a}, \dots$ = the \textbf{entry point} of the $\rho$ that starts at $a$ \\
Size & $|V| = |E| = n \le 1000$; one walk has at most $n+1$ steps, so even $n$ walks ($10^6$ steps) are cheap \\
\end{tabularx}
\end{modelbox}

\begin{diagrambox}
\begin{tikzpicture}[x=1.45cm]
    \foreach \v/\i/\s in {1/0/gvsrc, 2/1/gvdst, 3/2/gvseen}{
        \node[\s] (w\i) at (\i,0) {\v};
        \node[gdist, below=1mm of w\i] {hole 1};
    }
    \node[gvdst, dashed] (w3) at (3,0) {2};
    \node[gdist, below=1mm of w3, text=warningred] {hole 2!};
    \foreach \i [evaluate=\i as \j using int(\i+1)] in {0,...,2} \draw[gde] (w\i) -- (w\j);
    \draw[gback] (w3.north) to[bend right=35] node[lbl, above]{already has a hole} (w1.north);
    \node[note, text width=5.6cm, align=left, anchor=west] at (3.6,0) {Sample 1, $a = 1$. The teacher's route is the $\rho$ of vertex 1; the student with two holes is the \textbf{first vertex that repeats} = the entry point, here $2$.};
\end{tikzpicture}
\end{diagrambox}

\subsubsection*{3. Key observations}

\begin{itemize}[leftmargin=*]
    \item \textbf{The walk is forced}, so for a fixed $a$ the answer is simply ``simulate until a repeat''. The process always stops within $n+1$ visits (pigeonhole).
    \item \textbf{The answer is the entry point.} Tail vertices are visited once, then the walk enters the cycle at $e$ and goes around; the first vertex it sees again is $e$ (shape theorem, \S\ref{ssec:ch1-p5-theory}).
    \item \textbf{A cycle vertex is its own answer.} If $a$ lies on a cycle, the walk goes once around and returns to $a$ first: $\text{ans}[a] = a$.
    \item \textbf{A tail vertex inherits the answer of its successor.} If $a$ is not on a cycle, the walks from $a$ and from $p_a$ enter the cycle at the same vertex: $\text{ans}[a] = \text{ans}[p_a]$. This is what makes an \bigO{n} solution possible.
\end{itemize}

\begin{warningbox}[Not ``the vertex with the largest in-degree'']
In Sample 1 vertex $2$ has in-degree 2 and is the answer for $a = 1, 2$ --- but for $a = 3$ the answer is $3$ (in-degree 1). The answer depends on where you \emph{enter} the cycle, not on popularity.
\end{warningbox}

\subsubsection*{4. Approaches}

\begin{bruteforcebox}[Approach 1: Direct Simulation (already accepted here)]
For each $a$, use a fresh array \texttt{hole[]} and walk until the current student already has a hole.

\begin{lstlisting}[style=cpp]
#include <bits/stdc++.h>
using namespace std;

int main() {
    int n;
    cin >> n;
    vector<int> p(n + 1);
    for (int i = 1; i <= n; i++) cin >> p[i];
    for (int a = 1; a <= n; a++) {
        vector<bool> hole(n + 1, false);  // fresh badges for every experiment
        int v = a;
        while (!hole[v]) {                // first visit: one hole, move on
            hole[v] = true;
            v = p[v];
        }
        cout << v << (a == n ? '\n' : ' ');  // v is the first student visited twice
    }
}
\end{lstlisting}

\begin{complexitybox}
\textbf{Time:} \bigO{n^2} --- at most $n+1$ steps per start, about $10^6$ for $n = 1000$. \textbf{Space:} \bigO{n}. \textbf{Verdict:} \ac{} for this problem (it is the intended solution), but \tle{} if $n$ were $2\cdot10^5$ (about $2\cdot10^{10}$ steps on one long chain).
\end{complexitybox}
\end{bruteforcebox}

\begin{algorithmbox}[Approach 2: Linear Time --- Find the Cycles, Then Fill the Tails]
Use Algorithm 2 of the theory (three states). When a walk closes a \emph{new} cycle, every cycle vertex is its own answer. Then the tail vertices of this walk are answered \textbf{backwards}, from the cycle outwards, with $\text{ans}[v] = \text{ans}[p_v]$. When a walk instead runs into a finished vertex, that vertex already has its answer, and again the path is filled backwards.

\begin{lstlisting}[style=cpp]
#include <bits/stdc++.h>
using namespace std;

int main() {
    int n;
    cin >> n;
    vector<int> p(n + 1);
    for (int i = 1; i <= n; i++) cin >> p[i];

    vector<int> state(n + 1, 0);          // 0 = new, 1 = on the current walk, 2 = answered
    vector<int> ans(n + 1, 0);
    for (int s = 1; s <= n; s++) {
        if (state[s] != 0) continue;
        vector<int> path;
        int v = s;
        while (state[v] == 0) {           // walk until we meet something already seen
            state[v] = 1;
            path.push_back(v);
            v = p[v];
        }
        int k = path.size();              // path[0..k-1] still need an answer
        if (state[v] == 1) {              // NEW cycle at v: cycle vertices answer themselves
            while (path[k - 1] != v) { ans[path[k - 1]] = path[k - 1]; k--; }
            ans[v] = v; k--;
        }
        for (int i = k - 1; i >= 0; i--)  // tail vertices, from the cycle outwards
            ans[path[i]] = ans[p[path[i]]];
        for (int u : path) state[u] = 2;
    }
    for (int a = 1; a <= n; a++) cout << ans[a] << (a == n ? '\n' : ' ');
}
\end{lstlisting}

\begin{complexitybox}
\textbf{Time:} \bigO{n} --- every vertex enters a \texttt{path} once and is answered once. \textbf{Space:} \bigO{n}. \textbf{Verdict:} \ac{}, and it would still be fast for $n = 10^6$.
\end{complexitybox}
\end{algorithmbox}

\begin{invariantbox}[Why it is correct]
\textbf{Claim:} when a vertex is answered, its answer is the entry point of its $\rho$.

\emph{Cycle vertices.} The walk stopped at a state-1 vertex $v$, so $v, p_v, \dots$ up to the end of \texttt{path} is a cycle (theory, Algorithm 2). Starting on a cycle, the first vertex seen twice is the start itself: $\text{ans}[u] = u$ is right.

\emph{Tail vertices.} For a vertex $u$ not on a cycle, the walk from $u$ is $u$ followed by the walk from $p_u$; $u$ itself is never visited again (it is not on a cycle), so the first repeated vertex of both walks is the same: $\text{ans}[u] = \text{ans}[p_u]$. Filling \texttt{path} \emph{backwards} guarantees that $p_u$ (the next element of \texttt{path}, or the vertex $v$ where the walk stopped) is already answered when $u$ is processed. \hfill$\blacksquare$
\end{invariantbox}

\begin{alternativebox}[Approach 3: Peeling (Kahn) --- Cycles First, Then Reverse Order]
Delete in-degree-0 vertices repeatedly (theory, Algorithm 3). The survivors are the cycle vertices, $\text{ans}[v] = v$. A deleted vertex always comes out \emph{before} its successor, so processing the deletion order \textbf{backwards} sees $p_v$ before $v$, and $\text{ans}[v] = \text{ans}[p_v]$ can be filled directly.

\begin{lstlisting}[style=cpp]
#include <bits/stdc++.h>
using namespace std;

int main() {
    int n;
    cin >> n;
    vector<int> p(n + 1), indeg(n + 1, 0);
    for (int i = 1; i <= n; i++) { cin >> p[i]; indeg[p[i]]++; }

    vector<int> order;                    // vertices in the order they are peeled
    queue<int> q;
    for (int v = 1; v <= n; v++) if (indeg[v] == 0) q.push(v);
    while (!q.empty()) {
        int v = q.front(); q.pop();
        order.push_back(v);
        if (--indeg[p[v]] == 0) q.push(p[v]);
    }
    vector<int> ans(n + 1, 0);
    for (int v = 1; v <= n; v++)
        if (indeg[v] > 0) ans[v] = v;     // survivors lie on cycles
    for (int i = (int)order.size() - 1; i >= 0; i--)  // p[v] is peeled after v,
        ans[order[i]] = ans[p[order[i]]];             // so it is answered before v
    for (int a = 1; a <= n; a++) cout << ans[a] << (a == n ? '\n' : ' ');
}
\end{lstlisting}

\begin{diagrambox}
\begin{tikzpicture}
    \foreach \v/\x/\y/\s in {1/-1.4/0.8/gvdead, 2/0/0.8/gvsel, 3/1.3/0/gvsel, 4/0/-0.8/gvsel, 5/-1.4/-0.8/gvdead, 9/-2.8/-0.8/gvdead,
                             6/3.2/0.8/gvsel, 7/4.5/0/gvsel, 8/3.2/-0.8/gvsel}
        \node[\s] (k\v) at (\x,\y) {\v};
    \foreach \u/\w in {2/3, 3/4, 4/2, 6/7, 7/8, 8/6} \draw[gde] (k\u) -- (k\w);
    \foreach \u/\w in {1/2, 5/4, 9/5} \draw[gfaded, dashed, -{Stealth[length=2mm]}] (k\u) -- (k\w);
    \node[gdist, above=1mm of k1] {peeled 1st};
    \node[gdist, above=1mm of k9] {peeled 2nd};
    \node[gdist, below=1mm of k5] {peeled 3rd};
    \node[note, text width=4.4cm, align=left, anchor=west] at (5.2,0) {Peeling order $1, 9, 5$. Backwards: $\text{ans}[5] = \text{ans}[4] = 4$, then $\text{ans}[9] = \text{ans}[5] = 4$, then $\text{ans}[1] = \text{ans}[2] = 2$.};
\end{tikzpicture}
\figcap{Approach 3 on the trace input $p = [2,3,4,2,4,7,8,6,5]$ of Section 5. Purple: survivors = cycle vertices, each its own answer.}
\end{diagrambox}

\begin{complexitybox}
\textbf{Time:} \bigO{n}. \textbf{Space:} \bigO{n}. \textbf{Verdict:} \ac{}. Same answers, different order of work: Approach 2 discovers cycles while walking, Approach 3 first removes everything that is not on a cycle.
\end{complexitybox}
\end{alternativebox}

\begin{tricksbox}[Approach 4: Implementation Notes]
\begin{itemize}[leftmargin=*]
    \item \textbf{No recursion.} The tempting \texttt{solve(v)} that returns \texttt{ans[v] = solve(p[v])} recurses along a tail of length up to $n$; with $n = 1000$ that is safe, with $n = 10^6$ it overflows the stack. The explicit \texttt{path} vector is the iterative version of that recursion.
    \item \textbf{1-indexing.} Students are $1..n$; arrays of size $n+1$ avoid every $\pm 1$.
    \item \textbf{Output.} The $n$ answers go on one line separated by spaces; the code prints a newline after the last one instead of a space.
\end{itemize}
\end{tricksbox}

\subsubsection*{5. Trace}

We trace Approach 2 on a bigger input with two cycles, a tail of length 2 and two tails entering the same cycle at different vertices:\quad $n = 9$, \;$p = [2,3,4,2,4,7,8,6,5]$. The program prints \mbox{\texttt{2 2 3 4 4 6 7 8 4}}.

\begin{diagrambox}
\begin{tikzpicture}
    \foreach \v/\x/\y in {1/-1.4/0.8, 2/0/0.8, 3/1.3/0, 4/0/-0.8, 5/-1.4/-0.8, 9/-2.8/-0.8, 6/3.2/0.8, 7/4.5/0, 8/3.2/-0.8}
        \node[gv] (g\v) at (\x,\y) {\v};
    \foreach \u/\w in {1/2, 2/3, 3/4, 4/2, 5/4, 9/5, 6/7, 7/8, 8/6} \draw[gde] (g\u) -- (g\w);
\end{tikzpicture}
\figcap{The trace input: cycle $2 \to 3 \to 4 \to 2$ with tails $1 \to 2$ and $9 \to 5 \to 4$; cycle $6 \to 7 \to 8 \to 6$.}
\end{diagrambox}

\begin{tracebox}[{Approach 2 on $p = [2,3,4,2,4,7,8,6,5]$}]
\centering
\begin{tikzpicture}
    \begin{scope}
        \foreach \v/\x/\y/\s in {1/-1.4/0.8/gvseen, 2/0/0.8/gvseen, 3/1.3/0/gvseen, 4/0/-0.8/gvseen, 5/-1.4/-0.8/gv, 9/-2.8/-0.8/gv,
                                 6/3.2/0.8/gv, 7/4.5/0/gv, 8/3.2/-0.8/gv}
            \node[\s] (a\v) at (\x,\y) {\v};
        \foreach \u/\w in {1/2, 2/3, 3/4, 5/4, 9/5, 6/7, 7/8, 8/6} \draw[gde] (a\u) -- (a\w);
        \draw[gback] (a4) -- (a2);
        \node[note, text width=5.2cm, align=left, anchor=west] at (5.2,0) {\textbf{$s=1$:} \texttt{path} $= 1,2,3,4$, stops at $2$ (state 1). New cycle $2,3,4$: they are their own answers; then the tail: $\text{ans}[1] = \text{ans}[2] = 2$.};
    \end{scope}
    \begin{scope}[yshift=-2.75cm]
        \foreach \v/\x/\y/\s in {1/-1.4/0.8/gvdone, 2/0/0.8/gvdone, 3/1.3/0/gvdone, 4/0/-0.8/gvdone, 5/-1.4/-0.8/gvseen, 9/-2.8/-0.8/gv,
                                 6/3.2/0.8/gv, 7/4.5/0/gv, 8/3.2/-0.8/gv}
            \node[\s] (b\v) at (\x,\y) {\v};
        \foreach \u/\w in {1/2, 2/3, 3/4, 4/2, 9/5, 6/7, 7/8, 8/6} \draw[gde] (b\u) -- (b\w);
        \draw[gcross] (b5) -- (b4);
        \node[note, text width=5.2cm, align=left, anchor=west] at (5.2,0) {\textbf{$s=5$} ($2,3,4$ skipped: finished): \texttt{path} $= 5$, stops at $4$ (state 2, answered). No new cycle: $\text{ans}[5] = \text{ans}[4] = 4$.};
    \end{scope}
    \begin{scope}[yshift=-5.5cm]
        \foreach \v/\x/\y/\s in {1/-1.4/0.8/gvdone, 2/0/0.8/gvdone, 3/1.3/0/gvdone, 4/0/-0.8/gvdone, 5/-1.4/-0.8/gvdone, 9/-2.8/-0.8/gv,
                                 6/3.2/0.8/gvseen, 7/4.5/0/gvseen, 8/3.2/-0.8/gvseen}
            \node[\s] (c\v) at (\x,\y) {\v};
        \foreach \u/\w in {1/2, 2/3, 3/4, 4/2, 5/4, 9/5, 6/7, 7/8} \draw[gde] (c\u) -- (c\w);
        \draw[gback] (c8) -- (c6);
        \node[note, text width=5.2cm, align=left, anchor=west] at (5.2,0) {\textbf{$s=6$:} \texttt{path} $= 6,7,8$, stops at $6$ (state 1). A new cycle with no tail: $\text{ans}[8]=8$, $\text{ans}[7]=7$, $\text{ans}[6]=6$.};
    \end{scope}
    \begin{scope}[yshift=-8.25cm]
        \foreach \v/\x/\y/\s in {1/-1.4/0.8/gvdone, 2/0/0.8/gvdone, 3/1.3/0/gvdone, 4/0/-0.8/gvdone, 5/-1.4/-0.8/gvdone, 9/-2.8/-0.8/gvseen,
                                 6/3.2/0.8/gvdone, 7/4.5/0/gvdone, 8/3.2/-0.8/gvdone}
            \node[\s] (d\v) at (\x,\y) {\v};
        \foreach \u/\w in {1/2, 2/3, 3/4, 4/2, 5/4, 6/7, 7/8, 8/6} \draw[gde] (d\u) -- (d\w);
        \draw[gcross] (d9) -- (d5);
        \node[note, text width=5.2cm, align=left, anchor=west] at (5.2,0) {\textbf{$s=9$} ($7, 8$ skipped): \texttt{path} $= 9$, stops at $5$ (state 2): $\text{ans}[9] = \text{ans}[5] = 4$. Done.};
    \end{scope}
\end{tikzpicture}

\medskip
\small
\begin{tabular}{@{}c|ccccccccc@{}}
\toprule
$a$ & 1 & 2 & 3 & 4 & 5 & 6 & 7 & 8 & 9 \\
$p_a$ & 2 & 3 & 4 & 2 & 4 & 7 & 8 & 6 & 5 \\
\midrule
answered in walk & $s{=}1$ & $s{=}1$ & $s{=}1$ & $s{=}1$ & $s{=}5$ & $s{=}6$ & $s{=}6$ & $s{=}6$ & $s{=}9$ \\
\textbf{ans}$[a]$ & \textbf{2} & \textbf{2} & \textbf{3} & \textbf{4} & \textbf{4} & \textbf{6} & \textbf{7} & \textbf{8} & \textbf{4} \\
\bottomrule
\end{tabular}
\end{tracebox}

\begin{diagrambox}
\begin{tikzpicture}
    \foreach \v/\x/\y/\s/\a in {1/-1.4/0.8/gv/2, 2/0/0.8/gvdst/2, 3/1.3/0/gvdst/3, 4/0/-0.8/gvdst/4, 5/-1.4/-0.8/gv/4, 9/-2.8/-0.8/gv/4,
                                6/3.2/0.8/gvdst/6, 7/4.5/0/gvdst/7, 8/3.2/-0.8/gvdst/8}{
        \node[\s] (z\v) at (\x,\y) {\v};
        \node[gdist, text=accentorange, above right=-1.5mm and -1mm of z\v] {\a};
    }
    \foreach \u/\w in {1/2, 2/3, 3/4, 4/2, 5/4, 9/5, 6/7, 7/8, 8/6} \draw[gde] (z\u) -- (z\w);
    \node[note, text width=4.6cm, align=left, anchor=west] at (5.4,0) {Final answers (small orange numbers). Cycle vertices (orange) answer themselves; the tail $9 \to 5$ inherits $4$, the tail $1$ inherits $2$: the two tails enter the same cycle at different points.};
\end{tikzpicture}
\end{diagrambox}

Check $a = 9$ by hand: the teacher visits $9, 5, 4, 2, 3$ and then $4$ again --- student $4$ gets two holes, although the walk from $1$ enters the same cycle at $2$.

\subsubsection*{6. Test suite}

\begin{testbox}
\centering\small
\begin{tabular}{@{}lll@{}}
\toprule
\textbf{Input ($n$ / $p$)} & \textbf{Expected} & \textbf{Why this test} \\
\midrule
\texttt{3 / 2 3 2} & \texttt{2 2 3} & official sample 1: tail + 2-cycle \\
\texttt{3 / 1 2 3} & \texttt{1 2 3} & official sample 2: only self-loops \\
\texttt{1 / 1} & \texttt{1} & smallest input \\
\texttt{9 / 2 3 4 2 4 7 8 6 5} & \texttt{2 2 3 4 4 6 7 8 4} & two pieces, tail of length 2 (the trace) \\
\texttt{4 / 2 3 4 4} & \texttt{4 4 4 4} & a chain ending in a self-loop \\
$n = 1000$, $p_i = i+1$, $p_n = n$ & $n$ times \texttt{1000} & longest tail: brute force does $\approx 5\cdot10^5$ steps \\
\bottomrule
\end{tabular}
\end{testbox}

All three programs were also compared on 500 random inputs with $n \le 9$: identical outputs.

\subsubsection*{7. Common mistakes}

\begin{warningbox}[Mistake 1 --- forgetting to reset the holes between experiments]
In Approach 1 the array \texttt{hole[]} must be fresh for every $a$ (each $a$ is a separate story). Reusing it makes the second experiment stop immediately at a student who got his hole in the first one.
\end{warningbox}

\begin{warningbox}[Mistake 2 --- answering with the last new student instead of the repeated one]
\begin{lstlisting}[style=snippet]
while (!hole[p[v]]) { hole[v] = true; v = p[v]; }   // WRONG: stops one student too early
while (!hole[v]) { hole[v] = true; v = p[v]; }      // RIGHT: stop ON the repeated student
\end{lstlisting}
The answer is the student who receives the \emph{second} hole --- the vertex where the walk stops, not the one before it.
\end{warningbox}

\begin{warningbox}[Mistake 3 --- giving every vertex of the path the entry point]
In Approach 2, setting $\text{ans}[u] = v$ for the whole \texttt{path} is wrong for cycle vertices: on the cycle each vertex is its own answer ($\text{ans}[3] = 3$, not $2$). Only tail vertices inherit.
\end{warningbox}

\subsubsection*{8. Takeaways}

\begin{itemize}[leftmargin=*]
    \item ``Follow until something repeats'' $\Rightarrow$ functional graph; the first repeated vertex is the \textbf{entry point} of the $\rho$.
    \item Answers per vertex: cycle vertices are special (here: their own answer), tail vertices \textbf{inherit from their successor}.
    \item Fill tails backwards along the walk (or in reverse peeling order) so that the successor is always answered first.
    \item Small $n$: plain simulation is a perfectly good solution --- but know the \bigO{n} version for when $n$ grows.
\end{itemize}

\begin{variantsbox}
\variant{Longest cycle}{\varin{https://leetcode.com/problems/longest-cycle-in-a-graph/}}
  {instead of an entry point per start, report the length of the longest cycle; some vertices have no arrow.}
  {the same functional graph with \texttt{-1} for missing arrows; only cycles matter, tails are ignored.}
  {three-state walks (or visit times); a walk that stops inside itself gives a cycle length; keep the maximum.}
\variant{Where after $K$ visits?}{\varin{https://atcoder.jp/contests/abc167/tasks/abc167_d}}
  {one start only, but which vertex is reached after exactly $K \le 10^{18}$ steps?}
  {the $\rho$ of the start: tail $\mu$, cycle $\lambda$, entry point as here.}
  {walk until the first repeat, then $K \to \mu + (K - \mu) \bmod \lambda$ when $K \ge \mu$.}
\variant{Number of holed students}{\varnone}
  {for every $a$, how many \emph{different} students get at least one hole?}
  {per vertex: distance $d[a]$ to its cycle and length $\lambda[a]$ of that cycle; the answer is $d[a] + \lambda[a] = \mu + \lambda$.}
  {cycle vertices: $d = 0$; tails filled from the cycle outwards: $d[a] = d[p_a] + 1$, $\lambda[a] = \lambda[p_a]$ --- the skeleton of Approach 2 (more in Chapter~11).}
\variant{Trap the student cheaply}{\varin{https://codeforces.com/problemset/problem/1027/D}}
  {each vertex has a price; choose vertices of minimum total price so that \emph{every} walk eventually hits a chosen one.}
  {each walk ends in its cycle, so every cycle needs one chosen vertex; tails need none.}
  {for every new cycle found by the three-state walk, add its cheapest vertex.}
\end{variantsbox}

%============================================================
% CH1 · P5 · PROBLEM 2 — LONGEST CYCLE IN A GRAPH (LeetCode 2360)
%============================================================
\subsection{Problem 1.5.2: Longest Cycle in a Graph}
\problemlabel{1.5.2}{ssec:ch1-longest-cycle}

\problemcard{1.5.2}{Longest Cycle in a Graph}{LeetCode 2360}{https://leetcode.com/problems/longest-cycle-in-a-graph/}{\exhard\ \lc{Hard}}{Out-degree $\le 1$: walk from every unvisited vertex with global visit times; a walk that returns into itself measures a cycle}

\subsubsection*{1. Problem statement}

A directed graph has $n$ vertices $0..n-1$, and every vertex has \textbf{at most one} outgoing edge: \texttt{edges[i]} is the target of the edge leaving $i$, or $-1$ if there is none. Return the length of the \textbf{longest cycle}, or $-1$ if the graph has no cycle.

\textbf{Constraints:} $2 \le n \le 10^5$, \; $-1 \le \texttt{edges}[i] < n$, \; $\texttt{edges}[i] \ne i$ (no self-loops).

\begin{table}[H]
\centering\small
\begin{tabular}{@{}lcl@{}}
\toprule
\textbf{edges} & \textbf{Output} & \textbf{Explanation} \\
\midrule
$[3,3,4,2,3]$ & 3 & the only cycle is $2 \to 4 \to 3 \to 2$ \\
$[2,-1,3,1]$ & $-1$ & $0 \to 2 \to 3 \to 1 \to$ (nothing): no cycle \\
\bottomrule
\end{tabular}
\end{table}

\begin{diagrambox}
\begin{tikzpicture}
    \begin{scope}
        \node[font=\small\bfseries] at (1.2,1.7) {Example 1: $[3,3,4,2,3]$};
        \foreach \v/\x/\y/\s in {0/-0.6/1/gv, 1/-0.6/-0.6/gv, 3/0.8/0.2/gvseen, 2/2.2/1/gvseen, 4/2.2/-0.6/gvseen}
            \node[\s] (a\v) at (\x,\y) {\v};
        \draw[gde] (a0) -- (a3); \draw[gde] (a1) -- (a3);
        \draw[gpath, -{Stealth[length=3mm]}] (a3) -- (a2);
        \draw[gpath, -{Stealth[length=3mm]}] (a2) -- (a4);
        \draw[gpath, -{Stealth[length=3mm]}] (a4) -- (a3);
    \end{scope}
    \begin{scope}[xshift=6.8cm]
        \node[font=\small\bfseries] at (2,1.7) {Example 2: $[2,-1,3,1]$};
        \foreach \v/\x in {0/0, 2/1.3, 3/2.6, 1/3.9}
            \node[gv] (b\v) at (\x,0.2) {\v};
        \draw[gde] (b0) -- (b2); \draw[gde] (b2) -- (b3); \draw[gde] (b3) -- (b1);
        \draw[-{Stealth[length=2mm]}, warningred, dashed] (b1) -- (4.8,0.2);
        \node[font=\scriptsize\bfseries, text=warningred, anchor=west] at (4.8,0.2) {$-1$};
    \end{scope}
\end{tikzpicture}
\figcap{Left: a $\rho$-shape; the cycle (orange) has length 3. Right: a single chain that ends in a vertex without an outgoing edge --- no cycle at all.}
\end{diagrambox}

\subsubsection*{2. The graph model}

\begin{modelbox}[Graph Model --- Longest Cycle in a Graph]
\renewcommand{\arraystretch}{1.35}
\begin{tabularx}{\linewidth}{@{}>{\bfseries}l X@{}}
Vertices & $0..n-1$, given directly \\
Edges & $i \to \texttt{edges}[i]$ when $\texttt{edges}[i] \ne -1$; out-degree $\le 1$, stored in the input array (no adjacency list) \\
Weights & none --- a cycle's length is its number of edges = number of vertices \\
Question & over all pieces of this ``functional graph with holes'': the maximum cycle length, $-1$ if no piece has a cycle \\
Size & $|V| \le 10^5$, $|E| \le 10^5$: we need \bigO{n} or \bigO{n \log n}; \bigO{n^2} = $10^{10}$ is far too slow \\
\end{tabularx}
\end{modelbox}

\begin{diagrambox}
\begin{tikzpicture}
    \begin{scope}
        \node[gv] (p0) at (0,0.6) {};
        \node[gv] (p1) at (1.1,0.6) {};
        \node[gvseen] (p2) at (2.2,0.6) {};
        \node[gvseen] (p3) at (3.1,1.2) {};
        \node[gvseen] (p4) at (3.1,0) {};
        \draw[gde] (p0) -- (p1); \draw[gde] (p1) -- (p2);
        \draw[gde] (p2) -- (p3); \draw[gde] (p3) -- (p4); \draw[gde] (p4) -- (p2);
        \node[note] at (1.6,-0.7) {piece with a cycle ($\rho$): \textbf{counts}};
    \end{scope}
    \begin{scope}[xshift=5.6cm]
        \node[gv] (q0) at (0,0.6) {};
        \node[gv] (q1) at (1.1,1.2) {};
        \node[gv] (q2) at (1.1,0) {};
        \node[gvdead] (q3) at (2.2,0.6) {};
        \draw[gde] (q1) -- (q0); \draw[gde] (q2) -- (q0);
        \draw[gde] (q3) -- (q1);
        \draw[-{Stealth[length=2mm]}, warningred, dashed] (q0) -- (-1,0.6);
        \node[note] at (0.6,-0.7) {tree ending in a sink: \textbf{no cycle}};
    \end{scope}
    \node[note, text width=4.2cm, align=left, anchor=west] at (9.3,0.6) {Every piece is one of these two kinds (theory, ``three relatives''). We need the longest cycle over all $\rho$-pieces.};
\end{tikzpicture}
\end{diagrambox}

\subsubsection*{3. Key observations}

\begin{itemize}[leftmargin=*]
    \item \textbf{At most one cycle per piece, and cycles are disjoint.} A vertex has only one outgoing edge, so it can lie on at most one cycle. Hence the total length of all cycles is $\le n$ --- measuring \emph{every} cycle is affordable.
    \item \textbf{A walk finds a cycle iff it returns into itself.} Walking from $s$, we either fall off at $-1$ (no cycle on this walk), reach a vertex seen on an \emph{earlier} walk (its cycle, if any, was measured then), or reach a vertex of the \emph{current} walk (a new cycle).
    \item \textbf{Visit times measure lengths for free.} If vertex $v$ got time $t_v$ when first visited and the walk comes back to $v$ when the clock shows $T$, the cycle has exactly $T - t_v$ vertices.
    \item \textbf{Why not DFS from every vertex?} Each walk can be $\Theta(n)$ long (one long chain into a cycle), so restarting from every vertex costs $\Theta(n^2)$.
\end{itemize}

\subsubsection*{4. Approaches}

\begin{bruteforcebox}
For every vertex $s$, walk at most $n$ steps and check whether we come back to $s$; if we do after \texttt{len} steps, $s$ is on a cycle of length \texttt{len}.

\begin{lstlisting}[style=cpp]
#include <bits/stdc++.h>
using namespace std;

class Solution {
public:
    int longestCycle(vector<int>& edges) {
        int n = edges.size(), best = -1;
        for (int s = 0; s < n; s++) {     // is s on a cycle? walk at most n steps
            int v = edges[s], len = 1;
            while (v != -1 && v != s && len <= n) {
                v = edges[v];
                len++;
            }
            if (v == s) best = max(best, len);   // came back to s after len steps
        }
        return best;
    }
};
\end{lstlisting}

\begin{complexitybox}
\begin{diagrambox}
\begin{tikzpicture}[x=1.15cm]
    \foreach \i/\l in {0/0, 1/1, 2/2, 3/{\dots}, 4/{n{-}3}}
        \node[gv, minimum size=8.5mm, font=\scriptsize\bfseries] (c\i) at (\i,0) {$\l$};
    \node[gvseen, minimum size=8.5mm, font=\scriptsize\bfseries] (c5) at (5.3,0) {$n{-}2$};
    \node[gvseen, minimum size=8.5mm, font=\scriptsize\bfseries] (c6) at (6.6,0) {$n{-}1$};
    \foreach \i [evaluate=\i as \j using int(\i+1)] in {0,...,4} \draw[gde] (c\i) -- (c\j);
    \draw[gde] (c5) to[bend left=35] (c6); \draw[gde] (c6) to[bend left=35] (c5);
    \draw[gback] (c0.south) to[bend right=12] node[lbl, below]{start $0$: about $n$ steps, never returns to $0$} (c6.south);
\end{tikzpicture}
\figcap{The worst case for the brute force: every chain vertex walks to the 2-cycle and around it until the step limit $n$, without ever coming back to itself.}
\end{diagrambox}

\textbf{Time:} \bigO{n^2}: on a chain $0 \to 1 \to \dots \to n{-}3$ ending in a 2-cycle, every start walks $\approx n$ steps before giving up --- $10^{10}$ steps for $n = 10^5$. \textbf{Space:} \bigO{1}. \textbf{Verdict:} \tle{} (used below as the reference on small random graphs).
\end{complexitybox}
\end{bruteforcebox}

\begin{algorithmbox}[Approach 2: One Clock for All Walks]
Keep one global clock and one array \texttt{when[v]} ($0$ = never visited). For every unvisited $s$, remember $\texttt{start} = $ the clock value, then walk, stamping each new vertex with the current time. Stop at $-1$ or at a visited vertex $v$:
\begin{itemize}[leftmargin=*]
    \item $\texttt{when}[v] \ge \texttt{start}$: $v$ was stamped during \emph{this} walk $\Rightarrow$ new cycle of length $\texttt{time} - \texttt{when}[v]$;
    \item $\texttt{when}[v] < \texttt{start}$: an older walk $\Rightarrow$ nothing new.
\end{itemize}
The comparison with \texttt{start} plays exactly the role of ``state 1 versus state 2'' in the theory's Algorithm 2 --- one integer array instead of colours plus a \texttt{path}.

\begin{lstlisting}[style=cpp]
#include <bits/stdc++.h>
using namespace std;

class Solution {
public:
    int longestCycle(vector<int>& edges) {
        int n = edges.size();
        vector<int> when(n, 0);           // 0 = never visited, otherwise the visit time
        int time = 1, best = -1;
        for (int s = 0; s < n; s++) {
            if (when[s] != 0) continue;
            int start = time;             // every vertex of THIS walk gets a time >= start
            int v = s;
            while (v != -1 && when[v] == 0) {
                when[v] = time++;
                v = edges[v];
            }
            if (v != -1 && when[v] >= start)      // closed a cycle inside this walk
                best = max(best, time - when[v]);
        }
        return best;
    }
};
\end{lstlisting}

\begin{complexitybox}
\textbf{Time:} \bigO{n} --- each vertex is stamped once; the \texttt{while} loop runs $n$ times in total over all walks. \textbf{Space:} \bigO{n} for \texttt{when}. \textbf{Verdict:} \ac{}.
\end{complexitybox}
\end{algorithmbox}

\begin{invariantbox}[Why it is correct]
\textbf{Invariant:} during the walk from $s$, the vertices with $\texttt{when} \ge \texttt{start}$ are exactly the vertices visited by this walk, and their times are consecutive: $\texttt{start}, \texttt{start}+1, \dots, \texttt{time}-1$ in walk order.

\emph{New cycle found.} If the walk reaches $v$ with $\texttt{when}[v] \ge \texttt{start}$, then the vertices stamped $\texttt{when}[v], \dots, \texttt{time}-1$ are the walk from $v$ back to (just before) $v$ --- a cycle with $\texttt{time} - \texttt{when}[v]$ vertices.

\emph{No cycle missed.} Every cycle vertex is stamped by some walk. Take the first walk that touches a cycle $C$: it enters $C$ at some vertex and, because no vertex of $C$ was stamped before, it goes all the way round and returns into itself --- so $C$ is measured. Later walks reaching $C$ see old stamps and correctly do nothing. Cycles are disjoint, so each is measured exactly once. \hfill$\blacksquare$
\end{invariantbox}

\begin{alternativebox}[Approach 3: Peel the Trees, Measure What Is Left]
Repeatedly delete vertices of in-degree 0 (Kahn, theory Algorithm 3). Every tree vertex --- including whole sink-trees with no cycle --- is deleted; the survivors are disjoint cycles. Walk once around each survivor cycle, switching its vertices off so that every cycle is measured once.

\begin{lstlisting}[style=cpp]
#include <bits/stdc++.h>
using namespace std;

class Solution {
public:
    int longestCycle(vector<int>& edges) {
        int n = edges.size();
        vector<int> indeg(n, 0);
        for (int v = 0; v < n; v++)
            if (edges[v] != -1) indeg[edges[v]]++;
        queue<int> q;
        for (int v = 0; v < n; v++)
            if (indeg[v] == 0) q.push(v);
        while (!q.empty()) {              // peel vertices nobody points to
            int v = q.front(); q.pop();
            if (edges[v] != -1 && --indeg[edges[v]] == 0) q.push(edges[v]);
        }
        int best = -1;                    // survivors (indeg > 0) form disjoint cycles
        for (int s = 0; s < n; s++) {
            if (indeg[s] == 0) continue;
            int len = 0, v = s;
            do {                          // walk once around, switching vertices off
                indeg[v] = 0;
                v = edges[v];
                len++;
            } while (v != s);
            best = max(best, len);
        }
        return best;
    }
};
\end{lstlisting}

\begin{diagrambox}
\begin{tikzpicture}
    \foreach \v/\x/\y/\s in {0/0/0.8/gvsel, 1/1.3/0/gvsel, 2/0/-0.8/gvsel, 3/3/0.8/gvsel, 4/4.4/0.8/gvsel, 5/4.4/-0.8/gvsel, 6/3/-0.8/gvsel,
                             7/3/2.1/gvdead, 8/6.2/0/gvdead, 9/7.5/0/gvdead}
        \node[\s] (k\v) at (\x,\y) {\v};
    \foreach \u/\w in {0/1, 1/2, 2/0, 3/4, 4/5, 5/6, 6/3} \draw[gde] (k\u) -- (k\w);
    \foreach \u/\w in {7/3, 8/9} \draw[gfaded, dashed, -{Stealth[length=2mm]}] (k\u) -- (k\w);
    \node[gdist, right=1mm of k7] {peeled 1st};
    \node[gdist, above=1mm of k8] {peeled 2nd};
    \node[gdist, below=1mm of k9] {peeled 3rd};
    \node[note, text width=3.9cm, align=left, anchor=west] at (8.3,0.4) {On the trace input (Section 5): peel $7, 8, 9$; the survivors are the cycles $\{0,1,2\}$ (length 3) and $\{3,4,5,6\}$ (length 4) $\Rightarrow$ answer 4.};
\end{tikzpicture}
\end{diagrambox}

\begin{complexitybox}
\textbf{Time:} \bigO{n}. \textbf{Space:} \bigO{n} plus the queue. \textbf{Verdict:} \ac{}. It needs two passes, but it separates the two questions cleanly: ``who is on a cycle?'' and ``how long is each cycle?''.
\end{complexitybox}
\end{alternativebox}

\begin{tricksbox}[Approach 4: Implementation Refinements]
\begin{itemize}[leftmargin=*]
    \item \textbf{Timestamps replace three colours.} The single test $\texttt{when}[v] \ge \texttt{start}$ answers both ``visited?'' and ``visited by this walk?'' --- no \texttt{path} vector, no second colouring pass.
    \item \textbf{Iterative, never recursive.} A recursive DFS along the arrows would go $10^5$ levels deep on a long chain; the \texttt{while} loop has no depth at all.
    \item \textbf{Guard $-1$ first.} Both conditions are written \texttt{v != -1 \&\& \ldots}; with the order reversed, \texttt{when[-1]} is read before the test.
    \item \textbf{No overflow worries:} the clock never exceeds $n + 1$.
\end{itemize}
\end{tricksbox}

\subsubsection*{5. Trace}

Our own input with two cycles of different lengths, a tail into the longer one and a chain without a cycle: \;$\texttt{edges} = [1,2,0,4,5,6,3,3,9,-1]$. Answer: $4$.

\begin{diagrambox}
\begin{tikzpicture}
    \foreach \v/\x/\y in {0/0/0.8, 1/1.3/0, 2/0/-0.8, 3/3/0.8, 4/4.4/0.8, 5/4.4/-0.8, 6/3/-0.8, 7/3/2.1, 8/6.2/0, 9/7.5/0}
        \node[gv] (g\v) at (\x,\y) {\v};
    \foreach \u/\w in {0/1, 1/2, 2/0, 3/4, 4/5, 5/6, 6/3, 7/3, 8/9} \draw[gde] (g\u) -- (g\w);
    \draw[-{Stealth[length=2mm]}, warningred, dashed] (g9) -- (8.5,0);
    \node[font=\scriptsize\bfseries, text=warningred, anchor=west] at (8.5,0) {$-1$};
\end{tikzpicture}
\figcap{The trace graph: cycle $0\to1\to2\to0$, cycle $3\to4\to5\to6\to3$ with tail $7 \to 3$, and the chain $8 \to 9 \to -1$.}
\end{diagrambox}

\begin{tracebox}[{Approach 2 on $[1,2,0,4,5,6,3,3,9,-1]$}]
\centering
\begin{tikzpicture}
    \begin{scope}
        \foreach \v/\x/\y/\s/\t in {0/0/0.8/gvseen/1, 1/1.3/0/gvseen/2, 2/0/-0.8/gvseen/3, 3/3/0.8/gv/{}, 4/4.4/0.8/gv/{}, 5/4.4/-0.8/gv/{}, 6/3/-0.8/gv/{}, 7/3/2.1/gv/{}, 8/6.2/0/gv/{}, 9/7.5/0/gv/{}}{
            \node[\s] (a\v) at (\x,\y) {\v};
            \node[gdist, above left=0mm and 0mm of a\v] {\t};
        }
        \foreach \u/\w in {0/1, 1/2, 3/4, 4/5, 5/6, 6/3, 7/3, 8/9} \draw[gde] (a\u) -- (a\w);
        \draw[gback] (a2) -- (a0);
        \node[note, text width=4.3cm, align=left, anchor=west] at (8.3,0.4) {\textbf{$s=0$}, start $= 1$: stamps $0{:}1$, $1{:}2$, $2{:}3$; clock $= 4$; back at $0$ with $\texttt{when}[0] = 1 \ge 1$ $\Rightarrow$ cycle $4 - 1 = 3$. best $= 3$.};
    \end{scope}
    \begin{scope}[yshift=-3.9cm]
        \foreach \v/\x/\y/\s/\t in {0/0/0.8/gvdone/1, 1/1.3/0/gvdone/2, 2/0/-0.8/gvdone/3, 3/3/0.8/gvseen/4, 4/4.4/0.8/gvseen/5, 5/4.4/-0.8/gvseen/6, 6/3/-0.8/gvseen/7, 7/3/2.1/gv/{}, 8/6.2/0/gv/{}, 9/7.5/0/gv/{}}{
            \node[\s] (b\v) at (\x,\y) {\v};
            \node[gdist, above left=0mm and 0mm of b\v] {\t};
        }
        \foreach \u/\w in {0/1, 1/2, 2/0, 3/4, 4/5, 5/6, 7/3, 8/9} \draw[gde] (b\u) -- (b\w);
        \draw[gback] (b6) -- (b3);
        \node[note, text width=4.3cm, align=left, anchor=west] at (8.3,0.4) {\textbf{$s=3$} ($1, 2$ skipped), start $= 4$: stamps $4..7$; clock $= 8$; back at $3$ with $\texttt{when}[3] = 4 \ge 4$ $\Rightarrow$ cycle $8 - 4 = 4$. best $= 4$.};
    \end{scope}
    \begin{scope}[yshift=-7.8cm]
        \foreach \v/\x/\y/\s/\t in {0/0/0.8/gvdone/1, 1/1.3/0/gvdone/2, 2/0/-0.8/gvdone/3, 3/3/0.8/gvdone/4, 4/4.4/0.8/gvdone/5, 5/4.4/-0.8/gvdone/6, 6/3/-0.8/gvdone/7, 7/3/2.1/gvseen/8, 8/6.2/0/gvseen/9, 9/7.5/0/gvseen/10}{
            \node[\s] (c\v) at (\x,\y) {\v};
            \node[gdist, above left=0mm and 0mm of c\v] {\t};
        }
        \foreach \u/\w in {0/1, 1/2, 2/0, 3/4, 4/5, 5/6, 6/3, 8/9} \draw[gde] (c\u) -- (c\w);
        \draw[gcross] (c7) -- (c3);
        \node[note, text width=4.3cm, align=left, anchor=west] at (8.3,0.4) {\textbf{$s=7$}, start $= 8$: stamp $7{:}8$, reach $3$ with $\texttt{when}[3] = 4 < 8$: old walk, nothing. \textbf{$s=8$}, start $= 9$: stamps $8{:}9$, $9{:}10$, then $-1$: no cycle.};
    \end{scope}
\end{tikzpicture}

\medskip
\small
\begin{tabular}{@{}c|cccccccccc@{}}
\toprule
$v$ & 0 & 1 & 2 & 3 & 4 & 5 & 6 & 7 & 8 & 9 \\
\texttt{edges}$[v]$ & 1 & 2 & 0 & 4 & 5 & 6 & 3 & 3 & 9 & $-1$ \\
\texttt{when}$[v]$ & 1 & 2 & 3 & 4 & 5 & 6 & 7 & 8 & 9 & 10 \\
\midrule
walk ($s$ / start) & 0 / 1 & 0 / 1 & 0 / 1 & 3 / 4 & 3 / 4 & 3 / 4 & 3 / 4 & 7 / 8 & 8 / 9 & 8 / 9 \\
\bottomrule
\end{tabular}
\par\smallskip
Result: cycles of length $3$ and $4$ found, \textbf{answer $4$}.
\end{tracebox}

The two official examples run the same way: in $[3,3,4,2,3]$ the walk from $0$ stamps $0{:}1, 3{:}2, 2{:}3, 4{:}4$ and returns to $3$ (time $2 \ge 1$), giving $5 - 2 = 3$; in $[2,-1,3,1]$ the single walk $0,2,3,1$ falls off at $-1$, so the answer stays $-1$.

\subsubsection*{6. Test suite}

\begin{testbox}
\centering\small
\begin{tabular}{@{}lll@{}}
\toprule
\textbf{edges} & \textbf{Expected} & \textbf{Why this test} \\
\midrule
$[3,3,4,2,3]$ & 3 & official example 1: cycle reached through tails \\
$[2,-1,3,1]$ & $-1$ & official example 2: no cycle at all \\
$[1,0]$ & 2 & smallest cycle (self-loops are forbidden) \\
$[-1,-1]$ & $-1$ & no edges \\
$[1,2,0,4,5,6,3,3,9,-1]$ & 4 & two cycles, tail into the longer one (the trace) \\
$\texttt{edges}[i] = (i+1) \bmod n$, $n = 10^5$ & 100000 & one huge cycle: the maximum answer \\
\bottomrule
\end{tabular}
\end{testbox}

All three programs agree on the examples above and on $200\,000$ random graphs with $2 \le n \le 10$.

\subsubsection*{7. Common mistakes}

\begin{warningbox}[Mistake 1 --- a plain \texttt{visited} array]
\begin{lstlisting}[style=snippet]
if (v != -1 && visited[v]) best = max(best, time - when[v]);   // WRONG
if (v != -1 && when[v] >= start) best = max(best, time - when[v]);   // RIGHT
\end{lstlisting}
With only ``visited / not visited'', the walk from $7$ in the trace hits $3$ (visited earlier) and reports a bogus cycle of length $9 - 4 = 5$. Only a vertex of the \emph{current} walk closes a cycle.
\end{warningbox}

\begin{warningbox}[Mistake 2 --- restarting from every vertex]
Running ``walk and detect'' from all $n$ vertices with a fresh visited set each time is the \bigO{n^2} brute force in disguise. One shared \texttt{when} array makes the whole algorithm \bigO{n}.
\end{warningbox}

\begin{warningbox}[Mistake 3 --- counting the tail into the cycle length]
The cycle length is $\texttt{time} - \texttt{when}[v]$ where $v$ is the \emph{repeated} vertex, not $\texttt{time} - \texttt{start}$: in Example 1 the walk from $0$ has $4$ vertices, but vertex $0$ is a tail vertex and the cycle has only $3$.
\end{warningbox}

\begin{warningbox}[{Mistake 4 --- returning $0$ instead of $-1$}]
Initialise \texttt{best} with $-1$: a graph without cycles must return $-1$, and every real cycle has length $\ge 2$, so $-1$ is never confused with an answer.
\end{warningbox}

\subsubsection*{8. Takeaways}

\begin{itemize}[leftmargin=*]
    \item Out-degree $\le 1$ $\Rightarrow$ cycles are vertex-disjoint and their total length is $\le n$: measure them all.
    \item One global clock: ``visited during this walk'' $\Leftrightarrow$ $\texttt{when}[v] \ge \texttt{start}$; cycle length $= \texttt{time} - \texttt{when}[v]$.
    \item Peeling in-degree-0 vertices removes trees \emph{and} sink-chains, leaving only cycles --- a second, equally linear view.
    \item With arbitrary out-degrees this problem changes character completely (see V5): the ``one arrow'' structure is what makes it easy.
\end{itemize}

\begin{variantsbox}
\variant{Permutation: every vertex on a cycle}{\varin{https://leetcode.com/problems/array-nesting/}}
  {\texttt{nums} is a permutation, so there are no tails and no $-1$.}
  {every walk from an unvisited vertex is exactly one whole cycle.}
  {count the length of each walk with a plain \texttt{seen} array; answer the maximum.}
\variant{First repeated vertex from every start}{\varin{https://codeforces.com/problemset/problem/1020/B}}
  {out-degree exactly 1; for every start output where the walk first repeats.}
  {same walks; the answer is the entry point of each start's $\rho$.}
  {cycle vertices answer themselves, tail vertices copy the answer of their successor (Problem 1.5.1).}
\variant{Number of cycles}{\varnone}
  {count the cycles instead of measuring the longest.}
  {same graph and same walks.}
  {increment a counter whenever $\texttt{when}[v] \ge \texttt{start}$; or count the survivor cycles after peeling.}
\variant{Round table of favourites}{\varin{https://leetcode.com/problems/maximum-employees-to-be-invited-to-a-meeting/}}
  {out-degree exactly 1; seat as many people as possible next to their favourite.}
  {either one long cycle alone, or all 2-cycles together, each extended by the longest chain into both ends.}
  {peeling with ``longest chain ending here'' values; answer $\max(\text{longest cycle},\ \sum_{2\text{-cycles}} \text{chains})$ (E5.7).}
\variant{Any number of outgoing edges}{\varnone}
  {a general directed graph: a vertex may have many outgoing edges.}
  {cycles now overlap and walks are no longer forced.}
  {detecting \emph{a} cycle needs DFS with colours (Chapter~2 (DFS)); the \emph{longest} simple cycle is NP-hard --- no efficient algorithm is known.}
\end{variantsbox}

%============================================================
% CH1 · P5 · PROBLEM 3 — FIND THE DUPLICATE NUMBER (LeetCode 287)
%============================================================
\subsection{Problem 1.5.3: Find the Duplicate Number}
\problemlabel{1.5.3}{ssec:ch1-find-duplicate}

\problemcard{1.5.3}{Find the Duplicate Number}{LeetCode 287}{https://leetcode.com/problems/find-the-duplicate-number/}{\exmedium\ \lc{Medium}}{Read the array as arrows $i \to a[i]$; the walk from $0$ is a $\rho$ whose entry point is the duplicate; Floyd's tortoise and hare finds it with \bigO{1} memory}

\textbf{What this problem adds.} Problems~1.5.1 and~1.5.2 marked vertices in a \texttt{visited} or \texttt{when} array. Here that is forbidden: the array may not be modified and only \bigO{1} extra memory is allowed, so Algorithm~4 of the theory (Floyd's tortoise and hare) is no longer an optional curiosity but the whole point. The graph is hidden in an array of numbers, and the answer is the \textbf{entry point} of the $\rho$ that starts at vertex $0$ --- this time we also prove carefully why the second phase of Floyd's method lands exactly there.

\subsubsection*{1. Problem statement}

An array \texttt{nums} has $n+1$ integers, each in the range $[1, n]$. Exactly one value occurs more than once (it may occur two or more times); return that value. You must \textbf{not modify} \texttt{nums} and may use only \textbf{constant extra space}. Follow-up questions: prove that a repeated value must exist, and solve the problem in less than \bigO{n^2} time.

\textbf{Constraints:} $1 \le n \le 10^5$, \; $\texttt{nums.length} = n+1$, \; $1 \le \texttt{nums}[i] \le n$, \; all values except one appear at most once.

\begin{table}[H]
\centering\small
\begin{tabular}{@{}lcl@{}}
\toprule
\textbf{nums} & \textbf{Output} & \textbf{Explanation} \\
\midrule
$[1,3,4,2,2]$ & 2 & $2$ appears at indices $3$ and $4$ \\
$[3,1,3,4,2]$ & 3 & $3$ appears at indices $0$ and $2$ \\
$[3,3,3,3,3]$ & 3 & one value five times; $1, 2, 4$ are missing \\
\bottomrule
\end{tabular}
\end{table}

\begin{diagrambox}
\begin{tikzpicture}
    \begin{scope}
        \node[font=\small\bfseries] at (1.7,0.9) {Example 1: $[1,3,4,2,2]$};
        \node[gvsrc] (a0) at (0,0) {0};
        \node[gv] (a1) at (1.1,0) {1};
        \node[gv] (a3) at (2.2,0) {3};
        \node[gvdst] (a2) at (3.3,0) {2};
        \node[gvseen] (a4) at (3.3,-1.3) {4};
        \draw[gde] (a0) -- (a1); \draw[gde] (a1) -- (a3); \draw[gde] (a3) -- (a2);
        \draw[gde] (a2) to[bend left=35] (a4); \draw[gde] (a4) to[bend left=35] (a2);
    \end{scope}
    \begin{scope}[xshift=5.1cm]
        \node[font=\small\bfseries] at (1.2,0.9) {Example 2: $[3,1,3,4,2]$};
        \node[gvsrc] (b0) at (0,0) {0};
        \node[gvdst] (b3) at (1.2,0) {3};
        \node[gvseen] (b4) at (2.5,0) {4};
        \node[gvseen] (b2) at (1.85,-1.1) {2};
        \node[gvdead] (b1) at (0,-1.1) {1};
        \draw[gde] (b0) -- (b3); \draw[gde] (b3) -- (b4); \draw[gde] (b4) -- (b2); \draw[gde] (b2) -- (b3);
        \draw[gde] (b1) to[out=200, in=250, looseness=6] (b1);
    \end{scope}
    \begin{scope}[xshift=9.9cm]
        \node[font=\small\bfseries] at (1.2,0.9) {Example 3: $[3,3,3,3,3]$};
        \node[gvsrc] (c0) at (0,0) {0};
        \node[gv] (c1) at (0,-1.1) {1};
        \node[gv] (c2) at (2.4,0) {2};
        \node[gv] (c4) at (2.4,-1.1) {4};
        \node[gvdst] (c3) at (1.2,-0.55) {3};
        \foreach \v in {0,1,2,4} \draw[gde] (c\v) -- (c3);
        \draw[gde] (c3) to[out=240, in=300, looseness=6] (c3);
    \end{scope}
\end{tikzpicture}
\figcap{Each index $i$ is a vertex with one arrow to \texttt{nums[i]}. Green: vertex $0$, which no arrow can enter (all values are $\ge 1$). Orange: the duplicate --- in every example it is the vertex where the walk from $0$ enters its cycle. In Example 2 the self-loop at $1$ (grey) is a separate piece that the walk from $0$ never sees.}
\end{diagrambox}

\subsubsection*{2. The graph model}

\begin{modelbox}[Graph Model --- Find the Duplicate Number]
\renewcommand{\arraystretch}{1.35}
\begin{tabularx}{\linewidth}{@{}>{\bfseries}l X@{}}
Vertices & the $n+1$ indices $0..n$ \\
Edges & $i \to \texttt{nums}[i]$; every value is a valid index ($1 \le \texttt{nums}[i] \le n$), so out-degree is exactly 1: a \textbf{functional graph}, stored in the read-only input array \\
Weights & none \\
Question & the entry point of the $\rho$ that starts at vertex $0$ (= the vertex with two arrows into it on that walk) \\
Size & $|V| = n+1 \le 10^5+1$, $|E| = |V|$; memory \bigO{1} besides the input, input read-only; time below \bigO{n^2} \\
\end{tabularx}
\end{modelbox}

\begin{diagrambox}
\begin{tikzpicture}
    \dpindex{ix}{0.75}{$i$}{0}{4}
    \dparray{va}{0}{$a[i]$}{1/cell, 3/cell, 4/cell, 2/anscell, 2/anscell}
    \begin{scope}[xshift=6.2cm]
        \node[gvsrc] (m0) at (0,0) {0};
        \node[gv] (m1) at (1.25,0) {1};
        \node[gv] (m3) at (2.5,0) {3};
        \node[gvdst] (m2) at (3.75,0) {2};
        \node[gvseen] (m4) at (5.25,0) {4};
        \draw[gde] (m0) -- (m1); \draw[gde] (m1) -- (m3);
        \draw[gpath, -{Stealth[length=3mm]}] (m3) -- node[lbl, below=4mm]{$a[3]{=}2$} (m2);
        \draw[gde] (m2) to[bend left=40] (m4);
        \draw[gpath, -{Stealth[length=3mm]}] (m4) to[bend left=40] node[lbl, below]{$a[4]{=}2$} (m2);
        \draw[decorate, decoration={brace, amplitude=4pt}, darkgray] (-0.3,0.55) -- (2.8,0.55)
            node[midway, above=4pt, font=\scriptsize]{tail $\mu = 3$};
        \draw[decorate, decoration={brace, amplitude=4pt}, darkgray] (3.45,0.75) -- (5.55,0.75)
            node[midway, above=4pt, font=\scriptsize]{cycle $\lambda = 2$};
    \end{scope}
\end{tikzpicture}
\figcap{Example 1 as an array and as a $\rho$. The entry point $2$ receives one arrow from the tail ($a[3] = 2$) and one from the cycle ($a[4] = 2$): two different indices hold the value $2$ --- that \emph{is} the duplicate.}
\end{diagrambox}

\subsubsection*{3. Key observations}

\begin{itemize}[leftmargin=*]
    \item \textbf{A duplicate exists (pigeonhole).} $n+1$ values in the $n$ boxes $1..n$: some box gets two. This answers the first follow-up.
    \item \textbf{Vertex $0$ is never on a cycle.} Every cycle vertex has an incoming arrow from its cycle predecessor, but no value equals $0$. So the walk from $0$ has a tail of length $\mu \ge 1$.
    \item \textbf{The entry point has two incoming arrows.} The entry $e$ is entered from the last tail vertex \emph{and} from the last cycle vertex; these are two different indices with value $e$. Since only one value repeats, $e$ is the answer.
    \item \textbf{The start matters.} From a vertex that lies on a cycle (index $1$ in Example 2) the walk has no tail, its first repeated vertex is the start itself, and that value need not be a duplicate. Vertex $0$ is the one start that is \emph{guaranteed} to have a tail.
    \item \textbf{The constraints forbid the earlier tools.} \texttt{first[]}, \texttt{state[]}, \texttt{when[]} all cost \bigO{n} memory, and marking inside \texttt{nums} (e.g.\ negating entries) modifies the input. Only a constant number of vertex variables may be kept: pointers that walk.
\end{itemize}

\begin{diagrambox}
\begin{tikzpicture}
    \begin{scope}
        \node[font=\small\bfseries] at (1.4,1.0) {walk from $0$: tail, then cycle};
        \node[gvsrc] (p0) at (0,0) {0};
        \node[gvdst] (p3) at (1.3,0) {3};
        \node[gvseen] (p4) at (2.7,0.55) {4};
        \node[gvseen] (p2) at (2.7,-0.55) {2};
        \draw[gpath, -{Stealth[length=3mm]}] (p0) -- (p3);
        \draw[gde] (p3) -- (p4); \draw[gde] (p4) -- (p2);
        \draw[gpath, -{Stealth[length=3mm]}] (p2) -- (p3);
        \node[note, text width=4.4cm] at (1.4,-1.4) {$3$ is entered from $0$ (tail) and from $2$ (cycle): $a[0] = a[2] = 3$, \textbf{a duplicate}};
    \end{scope}
    \begin{scope}[xshift=7.6cm]
        \node[font=\small\bfseries] at (0.6,1.0) {walk from $1$: no tail};
        \node[gvdead] (q1) at (0,0) {1};
        \draw[gde] (q1) to[out=30, in=-30, looseness=6] (q1);
        \node[note, text width=4.6cm] at (0.6,-1.4) {the first repeated vertex is the start $1$ itself, entered only by its own arrow: $1$ appears once, \textbf{not} a duplicate};
    \end{scope}
\end{tikzpicture}
\figcap{Example 2 seen from two starts. Only a start outside every cycle --- and $0$ is always one --- produces an entry point with two incoming arrows.}
\end{diagrambox}

\subsubsection*{4. Approaches}

\begin{bruteforcebox}[Approach 1: Compare All Pairs]
Respect both memory rules in the most naive way: compare every pair of positions.

\begin{lstlisting}[style=cpp]
#include <bits/stdc++.h>
using namespace std;

class Solution {
public:
    int findDuplicate(vector<int>& nums) {
        int m = nums.size();              // m = n + 1
        for (int i = 0; i < m; i++)
            for (int j = i + 1; j < m; j++)
                if (nums[i] == nums[j]) return nums[i];   // read-only, O(1) memory
        return -1;                        // unreachable: a duplicate always exists
    }
};
\end{lstlisting}

\begin{complexitybox}
\textbf{Time:} \bigO{n^2} --- about $5\cdot10^9$ comparisons for $n = 10^5$ when the duplicate pair sits at the end. \textbf{Space:} \bigO{1}, input untouched. \textbf{Verdict:} \tle{}; it obeys the memory rules but fails the time follow-up. It serves as the reference for random tests.
\end{complexitybox}
\end{bruteforcebox}

\begin{algorithmbox}[Approach 2: Floyd's Tortoise and Hare from Vertex 0]
Use template (5) of the theory with $s = 0$ and $f(i) = \texttt{nums}[i]$, but return the entry point instead of $\mu$ and $\lambda$.
\begin{enumerate}[leftmargin=*]
    \item \textbf{Phase 1:} tortoise $t$ moves 1 step, hare $h$ moves 2 steps per round, until $t = h$ (a vertex on the cycle).
    \item \textbf{Phase 2:} put $t$ back on $0$; both move 1 step per round; where they meet is the entry point $=$ the duplicate.
\end{enumerate}

\begin{lstlisting}[style=cpp]
#include <bits/stdc++.h>
using namespace std;

class Solution {
public:
    int findDuplicate(vector<int>& nums) {
        // functional graph i -> nums[i] on 0..n; no arrow enters 0
        int t = nums[0], h = nums[nums[0]];   // round 1: tortoise 1 step, hare 2 steps
        while (t != h) {                      // phase 1: meet somewhere on the cycle
            t = nums[t];
            h = nums[nums[h]];
        }
        t = 0;                                // phase 2: tortoise back to the start
        while (t != h) {                      // both 1 step: they meet at the entry
            t = nums[t];
            h = nums[h];
        }
        return t;                             // entry of the cycle = the duplicate value
    }
};
\end{lstlisting}

\begin{complexitybox}
\textbf{Time:} \bigO{n} --- phase 1 takes fewer than $\mu + \lambda \le n+1$ rounds, phase 2 exactly $\mu$ steps (proof below). \textbf{Space:} \bigO{1}: two integers; \texttt{nums} is only read. \textbf{Verdict:} \ac{}, and it meets every requirement of the statement.
\end{complexitybox}
\end{algorithmbox}

\begin{invariantbox}[Why it is correct --- the $\mu$, $\lambda$ argument]
Write $x_k = f^k(0)$ for the vertex after $k$ steps. By the shape theorem $x_0, \dots, x_{\mu+\lambda-1}$ are all different, $x_\mu = e$ is the entry point, and $x_{k+\lambda} = x_k$ for every $k \ge \mu$. A tail vertex ($k < \mu$) appears only once in the whole sequence.

\textbf{(a) The entry point is the duplicate.} $\mu \ge 1$ because $0$ is not on a cycle. The vertex $e$ is the successor of $x_{\mu-1}$ (a tail vertex) and of $x_{\mu+\lambda-1}$ (a cycle vertex; for $\lambda = 1$ it is $e$ itself). These indices differ, and both hold the value $e$. Exactly one value repeats, so it is $e$.

\textbf{(b) Phase 1 meets.} After $t$ rounds the tortoise is at $x_t$ and the hare at $x_{2t}$. For $t < \mu$, $x_t$ is a tail vertex that never appears again, so $x_{2t} \ne x_t$. For $t \ge \mu$ both are on the cycle, and $x_{2t} = x_t$ iff $\lambda$ divides $t$. So the first meeting happens after $t^\ast$ rounds, the smallest positive multiple of $\lambda$ with $t^\ast \ge \mu$; since $\mu \ge 1$, $t^\ast \le \mu + \lambda - 1 \le n$.

\textbf{(c) Phase 2 meets at $e$, and nowhere earlier.} After $k$ more steps the restarted pointer is at $x_k$ and the other one at $x_{t^\ast + k}$. For $k = \mu$: $x_{t^\ast + \mu} = x_\mu$, because $t^\ast$ is a multiple of $\lambda$ and whole laps after the tail change nothing. For $k < \mu$: $x_k$ is on the tail, but $x_{t^\ast + k}$ has index $\ge t^\ast \ge \mu$, so it is on the cycle; they differ. Hence the pointers first coincide after exactly $\mu$ steps, at $x_\mu = e$. \hfill$\blacksquare$
\end{invariantbox}

\begin{diagrambox}
\begin{tikzpicture}
    \node[gvsrc] (z0) at (0,0) {0};
    \node[gv] (z1) at (1.2,0) {};
    \node[gv] (z2) at (2.4,0) {};
    \draw[gde] (z0) -- (z1); \draw[gde, dashed] (z1) -- (z2); \draw[gde] (z2) -- (3.25,0);
    \draw[gpath, -{Stealth[length=3mm]}, shorten >=3.6mm] (3.6,0) arc[start angle=180, end angle=300, radius=1.3];
    \draw[softpurple, line width=1.4pt, -{Stealth[length=3mm]}, shorten >=3.6mm] ($(4.9,0)+(300:1.3)$) arc[start angle=300, end angle=540, radius=1.3];
    \node[gvdst] (ze) at (3.6,0) {$e$};
    \node[gvcur] (zm) at ($(4.9,0)+(300:1.3)$) {$m$};
    \node[lbl, text=accentorange] at ($(4.9,0)+(235:1.75)$) {$d$};
    \node[lbl, text=softpurple] at ($(4.9,0)+(60:1.75)$) {$\lambda - d$};
    \draw[decorate, decoration={brace, amplitude=5pt, mirror}, darkgray] (-0.3,-0.6) -- (3.3,-0.6)
        node[midway, below=5pt, font=\scriptsize]{tail: $\mu$ steps from $0$ to $e$};
    \node[note, text width=6.4cm, align=left, anchor=west] at (6.9,0.1)
        {Phase 1 ends at $m = x_{t^\ast}$, which is $d = t^\ast - \mu$ steps after $e$ on the cycle. Since $\lambda \mid t^\ast$, $\mu \equiv -d \equiv \lambda - d \pmod{\lambda}$: walking $\mu$ steps from $m$ means $\lambda - d$ steps back to $e$ plus whole laps --- and the pointer from $0$ needs the same $\mu$ steps to reach $e$.\\[2pt]
        Trace input of Section 5: $\mu = 4$, $\lambda = 3$, $t^\ast = 6$, $d = 2$, and $4 = (3-2) + 1\cdot 3$.};
\end{tikzpicture}
\figcap{The picture behind phase 2. Orange arc: the $d$ steps from $e$ to the meeting point $m$. Purple arc: the $\lambda - d$ steps from $m$ back to $e$; the remaining $\mu - (\lambda - d)$ steps are whole laps.}
\end{diagrambox}

\begin{alternativebox}[Approach 3: Binary Search on the Value (Counting)]
Forget the graph and apply the pigeonhole principle to a \emph{range of values}. Let $\text{cnt}(x)$ be the number of entries $\le x$ and $d$ the duplicate. Then $\text{cnt}(x) > x$ \textbf{iff} $d \le x$ (proof below), so binary search for the smallest $x$ with $\text{cnt}(x) > x$; each test is one read-only pass over the array.

\begin{lstlisting}[style=cpp]
#include <bits/stdc++.h>
using namespace std;

class Solution {
public:
    int findDuplicate(vector<int>& nums) {
        int n = nums.size() - 1, lo = 1, hi = n;   // the answer lies in [lo, hi]
        while (lo < hi) {
            int mid = (lo + hi) / 2, cnt = 0;
            for (int x : nums) if (x <= mid) cnt++; // read-only counting pass
            if (cnt > mid) hi = mid;                // too many small values: d <= mid
            else lo = mid + 1;                      // otherwise d > mid
        }
        return lo;
    }
};
\end{lstlisting}

\textbf{Why the test is monotone.} If $\text{cnt}(\textit{mid}) > \textit{mid}$, two of those entries share a value in $[1, \textit{mid}]$ (pigeonhole), so $d \le \textit{mid}$. If $\text{cnt}(\textit{mid}) \le \textit{mid}$, the other $n + 1 - \text{cnt} \ge n + 1 - \textit{mid}$ entries lie in $[\textit{mid}+1, n]$, which has only $n - \textit{mid}$ values: two of them collide, so $d > \textit{mid}$. Each step keeps $d$ inside $[lo, hi]$.

\begin{diagrambox}
\begin{tikzpicture}
    \dpindex{bx}{0.75}{$x$}{1}{9}
    \dparray{bc}{0}{$\text{cnt}(x)$}{1/cell, 2/cell, 3/cell, 4/cell, 5/curcell, 6/curcell, 8/anscell, 9/cell, 10/cell}
    \node[note, text width=3.6cm, align=left, anchor=west] at (8.4,0.35) {mid $= 5$: $5 \not> 5$, so $lo = 6$\\ mid $= 7$: $8 > 7$, so $hi = 7$\\ mid $= 6$: $6 \not> 6$, so $lo = 7$};
\end{tikzpicture}
\figcap{Counting on the trace input $[5,3,7,6,9,8,1,4,2,7]$. $\text{cnt}(x) > x$ holds exactly for $x \ge 7$; the three probed values are highlighted, and the search stops at $lo = hi = 7$.}
\end{diagrambox}

\begin{complexitybox}
\textbf{Time:} \bigO{n \log n} --- $\lceil \log_2 n \rceil \le 17$ passes of $n+1$ reads. \textbf{Space:} \bigO{1}, input untouched. \textbf{Verdict:} \ac{}. It satisfies all the rules too, but it is slower than Floyd and uses no graph structure at all; it is the solution to reach for if you do not spot the functional graph.
\end{complexitybox}
\end{alternativebox}

\begin{tricksbox}[Approach 4: Approaches That Break a Rule, and Brent's Variant]
\renewcommand{\arraystretch}{1.25}
\begin{tabularx}{\linewidth}{@{}>{\bfseries}l l X@{}}
Idea & Cost & Which rule it breaks \\
\midrule
\texttt{seen[]} / hash set & \bigO{n} time and memory & \bigO{n} extra memory (Algorithm~1 of the theory from $0$, or a plain ``seen before?'' scan) \\
sort, then neighbours & \bigO{n \log n} & sorting \texttt{nums} modifies it; sorting a copy needs \bigO{n} memory \\
negate \texttt{nums[|v|]} & \bigO{n}, \bigO{1} memory & writes into the input (the visited marks live in the signs) \\
\end{tabularx}

\medskip
\textbf{Brent's algorithm} (in words) is a second \bigO{1}-memory cycle finder. The hare walks one step at a time; the tortoise stays put and \emph{teleports} to the hare's position whenever the number of steps since the last teleport reaches a power of two ($1, 2, 4, \dots$). As soon as the hare lands on the tortoise, the number of steps since the last teleport is exactly $\lambda$. Then start two pointers at $0$, advance one of them $\lambda$ steps, and move both one step at a time: they meet at the entry after $\mu$ steps (the gap of $\lambda$ is one full lap). It usually evaluates $f$ fewer times than Floyd; on this problem it returns the same value (checked on the random tests of Section 6).
\end{tricksbox}

\subsubsection*{5. Trace}

Our own input with a longer tail and a separate cycle: $n = 9$, \;$\texttt{nums} = [5,3,7,6,9,8,1,4,2,7]$. The value $7$ is at indices $2$ and $9$; the walk from $0$ is $0,5,8,2,7,4,9,7,\dots$ with $\mu = 4$ and $\lambda = 3$. The program returns $7$.

\begin{diagrambox}
\begin{tikzpicture}
    \foreach \v/\x/\y in {0/0/0, 5/1.2/0, 8/2.4/0, 2/3.6/0, 7/4.8/0, 4/5.9/0.8, 9/5.9/-0.8, 1/8.4/0.8, 3/9.7/0, 6/8.4/-0.8}
        \node[gv] (g\v) at (\x,\y) {\v};
    \foreach \u/\w in {0/5, 5/8, 8/2, 2/7, 7/4, 4/9, 9/7, 1/3, 3/6, 6/1} \draw[gde] (g\u) -- (g\w);
\end{tikzpicture}
\figcap{The trace graph: tail $0 \to 5 \to 8 \to 2$, cycle $7 \to 4 \to 9 \to 7$ entered at $7$ (from $2$ and from $9$), and a separate cycle $1 \to 3 \to 6 \to 1$ that the walk from $0$ never reaches.}
\end{diagrambox}

\begin{tracebox}[{Phase 1 on $[5,3,7,6,9,8,1,4,2,7]$: T moves 1 step, H moves 2}]
\centering
\begin{tikzpicture}
    \begin{scope}
        \foreach \v/\x/\y/\s in {0/0/0/gvsrc, 5/1.1/0/gv, 8/2.2/0/gvcur, 2/3.3/0/gv, 7/4.4/0/gvseen, 4/5.3/0.75/gv, 9/5.3/-0.75/gv}
            \node[\s] (a\v) at (\x,\y) {\v};
        \foreach \u/\w in {0/5, 5/8, 8/2, 2/7, 7/4, 4/9, 9/7} \draw[gde] (a\u) -- (a\w);
        \node[gdist, above=1mm of a8] {T};
        \node[gdist, text=recursionred, above=1mm of a7] {H};
        \node[note, text width=5.6cm, align=left, anchor=west] at (6.2,0) {\textbf{Round 2:} T $= x_2 = 8$, H $= x_4 = 7$. The hare is already on the cycle, the tortoise still on the tail.};
    \end{scope}
    \begin{scope}[yshift=-2.8cm]
        \foreach \v/\x/\y/\s in {0/0/0/gvsrc, 5/1.1/0/gv, 8/2.2/0/gv, 2/3.3/0/gv, 7/4.4/0/gvcur, 4/5.3/0.75/gvseen, 9/5.3/-0.75/gv}
            \node[\s] (b\v) at (\x,\y) {\v};
        \foreach \u/\w in {0/5, 5/8, 8/2, 2/7, 7/4, 4/9, 9/7} \draw[gde] (b\u) -- (b\w);
        \node[gdist, above=1mm of b7] {T};
        \node[gdist, text=recursionred, right=1mm of b4] {H};
        \node[note, text width=5.6cm, align=left, anchor=west] at (6.2,0) {\textbf{Round 4:} T $= x_4 = 7$ reaches the entry ($\mu = 4$), H $= x_8 = 4$. Both are on the cycle now; the hare closes the gap by one step per round.};
    \end{scope}
    \begin{scope}[yshift=-5.6cm]
        \foreach \v/\x/\y/\s in {0/0/0/gvsrc, 5/1.1/0/gv, 8/2.2/0/gv, 2/3.3/0/gv, 7/4.4/0/gv, 4/5.3/0.75/gv, 9/5.3/-0.75/gvdst}
            \node[\s] (c\v) at (\x,\y) {\v};
        \foreach \u/\w in {0/5, 5/8, 8/2, 2/7, 7/4, 4/9, 9/7} \draw[gde] (c\u) -- (c\w);
        \node[gdist, text=accentorange, right=1mm of c9] {T = H};
        \node[note, text width=5.6cm, align=left, anchor=west] at (6.2,0) {\textbf{Round 6:} T $= x_6 = 9$, H $= x_{12} = 9$: they meet at $m = 9$. Indeed $t^\ast = 6$ is the smallest multiple of $\lambda = 3$ that is $\ge \mu = 4$.};
    \end{scope}
\end{tikzpicture}

\medskip
\small
\begin{tabular}{@{}l|cccccc@{}}
\toprule
round $t$ & 1 & 2 & 3 & 4 & 5 & 6 \\
\midrule
tortoise T $= x_t$ & 5 & 8 & 2 & 7 & 4 & \textbf{9} \\
hare H $= x_{2t}$ & 8 & 7 & 9 & 4 & 7 & \textbf{9} \\
\bottomrule
\end{tabular}
\end{tracebox}

\begin{tracebox}[{Phase 2: T restarts at 0, both move 1 step}]
\centering
\begin{tikzpicture}
    \begin{scope}
        \foreach \v/\x/\y/\s in {0/0/0/gvcur, 5/1.1/0/gv, 8/2.2/0/gv, 2/3.3/0/gv, 7/4.4/0/gv, 4/5.3/0.75/gv, 9/5.3/-0.75/gvseen}
            \node[\s] (d\v) at (\x,\y) {\v};
        \foreach \u/\w in {0/5, 5/8, 8/2, 2/7, 7/4, 4/9, 9/7} \draw[gde] (d\u) -- (d\w);
        \node[gdist, above=1mm of d0] {T};
        \node[gdist, text=recursionred, right=1mm of d9] {H};
        \node[note, text width=5.6cm, align=left, anchor=west] at (6.2,0) {\textbf{Step 0:} T $= 0$, H $= m = 9$. T is $\mu = 4$ steps from the entry; H is $\lambda - d = 1$ step from it, plus one lap of $3$.};
    \end{scope}
    \begin{scope}[yshift=-2.8cm]
        \foreach \v/\x/\y/\s in {0/0/0/gvsrc, 5/1.1/0/gv, 8/2.2/0/gvcur, 2/3.3/0/gv, 7/4.4/0/gv, 4/5.3/0.75/gvseen, 9/5.3/-0.75/gv}
            \node[\s] (e\v) at (\x,\y) {\v};
        \foreach \u/\w in {0/5, 5/8, 8/2, 2/7, 7/4, 4/9, 9/7} \draw[gde] (e\u) -- (e\w);
        \node[gdist, above=1mm of e8] {T};
        \node[gdist, text=recursionred, right=1mm of e4] {H};
        \node[note, text width=5.6cm, align=left, anchor=west] at (6.2,0) {\textbf{Step 2:} T $= 8$ (tail), H $= 4$ (cycle). A tail vertex never equals a cycle vertex, so no false meeting can happen before step $\mu$.};
    \end{scope}
    \begin{scope}[yshift=-5.6cm]
        \foreach \v/\x/\y/\s in {0/0/0/gvsrc, 5/1.1/0/gv, 8/2.2/0/gv, 2/3.3/0/gv, 7/4.4/0/gvdst, 4/5.3/0.75/gv, 9/5.3/-0.75/gv}
            \node[\s] (f\v) at (\x,\y) {\v};
        \foreach \u/\w in {0/5, 5/8, 8/2, 2/7, 7/4, 4/9, 9/7} \draw[gde] (f\u) -- (f\w);
        \node[gdist, text=accentorange, above=1mm of f7] {T = H};
        \node[note, text width=5.6cm, align=left, anchor=west] at (6.2,0) {\textbf{Step 4 $= \mu$:} both at $7$, the entry point. Return \textbf{7} --- indeed $\texttt{nums}[2] = \texttt{nums}[9] = 7$.};
    \end{scope}
\end{tikzpicture}

\medskip
\small
\begin{tabular}{@{}l|ccccc@{}}
\toprule
step $k$ & 0 & 1 & 2 & 3 & 4 \\
\midrule
T $= x_k$ & 0 & 5 & 8 & 2 & \textbf{7} \\
H $= x_{6+k}$ & 9 & 7 & 4 & 9 & \textbf{7} \\
\bottomrule
\end{tabular}
\end{tracebox}

The official examples run the same way. Example 1: phase 1 visits $(t,h) = (1,3), (3,4), (2,4), (4,4)$ and meets at $4$ after $t^\ast = 4$ rounds ($\mu = 3$, $\lambda = 2$); phase 2 gives $(0,4), (1,2), (3,4), (2,2)$: answer $2$. Example 2: $(3,4), (4,3), (2,2)$, then $(0,2), (3,3)$: answer $3$. Example 3: the first round already gives $(3,3)$, then $(0,3), (3,3)$: answer $3$. The counting search probes $\textit{mid} = 2$ ($\text{cnt} = 3$) and $\textit{mid} = 1$ ($\text{cnt} = 1$) in Example 1 and also returns $2$.

\subsubsection*{6. Test suite}

\begin{testbox}
\centering\small
\begin{tabular}{@{}lll@{}}
\toprule
\textbf{nums} & \textbf{Expected} & \textbf{Why this test} \\
\midrule
$[1,3,4,2,2]$ & 2 & official example 1: $\mu = 3$, 2-cycle \\
$[3,1,3,4,2]$ & 3 & official example 2: a separate self-loop at $1$ \\
$[3,3,3,3,3]$ & 3 & official example 3: value repeated 5 times \\
$[1,1]$ & 1 & smallest input ($n = 1$) \\
$[5,3,7,6,9,8,1,4,2,7]$ & 7 & the trace: $\mu = 4$, $\lambda = 3$, $t^\ast = 6$ \\
$a[i] = i+1$ ($i < n$), $a[n] = n$, $n = 10^5$ & 100000 & longest tail: $\mu = 10^5$, $\lambda = 1$ \\
$a[i] = i+1$ ($i < n$), $a[n] = 1$, $n = 10^5$ & 1 & longest cycle: $\mu = 1$, $\lambda = 10^5$ \\
\bottomrule
\end{tabular}
\end{testbox}

Floyd, the counting search and Brent's variant agree with the all-pairs brute force on the examples above and on $200\,000$ random arrays with $n \le 9$ (one value repeated $2..n+1$ times, the others distinct).

\subsubsection*{7. Common mistakes}

\begin{warningbox}[Mistake 1 --- starting the walk at a vertex that may lie on a cycle]
\begin{lstlisting}[style=snippet]
int t = nums[1], h = nums[nums[1]];  ...  t = 1;   // WRONG: vertex 1 may be on a cycle
int t = nums[0], h = nums[nums[0]];  ...  t = 0;   // RIGHT: nothing points to 0
\end{lstlisting}
Starting at $1$ returns $1$ for Example 2 (the self-loop) and $1$ for the trace input (cycle $1 \to 3 \to 6$) --- both wrong. Only vertex $0$ is guaranteed to have a tail, and only an entry point reached from a tail is a duplicate.
\end{warningbox}

\begin{warningbox}[Mistake 2 --- phase 1 and phase 2 from different starts]
\begin{lstlisting}[style=snippet]
int t = 0, h = 0;
do { t = nums[t]; h = nums[nums[h]]; } while (t != h);   // phase 1 from 0
t = nums[0];                                            // WRONG: phase 2 from nums[0]
\end{lstlisting}
The proof needs the same start in both phases. Here the restarted pointer is one step ahead, the two pointers circle the cycle at a constant offset and \textbf{never meet}: the loop runs forever on Examples 1 and 2 and on the trace input. Restart at $0$ --- or start \emph{both} phases at \texttt{nums[0]}, which also works: that is the $\rho$ of vertex \texttt{nums[0]}, with the same entry point.
\end{warningbox}

\begin{warningbox}[Mistake 3 --- \texttt{while} with both pointers at the start]
Writing \texttt{t = h = 0; while (t != h) \{\ldots\}} never enters phase 1, and phase 2 then returns $0$, which is not even a possible value. Either do one round before testing (as in Approach 2) or use \texttt{do \{\ldots\} while}.
\end{warningbox}

\begin{warningbox}[Mistake 4 --- binary search on positions instead of values]
\texttt{nums} is not sorted, so ``look at \texttt{nums[mid]}'' means nothing. Approach 3 searches the \emph{value range} $[1, n]$ and inspects the whole array at every step; the test is $\text{cnt}(\textit{mid}) > \textit{mid}$, with $\le$ inside the count.
\end{warningbox}

\subsubsection*{8. Takeaways}

\begin{itemize}[leftmargin=*]
    \item ``$n+1$ values in $[1, n]$'' $\Rightarrow$ the array is a functional graph on $0..n$ in which $0$ has no incoming arrow: the walk from $0$ always has a tail.
    \item The repeated value is the vertex with two incoming arrows on that walk --- the \textbf{entry point} of the $\rho$.
    \item Floyd: meet after $t^\ast$ rounds with $\lambda \mid t^\ast$; restart one pointer at the \emph{same} start, step both by one, and they meet after exactly $\mu$ steps at the entry. \bigO{n} time, \bigO{1} memory, read-only.
    \item When memory is no issue, the \texttt{first[]} array of Algorithm~1 is simpler; when the input is read-only and memory is \bigO{1}, pointers that walk are the only tool left.
\end{itemize}

\begin{variantsbox}
\variant{Entry of a cycle in a linked list}{\exmedium\ \lc{Medium}\ \varlink{Linked List Cycle II}{LeetCode 142}{https://leetcode.com/problems/linked-list-cycle-ii/}}
  {nodes with \texttt{next} pointers instead of an array; return the node where the cycle begins, or \texttt{null}.}
  {a functional graph with out-degree $\le 1$ (\texttt{null} = no arrow); the start is \texttt{head}, which may itself be on the cycle.}
  {the same two phases from \texttt{head}; phase 1 stops early when the hare reaches \texttt{null} (no cycle). The examples give the nodes at positions $1$, $0$ and \texttt{null}.}
\variant{Happy numbers}{\exeasy\ \lc{Easy}\ \varlink{Happy Number}{LeetCode 202}{https://leetcode.com/problems/happy-number/}}
  {no array at all: $x \mapsto$ sum of the squares of the digits of $x$; is $1$ eventually reached?}
  {an implicit functional graph on the integers; $1$ is a self-loop, so every walk ends in the cycle $\{1\}$ or in another cycle.}
  {Floyd on $f$ with \bigO{1} memory; the answer is ``meeting point $= 1$''. Examples: $19 \to$ true, $2 \to$ false; matches a hash-set walk for all $x \le 10^5$.}
\variant{Cycles in a circular array}{\exmedium\ \lc{Medium}\ \varlink{Circular Array Loop}{LeetCode 457}{https://leetcode.com/problems/circular-array-loop/}}
  {$i \to (i + \texttt{nums}[i]) \bmod n$; is there a cycle of length $> 1$ whose moves all go in the same direction?}
  {a functional graph; walks are cut as soon as the sign changes, so the usable graph has out-degree $\le 1$.}
  {Floyd from every start along same-sign arrows, rejecting self-loops; zero out each explored walk so the total stays \bigO{n}. Examples: true, false, true; matches a brute force on $10^5$ random arrays.}
\variant{Several values may repeat}{\varnone}
  {same array shape, but any number of values may repeat; return \emph{one} repeated value.}
  {unchanged: the walk from $0$ still has a tail, and its entry point still has two incoming arrows.}
  {Floyd still returns a repeated value (the entry point of $0$'s $\rho$); the counting search also still works, because its two pigeonhole arguments never use uniqueness. Both checked on $300\,000$ random arrays.}
\end{variantsbox}

%============================================================
% CH1 · PATTERN 5 — EXERCISES, HINTS, ANSWERS
%============================================================
\subsection{Pattern 5: Exercises and Going Further}
\label{ssec:ch1-p5-ex}

\begin{patternbox}{Pattern 5 Summary Card}
\renewcommand{\arraystretch}{1.25}
\begin{tabularx}{\linewidth}{@{}>{\bfseries}l X@{}}
Recognise & one array \texttt{a[i]} = ``where $i$ goes next''; a deterministic process; ``eventually'', ``forever'', ``after $K$ steps'' \\
Shape & every walk is a $\rho$: tail $\mu$ + cycle $\lambda$, $\mu + \lambda \le n$; every piece has exactly one cycle (none if some \texttt{a[i] = -1}) \\
One start & walk with \texttt{first[v]} = step number; the first repeat gives $\mu$, $\lambda$ and the entry point; with \bigO{1} memory: Floyd's tortoise and hare (\hyperref[ssec:ch1-find-duplicate]{Problem~1.5.3}) \\
All vertices & three states (new / this walk / finished) or visit times $\ge$ start; or peel in-degree-0 vertices \\
Per vertex & cycle vertices first, then tails from the cycle outwards: $\text{ans}[v]$ from $\text{ans}[f(v)]$ \\
Top traps & fresh \texttt{visited} per start ($\bigO{n^2}$); old walk mistaken for a new cycle; \texttt{-1}; huge $K$ in \texttt{int}; deep recursion \\
\end{tabularx}
\end{patternbox}

\exlegend
\begin{enumerate}[leftmargin=*, label=\textbf{E5.\arabic*}, itemsep=6pt]
    \item \exeasy\ \pconcept \exprob{Array Nesting}{LeetCode 565}{https://leetcode.com/problems/array-nesting/}\ \lc{Medium}. \texttt{nums} is a permutation of $0..n-1$. The set $s[k] = \{\texttt{nums}[k], \texttt{nums}[\texttt{nums}[k]], \dots\}$ grows until an element would repeat. Return the largest $|s[k]|$. ($n \le 10^5$.)
    \item \exeasy\ \pconcept \exown{The $\rho$ of a generator.} The sequence $x_0 = 2$, $x_{i+1} = (x_i^2 + 1) \bmod 1000$ must eventually repeat. Find the tail length $\mu$ and the period $\lambda$ without storing more than one array of size $1000$.
    \item \exmedium\ \ptricky \exprob{Find Closest Node to Given Two Nodes}{LeetCode 2359}{https://leetcode.com/problems/find-closest-node-to-given-two-nodes/}\ \lc{Medium}. Out-degree $\le 1$ graph (\texttt{edges[i] = -1} allowed), two start vertices. Among the vertices reachable from both, return the one minimising the larger of the two distances (smallest index on ties), or $-1$.
    \item \exmedium\ \pclassic \exprob{Teleporter}{AtCoder ABC167 D}{https://atcoder.jp/contests/abc167/tasks/abc167_d}. Town $i$'s teleporter sends you to town $A_i$. Starting in town 1, where are you after exactly $K$ teleports? ($N \le 2\cdot10^5$, $K \le 10^{18}$.)
    \item \exmedium\ \pconcept \exown{How many shuffles?} A deck of $n$ cards is shuffled by always moving the card at position $i$ to position $p_i$ ($p$ is a permutation). After how many shuffles is the deck in its original order for the first time? ($p = [1,2,0,4,3] \to 6$.)
    \item \exhard\ \pclassic \cf{1700}\ \exprob{Mouse Hunt}{Codeforces 1027D}{https://codeforces.com/problemset/problem/1027/D}. A mouse in room $i$ moves to room $a_i$ every second; its starting room is unknown. A trap in room $i$ costs $c_i$. Find the minimum total cost that guarantees the mouse is eventually caught. ($n \le 2\cdot10^5$, $c_i \le 10^4$.)
    \item \exhard\ \ptricky \exprob{Maximum Employees to Be Invited to a Meeting}{LeetCode 2127}{https://leetcode.com/problems/maximum-employees-to-be-invited-to-a-meeting/}\ \lc{Hard}. Employee $i$ attends a round-table meeting only if seated next to \texttt{favorite[i]} ($\ne i$). Return the maximum number of employees that can be invited. ($n \le 10^5$.)
    \item \exhard\ \precent \cf{1700}\ \exprob{Medium Demon Problem (easy version)}{Codeforces 2044G1}{https://codeforces.com/problemset/problem/2044/G1}. Each of $n$ spiders starts with one plushie; every year each spider that has a plushie gives one to spider $r_i$ ($r_i \ne i$), all at the same time, and nobody keeps more than one. Year $y$ is stable if the distribution before year $y$ equals the one before year $y-1$. Find the first stable year. ($\sum n \le 2\cdot10^5$.)
\end{enumerate}

\subsubsection*{Hints}
\begin{enumerate}[leftmargin=*, label=\textbf{E5.\arabic*}]
    \item In a permutation every vertex also has in-degree 1, so there are no tails: $s[k]$ is exactly the cycle through $k$.
    \item The function $x \mapsto (x^2+1) \bmod 1000$ is a functional graph on $1000$ implicit vertices. Use Algorithm~1 (template \texttt{rho}).
    \item From each start, the reachable vertices form one path (no branching!). Record the distance along it, stopping at $-1$ or at a repeat.
    \item Store the walk from town 1 until the first repeat. Then $K$ is either inside the tail or can be reduced modulo the cycle length.
    \item Each cycle of $p$ returns to its start after exactly its length many shuffles, and only after multiples of it.
    \item Wherever the mouse starts, it ends up running around one cycle forever. Which rooms really need a trap?
    \item A seating is either one cycle of favourites on its own, or a collection of mutual pairs ($a \leftrightarrows b$), each extended by a chain of fans on both sides.
    \item Who loses the plushie first? A spider nobody gives to. Who next? Think of peeling the graph layer by layer.
\end{enumerate}

\subsubsection*{Answers}

\begin{solutionbox}[E5.1 Array Nesting]
Each walk from an unseen index is a complete cycle; a shared \texttt{seen} array makes the total work \bigO{n}.
\begin{lstlisting}[style=cpp]
#include <bits/stdc++.h>
using namespace std;

class Solution {
public:
    int arrayNesting(vector<int>& nums) {
        int n = nums.size(), best = 0;
        vector<bool> seen(n, false);
        for (int s = 0; s < n; s++) {
            if (seen[s]) continue;        // s already belongs to a measured cycle
            int len = 0, v = s;
            while (!seen[v]) {            // a permutation: the walk from s IS s's cycle
                seen[v] = true;
                v = nums[v];
                len++;
            }
            best = max(best, len);
        }
        return best;
    }
};
\end{lstlisting}
$[5,4,0,3,1,6,2]$: cycles $\{0,5,6,2\}$, $\{1,4\}$, $\{3\}$ $\Rightarrow$ \textbf{4}; \;$[0,1,2]$: three self-loops $\Rightarrow$ \textbf{1}.
\end{solutionbox}

\begin{solutionbox}[E5.2 The $\rho$ of a generator]
The sequence is $2, 5, 26, 677, 330, 901, 802, 205, 26, \dots$: the value $26$ first appears at step 2 and again at step 8, so $\mu = 2$ and $\lambda = 8 - 2 = 6$.
\begin{center}
\begin{tikzpicture}[x=1.25cm]
    \foreach \v/\i/\s in {2/0/gvsrc, 5/1/gv, 26/2/gvdst, 677/3/gvseen, 330/4/gvseen, 901/5/gvseen, 802/6/gvseen, 205/7/gvseen}{
        \node[\s, minimum size=9mm, font=\scriptsize\bfseries] (x\i) at (\i,0) {\v};
        \node[gdist, below=1mm of x\i] {step \i};
    }
    \foreach \i [evaluate=\i as \j using int(\i+1)] in {0,...,6} \draw[gde] (x\i) -- (x\j);
    \draw[gback] (x7.north) to[bend right=30] node[lbl, above]{$205^2 + 1 = 42026 \equiv 26$} (x2.north);
\end{tikzpicture}
\end{center}
\begin{lstlisting}[style=cpp]
#include <bits/stdc++.h>
using namespace std;
typedef long long ll;

int main() {
    const int m = 1000;
    vector<int> first(m, -1);             // implicit vertices 0..m-1, f(x) = (x*x + 1) mod m
    int x = 2, step = 0;
    while (first[x] == -1) {
        first[x] = step++;
        x = (int)(((ll)x * x + 1) % m);
    }
    cout << first[x] << " " << step - first[x] << "\n";   // 2 6
}
\end{lstlisting}
This is the ``$\rho$'' that gives Pollard's rho factorisation method its name. If even $m$ entries are too many, Floyd's tortoise and hare (template 5, used in \hyperref[ssec:ch1-find-duplicate]{Problem~1.5.3}) finds $\mu$ and $\lambda$ with \bigO{1} memory.
\end{solutionbox}

\begin{solutionbox}[E5.3 Find Closest Node to Given Two Nodes]
Out-degree $\le 1$ means the vertices reachable from a start lie on \emph{one} path, so ``BFS'' is just a walk. Compute $d_1$ and $d_2$ by walking (stop at $-1$ or a repeat --- the walk may enter a cycle), then scan $v = 0..n-1$ in increasing order.
\begin{center}
\begin{tikzpicture}
    \node[gvsrc] (a0) at (0,0.6) {0};
    \node[gvsrc] (a1) at (0,-0.6) {1};
    \node[gvdst] (a2) at (1.5,0) {2};
    \node[gv] (a3) at (3,0) {3};
    \draw[gde] (a0) -- (a2); \draw[gde] (a1) -- (a2); \draw[gde] (a2) -- (a3);
    \draw[-{Stealth[length=2mm]}, warningred, dashed] (a3) -- (4,0);
    \node[gdist, above=1mm of a2] {$d_1 = 1, d_2 = 1$};
    \node[gdist, above=1mm of a3] {$d_1 = 2, d_2 = 2$};
\end{tikzpicture}
\end{center}
\begin{lstlisting}[style=cpp]
#include <bits/stdc++.h>
using namespace std;

class Solution {
    vector<int> distFrom(const vector<int>& edges, int s) {
        vector<int> d(edges.size(), -1);  // -1 = not reachable from s
        for (int v = s, k = 0; v != -1 && d[v] == -1; v = edges[v], k++)
            d[v] = k;                     // stop at a dead end or when the walk repeats
        return d;
    }
public:
    int closestMeetingNode(vector<int>& edges, int node1, int node2) {
        vector<int> d1 = distFrom(edges, node1), d2 = distFrom(edges, node2);
        int best = -1, bestVal = INT_MAX;
        for (int v = 0; v < (int)edges.size(); v++) {  // increasing v: ties keep the smaller index
            if (d1[v] == -1 || d2[v] == -1) continue;
            int m = max(d1[v], d2[v]);
            if (m < bestVal) { bestVal = m; best = v; }
        }
        return best;
    }
};
\end{lstlisting}
$[2,2,3,-1]$, $0$, $1$ $\to$ \textbf{2} (distances 1 and 1); \;$[1,2,-1]$, $0$, $2$ $\to$ \textbf{2} (distances 2 and 0). Also checked against a brute force on $20\,000$ random graphs. Without the test \texttt{d[v] == -1} the walk would loop forever on a cycle.
\end{solutionbox}

\begin{solutionbox}[E5.4 Teleporter]
Record \texttt{order[t]} = town after $t$ teleports until the first repeat; with $\mu = \texttt{first}[v]$ and $\lambda = |\texttt{order}| - \mu$, the answer is \texttt{order[K]} if $K < \mu$, else \texttt{order}$[\mu + (K-\mu) \bmod \lambda]$.
\begin{lstlisting}[style=cpp]
#include <bits/stdc++.h>
using namespace std;
typedef long long ll;

int main() {
    int n;
    ll k;
    cin >> n >> k;
    vector<int> a(n + 1);
    for (int i = 1; i <= n; i++) cin >> a[i];

    vector<int> first(n + 1, -1);         // first[v] = step when v was first reached
    vector<int> order;                    // order[t] = town after t teleports
    int v = 1;
    while (first[v] == -1) {
        first[v] = order.size();
        order.push_back(v);
        v = a[v];
    }
    ll mu = first[v];                     // tail length (v = cycle entry)
    ll lam = (ll)order.size() - mu;       // cycle length
    if (k < mu) cout << order[k] << "\n";
    else        cout << order[mu + (k - mu) % lam] << "\n";
}
\end{lstlisting}
Sample 1 ($A = [3,2,4,1]$, $K = 5$): walk $1,3,4$ then $1$ again, so $\mu = 0$, $\lambda = 3$, answer \texttt{order}$[5 \bmod 3] = $ \textbf{4}. Sample 2 ($K \approx 7.3\cdot10^{17}$) prints \textbf{2}. Matches a step-by-step simulation on $5000$ random inputs with small $K$.
\end{solutionbox}

\begin{solutionbox}[E5.5 How many shuffles?]
A card on a cycle of length $L$ is back home after $t$ shuffles iff $L$ divides $t$. All cards are home iff every cycle length divides $t$: the answer is $\operatorname{lcm}$ of the cycle lengths.
\begin{lstlisting}[style=cpp]
#include <bits/stdc++.h>
using namespace std;
typedef long long ll;

ll shufflesToReturn(const vector<int>& p) {
    int n = p.size();
    vector<bool> seen(n, false);
    ll ans = 1;
    for (int s = 0; s < n; s++) {
        if (seen[s]) continue;
        ll len = 0;
        for (int v = s; !seen[v]; v = p[v]) { seen[v] = true; len++; }
        ans = ans / __gcd(ans, len) * len;   // lcm (may overflow for huge n: then work modulo primes)
    }
    return ans;
}

int main() {
    cout << shufflesToReturn({1, 2, 0, 4, 3}) << " " << shufflesToReturn({1, 2, 3, 0, 5, 6, 4}) << "\n";  // 6 12
}
\end{lstlisting}
$[1,2,0,4,3]$ has cycles of lengths 3 and 2 $\Rightarrow$ \textbf{6}; $[1,2,3,0,5,6,4]$ has 4 and 3 $\Rightarrow$ \textbf{12}. Note that the answer can be astronomically large for large $n$ --- real problems then ask for it modulo a prime, and the lcm is assembled from prime powers.
\end{solutionbox}

\begin{solutionbox}[E5.6 Mouse Hunt]
Wherever the mouse starts, it eventually runs around the cycle of its piece forever. A trap on a tail is useless for a mouse that starts on the cycle, while one trap anywhere on the cycle catches every mouse of that piece. So: \textbf{for each cycle, pay its cheapest room}; the answer is the sum. Use the three-state walk (theory, Algorithm 2) and walk once around each new cycle.
\begin{lstlisting}[style=cpp]
#include <bits/stdc++.h>
using namespace std;
typedef long long ll;

int main() {
    ios::sync_with_stdio(false);
    cin.tie(nullptr);
    int n;
    cin >> n;
    vector<int> c(n + 1), a(n + 1);
    for (int i = 1; i <= n; i++) cin >> c[i];
    for (int i = 1; i <= n; i++) cin >> a[i];

    vector<int> state(n + 1, 0);          // 0 new, 1 on the current walk, 2 finished
    ll total = 0;                         // up to 2e5 * 1e4 = 2e9: needs long long
    for (int s = 1; s <= n; s++) {
        if (state[s] != 0) continue;
        vector<int> path;
        int v = s;
        while (state[v] == 0) { state[v] = 1; path.push_back(v); v = a[v]; }
        if (state[v] == 1) {              // new cycle through v: trap its cheapest room
            int best = c[v];
            for (int u = a[v]; u != v; u = a[u]) best = min(best, c[u]);
            total += best;
        }
        for (int u : path) state[u] = 2;
    }
    cout << total << "\n";
}
\end{lstlisting}
The three official samples give \textbf{3}, \textbf{10} and \textbf{2}; the program also matches an exhaustive search over all trap sets on $5000$ random inputs with $n \le 8$. Note the \texttt{long long}: if every room points to itself, there are $2\cdot10^5$ cycles of cost up to $10^4$ each, a total of $2\cdot10^9 > 2^{31}-1$.
\end{solutionbox}

\begin{solutionbox}[E5.7 Maximum Employees to Be Invited to a Meeting]
The favourites form a functional graph. Two kinds of seatings work:
\begin{itemize}[leftmargin=*]
    \item \textbf{one cycle of length $\ge 3$}, seated in cycle order --- nobody else fits (each seat's neighbours are taken);
    \item \textbf{all 2-cycles at once}: for a mutual pair $a \leftrightarrows b$ seat the longest chain of fans ending at $a$, then $a, b$, then the longest chain ending at $b$; such blocks can be placed one after another around the table.
\end{itemize}
Answer $= \max(\text{longest cycle},\ \sum_{\text{2-cycles}} (\text{chain}[a] + \text{chain}[b]))$, where $\text{chain}[v]$ = vertices on the longest arrow chain ending at $v$ (including $v$). Peeling computes $\text{chain}$ because a vertex is always peeled before its successor.
\begin{center}
\begin{tikzpicture}
    \begin{scope}
        \foreach \v/\x/\y/\s in {0/0/0.7/gvsel, 3/1.4/0.7/gvsel, 4/1.4/-0.7/gvsel, 1/0/-0.7/gvsel, 2/-1.4/-0.7/gvdead}
            \node[\s] (e\v) at (\x,\y) {\v};
        \draw[gde] (e0) -- (e3); \draw[gde] (e3) -- (e4); \draw[gde] (e4) -- (e1); \draw[gde] (e1) -- (e0);
        \draw[gde] (e2) -- (e1);
        \node[note, text width=4.6cm] at (0,-1.7) {Example 3, $[3,0,1,4,1]$: one cycle $0 \to 3 \to 4 \to 1 \to 0$ of length 4 fills the table; $2$ cannot also sit next to $1$ $\Rightarrow$ \textbf{4}.};
    \end{scope}
    \begin{scope}[xshift=7.2cm]
        \foreach \v/\x/\y/\s in {4/-1.4/0.7/gvseen, 0/0/0.7/gvsel, 1/1.4/0.7/gvsel, 5/-1.4/-0.7/gvseen, 2/0/-0.7/gvsel, 3/1.4/-0.7/gvsel}
            \node[\s] (h\v) at (\x,\y) {\v};
        \draw[gde] (h4) -- (h0); \draw[gde] (h5) -- (h2);
        \draw[gde] (h0) to[bend left=30] (h1); \draw[gde] (h1) to[bend left=30] (h0);
        \draw[gde] (h2) to[bend left=30] (h3); \draw[gde] (h3) to[bend left=30] (h2);
        \node[note, text width=4.6cm] at (0,-1.7) {Own example $[1,0,3,2,0,2]$: two pairs, each with a fan chain of length 2 on one side: blocks $4,0,1$ and $5,2,3$ $\Rightarrow$ $3 + 3 =$ \textbf{6}.};
    \end{scope}
\end{tikzpicture}
\end{center}
\begin{lstlisting}[style=cpp]
#include <bits/stdc++.h>
using namespace std;

class Solution {
public:
    int maximumInvitations(vector<int>& favorite) {
        int n = favorite.size();
        vector<int> indeg(n, 0), chain(n, 1);  // chain[v]: longest chain ending at v
        for (int v = 0; v < n; v++) indeg[favorite[v]]++;
        queue<int> q;
        for (int v = 0; v < n; v++) if (indeg[v] == 0) q.push(v);
        while (!q.empty()) {                    // peel trees, push chains on
            int v = q.front(); q.pop();
            int w = favorite[v];
            chain[w] = max(chain[w], chain[v] + 1);
            if (--indeg[w] == 0) q.push(w);
        }
        int bigCycle = 0, pairs = 0;
        for (int s = 0; s < n; s++) {
            if (indeg[s] == 0) continue;        // not on a cycle, or already measured
            int len = 0, v = s;
            do { indeg[v] = 0; v = favorite[v]; len++; } while (v != s);
            if (len == 2) pairs += chain[s] + chain[favorite[s]];  // two mutual fans + their tails
            else bigCycle = max(bigCycle, len);
        }
        return max(bigCycle, pairs);
    }
};
\end{lstlisting}
The examples give \textbf{3}, \textbf{3}, \textbf{4}; an exhaustive search over all subsets and circular orders ($n \le 7$, $3000$ random inputs) agrees; the own example $[1,0,3,2,0,2]$ gives \textbf{6}. Why only \emph{one} long cycle but \emph{all} pairs? A cycle of length $\ge 3$ occupies both neighbours of each of its members, so the table is closed; a pair $a \leftrightarrows b$ is already happy with each other, leaving the outer seats free for fans --- and for the next block. This is the hardest exercise of the chapter; the peeling trick (``DP flowing along the trees towards the cycle'') returns in Chapter~11.
\end{solutionbox}

\begin{solutionbox}[E5.8 Medium Demon Problem (easy version)]
A spider has a plushie before year $y \ge 2$ iff someone who had one before year $y-1$ points to it. Cycle spiders always receive from their cycle predecessor, so they never lose theirs. A spider with in-degree 0 loses its plushie in year 1; a spider all of whose givers are gone after $k$ years loses it one year later. This is \textbf{peeling by layers} (theory, Algorithm 3): the spiders of layer $k$ hold a plushie before years $1..k$ and none before year $k+1$. So the distribution changes before every year up to $L+1$, where $L$ is the number of non-empty layers, and the first stable year is $\mathbf{L + 2}$. (Without tails, $L = 0$ and the answer is $2$: year 1 can never be stable.)
\begin{center}
\begin{tikzpicture}
    \foreach \v/\x/\y/\s/\t in {1/0/0.7/gvsel/{}, 2/0/-0.7/gvsel/{}, 4/1.6/-0.7/gvdead/{layer 3}, 3/3.2/-0.7/gvdead/{layer 2}, 5/4.8/-0.7/gvdead/{layer 1}}{
        \node[\s] (m\v) at (\x,\y) {\v};
        \node[gdist, below=1mm of m\v] {\t};
    }
    \draw[gde] (m1) to[bend left=30] (m2); \draw[gde] (m2) to[bend left=30] (m1);
    \draw[gde] (m5) -- (m3); \draw[gde] (m3) -- (m4); \draw[gde] (m4) -- (m2);
    \node[note, text width=4.8cm, align=left, anchor=west] at (5.6,0) {Sample test 3, $r = [2,1,4,2,3]$: three layers $\{5\}, \{3\}, \{4\}$, so the answer is $3 + 2 = 5$, as in the official explanation.};
\end{tikzpicture}
\end{center}
\begin{lstlisting}[style=cpp]
#include <bits/stdc++.h>
using namespace std;

int main() {
    int t;
    cin >> t;
    while (t--) {
        int n;
        cin >> n;
        vector<int> r(n + 1), indeg(n + 1, 0);
        for (int i = 1; i <= n; i++) { cin >> r[i]; indeg[r[i]]++; }
        vector<int> layer;                // current peeling layer
        for (int v = 1; v <= n; v++) if (indeg[v] == 0) layer.push_back(v);
        int layers = 0;
        while (!layer.empty()) {          // layer k loses its plushie in year k
            layers++;
            vector<int> next;
            for (int v : layer)
                if (--indeg[r[v]] == 0) next.push_back(r[v]);
            layer.swap(next);
        }
        cout << layers + 2 << "\n";       // last change in year layers + 1
    }
}
\end{lstlisting}
The official sample (five tests) prints \textbf{2, 2, 5, 4, 5}; a year-by-year simulation agrees on $20\,000$ random inputs with $n \le 9$. The \emph{tricky} part is the model, not the code: the statement talks about plushies and stability, but the answer is just ``number of peeling layers $+\,2$''.
\end{solutionbox}


%============================================================
\clearpage
\section{Chapter Summary \& What Comes Next}
\label{sec:ch1-summary}
%============================================================
%============================================================
% CH1 — CHAPTER SUMMARY & WHAT COMES NEXT   (tag g1s)
%   cheat sheet, the five patterns in pictures, storage picture,
%   "which pattern?" recognition table, complexity tables, mind map,
%   pattern growth map, self-check quiz (answers at the end of the chapter),
%   readiness checklist, final words
%============================================================

\subsection{The Chapter on One Page}
\label{ssec:ch1-summary-main}

\begin{chaptertheorybox}{Cheat Sheet --- Chapter 1}
\begin{center}
\small
\renewcommand{\arraystretch}{1.45}
\begin{tabularx}{\linewidth}{@{}>{\bfseries}l >{\raggedright\arraybackslash}p{3.3cm} >{\raggedright\arraybackslash}X l@{}}
\toprule
Pattern & Graph model & Core move & Cost \\
\midrule
P1 Degrees & edge list, one counter per vertex & \texttt{deg[u]++}, \texttt{deg[v]++} (or in/out for arrows), then a formula on the counters & \bigO{n+m} \\
P2 Reachability & adjacency list & DFS or BFS from the source; mark a vertex \emph{when it is discovered} & \bigO{n+m} \\
P3 Grids & cell $(r,c)$, 4 neighbours & flood fill from every free, unvisited cell; one fill $=$ one region & \bigO{RC} \\
P4 Implicit states & state $=$ number / tuple, edge $=$ one move & BFS over states, neighbours generated on the fly & \bigO{S \cdot b} \\
P5 Functional & \texttt{nxt[v]}, out-degree $1$ & follow arrows; colours $0/1/2$ (or timestamps) reveal the cycle; Floyd's two pointers when memory is \bigO{1} & \bigO{n} \\
\bottomrule
\end{tabularx}
\end{center}

\textbf{The vocabulary in one breath:} $\sum_v \deg(v) = 2|E|$, so the number of odd-degree vertices is even; a tree is connected with $n-1$ edges; a DAG has no directed cycle; every component of a functional graph is one cycle with trees hanging into it.

\textbf{The storage rule:} adjacency list by default; an adjacency matrix only when $n \lesssim 5000$; an undirected edge is stored as two directed edges; reset everything between test cases.

\textbf{The visited rule:} every vertex enters the stack or queue at most once, every edge is looked at a constant number of times $\Rightarrow$ \bigO{|V|+|E|}.

\textbf{The golden rule:} before writing code, fill in the Graph Model --- what is a vertex, what is an edge, directed or not, weighted or not, what is the question in graph language, and how large are $|V|$ and $|E|$?
\end{chaptertheorybox}

\subsection{The Five Patterns in Pictures}

Each panel shows the smallest picture that still explains its pattern. If you can redraw these six panels from memory, you remember the chapter.

\begin{diagrambox}
\begin{tikzpicture}[
    sv/.style={minimum size=6mm, font=\scriptsize\bfseries},
    pnl/.style={draw=patternborder, fill=patternbg!45, rounded corners=3pt, line width=0.8pt},
    ptitle/.style={font=\scriptsize\bfseries\color{patternborder!80!black}},
    pcap/.style={font=\scriptsize, text width=4.5cm, align=center, text=darkgray},
    mv/.style={font=\tiny\bfseries, fill=white, inner sep=0.8pt, text=deepblue}]
    % ---------------- P1 degrees ----------------
    \begin{scope}[shift={(0,4.6)}]
        \draw[pnl] (0,0) rectangle (4.8,4.2);
        \node[ptitle] at (2.4,3.9) {P1 Degrees \& Counting};
        \node[gv, sv] (a) at (0.8,3.15) {1};
        \node[gv, sv] (b) at (4.0,3.15) {2};
        \node[gv, sv] (c) at (0.8,1.65) {4};
        \node[gvdst, sv] (j) at (2.4,2.3) {3};
        \draw[gde] (a) -- (j); \draw[gde] (b) -- (j); \draw[gde] (c) -- (j); \draw[gde] (a) -- (b);
        \node[gdist, anchor=west, text=accentorange] at (2.8,2.0) {in $3$, out $0$};
        \node[pcap] at (2.4,0.55) {One pass counts in/out-degrees;\\judge: in $n{-}1$, out $0$.};
    \end{scope}
    % ---------------- P2 reachability ----------------
    \begin{scope}[shift={(5.1,4.6)}]
        \draw[pnl] (0,0) rectangle (4.8,4.2);
        \node[ptitle] at (2.4,3.9) {P2 Reachability};
        \node[gvsrc, sv] (s) at (0.7,2.9) {0};
        \node[gvdone, sv] (u1) at (1.9,3.15) {1};
        \node[gvdone, sv] (u2) at (1.2,1.7) {2};
        \node[gvdone, sv] (u3) at (2.6,2.1) {3};
        \node[gvdead, sv] (x1) at (3.9,3.0) {4};
        \node[gvdead, sv] (x2) at (4.0,1.7) {5};
        \draw[gtree] (s) -- (u1); \draw[gtree] (s) -- (u2); \draw[gtree] (u1) -- (u3);
        \draw[ge] (u2) -- (u3); \draw[gfaded] (x1) -- (x2);
        \node[pcap] at (2.4,0.55) {DFS/BFS from $0$ with \texttt{visited}:\\$\{0,1,2,3\}$ reached, $\{4,5\}$ not.};
    \end{scope}
    % ---------------- P3 grids ----------------
    \begin{scope}[shift={(10.2,4.6)}]
        \draw[pnl] (0,0) rectangle (4.8,4.2);
        \node[ptitle] at (2.4,3.9) {P3 Grids as Graphs};
        % cell (column c, row r) has its lower-left corner at (0.75 + 0.55c, 2.75 - 0.55r)
        \foreach \c/\r in {0/0,3/0,0/1,2/1,3/1,5/1,2/2,5/2,0/3,1/3,2/3,4/3,5/3}
            \fill[blockedfill] (0.75+0.55*\c, 2.75-0.55*\r) rectangle ++(0.55,0.55);
        \foreach \c/\r in {1/0,2/0,1/1,0/2,1/2}
            \fill[primaryblue!22] (0.75+0.55*\c, 2.75-0.55*\r) rectangle ++(0.55,0.55);
        \foreach \c/\r in {4/0,5/0,4/1,3/2,4/2,3/3}
            \fill[dupfill] (0.75+0.55*\c, 2.75-0.55*\r) rectangle ++(0.55,0.55);
        \foreach \k in {0,...,6} \draw[darkgray!50] (0.75+0.55*\k, 1.1) -- (0.75+0.55*\k, 3.3);
        \foreach \k in {0,...,4} \draw[darkgray!50] (0.75, 1.1+0.55*\k) -- (4.05, 1.1+0.55*\k);
        \node[font=\scriptsize\bfseries, text=primaryblue] at (0.75+0.55*1.5, 2.75-0.55*0.5) {A};
        \node[font=\scriptsize\bfseries, text=recursionred] at (0.75+0.55*4.5, 2.75-0.55*0.5) {B};
        \node[pcap] at (2.4,0.55) {Cells are vertices, sides are edges;\\one flood fill per room: $2$ rooms.};
    \end{scope}
    % ---------------- P4 implicit states ----------------
    \begin{scope}[shift={(0,0)}]
        \draw[pnl] (0,0) rectangle (4.8,4.2);
        \node[ptitle] at (2.4,3.9) {P4 Implicit State Graphs};
        \node[gvsrc, sv] (n4) at (2.4,3.25) {4};
        \node[gvseen, sv] (n8) at (1.3,2.4) {8};
        \node[gvseen, sv] (n3) at (3.5,2.4) {3};
        \node[gv, sv] (n16) at (0.6,1.55) {16};
        \node[gv, sv] (n7) at (1.8,1.55) {7};
        \node[gvdst, sv] (n6) at (3.0,1.55) {6};
        \node[gv, sv] (n2) at (4.2,1.55) {2};
        \draw[gde] (n4) -- node[mv, pos=0.45] {$\times2$} (n8);
        \draw[gpath, -{Stealth[length=2.4mm]}] (n4) -- node[mv, pos=0.45] {$-1$} (n3);
        \draw[gde] (n8) -- (n16); \draw[gde] (n8) -- (n7);
        \draw[gpath, -{Stealth[length=2.4mm]}] (n3) -- node[mv, pos=0.45] {$\times2$} (n6);
        \draw[gde] (n3) -- (n2);
        \node[pcap] at (2.4,0.55) {Numbers $=$ vertices, presses $=$ edges;\\BFS: $4 \to 3 \to 6$, $2$ presses.};
    \end{scope}
    % ---------------- P5 functional ----------------
    \begin{scope}[shift={(5.1,0)}]
        \draw[pnl] (0,0) rectangle (4.8,4.2);
        \node[ptitle] at (2.4,3.9) {P5 Functional Graphs};
        \node[gv, sv] (f1) at (0.6,3.05) {1};
        \node[gv, sv] (f6) at (0.6,1.65) {6};
        \node[gv, sv] (f2) at (1.55,2.25) {2};
        \node[gvsel, sv] (f3) at (2.65,2.25) {3};
        \node[gvsel, sv] (f4) at (3.85,3.0) {4};
        \node[gvsel, sv] (f5) at (3.85,1.6) {5};
        \draw[gde] (f1) -- (f2); \draw[gde] (f6) -- (f2); \draw[gde] (f2) -- (f3);
        \draw[gde, softpurple] (f3) -- (f4); \draw[gde, softpurple] (f4) -- (f5); \draw[gde, softpurple] (f5) -- (f3);
        \node[pcap] at (2.4,0.55) {One arrow leaves every vertex;\\every walk ends in the cycle $3,4,5$.};
    \end{scope}
    % ---------------- the Graph Model ----------------
    \begin{scope}[shift={(10.2,0)}]
        \draw[pnl, draw=definitionborder, fill=definitionbg] (0,0) rectangle (4.8,4.2);
        \node[ptitle, text=definitionborder] at (2.4,3.9) {Before Any Code: the Graph Model};
        \node[font=\scriptsize, text width=4.3cm, align=left, anchor=north west] at (0.25,3.6) {%
            \textbf{1.} What is a vertex?\\
            \textbf{2.} What is an edge? Directed?\\
            \textbf{3.} Weighted? Parallel edges?\\
            \textbf{4.} The question in graph words\\
            \textbf{5.} $|V|$ and $|E|$ from the limits\\
            \textbf{6.} Shape: tree, DAG, grid,\\
            \phantom{\textbf{6.} }out-degree $1$?};
        \node[pcap] at (2.4,0.55) {Fill it in first;\\the pattern usually follows.};
    \end{scope}
\end{tikzpicture}
\figcap{The five patterns of the chapter and the one habit that connects them. Colours: green border $=$ source, orange $=$ target, blue fill $=$ finished, grey $=$ unreachable, purple $=$ on the cycle, thick orange $=$ chosen path.}
\end{diagrambox}

\subsection{Three Ways to Store the Same Graph}

\begin{diagrambox}
\begin{tikzpicture}[
    sv/.style={minimum size=6.5mm, font=\scriptsize\bfseries},
    mc/.style={draw=darkgray!60, minimum size=5.5mm, inner sep=0pt, font=\scriptsize},
    mh/.style={minimum size=5.5mm, inner sep=0pt, font=\scriptsize\bfseries\color{deepblue}},
    hd/.style={font=\scriptsize\bfseries\color{deepblue}}]
    % the graph
    \node[hd] at (1.1,2.75) {the graph};
    \node[gv, sv] (v0) at (0.3,1.9) {0};
    \node[gv, sv] (v1) at (1.9,1.9) {1};
    \node[gv, sv] (v2) at (1.1,0.6) {2};
    \node[gv, sv] (v3) at (2.4,-0.3) {3};
    \draw[ge] (v0) -- (v1); \draw[ge] (v0) -- (v2); \draw[ge] (v1) -- (v2); \draw[ge] (v2) -- (v3);
    \node[font=\tiny, text=darkgray] at (1.1,-0.85) {degrees $2,2,3,1$; sum $8 = 2 \cdot 4$};
    % edge list
    \node[hd] at (4.1,2.75) {edge list};
    \foreach \e [count=\k from 0] in {{(0,1)},{(0,2)},{(1,2)},{(2,3)}}
        \node[mc, minimum width=12mm] at (4.1,1.9-0.6*\k) {\e};
    \node[font=\tiny, text=darkgray] at (4.1,-0.85) {$m$ pairs};
    % adjacency matrix
    \node[hd] at (7.25,2.75) {adjacency matrix};
    \foreach \j in {0,...,3} \node[mh] at (6.35+0.55*\j,2.25) {\j};
    \foreach \i in {0,...,3} \node[mh] at (5.75,1.7-0.55*\i) {\i};
    \foreach \row [count=\i from 0] in {{0,1,1,0},{1,0,1,0},{1,1,0,1},{0,0,1,0}}
        \foreach \x [count=\j from 0] in \row {
            \ifnum\x=1 \node[mc, fill=primaryblue!18, font=\scriptsize\bfseries] at (6.35+0.55*\j,1.7-0.55*\i) {1};
            \else \node[mc, text=midgray] at (6.35+0.55*\j,1.7-0.55*\i) {0}; \fi}
    \node[font=\tiny, text=darkgray] at (7.2,-0.85) {$n \times n$ cells, symmetric};
    % adjacency list
    \node[hd] at (11.6,2.75) {adjacency list};
    \foreach \v/\lst [count=\i from 0] in {0/{1,2}, 1/{0,2}, 2/{0,1,3}, 3/{2}} {
        \node[mh] at (10.1,1.9-0.6*\i) {\v};
        \foreach \w [count=\k from 0] in \lst
            \node[mc, fill=computedfill] at (10.75+0.6*\k,1.9-0.6*\i) {\w};
        \draw[flow] (10.35,1.9-0.6*\i) -- (10.47,1.9-0.6*\i);
    }
    \node[font=\tiny, text=darkgray] at (11.6,-0.85) {$2m = 8$ entries in total};
\end{tikzpicture}
\figcap{One undirected graph, three representations (Lesson~\ref{ssec:ch1-storing}). The adjacency list stores every undirected edge twice, which is the handshake lemma in memory form.}
\end{diagrambox}

\subsection{Which Pattern? A Recognition Table}

Read the statement, find the row whose left column sounds like it, and start from the model in the middle column. The rows are phrased the way contest statements are phrased, not the way textbooks are.

\begin{center}
\small
\renewcommand{\arraystretch}{1.35}
\begin{tabularx}{\linewidth}{@{}>{\raggedright\arraybackslash}p{4.7cm} >{\raggedright\arraybackslash}X c >{\raggedright\arraybackslash}p{2.6cm}@{}}
\toprule
\textbf{The statement says \dots} & \textbf{The graph in disguise} & \textbf{Pattern} & \textbf{Seen in} \\
\midrule
``everybody trusts / beats / points to \dots'', ``find the special person'' & directed graph; the special vertex has a telltale in/out-degree & P1 & Town Judge \\
``rank / importance / score of a pair of cities'' & a formula in $\deg(u)$, $\deg(v)$ and ``is $u$--$v$ an edge?'' & P1 & Network Rank \\
``$n$ vertices and $n-1$ edges'', ``is it a star / a path?'' & a tree; shape tests are degree tests plus at most one traversal & P1 (+P2) & P1.03, P1.08 \\
``is there a route from $A$ to $B$?'', ``can every room be visited?'' & one DFS/BFS from the source, then read \texttt{visited} & P2 & Path Exists, Keys and Rooms \\
``to do $X$ you first need $Y$'', ``which items are needed?'' & reachability along prerequisite arrows & P2 & P1.12 \\
``map of \texttt{.} and \texttt{\#}'', ``islands'', ``rooms'', ``regions'' & cells $+$ 4 neighbours; count the flood fills & P3 & Counting Rooms \\
``largest / total / border of a region'' & a flood fill that also accumulates something & P3 & Max Area of Island \\
``minimum number of operations / presses to turn $x$ into $y$'' & values are vertices, operations are edges; BFS & P4 & Two Buttons \\
``lock'', ``puzzle'', ``configuration'', ``forbidden states'' & tuples encoded as integers; BFS with a banned set & P4 & Open the Lock \\
``each person names exactly one other'', ``$p_i$'', ``teleport $K$ times'' & functional graph: follow the arrows until a repeat & P5 & Badge \\
``at most one outgoing edge'', ``longest cycle'' & functional graph with dead ends; timestamps per walk & P5 & Longest Cycle \\
``$n+1$ values in $[1,n]$'', ``do not modify the array, \bigO{1} extra space'' & functional graph $i \to a_i$ from vertex $0$; Floyd's tortoise and hare finds the entry point & P5 & Find the Duplicate \\
``a permutation applied again and again'' & disjoint cycles (a functional graph without tails) & P5 & P1.21 \\
\bottomrule
\end{tabularx}
\end{center}

\begin{warningbox}[When the table points to the wrong row]
The words in the statement are only a hint; the constraints decide. ``Minimum number of moves'' with $10^{18}$ possible states is not a BFS (look for a formula or a greedy, as in Two Buttons read backwards), and ``islands'' on a $10^5 \times 10^5$ map are not a flood fill. Always finish the Graph Model, including $|V|$ and $|E|$, before trusting a row.
\end{warningbox}

\subsection{Complexity at a Glance}

\begin{center}
\small
\renewcommand{\arraystretch}{1.35}
\begin{tabularx}{\linewidth}{@{}>{\raggedright\arraybackslash}X c c c c@{}}
\toprule
\textbf{Representation} & \textbf{Memory} & \textbf{Neighbours of $v$} & \textbf{Is $(u,v)$ an edge?} & \textbf{OK for $n = 2\cdot 10^5$?} \\
\midrule
Edge list & \bigO{m} & \bigO{m} & \bigO{m} & \yes\ (for P1 counting) \\
Adjacency matrix & \bigO{n^2} & \bigO{n} & \bigO{1} & \no\ ($4\cdot 10^{10}$ cells) \\
Adjacency list & \bigO{n+m} & \bigO{\deg v} & \bigO{\deg v} & \yes \\
\bottomrule
\end{tabularx}

\medskip
\begin{tabularx}{\linewidth}{@{}l >{\raggedright\arraybackslash}X c c@{}}
\toprule
\textbf{Pattern} & \textbf{Algorithm} & \textbf{Time} & \textbf{Memory} \\
\midrule
P1 Degrees & one pass over the edges, then one pass over the vertices & \bigO{n+m} & \bigO{n} \\
P2 Reachability & DFS or BFS with a \texttt{visited} array & \bigO{n+m} & \bigO{n+m} \\
P3 Grids & flood fill over an $R \times C$ grid ($< 2RC$ edges) & \bigO{RC} & \bigO{RC} \\
P4 Implicit states & BFS over $S$ states with $b$ moves each & \bigO{S \cdot b} & \bigO{S} \\
P5 Functional & follow \texttt{nxt}, each vertex finished once & \bigO{n} & \bigO{n} \\
P5, one start & Floyd's tortoise and hare (Find the Duplicate) & \bigO{n} & \bigO{1} \\
\bottomrule
\end{tabularx}
\end{center}

\noindent Rule of thumb: about $10^8$ simple operations per second. For P4, estimate $S$ \emph{before} coding: $10^4$ lock states are nothing, $10^7$ states are fine with an \texttt{int} array, $10^{12}$ states mean BFS is the wrong tool.

\subsection{Mind Map}

\begin{diagrambox}
\begin{tikzpicture}[
    mindmap, grow cyclic, every node/.style={concept, execute at begin node=\hskip0pt},
    concept color=primaryblue!70, text=white, font=\scriptsize\bfseries,
    level 1/.append style={level distance=4.1cm, sibling angle=72, font=\scriptsize\bfseries},
    level 2/.append style={level distance=2.3cm, sibling angle=40, font=\tiny}]
\node[font=\small\bfseries, minimum size=2.6cm] {Graphs}
    child[concept color=successgreen!80!black] { node {Vocabulary}
        child { node {degree, handshake} }
        child { node {path, cycle, component} }
        child { node {tree, DAG, bipartite} } }
    child[concept color=accentorange] { node {Storage}
        child { node {edge list} }
        child { node {matrix} }
        child { node {adjacency list} } }
    child[concept color=softpurple] { node {Traversal}
        child { node {DFS (stack)} }
        child { node {BFS (queue)} }
        child { node {visited $\Rightarrow$ $\bigO{n{+}m}$} } }
    child[concept color=patternborder] { node {5 Patterns}
        child { node {degrees, reach} }
        child { node {grids, states} }
        child { node {functional} } }
    child[concept color=warningred!80!black] { node {Modelling}
        child { node {what is a vertex?} }
        child { node {what is an edge?} }
        child { node {$|V|$, $|E|$ from limits} } };
\end{tikzpicture}
\end{diagrambox}

\subsection{How Each Pattern Grows Through the Book}

Nothing in this chapter is thrown away. Every later technique is one of these five ideas with a weight, a priority queue, a cleverer order or a richer state.

\begin{diagrambox}
\begin{tikzpicture}[x=3.15cm, y=1.25cm,
    pat/.style={rectangle, rounded corners=3pt, draw=patternborder, fill=patternbg, line width=1pt,
                font=\scriptsize\bfseries, align=center, text width=2.2cm, minimum height=9mm},
    lvl/.style={rectangle, rounded corners=2pt, draw=#1, fill=#1!10, font=\tiny, align=center,
                text width=2.45cm, minimum height=8mm}]
    \node[font=\scriptsize\bfseries, color=deepblue] at (0,0.8) {Chapter 1};
    \node[font=\scriptsize\bfseries, color=successgreen] at (1,0.8) {Part I (Ch.~2--6)};
    \node[font=\scriptsize\bfseries, color=softpurple] at (2,0.8) {Part II--III (Ch.~7--17)};
    \node[font=\scriptsize\bfseries, color=warningred] at (3,0.8) {Part IV--VI (Ch.~18--35)};
    \node[pat] (p1) at (0,0)  {P1 Degrees};
    \node[pat] (p2) at (0,-1) {P2 Reachability};
    \node[pat] (p3) at (0,-2) {P3 Grids};
    \node[pat] (p4) at (0,-3) {P4 State graphs};
    \node[pat] (p5) at (0,-4) {P5 Functional};
    \node[lvl=successgreen] (a1) at (1,0)  {in-degrees and Kahn's order (Ch.~5)};
    \node[lvl=successgreen] (a2) at (1,-1) {DFS, components, DSU (Ch.~2, 4, 6)};
    \node[lvl=successgreen] (a3) at (1,-2) {BFS distances on grids (Ch.~3)};
    \node[lvl=successgreen] (a4) at (1,-3) {BFS on states, multi-source BFS (Ch.~3)};
    \node[lvl=successgreen] (a5) at (1,-4) {cycle detection in DFS (Ch.~2)};
    \node[lvl=softpurple] (b1) at (2,0)  {leaves and degrees in trees (Ch.~13)};
    \node[lvl=softpurple] (b2) at (2,-1) {weighted reachability: Dijkstra (Ch.~7)};
    \node[lvl=softpurple] (b3) at (2,-2) {0-1 BFS on grids (Ch.~8)};
    \node[lvl=softpurple] (b4) at (2,-3) {Dijkstra on states, graph construction (Ch.~7, 12)};
    \node[lvl=softpurple] (b5) at (2,-4) {binary lifting, $K$-th successor (Ch.~11)};
    \node[lvl=warningred] (c1) at (3,0)  {Euler parity, degree tricks (Ch.~21, 22)};
    \node[lvl=warningred] (c2) at (3,-1) {strongly connected components (Ch.~19)};
    \node[lvl=warningred] (c3) at (3,-2) {min cut on grids (Ch.~24)};
    \node[lvl=warningred] (c4) at (3,-3) {flow modelling of states (Ch.~24)};
    \node[lvl=warningred] (c5) at (3,-4) {condensation of SCCs (Ch.~19)};
    \foreach \i in {1,...,5}{\draw[flow] (p\i) -- (a\i); \draw[flow] (a\i) -- (b\i); \draw[flow] (b\i) -- (c\i);}
\end{tikzpicture}
\figcap{Roughly: the further right, the higher the rating. Each arrow means ``the same idea with a richer graph or a cleverer order of visiting it''.}
\end{diagrambox}

\subsection{Self-Check Quiz}

Answer without looking back.
\begin{enumerate}[leftmargin=*]
    \item Can a simple undirected graph have the degree sequence $3, 3, 3, 1, 1$? Why or why not?
    \item An undirected graph with $m$ edges is stored as adjacency lists. How many entries do the lists hold in total?
    \item Why is an adjacency matrix out of the question for $n = 2 \cdot 10^5$, and what do you use instead?
    \item Why does DFS or BFS with a \texttt{visited} array run in \bigO{n+m} and not, say, \bigO{n \cdot m}?
    \item In BFS, why is a vertex marked visited when it is \emph{pushed} into the queue rather than when it is popped?
    \item A recursive DFS crashes on a path of $10^6$ vertices although the code is correct. What happened, and what is the fix?
    \item A connected graph has $n$ vertices and $n-1$ edges. What is it? And with exactly $n$ edges?
    \item In Find the Town Judge, why does counting in- and out-degrees suffice --- no traversal at all?
    \item An $R \times C$ grid with 4-neighbour moves: how many vertices and how many edges does its graph have when no cell is blocked?
    \item Open the Lock: how many states are there, and how many neighbours does each state have?
    \item In Two Buttons ($x \to 2x$ or $x \to x-1$, from $n$ to $m$), why does a BFS never need values above $\max(n, 2m)$?
    \item In a functional graph, what does every connected component look like, and why must every walk eventually repeat a vertex?
\end{enumerate}

\noindent\textit{Answers are at the end of the chapter, after the solution sketches (page~\pageref{ssec:ch1-quiz-answers}). Try every question first.}

\subsection{Ready for Chapter 2?}

\begin{checklistbox}[Move on when every box is ticked]
\begin{itemize}[leftmargin=*]
    \item[$\square$] I can define degree, path, cycle, component, tree, DAG, bipartite and functional graph, and draw a small example of each.
    \item[$\square$] I can read an edge list (0- or 1-indexed, directed or undirected, several test cases) into an adjacency list without looking anything up.
    \item[$\square$] I can write recursive DFS, iterative DFS and BFS with a \texttt{visited} array from memory, each in under 5 minutes.
    \item[$\square$] For every problem of this chapter I can fill in the Graph Model: vertices, edges, direction, the question, $|V|$ and $|E|$.
    \item[$\square$] I can turn a grid, a number game or a lock into a graph and bound its number of states before coding.
    \item[$\square$] I can find the cycle that a walk in a functional graph falls into, and the length of that cycle.
    \item[$\square$] I have solved at least 18 of the 26 practice problems of Tiers 1--3, including at least 3 from Tier 3.
\end{itemize}
\end{checklistbox}

\subsection{What Comes Next}

A graph problem is solved in two steps: \emph{see} the graph, then \emph{walk} it. This chapter trained the first step --- vertices and edges hide in friendships, grids, numbers, locks and arrays of pointers --- and gave you the simplest walks: counting degrees, flooding everything reachable, and following single arrows. Chapter~2 (DFS) takes the walk seriously: entry and exit times, the DFS tree, back edges and cycle detection, and how to avoid the recursion-depth trap for good. Chapter~3 (BFS) turns ``can I get there?'' into ``how fast can I get there?'', with distances, multi-source searches and shortest paths on grids and state graphs. Chapter~4 (components, bipartite), Chapter~5 (DAGs) and Chapter~6 (DSU) then complete the foundations on which every later part of the book is built.


%============================================================
\clearpage
\section{Practice Problem Set}
\label{sec:ch1-problems}
%============================================================
%============================================================
% CH1 — PRACTICE PROBLEM SET (28 problems: 3 tiers + Stretch)   (tag g1s)
% Every problem carries a level badge and exactly one kind tag (AUTHORING.md 6b).
% None of these is a solved problem or an exercise anywhere else in the book, and
% no two of them share the same model + algorithm. Every sketch in
% problem_set_solutions.tex was implemented and checked on the official samples.
%============================================================
\subsection{Practice Problems}
\label{ssec:ch1-problem-set}

\begin{insightbox}[How to use this set]
\begin{itemize}[leftmargin=*]
    \item Work tier by tier: \exeasy{} problems should feel routine before you start the \exmedium{} ones; \exhard{} problems (up to \cf{1400}, AtCoder difficulty up to about $1100$) need one non-obvious observation on top of a pattern. The \exstretch{} group is \emph{above} this chapter's level: optional, never needed to continue.
    \item Every tier mixes well-known classics, problems that isolate one idea of the chapter, recent contest problems and problems with a trap --- an edge case, a tempting wrong idea or a twist in the model.
    \item For each problem, fill in the Graph Model (vertices, edges, direction, the question, $|V|$ and $|E|$ --- Lesson~\ref{ssec:ch1-modelling}) \emph{before} coding. Half of these problems never say the word ``graph''.
    \item The hint names the pattern and the key idea. Cover it first; peek only after 20 minutes.
    \item Stuck after an hour? Draw the samples as graphs by hand, write the slowest correct solution, and compare it with your fast one on small random inputs.
    \item Design sketches for all 28 problems follow the set (page~\pageref{ssec:ch1-sketches}) --- compare them with your Graph Model \emph{after} a real attempt.
\end{itemize}
\end{insightbox}

\subsection*{Tier 1 --- Warm-up \exeasy}

\problembox{P1.01}{\exeasy\ \pconcept Counting Roads --- AtCoder ABC061 B}{P1 in its purest form: print every degree. Parallel roads count separately.}{https://atcoder.jp/contests/abc061/tasks/abc061_b}
\problembox{P1.02}{\exeasy\ \precent Adjacency Matrix --- AtCoder ABC343 B}{Storage: turn an adjacency matrix into adjacency lists (print the neighbours of every vertex in increasing order).}{https://atcoder.jp/contests/abc343/tasks/abc343_b}
\problembox{P1.03}{\exeasy\ \ptricky Path Graph? --- AtCoder ABC287 C}{P1 + P2: $M = N-1$, all degrees $\le 2$, and connected. Find a graph that passes the first two tests but is not a path.}{https://atcoder.jp/contests/abc287/tasks/abc287_c}
\problembox{P1.04}{\exeasy\ \pclassic Love Triangle --- Codeforces 939A \cf{800}}{P5: every plane likes exactly one other plane; a love triangle is a cycle of length $3$, i.e.\ $f(f(f(i))) = i$.}{https://codeforces.com/problemset/problem/939/A}
\problembox{P1.05}{\exeasy\ \pconcept Computer Game --- Codeforces 1598A \cf{800}}{P3 with 8 neighbours on a $2 \times n$ grid: the only way to be blocked is a column with two traps.}{https://codeforces.com/problemset/problem/1598/A}
\problembox{P1.06}{\exeasy\ \ptricky Li Hua and Maze --- Codeforces 1797A \cf{800}}{P1 on a grid graph, no search at all: the answer is the smaller of the two cells' degrees ($2$ in a corner, $3$ on a side, $4$ inside).}{https://codeforces.com/problemset/problem/1797/A}
\problembox{P1.07}{\exeasy\ \pclassic Party --- Codeforces 115A \cf{900}}{P5 flavour: ``immediate manager'' arrays are parent pointers of a forest; the answer is the largest depth.}{https://codeforces.com/problemset/problem/115/A}
\problembox{P1.08}{\exeasy\ \precent Begginer's Zelda --- Codeforces 1905B \cf{1100}}{P1 on a tree: one operation can remove at most two leaves; count the leaves.}{https://codeforces.com/problemset/problem/1905/B}

\subsection*{Tier 2 --- Core \exmedium}

\problembox{P1.09}{\exmedium\ \ptricky Validate Binary Tree Nodes --- LeetCode 1361 \lc{Medium}}{P1 + P2: every in-degree $\le 1$, exactly one vertex of in-degree $0$, and that root reaches all $n$ vertices.}{https://leetcode.com/problems/validate-binary-tree-nodes/}
\problembox{P1.10}{\exmedium\ \ptricky Coloring A Border --- LeetCode 1034 \lc{Medium}}{P3: flood fill the component of the start cell, collect its border cells, and recolour only \emph{after} the search.}{https://leetcode.com/problems/coloring-a-border/}
\problembox{P1.11}{\exmedium\ \precent Remove Methods From Project --- LeetCode 3310 \lc{Medium}}{P2: the suspicious methods are those reachable from $k$; they may be removed only if no outside method calls into them.}{https://leetcode.com/problems/remove-methods-from-project/}
\problembox{P1.12}{\exmedium\ \pconcept Martial artist --- AtCoder ABC226 C}{P2 along prerequisite arrows: the moves you must learn are exactly those reachable from move $N$.}{https://atcoder.jp/contests/abc226/tasks/abc226_c}
\problembox{P1.13}{\exmedium\ \pconcept Ladder Takahashi --- AtCoder ABC277 C}{P2 with huge vertex labels: floors up to $10^9$, but only floors touched by a ladder matter --- store the graph in a \texttt{map}.}{https://atcoder.jp/contests/abc277/tasks/abc277_c}
\problembox{P1.14}{\exmedium\ \pconcept Little Pigs and Wolves --- Codeforces 116B \cf{1100}}{P1 on a grid: every pig has at most one wolf neighbour, so count the wolves that have at least one pig neighbour.}{https://codeforces.com/problemset/problem/116/B}
\problembox{P1.15}{\exmedium\ \pclassic Students and Shoelaces --- Codeforces 129B \cf{1200}}{P1 repeated: remove all degree-$1$ vertices \emph{simultaneously}, count the rounds.}{https://codeforces.com/problemset/problem/129/B}
\problembox{P1.16}{\exmedium\ \ptricky Mocha and Hiking --- Codeforces 1559C \cf{1200}}{Draw the graph: a spine $1 \to 2 \to \dots \to n$ plus one extra village. Insert the extra village where an arrow into it is followed by an arrow out of it. Is $-1$ ever the answer?}{https://codeforces.com/problemset/problem/1559/C}
\problembox{P1.17}{\exmedium\ \precent Arrow Path --- Codeforces 1948C \cf{1300}}{P3 + P4: a state is the cell where the robot stands after the forced arrow move; BFS over these at most $2n$ states.}{https://codeforces.com/problemset/problem/1948/C}

\subsection*{Tier 3 --- Challenge \exhard}

\problembox{P1.18}{\exhard\ \precent Path Existence Queries in a Graph I --- LeetCode 3532 \lc{Medium}}{P2 without building the graph: on sorted values, two vertices are connected iff no gap between neighbours on the way exceeds \texttt{maxDiff}.}{https://leetcode.com/problems/path-existence-queries-in-a-graph-i/}
\problembox{P1.19}{\exhard\ \ptricky Grid and Magnet --- AtCoder ABC351 D}{P3: cells next to a magnet can be entered but not left. Flood fill through the free cells; the dead-end cells may belong to several regions.}{https://atcoder.jp/contests/abc351/tasks/abc351_d}
\problembox{P1.20}{\exhard\ \pclassic Road Construction --- Codeforces 330B \cf{1300}}{P1: fewer than $n/2$ forbidden pairs touch fewer than $n$ cities, so some city is in none of them --- build a star around it.}{https://codeforces.com/problemset/problem/330/B}
\problembox{P1.21}{\exhard\ \precent Permutation Game --- Codeforces 1968D \cf{1300}}{P5 on a permutation: walk $i$ steps, then stay forever; only the first $\min(n,k)$ positions of the walk matter.}{https://codeforces.com/problemset/problem/1968/D}
\problembox{P1.22}{\exhard\ \precent Find the Permutation --- Codeforces 2056B \cf{1300}}{The adjacency matrix is a comparator: for two values it tells which one comes first in $p$. Sort with it.}{https://codeforces.com/problemset/problem/2056/B}
\problembox{P1.23}{\exhard\ \pclassic Divide by three, multiply by two --- Codeforces 977D \cf{1400}}{P4 + P5: numbers are vertices, $x \to x/3$ and $x \to 2x$ are arrows; every number has at most one successor, so follow the arrows from the start.}{https://codeforces.com/problemset/problem/977/D}
\problembox{P1.24}{\exhard\ \pclassic Ball in Berland --- Codeforces 1475C \cf{1400}}{P1 on a bipartite graph: for each pair (edge) count the edges that share no endpoint with it: $k - \deg(a) - \deg(b) + 1$.}{https://codeforces.com/problemset/problem/1475/C}
\problembox{P1.25}{\exhard\ \ptricky Web of Lies --- Codeforces 1548A \cf{1400}}{P1 with updates: a noble survives iff he has no friend of higher power; maintain that count per vertex.}{https://codeforces.com/problemset/problem/1548/A}
\problembox{P1.26}{\exhard\ \ptricky Trapped in the Witch's Labyrinth --- Codeforces 2034C \cf{1400}}{P3 backwards: compute the cells that are \emph{forced} to escape (reverse search from the border); every other cell can be trapped.}{https://codeforces.com/problemset/problem/2034/C}

\subsection*{Stretch --- Above This Chapter \exstretch}

\problembox{P1.27}{\exstretch\ \precent Reachability in Functional Graph --- AtCoder ABC357 E}{P5: a cycle vertex reaches exactly its cycle; a tail vertex reaches one more vertex than its successor.}{https://atcoder.jp/contests/abc357/tasks/abc357_e}
\problembox{P1.28}{\exstretch\ \pclassic Lunar New Year and a Wander --- Codeforces 1106D \cf{1500}}{P2 with a twist: replace BFS's queue by a min-heap; always step to the smallest discovered vertex.}{https://codeforces.com/problemset/problem/1106/D}

\begin{tcolorbox}[enhanced, breakable, colback=white, colframe=accentorange, boxrule=1pt, arc=2mm,
    title={Topic index of the practice set --- one idea per problem}, fonttitle=\bfseries, colbacktitle=accentorange]
\small
\renewcommand{\arraystretch}{1.2}
\begin{tabularx}{\linewidth}{@{}>{\bfseries}l X@{}}
P1 Degrees & 01 print all degrees $\cdot$ 02 matrix $\to$ lists $\cdot$ 03 path $=$ three tests $\cdot$ 06 cell degree is the cut $\cdot$ 08 $\lceil \text{leaves}/2 \rceil$ $\cdot$ 09 in-degree tests $+$ one search $\cdot$ 14 count wolves with a pig neighbour $\cdot$ 15 simultaneous leaf peeling $\cdot$ 20 a vertex in no forbidden pair $\to$ star $\cdot$ 24 disjoint edge pairs from degrees $\cdot$ 25 ``no stronger neighbour'' counter under updates \\
P2 Reachability & 11 reach set $+$ check of incoming edges $\cdot$ 12 reach backwards along prerequisites $\cdot$ 13 labels up to $10^9$ in a \texttt{map} $\cdot$ 16 Hamiltonian route by insertion $\cdot$ 18 components on a sorted line, no edges built $\cdot$ 22 matrix as a comparator $\cdot$ 28 min-heap frontier \\
P3 Grids & 05 8 neighbours, blocked column $\cdot$ 10 border of a component $\cdot$ 17 forced arrow moves $\cdot$ 19 dead-end cells shared by regions $\cdot$ 26 reverse search for forced escapes \\
P4 State graphs & 17 state $=$ cell after the arrow $\cdot$ 23 numbers as vertices, one path \\
P5 Functional & 04 a 3-cycle $\cdot$ 07 depth via parent pointers $\cdot$ 21 walk, then stay $\cdot$ 23 out-degree $\le 1$ chain $\cdot$ 27 reach $=$ tail $+$ cycle \\
\end{tabularx}
\end{tcolorbox}

\clearpage
%============================================================
% CH1 — SOLUTION SKETCHES FOR THE PRACTICE SET   (tag g1s)
% Graph model, algorithm and complexity for each problem.
% Sketches, not full code: the point is to check your MODEL.
% Every sketch was implemented and reproduces the official samples.
% Order and groups follow problem_set.tex (Tiers 1-3, then Stretch).
%============================================================
\subsection{Solution Sketches}
\label{ssec:ch1-sketches}

\begin{warningbox}[Read a sketch only after a real attempt]
Each sketch gives the \emph{design} --- the graph model, the algorithm, the complexity and the one trap of the problem --- in a few lines, so you can compare it with your own Graph Model. Writing and debugging the code is left to you. If your model differs but you can argue why it answers the question, it is probably also right.
\end{warningbox}

\begingroup
\providecommand{\sk}{}
\renewcommand{\sk}[2]{\par\smallskip\noindent\textbf{\textcolor{deepblue}{#1}}\quad #2\par}

\begin{solutionbox}[Tier 1 --- Warm-up]
\sk{P1.01 Counting Roads (ABC061 B, difficulty 108).}{For each road $(a,b)$: \texttt{deg[a]++}, \texttt{deg[b]++}; print \texttt{deg[1..N]}. Parallel roads are separate edges (sample~2 prints $5$ and $5$). \bigO{N+M}.}
\sk{P1.02 Adjacency Matrix (ABC343 B, difficulty $<0$).}{Row $i$ of the matrix is vertex $i$'s neighbour set: print every $j$ with $A_{ij} = 1$, in increasing order (an empty line for an isolated vertex, sample~2). This is exactly the conversion from a matrix to adjacency lists, \bigO{N^2} --- the input size.}
\sk{P1.03 Path Graph? (ABC287 C, difficulty 466).}{Yes iff $M = N-1$, every degree is $\le 2$, and one DFS from vertex $1$ reaches all $N$ vertices. \emph{Trap:} all three tests are needed --- a triangle plus an isolated vertex ($N = 4$, $M = 3$) passes the first two and fails the third. \bigO{N+M}.}
\sk{P1.04 Love Triangle (CF 939A).}{Functional graph $i \to f_i$ with $f_i \ne i$. A love triangle is a cycle of length exactly $3$; since $f_i \ne i$, the test $f(f(f(i))) = i$ for some $i$ cannot be fooled by a cycle of length $1$. \bigO{n}.}
\sk{P1.05 Computer Game (CF 1598A).}{Vertices: the safe cells of the $2 \times n$ grid; edges: the 8 neighbour moves. With diagonal moves, two consecutive columns that each have a safe cell are always connected; the path is cut only by a column whose two cells are both traps. Answer \texttt{NO} iff such a column exists (a BFS on the $2n$ cells gives the same answer, checked on random grids). \bigO{n}.}
\sk{P1.06 Li Hua and Maze (CF 1797A).}{Grid graph; $\deg(\text{cell}) = $ number of in-grid side neighbours ($2$, $3$ or $4$). Blocking all neighbours of the start (or of the end) cell disconnects them, so the answer is $\le \min(\deg s, \deg t)$; because $n, m \ge 4$ and the cells are not adjacent, there are that many vertex-disjoint routes, so fewer blocks never suffice. \emph{Trap:} the tempting BFS is impossible ($n, m \le 10^9$) and unnecessary. Answer $\min(\deg s, \deg t)$ in \bigO{1} --- checked against an exhaustive search on all $4 \times 4$ to $5 \times 5$ grids.}
\sk{P1.07 Party (CF 115A).}{Arrow from each employee to his manager ($-1$ for a root): a forest of parent pointers. Employees at the same depth are never superior to one another, and a chain of depth $d$ needs $d$ groups, so the answer is the maximum depth (counting the root as $1$). Walk up from every vertex (\bigO{n^2} with $n \le 2000$) or memoise depths for \bigO{n}. Sample: $3$.}
\sk{P1.08 Begginer's Zelda (CF 1905B).}{One operation lowers the number of leaves by at most two (only the ends of the path can stop being leaves), and contracting the path between two leaves lowers it by exactly two until at most two leaves are left. So with $L$ leaves the answer is $\lceil L/2 \rceil$. Count degree-$1$ vertices: \bigO{n}. Samples: $1, 3, 2, 2$; checked against an exhaustive search over all contraction sequences for $n \le 7$.}
\end{solutionbox}

\begin{solutionbox}[Tier 2 --- Core]
\sk{P1.09 Validate Binary Tree Nodes.}{Arrows from each node to its children. A valid binary tree needs: every in-degree $\le 1$; exactly one vertex with in-degree $0$ (the root); and the root reaches all $n$ vertices. \emph{Trap:} the degree tests alone accept a tree plus a separate cycle such as $2 \to 3 \to 2$, whose vertices all have in-degree $1$. One counting pass and one DFS. \bigO{n}. Samples: \texttt{true}, \texttt{false}, \texttt{false}.}
\sk{P1.10 Coloring A Border.}{Flood fill the component of \texttt{(row, col)} (same colour, 4 neighbours). A component cell is a \emph{border} cell if it lies on the grid edge or has a side neighbour of a different colour. \emph{Trap:} collect the border cells during the search and recolour them afterwards --- recolouring during the search changes the colours that later cells compare against. \bigO{mn}. All three samples reproduced.}
\sk{P1.11 Remove Methods From Project.}{Arrow $a \to b$: ``$a$ invokes $b$''. The suspicious set $S$ is everything reachable from $k$ (one DFS). Then scan all invocations once: if some $a \notin S$ invokes $b \in S$, nothing may be removed and the answer is all methods; otherwise the answer is the methods outside $S$. \bigO{n+m}. Samples: $[0,1,2,3]$, $[3,4]$, $[\,]$.}
\sk{P1.12 Martial artist (ABC226 C, difficulty 539).}{Arrow $i \to A_{i,j}$ (``move $i$ needs move $A_{i,j}$''). The moves that must be learned are exactly those reachable from $N$; learning them in increasing index order respects every prerequisite ($A_{i,j} < i$). Answer: $\sum T_i$ over the reachable set, in \texttt{long long} (up to $2 \cdot 10^{14}$). \bigO{N + \sum K_i}. Samples: $10$, $5000000000$.}
\sk{P1.13 Ladder Takahashi (ABC277 C, difficulty 540).}{Vertices: floor $1$ and every floor touched by a ladder (at most $2N+1$ of them); undirected edge per ladder. Labels go up to $10^9$, so use \texttt{map<int, vector<int>>} (or compress coordinates). DFS from floor $1$; answer the largest visited floor ($1$ if no ladder touches floor $1$). \bigO{N \log N}. Samples: $10$, $12$, $1$.}
\sk{P1.14 Little Pigs and Wolves (CF 116B).}{Bipartite graph wolves--pigs on the grid. Each pig has at most one wolf neighbour, so different wolves never compete for a pig: every wolf with at least one pig neighbour eats one. Answer: the number of wolves with a pig among their 4 neighbours. \bigO{nm}. Samples: $2$, $0$.}
\sk{P1.15 Students and Shoelaces (CF 129B).}{Repeat: collect \emph{all} vertices of degree $1$ first, then delete them and their edges; count the rounds until no vertex has degree $1$. Deleting while scanning would let a vertex that only became a leaf in this round leave too early. At most $n$ rounds of \bigO{n+m} each, fine for $n \le 100$. Samples: $0$, $2$, $1$.}
\sk{P1.16 Mocha and Hiking (CF 1559C).}{Graph: spine $1 \to 2 \to \dots \to n$ and village $n{+}1$ joined to each $i$ by $i \to n{+}1$ ($a_i = 0$) or $n{+}1 \to i$ ($a_i = 1$). If $a_1 = 1$: $n{+}1, 1, \dots, n$. Else if $a_n = 0$: $1, \dots, n, n{+}1$. Else $a$ starts with $0$ and ends with $1$, so some $i$ has $a_i = 0$, $a_{i+1} = 1$: $1, \dots, i, n{+}1, i{+}1, \dots, n$. \emph{Trap:} the statement allows $-1$, but a route always exists --- searching for a Hamiltonian path is never needed. \bigO{n}. Any valid route is accepted; ours was checked by a route validator on the samples and on random inputs.}
\sk{P1.17 Arrow Path (CF 1948C).}{State: the cell where the robot stands after a full second (move, then forced arrow step); start $(1,1)$. From a state, try the 4 neighbours inside the grid and apply the arrow of that neighbour (if it stays inside). BFS over the $\le 2n$ states; answer \texttt{YES} iff $(2,n)$ is reached. \bigO{n}. Samples: \texttt{YES YES NO YES}; counting $(2,n)$ also when it is only touched before the arrow step gives the same answers on random grids.}
\end{solutionbox}

\begin{solutionbox}[Tier 3 --- Challenge]
\sk{P1.18 Path Existence Queries in a Graph I.}{\texttt{nums} is sorted, and an edge joins $i$, $j$ iff $|nums_i - nums_j| \le \texttt{maxDiff}$. If $i < j$ are connected by any path, every gap between consecutive values in $[i, j]$ is $\le \texttt{maxDiff}$, and then consecutive indices are adjacent. So label components left to right: $c_0 = 0$, $c_i = c_{i-1} + [nums_i - nums_{i-1} > \texttt{maxDiff}]$, and answer each query by $c_u = c_v$. The graph may have $\Theta(n^2)$ edges --- never build it. \bigO{n + q}. Samples reproduced; checked against BFS on the explicit graph.}
\sk{P1.19 Grid and Magnet (ABC351 D, difficulty 974).}{Free cells are vertices. A free cell next to a magnet is a \emph{dead end}: it can be entered but not left. Flood fill from every free non-dead-end cell not yet in a region, expanding only non-dead-end cells, and count the cells reached. \emph{Trap:} one dead end can touch several regions, so mark visits with a region stamp, not with one global \texttt{visited} flag; and a lonely dead-end cell gives freedom $1$. \bigO{HW}. Samples: $9$, $1$; checked against a BFS from every cell.}
\sk{P1.20 Road Construction (CF 330B).}{Any valid network is connected, so it needs $\ge n-1$ roads. The $m < n/2$ forbidden pairs touch at most $2m < n$ cities, so some city $c$ appears in none of them. The star with centre $c$ uses $n-1$ allowed roads and any two cities are at distance $\le 2$. \bigO{n+m}. Sample: $3$ roads (a different star is also accepted).}
\sk{P1.21 Permutation Game (CF 1968D).}{A player at $P$ visits $v_0 = P, v_1 = p_{v_0}, \dots$ (a cycle of the permutation). Optimal play: move for $i$ turns, then stay on $v_i$ for the remaining $k - i$ turns: $\text{score}_i = \sum_{j<i} a_{v_j} + (k-i)\, a_{v_i}$. Positions repeat after at most $n$ moves, so take $\max$ over $i < \min(n, k)$, in \texttt{long long} (up to $10^{18}$). Compare the two players. \bigO{n} per test. All $10$ sample verdicts reproduced; also checked against trying all $2^k$ stay/move choices on small inputs.}
\sk{P1.22 Find the Permutation (CF 2056B).}{For values $a < b$, the edge $a$--$b$ exists iff $a$ stands before $b$ in $p$; without the edge, $b$ stands before $a$. So the matrix decides the order of every pair: sort the values $1..n$ with this comparator (it is consistent because it comes from a real permutation). \bigO{n^2}, the input size. Samples: $1$; $4\,2\,1\,3\,5$; $6\,5\,4\,3\,2\,1$; checked by building the graph of random permutations and recovering them.}
\sk{P1.23 Divide by three, multiply by two (CF 977D).}{Vertices: the numbers; arrows $x \to x/3$ and $x \to 2x$ when present. No number can have both successors in the list (reaching the other one later would need $2^a = 3^b$ with $a + b > 0$), so the graph is a single path: start at the number with in-degree $0$ and follow the arrows. \bigO{n^2} with $n \le 100$, or a \texttt{set} for \bigO{n \log n}; values reach $3 \cdot 10^{18}$, so \texttt{long long}. Equivalent shortcut: sort by (more factors $3$ first, then smaller value). Samples reproduced.}
\sk{P1.24 Ball in Berland (CF 1475C).}{Bipartite graph boys--girls with the $k$ pairs as edges. The edges that share an endpoint with edge $(a, b)$, itself included, number $\deg(a) + \deg(b) - 1$; the remaining $k - \deg(a) - \deg(b) + 1$ are compatible. Summing over all edges counts every compatible pair twice: answer $= \frac12 \sum_e (k - \deg a - \deg b + 1)$, in \texttt{long long}. \bigO{a + b + k}. Samples: $4$, $0$, $2$; checked against all pairs on random inputs.}
\sk{P1.25 Web of Lies (CF 1548A).}{Claim: the survivors are exactly the nobles with no friend of higher power. Such a noble is never vulnerable (his friends only disappear). A noble $v$ with a stronger friend $w$ dies: $w$ cannot die while $v$ is alive (then $w$ has a weaker living friend), and as long as some living noble has a stronger living friend, the weakest such noble has only stronger living friends, so the process does not stop. \emph{Trap:} simulating the process per query is far too slow. Maintain $\text{cnt}[v]$ = number of stronger friends: on adding $(u, v)$ increment $\text{cnt}[\min(u,v)]$, on removing decrement it, and keep the number of $v$ with $\text{cnt}[v] = 0$ up to date. \bigO{1} per query. Samples: $2, 1$ and $1$; checked against simulating the process on random small inputs.}
\sk{P1.26 Trapped in the Witch's Labyrinth (CF 2034C).}{Compute the set $E$ of cells that \emph{must} escape, by a reverse search from the border: a fixed arrow is in $E$ if it points outside or into $E$; a \texttt{?} is in $E$ only if all four neighbours are outside or in $E$. Every other cell can be trapped: point a \texttt{?} at a neighbour not in $E$ (two neighbouring \texttt{?} cells point at each other). Answer $nm - |E|$. \emph{Trap:} a lonely \texttt{?} surrounded by escaping cells cannot be saved. \bigO{nm}. Samples: $0$, $6$, $5$; checked against trying every assignment of the \texttt{?} cells on small grids.}
\end{solutionbox}

\begin{solutionbox}[Stretch --- above this chapter]
\sk{P1.27 Reachability in Functional Graph (ABC357 E, difficulty 1295).}{Each component is a cycle with trees hanging into it. A cycle vertex reaches exactly the $L$ vertices of its cycle; a tail vertex $u$ reaches itself plus everything $a_u$ reaches: $r(u) = 1 + r(a_u)$. Find the cycles with colours $0/1/2$ (a walk that meets a vertex of colour $1$ has closed a new cycle), then fill the tails by walking until a known $r$ and assigning backwards. Answer $\sum_u r(u)$, up to $4 \cdot 10^{10}$: \texttt{long long}. \bigO{N}. Samples: $8$, $14$, $41$; checked against brute force.}
\sk{P1.28 Lunar New Year and a Wander (CF 1106D).}{Bob may walk back through visited vertices, so at any moment he can record \emph{any} discovered but unrecorded vertex. Greedy: keep the frontier in a min-heap, always record the smallest, then push its undiscovered neighbours (marked on push, as in BFS). Exchange argument: recording anything else first makes the sequence lexicographically larger. \bigO{m \log n}. Samples reproduced.}
\end{solutionbox}

%------------------------------------------------------------
\subsection{Answers to the Self-Check Quiz}
\label{ssec:ch1-quiz-answers}

\begin{solutionbox}[Quiz answers]
\begin{enumerate}[leftmargin=*]
    \item No. The degrees sum to $11$, but every graph has $\sum_v \deg(v) = 2|E|$, an even number (equivalently: it would have three odd-degree vertices, an odd number).
    \item $2m$ entries: every undirected edge appears in the lists of both endpoints.
    \item The matrix has $n^2 = 4 \cdot 10^{10}$ cells --- tens of gigabytes even at one byte per cell, and \bigO{n^2} time just to clear it. Use adjacency lists: \bigO{n+m} memory.
    \item Thanks to \texttt{visited}, every vertex is pushed and processed at most once, and processing $v$ scans its list once; the scans add up to $\sum_v \deg(v) = 2m$ (or $m$ for a directed graph). Total \bigO{n+m}.
    \item Marking on pop lets the same vertex be pushed once per neighbour that sees it before it is popped, so the queue can hold \bigO{m} copies and work is repeated. Marking on push guarantees each vertex enters the queue exactly once (and, for distances, fixes its distance at the first, smallest value).
    \item Every recursive call takes a stack frame; $10^6$ nested calls exceed the default stack (a few megabytes), so the program dies with a stack overflow. Use an iterative DFS with an explicit \texttt{stack} (or BFS), or make sure the depth is small.
    \item With $n-1$ edges: a tree. With exactly $n$ edges: a tree plus one edge, i.e.\ a connected graph with exactly one cycle.
    \item The judge is defined by degrees only: trusted by all $n-1$ others (in-degree $n-1$) and trusting nobody (out-degree $0$). At most one vertex can satisfy both (two judges would have to trust each other), so one counting pass and one check per vertex suffice.
    \item $RC$ vertices; $R(C-1)$ horizontal plus $C(R-1)$ vertical edges, i.e.\ $2RC - R - C$ edges --- fewer than $2RC$.
    \item $10^4$ states (\texttt{0000} to \texttt{9999}); each has $8$ neighbours (4 wheels, each turned up or down).
    \item From a value $y \ge m$ at least $y - m$ presses are needed (a press lowers the value by at most $1$), and $y - m$ presses of the ``$-1$'' button achieve it, so an optimal sequence never doubles a value $\ge m$. Doubling therefore only happens to values $< m$, giving values $< 2m$, and otherwise values only go down from $n$. Hence every value on an optimal path is $\le \max(n, 2m)$.
    \item One directed cycle with trees hanging into it, all arrows of the trees pointing towards the cycle (the ``$\rho$'' shape). Every vertex has an outgoing arrow, so a walk never stops; with only $n$ vertices, by the pigeonhole principle some vertex repeats within $n+1$ steps, and from then on the walk goes round the cycle forever.
\end{enumerate}
\end{solutionbox}
\endgroup

